\xchapter{Questions and Answers on Genesis, I}

\emph{[Yonge's title, A Volume of Questions, and Solutions to Those Questions, Which Arise in Genesis. Yonge had access
only to Aucher's Latin translation of Aucher's Armenian version (J. B. Aucher's Philonis Judaei paralipomena Armena: libri
videlicet quatuor in Genesin, libri duo in Exodum, Venice, 1826). Apparently Yonge does not have access to Questions
and Answers on Exodus in either Latin or Armenian, although fragments of Questions and Answers on Exodus appear in Yonge's edition in ``Fragments, Extracted from the Parallels of John of Damascus.'']}

(1) Why does Moses, revolving and considering the creation of the world, say: ``This is the book
of the generation of heaven and earth, when they were created?'' (\#Ge 2:4). The expression,
``when they were created,'' indicates as it seems an indeterminate time not accurately described.
But this argument will confute those authors who calculate a certain number of years reduced to
one, from the time when it is possible that the world may have been created. And again, the
expression: ``This is the book of the generation,'' is as it were indicative of the book as it follows,
which contains an account of the creation of the world; in which it is intimated that what has been
related about the creation of the world is consistent with strict truth.

(2) What is the object of saying, ``And God made every green herb of the field, before it was upon
the earth, and every grass before it had sprung up?'' (\#Ge 2:5). He here by these expressions
intimates in enigmatical language the incorporeal species; since the expression, ``before it was
upon the earth,'' indicates the arriving at perfection of every herb, and of all seeds and trees. But
as to what he says, that ``before it had sprung up upon the earth,'' he had made every green herb,
and grass, et caetera, it is plain that the incorporeal species, as being indicative of the others,
were created first, in accordance with intellectual nature, which those things which are upon the
earth perceptible to the outward senses were to imitate.

(3) What is the meaning of saying: ``A fountain went up from the earth, and watered all the face of
the earth?'' (\#Ge 2:6). But here the question is how it could be that the whole earth was
watered by one fountain, not only on account of its size, but also because of the inequality of the
mountainous and champaign situations? Unless, indeed, just as the whole force of the king's
cavalry is called ``the horse,'' so the whole multitude of the veins of the earth which supply drinkable
water, may perhaps be called the fountain, inasmuch as they all bubble up like a fountain. And that
expression is peculiarly appropriate which says that the fountain watered, not the whole earth, but
its face; as in the living being it waters the chief and predominant part (the mind or the
countenance). Since that is the most important part of the earth which can be good and fertile and
productive, and that is the part which stands in need of the nourishment of fountains.

(4) What is the man who was created? And how is that man distinguished who was made after the
image of God? (\#Ge 2:7). This man was created as perceptible to the senses, and in the
similitude of a Being appreciable only by the intellect; but he who in respect of his form is
intellectual and incorporeal, is the similitude of the archetypal model as to appearance, and he is
the form of the principal character; but this is the word of God, the first beginning of all things, the
original species or the archetypal idea, the first measure of the universe. Moreover, that man who
was to be created as a vessel is formed by a potter, was formed out of dust and clay as far as his
body was concerned; but he received his soul by God breathing the breath of life into his face, so
that the temperament of his nature was combined of what was corruptible and of what was
incorruptible. But the other man, he who is only so in form, is found to be unalloyed without any
mixture proceeding from an invisible, simple, and transparent nature.

(5) Why is it said that God breathed into his face the breath of life? (\#Ge 2:7). In the first place
because life is the principal part of the body; for the rest was only made as a sort of foundation or
pedestal, and then life was put upon it as a statue. Besides, the sense is the fountain of the animal
form, and sense resides in the face. Secondly, man is created to be a partaker not only of a soul
but also of a rational soul; and the head is the temple of the reason, as some writers have called it.

(6) Why is God said to have planted a Paradise? And for whom? And what is meant by a
paradise? (\#Ge 2:8). The word paradise, if taken literally, has no need of any particular
explanation; for it means a place thickly crowded with every kind of tree; but symbolically taken, it
means wisdom, intelligence both divine and human, and the proper comprehension of the causes
of things; since it was proper, after the creation of the world, to establish a contemplative system of
life, in order that man, by the sight of the world and of the things which are contained in it, might be
able to attain to a correct notion of the praise due to the Father. And since it was not possible for
him to behold nature herself, nor properly to praise the Creator of the universe without wisdom,
therefore the Creator planted the outlines of it in the rational soul of the principal guide of man,
namely the mind, as he planted trees in the paradise. And when we are told that in the middle was
the tree of life, that means the knowledge not only of the creature, but also of the greater and
supreme cause of the universe; for if any one is able to arrive at a certain comprehension of that,
he will be fortunate and truly happy and immortal. Moreover, after the creation of the world human
wisdom was created, as also after the creation of the world the Paradise was planted; and so the
poets say that the chorus of musicians was established in order to praise the Creator and his
works; as Plato says, that the Creator was the first and greatest of causes, and that the world was
the most beautiful of all creatures.

(7) Why in Adin, or Eden, is God said to have planted the Paradise towards the east? (Genesis
2:8). This is said in the first place because the motion of the world proceeds from the rising of the
sun to its setting. And it first exists in that quarter from which it is moved; secondly, because that
part of the world which is in the region of the east is called the right side; and that which is in the
region of the west is called the left side of the world. Moreover the poet bears witness to this,
calling the birds from the east dexteras or right, and those on the west sinistras or left; \footnote{he is
referring here to Homer (as Pope translates it)--''Ye vagrants of the sky! your wings extend, / Or where the suns arise, or
where descend; / To right, to left, unheeded take your way, / While I the dictates of high Heaven obey.''} when he says,
whether they go to the right to the day and to the sun, or whether they go to the left towards the
dusky evening. But the name Eden, when rightly understood, is an indication of all kinds of delights,
and joys, and pleasures; since all good things and all blessings derive their beginning from the
place of the Lord. Thirdly, because wisdom itself is splendour and light.

(8) Why did God place man whom he had created in the Paradise, but not that man who is after his
own image? (\#Ge 2:15). Some persons have said, when they fancied that the Paradise was a
garden, that because the man who was created was endowed with senses, therefore he naturally
and properly proceeded into a sensible place; but the other man, who is made after God's own
image, being appreciable only by the intellect, and invisible, had all the incorporeal species for his
share; but I should rather say that the paradise was a symbol of wisdom, for that created man is a
kind of mixture, as having been compounded of soul and body, having work to do by learning and
discipline; desiring according to the law of philosophy that he may become happy; but he who is
according to God's own image is in need of nothing, being by himself a hearer, and being taught by
himself, and being found to be his own master by reason of his natural endowments.

(9) Why does Moses say that every tree in the Paradise was beautiful to look upon and good to
eat? (\#Ge 2:9). He says this because the virtue of trees is of a twofold nature, consisting in
bearing leaves and fruit, one of which qualities is referred to the pleasing of the sight, the other to
the gratification of the taste; but the word beautiful was not employed inappropriately. Indeed it is
very proper that the plants should be always green and flourishing perpetually, as belonging to a
divine Paradise, which as such must be everlasting; and it is fit too that they should never
degenerate so as to lose their leaves. But of the fruit he says, not that it was beautiful, but that it
was good, speaking in a very philosophical spirit; since men take food, not only because of the
pleasure which it affords, but also because of its use; and use is the flowing forth and imparting of
some good.

(10) What is meant by the tree of life, and why it was placed in the middle of the Paradise?
(\#Ge 2:9). Some people have believed that, if there were really plants of a corporeal and
deadly nature, there are also some which are causes of life and immortality, because, they say, life
and death are opposed to one another, and because some plants are ascertained to be
unwholesome, therefore of necessity there must be others from which health may be derived. But
what these are which are wholesome they know not; for generation, as the opinion of the wise has
it, is the beginning of corruption. But perhaps we ought to look on these things as spoken in an
allegorical sense; for some say the tree of life belongs to the earth, inasmuch as it is the earth
which produces everything which is of use for life, whether it be the life of mankind or of any other
animal; since God has appointed the situation in the centre for this plant, and the centre of the
universe is the earth. There are others who assert that what is meant by the tree of life is the
centre between the seven circles of heaven; but some affirm that it is the sun which is meant, as
that is nearly in the centre, between the different planets, and is likewise the cause of the four
seasons, and since it is owing to him that every thing which exists is called into existence. Others
again understand by the tree of life the direction of the soul, for this it is which renders the sense
nervous and solid, so as to produce actions corresponding to its nature, and to the community of
the parts of the body. But whatever is in the middle is in a manner the primary cause and beginning
of things, like the leader of a chorus. But still, the best and wisest authorities have considered that
by the tree of life is indicated the best of all the virtues of man, piety, by which alone the mind
attains to immortality.

(11) What is meant by the tree of the knowledge of good and evil? (\#Ge 2:9). This indeed
exhibits that meaning which is sought for in the letter of the Scriptures more clearly to the sight, as
it bears a manifest allegory on the face of it. What is meant then under this figure is prudence,
which is the comprehension of science, by which all things are known and distinguished from one
another, whether they be good and beautiful, or bad and unseemly, or in short every sort of
contrariety is discerned; since some things belong to the better class, and some to the worse.
Therefore the wisdom which exists in this world is not in truth God himself, but the work of God; that
it is which sees and thoroughly investigates every thing. But the wisdom which exists in man sees
in an incorrect and mixed manner with somewhat darkened eyes; for it is found to be incompetent
to see and comprehend clearly and without alloy each particular thing separately. Moreover, there
is a kind of deception mingled with human wisdom; since very often there are some shadows found
which hinder the eyes from contemplating a brilliant light; since what the eye is in the body, such
also is the mind and wisdom in the soul.

(12) What the river is which proceeded out of Adin by which the Paradise is watered, and from
which the four rivers proceed, the Phison, and the Gihon, and the Tigris, and the Euphrates?
(\#Ge 2:10). The sources of the Deglath and of the Arazania, that is to say of the Tigris and
Euphrates, are said to arise in the mountains of Armenia; but there is no paradise there at this day,
nor do both the sources of both these rivers remain there. Perhaps, therefore, we ought to
consider that the real situation of the Paradise is in a place at a distance from this part of the world
which we inhabit, and that it has a river running beneath the earth. which pours forth many veins of
the largest size; so that they, rising up together, pour themselves forth into other veins, which
receive them on account of their great size, and then those have been suppressed by the gulfs of
the waves; on which account, through the impulse given to them by the violence which is implanted
in them, they burst out on the face of the earth in other places, and also among others in the
mountains of Armenia. Therefore, those things which have been accounted the sources of the
rivers, are rather their flowing course; or again, they may be said to be correctly looked upon as
sources, because by all means we must consider the holy scriptures infallible, which point out the
fact of four rivers; for the river is the beginning, and not the spring. But perhaps this passage also
contains an allegorical meaning; for the four rivers are the signs of four virtues: Phison being the
sign of prudence, as deriving its name from parsimony; \footnote{from pheido—, ``parsimony.''} and Gihon being
the sign of sobriety, as having its employment in the regulation of meat and drink, and as
restraining the appetites of the belly, and of those parts which are blow the belly, as being earthly;
the Tigris again is the sign of fortitude, for this it is which regulates the raging commotion of anger
within us; and the Euphrates is the sign of justice, since there is nothing in which the thoughts of
men exult more than in justice.

(13) Why it is that he not only describes the situation of the Euphrates, but also says that the
Phison goes round all the land of Evilat, and that the Gihon goes round all the land of Ethiopia, and
that the Tigris goes toward Assyria? (\#Ge 2:14). The Tigris is a very cruel and mischievous
river, as the citizens of Babylon bear witness, and so do the magi, who have found it to be of a
character quite different from the nature of other rivers; however they might also have another
reason for looking on it with aversion. But the Euphrates is a gentler, and more salubrious, and
more nourishing stream. On which account, the wise men of the Hebrews and Assyrians speak of
it as one which increases and extends itself; and on this account it is not here characterised by its
connection with other things, as the other three rivers are, but by itself. My own opinion is, that
these expressions are all symbolical, for prudence is the virtue of the rational part of man; and it is
in this that wickedness is sometimes found. And fortitude is that portion of the human character
which is liable to degenerate into anger. And sobriety, again, may be impaired by the desires, but
anger and concupiscence are the characteristics of beasts; therefore the sacred historian has
here described those three rivers by the places which they flow round. But he has not described
the Euphrates in that manner, as being the symbol of justice, for there is no certain and limited
portion of it allotted to the soul, but a perfect harmony of the three parts of the soul and of the three
virtues is possessed by it.

(14) Why God placed man in the Paradise with a twofold object, namely, that he might both till it and
keep it, when the Paradise was in reality in need of no cultivation, because it was perfect in
everything, as having been planted by God; nor, again, did it require a keeper, for who was there to
ravage it? (\#Ge 2:15). These are the two objects which a cultivation of the land must attain to
and take care of, the cultivation of the land and the safe keeping of the things which are in it,
otherwise it will be spoiled by laziness or else by devastation. But although the Paradise did not
stand in need of these exertions, nevertheless it was proper that he who had the regulation and
care of it committed to him, namely, the first man, should be as it were a sort of pattern and law to
all workmen in future of everything which ought to be done by them. Moreover it was suitable that,
though all the Paradise was full of everything, it should still leave the cultivator some grounds for
care, and some means of displaying his industry; for instance, by digging around it, and tending it,
and softening it, and digging trenches, and irrigating it by water; and it was needful to attend to its
safety, although there was no one to lay it waste, because of the wild beasts, also more especially
in respect of the air and water; as, for instance, when a drought prevailed, to irrigate it with a
plenteous supply of water, and in moister weather to check the superabundance of moisture by
directing the course of the streams in other directions.

(15) Why, when God commanded the man to eat of every tree within the Paradise, he speaks in
the singular number, and says, ``Thou shalt eat;'' but when he commands that he shall abstain from
the tree which would give him the knowledge of good and evil, he speaks in the plural number, and
says, ``Ye shall not eat of it, for in the day in which ye eat of it ye shall surely die?'' (\#Ge 2:16).
In the first place he uses this language because one good was derived from many; and that also is
not unimportant in these principles, since he who has done anything which is of utility is one, and he
who attains to anything useful is also one; but when I say one, I am speaking not of that which in
point of number comes before duality, but of that one creative virtue by which many beings rightly
coalesce, and by their concord imitate singularity, as a flock, a herd, a troop, a chorus, an army, a
nation, a tribe, a family, a state; for all these things being many members form one community,
being united by affection as by a kiss; when things which are not combined, and which have no
principle of union by reason of their duality and multitude, fall into different divisions, for duality is
the beginning of discord. But two men living as if they were one, by the same philosophy, practise
an unalloyed and brilliant virtue, which is free from all taint of wickedness; but where good and evil
are mingled together the combination contains the principle of death.

(16) What is the meaning of the expression, ``Ye shall surely die?'' (\#Ge 2:17). The death of
the good is the beginning of another life; for life is a twofold thing, one life being in the body,
corruptible; the other without the body, incorruptible. Therefore one wicked man surely dies the
death, who while still breathing and among the living is in reality long since buried, so as to retain in
himself no single spark of real life, which is perfect virtue. But a good man, who deserves so high a
title, does not surely die, but has his life prolonged, and so attains to an eternal end.

(17) Why God says, ``It is not good for man to be alone; let us make him a help meet for him?''
(\#Ge 2:18). By these words God intimates that there is to be a communion, not with all men,
but with those who are willing to be assisted and in their turn to assist others, even though they may
scarcely have any power to do so; since love consists not more in utility than in the harmonious
concord of trustworthy and steadfast manners; so that every one who joins in a communion of love
may be entitled to utter the expression of Pythagoras, ``A friend is another I.''

(18) Why, when God had already said, ``Let us make a help for man,'' he creates beasts and
cattle? (\#Ge 2:19). Perhaps some gluttons and insatiably greedy persons may say that God
did this because beasts and flying things were, as it were, necessary food for man, and his
meetest helper; for that the eating of meat assists the belly so as to conduce to the health and
vigour of the body. But I should think that by reason of the evil implanted in them by nature animals
of all kinds, whether terrestrial or flying in the air, were in this age hostile to and contrary to man;
but that in the case of the first man, as one adorned with every imaginable virtue, they were, as it
were, allies, and a reinforcement in war, and familiar friends, as being tame and domestic by
nature, and this was the sole principle of their familiarity with man, for this it was fit that servants
should dwell with their lord.

(19) Why the creation of animals and flying creatures is mentioned a second time, when the
account of their creation had already been given in the history of the six days? (\#Ge 2:19).
Perhaps those things which were created in the six days were incorporeal angels, indicated under
these symbolical expressions, being the appearances of terrestrial and flying animals, but now they
were produced in reality, being the copies of what had been created before, images perceptible by
the outward senses of invisible models.

(20) Why did God bring every animal to man, that he might give them their names? (\#Ge 2:19).
He has here explained a great source of perplexity to the students of philosophy, admonishing
them that names proceed from having been given, and not from nature; for a natural nomenclature
is with peculiar fitness assigned to each creature when a man of wisdom and pre-eminent
knowledge appears; and, in fact, the office of assigning the names to animals is one which
particularly belongs to the mind of the wise man alone, and indeed to the first man born out of the
earth, since it was fitting that the first of the human race, and the sovereign of all the animals born
out of the earth, should have the dignity assigned to him. For inasmuch as he was the first person
to see the animals, and as he was the first person who deserved to govern them all as their chief,
so also it was fitting that he should be their first namer and the inventor of their names, since it
would have been inconsistent and mad to leave them without any names, or to allow them to
receive names from any one born at a later period, which would have been an insult to and a
derogation from the honour and glory due to the first born. But we may also adopt this idea, that the
giving of names to the different animals was so easily arranged that the very moment that Adam
gave the name the animal itself also heard it; being influenced by the name thus given to it as by a
familiar indication closely connected with it.

(21) Why does Moses say, ``He brought the animals to Adam, that he might see what he would call
them,'' when God can never entertain a doubt? (\#Ge 2:19). It is in truth inconsistent with the
nature of God to doubt; therefore it does not appear that he was in doubt on this occasion, but that
since he had given intellect to man as being the first man born out of the earth and endowed with a
great desire for virtue, by which he was made thoroughly wise as if he had been endowed with
wisdom by nature, so as to consider all things like the proper Ruler and Lord of all, God now
caused him to be influenced to display the proper performance of his task, and saw what was really
the most excellent point of his mind. Besides this, by this statement he evidently indicates the
perfect free-will existing in us, refuting those who affirm that everything exists by a certain
necessity. Or else because it belonged to man to employ the animals, therefore he also gave him
authority to give them names.

(22) What is the meaning of the expression, ``And whatever he called each living thing, that was the
name thereof?'' (\#Ge 2:19). We must consider that Adam gave names not only to all living
creatures, but also to plants, and to everything else which is inanimate, beginning with the more
excellent class; for the living creature is superior to that which has not life. Therefore the scripture
considers the mention of the better part sufficient, indicating by this mention to all who are not
utterly devoid of sense, that he in fact gave names to everything, since it was easy to fix names to
things without life, which were never likely to change their place, and which had no passions of the
soul to exercise, but the giving of proper appellations to living creatures was a more difficult task
on account of the motions of their bodies and the various impulses of their souls, in accordance
with the imagination and the variety of the outward senses, and the different agitations of the mind
from which the effects of their works proceed. Therefore the mind could give names to the more
difficult classes of living creatures. And on this account it was a very proper expression to employ,
that he gave them names as being easy to name, because they were near.

(23) What is the meaning of the expression: ``But for Adam there was not found a helper like to
him?'' (\#Ge 2:20). Every thing was helping and assisting the prince of the human race: the
earth, the rivers, the sea, the air, the light, the heaven. Moreover, every species of fruit and plant
co-operated with him, and every herd of cattle, and every beast which was not savage.
Nevertheless, of all these things which were assisting him, there was nothing like himself,
inasmuch as they were none of them human beings. Therefore, God gave a certain indication that
he might show that man ought to be an assistant to and co-operator with man, being endowed with
perfect similarity to one another in both body and soul.

(24) What is the meaning of the statement, ``And God sent a trance upon Adam, and caused him to
sleep?'' (\#Ge 2:21). How it is that man sleeps is a question which has caused an extraordinary
amount of perplexity to philosophers. But yet our prophet has distinctly explained this question; for
sleep is in itself properly a trance, not of that kind which is more nearly allied to insanity, but of that
which is in accordance with the dissolution of the senses and the absence of counsel; for then the
senses withdraw from those things which are their proper object, and the intellect withdraws from
the senses, not strengthening their nerves, nor giving any motion to those parts which have
received the power of action, inasmuch as they are withdrawn from the objects perceptible by the
outward senses.

(25) What the rib is which God took from the man whom he had formed out of the earth, and which
he made into a woman? (\#Ge 2:21—22). The letter of this statement is plain enough; for it is
expressed according to a symbol of the part, a half of the whole, each party, the man and the
woman, being as sections of nature co-equal for the production of that genus which is called man.
But with respect to the mind, man is understood in a symbolical manner, and his one rib is virtue,
proceeding from the senses; but woman, who is the sensation of counsel, will be more variable.
But some think that the rib means valour and vigour, on which account men call a boxer who as
strong loins eminently strong. Therefore, the lawgiver relates that the woman was formed out of the
rib of the man, indicating by that expression, that one half of the body of the man is woman. And
this is testified to by the formation of the body, by the way in which it is put together, by its motions
and vigour, by the force of the soul, and its strength; for all things are regarded as in a twofold light;
since, as the formation of the man is more perfect, and, if one may so say, more double than the
formation of the woman, so also it required half the time, that is to say forty days; when, for the
imperfect, and, if I may so call it, half section of the man, that is to say the woman, there was need
of a double allowance, that is to say, of eighty days, so that the doubling of the time required for the
nature of the man might be changed, in order to the formation of the peculiar properties of the
woman; for that body, and that soul, the nature of which is in a twofold ratio, the body and soul, that
is, of the man, require but half of the delineation and formation: but that body of which the nature
and construction is in the ratio of one half, namely, that of the woman, her formation and
delineation is in a twofold ratio.

(26) Why Moses calls the form of the woman a building? (\#Ge 2:22).\footnote{the margin of our Bible
points out that the Hebrew word translated ``made,'' means strictly ``builded.''} The union and plentitude of concord
formed by the man and woman is symbolically called a house; but every thing is altogether
imperfect and destitute of a home, which is deserted by a woman; for to the man the public affairs
of the state are committed, but the particular affairs of the house belong to the woman; and a want
of the woman will be the destruction of the house; but the actual presence of the woman shows the
regulation of the house.

(27) Why, as other animals and as man also was made, the woman was not also made out of the
earth, but out of the rib of the man? (\#Ge 2:21). This was so ordained in the first place, in
order that the woman might not be of equal dignity with the man. In the second place, that she might
not be of equal age with him, but younger; since those who marry wives more advanced in years
than themselves deserve blame, as having overturned the law of nature. Thirdly, the design of God
was, that the husband should take care of his wife, as of a necessary part of himself; but that the
woman should requite him in turn with service, as a portion of the universe. In the fourth place, he
admonishes man by this enigmatical intimation, that he should take care of his wife as of his
daughter; and he admonishes the woman that she should honour her husband as her father. And
very rightly, since the woman changes her habitation, passing from her own offspring to her
husband. On which account, it is altogether right and proper that he who has received should take
upon himself the liability in respect of what has been given; and that she who has been removed
should worthily give the same honour to her husband which she has previously given to her
parents; for the husband receives his wife from her parents, as a deposit which is entrusted to him;
and the woman receives her husband from the law.

(28) Why, when the man saw the woman who had been formed in this manner, he proceeded to
say: ``She'' (for ``this,'' touto) ``is now bone of my bone and flesh of my flesh: she shall be called
woman, because she has been taken out of man?'' (\#Ge 2:23).\footnote{woman, virgo, or virago, is here
looked upon as derived from vir, ``man;'' he also derives gyne— from gennao—.} He might have been amazed at
what he had seen, and have said in a negative manner: How can this exquisite and desirable
beauty have been derived from bones and from flesh which are endowed with neither beauty nor
elegance, being of a form so far more beautiful, and endowed with such excessive life and grace?
The matter is incredible because she is like; and yet it is credible, because God himself has been
her creator and painter. Again, he might have said affirmatively: Truly she is a living being, my bone
and my flesh, for she exists by having been taken from that bone and flesh of mine. But he makes
mention of his bone and flesh in a very natural manner; for the human or corporeal tabernacle is
the combination of bones, and flesh, and entrails, and veins, and nerves, and ligaments, and
blood-vessels, and breathing tubes, and blood. And she is called woman (in Greek gyne—) with
great correctness, as the power of producing with fertility, either because she becomes pregnant
through the reception of the seed, and so brings forth; or, as the prophet says, because she was
made out of man, not out of the earth, as he was; nor from seed, as all mankind after them; but of a
certain intermediate nature; and like a branch, brought out of one vine to produce another vine.

(29) Why he says, ``Therefore shall a man leave his father and his mother, and cleave to his wife,
and they shall be two in one flesh?'' (\#Ge 2:24). He here orders man to behave himself
towards his wife with such excess of affection in their intercourse, that he is willing to leave his
parents, not in order that by that means it may be more suitable, but as they would scarcely be a
motive for his fidelity to his wife. And we must remark, that it is very excellent and prudently done,
that he has avoided saying that the woman is to leave her parents and cleave to her husband,
since the character of the man is bolder than the nature of the woman; but he says that the man
ought to do this for the sake of the woman; for he is borne on by a cheerful and willing impulse to
the concord of knowledge, to which, becoming wholly devoted, he restrains and regulates his
desires, and clings to his wife alone like bird-lime. Especially because he himself, delighting in his
master-like authority, is to be respected for his pride: but the woman, being in the rank of a servant,
is praised, for assenting to a life of communion. And when it is said that the two are one flesh, that
indicates that the flesh is very tangible and fully endowed with outward senses, on which it depends
to be afflicted with pain and delighted with pleasure, so that both the man and woman may derive
pleasure and pain from the same sources, and may feel the same; aye, and may still more think
the same.

(30) Why both of them, the man and the woman, are said to have been naked, and not to have
been ashamed? (\#Ge 2:25). They were not ashamed, in the first place, because they were in
the neighbourhood of the world, and the different parts of the world are all naked, each of them
indicating some peculiar qualities, and having peculiar coverings of their own. In the second place,
on account of the sincerity and simplicity of their manners and of their natural disposition, which
had not taint of pride about it. For ambition had as yet no existence. Thirdly, because the climate
and the mildness of the atmosphere was a sufficient covering for them, so as to prevent either cold
or heat from hurting them. In the fourth place, because they, by reason of the relationship existing
between themselves and the world, could not receive injury from any part of it whatever, as being
related to them.

(31) Why does Moses say that the serpent was more cunning than all the beasts of the field?
(\#Ge 3:1). One may probably affirm with truth that the serpent in reality is more cunning than
any beast whatever. But the reason why he appears to me to be spoken of in these terms here is
on account of the natural proneness of mankind to vice, of which he is the symbol. And by vice I
mean concupiscence, inasmuch as those who are devoted to pleasure are more cunning, and are
the inventors of stratagems and means by which to indulge their passions. Being, forsooth, very
crafty in devising plans, both such as favour pleasure and also such as procure means of enjoying
it. But it appears to me that since that animal, so superior in wisdom, was about to seduce man, it
is not the whole race that is here meant to be spoken of as so exceedingly wise, but only that
single serpent, for the reason above mentioned.

(32) Did the serpent speak with a human voice? (\#Ge 3:2). In the first place, it may be the fact
that at the beginning of the world even the other animals besides man were not entirely destitute of
the power of articulate speech, but only that man excelled them in a greater fluency and perspicuity
of speech and language. In the second place, when anything very marvellous requires to be done,
God changes the subject natures by which he means to operate. Thirdly, because our soul is
entirely filled with many errors, and rendered deaf to all words except in one or two languages to
which it is accustomed; but the souls of those who were first created were rendered acute to
thoroughly understand every voice of every kind, in order that they might be pure from evil and
wholly unpolluted. Since we indeed are not endowed with senses in such perfection, for those
which we have received are in some degree depraved, just as the construction of our bodies too is
small; but the first created men, as they received bodies of vast size reaching to a gigantic height,
must also of necessity have received more accurate senses, and, what is more excellent still, a
power of examining into and hearing things in a philosophical manner. For some people think, and
perhaps with some reason, that they were endowed with such eyes as enabled them to behold
even those natures, and essences, and operations, which exist in heaven, as also ears by which
they could comprehend every kind of voice and language.

(33) Why did the serpent accost the woman, and not the man? (\#Ge 3:2). The serpent, having
formed his estimate of virtue, devised a treacherous stratagem against them, for the sake of
bringing mortality on them. But the woman was more accustomed to be deceived than the man. For
his counsels as well as his body are of a masculine sort, and competent to disentangle the notions
of seduction; but the mind of the woman is more effeminate, so that through her softness she
easily yields and is easily caught by the persuasions of falsehood, which imitate the resemblance
of truth. Since therefore, in his old age, the Serpent\footnote{the ancients believed that the serpent became young
again by casting his skin. Ovid says--Anguibus exuitur tenui cum pelle vetustas.} strips himself of his scales from
the top of his head to his tail, he, by his nakedness, reproaches man because he has exchanged
death for immortality. His nature is renewed by the beast, and made to resemble every time. The
woman, when she sees this, is deceived; when she ought rather to have looked upon him as an
example, who, while showing his ingenuity towards her, was full of devices, but she was led to
desire to acquire a life which should be free from old age, and from all decay.

(34) Why the serpent tells the woman lies, saying, ``God has said, Ye shall not eat of every tree in
the Paradise,'' when, on the contrary, what God really had said was, ``Ye shall eat of every tree in
the Paradise, except one?'' (\#Ge 3:4). It is the custom for contending arguers to speak falsely
in an artful manner, in order to produce ignorance of the real facts, as was done in this case, since
the man and woman had been commanded to eat of all the trees but one. But this insidious
prompter of wickedness coming in, says that the order which they had received was that they
should not eat of them all. He brought forward an ambiguous statement as a slippery
stumbling-block to cause the soul to trip. For this expression, ``Ye shall not eat of every tree,''
means in the first place either, not even of one, which is false; or, secondly, not of every one, as if
he intended to say, there are some of which you may not eat, which is true. Therefore he asserts
such a falsehood more explicitly.

(35) Why, when it was commanded them to avoid eating of one plant alone, the woman made also
a further addition to this injunction, saying, ``He said, Ye shall not eat of it, neither shall ye touch it?''
(\#Ge 3:3). In the first place she says this, because taste and every other sense after its kind
consists in the touch appropriate to it. In the second place she says it that it may seem to condemn
them themselves, who did what they had been forbidden. For if even the mere act of touching it
was prohibited, how could they who, besides touching the tree, presumed to eat of the fruit, and so
added a greater transgression to the lesser one, be anything but condemners and punishers of
themselves?

(36) What is the meaning of the expression, ``Ye shall be as gods, knowing good and evil?''
(\#Ge 3:5). Whence was it that the serpent found the plural word ``gods,'' when there is only one
true God, and when this is the first time that he names him? But perhaps this arises from there
having been in him a certain prescient wisdom, by which he now declared the notion of the
multitude of gods which was at a future time to prevail amongst men; and, perhaps, history now
relates this correctly at its first being advanced not by any rational being, nor by any creature of the
higher class, but as having derived its origin from the most virulent and vile of beasts and serpents,
since other similar creatures lie hid under the earth, and their lurking places are in the holes and
fissures of the earth. Moreover, it is the inseparable sign of a being endowed with reason to look
upon God as essentially one being, but it is the mark of a beast to imagine that there are many
gods, and these too devoid of reason, and who can scarcely be said with propriety to have any
existence at all. Moreover, the devil proceeds with great art, speaking by the mouth of the serpent.
For not only is there in the Divinity the knowledge of good and evil, but there is also an approval of
what is good and a repudiation of what is evil; but he does not speak of either of these feelings
because they were useful, but only suggested the mere knowledge of the two contrary things,
namely, of good and evil. In the second place, the expression, ``as gods,'' in the plural number, is in
this place not used inconsiderately, but in order to give the idea of there being both a bad and a
good God. And these are of a twofold quality. Therefore it is suitable to the notion of particular
gods to have a knowledge of contrary things; but the Supreme Cause is above all others.

(37) Why the woman first touched the tree and ate of its fruit, and the man afterwards, receiving it
from her? (\#Ge 3:6). The words used first of all, by their own intrinsic force, assert that it was
suitable that immortality and every good thing should be represented as under the power of the
man, and death and every evil under that of the woman. But with reference to the mind, the woman,
when understood symbolically, is sense, and the man is intellect. Moreover, the outward senses do
of necessity touch those things which are perceptible by them; but it is through the medium of the
outward senses that things are transmitted to the mind. For the outward senses are influenced by
the objects which are presented to them; and the intellect by the outward senses.

(38) What is the meaning of the expression, ``And she gave it to her husband to eat with her?''
(\#Ge 3:6). What has been just said bears on this point also, since the time is nearly one and
the same in which the outward senses are influenced by the object which is presented to them, and
the intellect has an impression made on it by the outward senses.

(39) What is the meaning of the expression, ``And the eyes of both of them were opened?''
(\#Ge 3:7). That they were not created blind is manifest even from this fact that as all other
things, both animals and plants, were created in perfection, so also man must have been adorned
with the things which are his most excellent parts, namely, eyes. And we may especially prove this,
because a little while before the earth-born Adam was giving names to all the animals on the earth.
Therefore it is perfectly plain that he saw them before doing so. Unless, indeed, Moses used the
expression ``eyes'' in a figurative sense for the vision of the soul, by which alone the perception of
good and evil, of what is elegant or unsightly, and, in fact, of all contrary natures, arise. But, if the
eye is to be taken separately as counsel, which is called the warning of the understanding, then
again there is a separate eye, which is a certain something devoid of sound reason, which is
called opinion.

(40) What is the meaning of the expression, ``Because they knew that they were naked?'' (Genesis
3:7). They first arrived at the knowledge of this fact, that is to say, of their nakedness, after they
had eaten of the forbidden fruit. Therefore, opinion was like the beginning of wickedness, when
they perceived that they had not as yet used any covering, inasmuch as all parts of the universe
are immortal and incorruptible; but they themselves immediately found themselves in need of some
corruptible coverings made with hands. But this knowledge was in the nakedness itself, not as
having been in itself the cause of any change, but because their mind now conceived a novelty
unlike the rest of the universal world.

(41) Why they sewed fig-leaves into girdles? (\#Ge 3:8). They did this in the first place,
because the fruit of the fig is very pleasant and agreeable to the taste. Therefore the sacred
historian here, by a symbolical expression, indicates those who sew together and join pleasures to
pleasures by every means and contrivance imaginable. Therefore they bind them around the place
where the parts of generation are seated, as that is the instrument of important transactions. And
they do this, secondly, because although the fruit of the fig-tree is, as I have already said, sweeter
than any other, yet its leaves are harder. And, therefore, Moses here wishes by this symbol to
intimate that the motions of pleasure are slippery and smooth in appearance, but that they,
nevertheless, are in reality hard, so that it is impossible that he who feels them should be delighted,
unless he was previously sorrowful, and he will again become sorrowful. For to be always
sorrowing is a melancholy thing between a double grief, the one being at its beginning, and the
other coming before the first is ended.

(42) What is meant by the statement that the sound was heard of God walking in the Paradise?
was it the sound of his voice, or of his feet? and can God be said to talk? (\#Ge 3:8). Those
gods who are in heaven, perceptible to our outward senses, walk in a ring, proceeding onwards by
a circuitous track; but the Supreme Cause is steadfast and immoveable, as the ancients have
decided. But the true God gives some indication also, as if he wished to give a sense of motion.
For in truth even without his uttering any words, the prophets hear him, by a certain virtue of some
diviner voice sounding in their ears, or perhaps being even articulately uttered. As therefore God is
heard without uttering any sound, so also he gives an idea of walking when he is not walking, nay,
though he is altogether immoveable. But do you not see that before they had tasted of wickedness,
as they were stable and constant, and immoveable and tranquil, and uniform, so also in an equal
manner must they have looked upon the Deity as immoveable, as in fact he is. But they once had
become endued with cunning, they, by judging from themselves, began to strip him of his attributes
of immobility and unchangeableness, and conjectured that he too was subject to variation and
change.

(43) Why while they are hiding themselves from the face of God, the woman is not mentioned first,
since she was the first to eat of the forbidden fruit: but why the man is spoken of in the first place;
for the sacred historian's words are, ``And Adam and his wife hid themselves?'' (\#Ge 3:9). The
woman, being imperfect and depraved by nature, made the beginning of sinning and prevaricating;
but the man, as being the more excellent and perfect creature, was the first to set the example of
blushing and of being ashamed, and indeed, of every good feeling and action.

(44) Why they did not hide themselves in some other place, but in the middle of the trees of the
Paradise? (\#Ge 3:9). Every thing is not done by sinners with wisdom and sagacity, but it often
happens that while thieves are watching for an opportunity of plunder, having no thoughts of the
Deity who presides over the world, the booty which is close to them and lying at their feet is by
some admirable management wrested from them without delay: and something of this kind took
place on the present occasion. For when they ought rather to have fled to a distance from the
garden in which their offence had taken place, they still were arrested in the middle of the Paradise
itself, in order that they might be convicted of their sin too clearly to find any refuge even in flight
itself. And this statement indicates in a figurative manner that every wicked man takes refuge in
wickedness, and that every man who is wholly devoted to his passions flies to those passions as to
an asylum.

(45) Why God asks Adam, ``Where art thou?'' when he knows everything: and why he does not
also put the same question to the woman? (\#Ge 3:10). The expression, ``Where art thou?''
does not here seem to be a mere interrogatory, but rather a threat and a conviction: ``Where art
thou now, O man? from how many good things art thou changed? having forsaken immortality and
a life of the most perfect happiness, you have become changed to death and misery in which you
are buried.'' But God did not condescend to put any question to the woman at all, looking upon her
as the cause of the evil which had occurred, and as the guide to her husband to a life of shame.
But there is an allegorical meaning in this passage, because the principal part is the man, his
guide, the mind, having in itself the masculine principle, when it gives ear to any one introduces
also the defect of the female part, namely that of the outward sense.

(46) Why the man says, ``The woman gave me of the tree, and I did eat;'' but the woman does not
say, ``The serpent gave to me,'' but, ``The serpent beguiled me and I did eat?'' (\#Ge 3:12— 13).
The literal expression here affords grounds for that probable opinion that woman is accustomed
rather to be deceived than to devise anything of importance out of her own head; but with the man
the case is just the contrary. But as regards the intellect, everything which is the object of the
outward senses beguiles and seduces each particular sense of every imperfect being to which it is
adapted. And the sense then, being vitiated by the object, infects the dominant and principal part,
the mind, with its own taint. Therefore the mind receives the impression from the outward sense,
giving it that which it has received itself. For the outward sense is deceived and beguiled by the
sensible object submitted to it, but the senses of the wise man are infallible, as are also the
cogitations of his mind.

(47) Why God curses the serpent first, then the woman, and the man last of all? (\#Ge 3:14).
The reason is that the order of the verses followed the order in which the offences were
committed. The first offence was the deceit practised by the serpent; the second was the sin of the
woman which was owing to him when she abandoned herself to his seduction; the third thing was
the guilt of the man in yielding rather to the inclination of the woman than to the commandment of
God. But this order is very admirable, containing within itself a perfect allegory; inasmuch as the
serpent is the emblem of desire, as is proved, and the woman of the outward sense; but the man is
the symbol of intellect. Therefore the infamous author of the sin is desire; and that first deceives
the outward sense, and then the outward sense captivates the mind.

(48) Why the curse is pronounced on the serpent in this manner, that he shall go on his breast and
on his belly, and eat dust, and be at enmity with the woman? (\#Ge 3:16). The words in
themselves are plain enough, and we have evidence of them in what we have seen. But the real
meaning contains an allegory concealed beneath it; since the serpent is the emblem of desire,
representing under a figure a man devoted to pleasure. For he creeps upon his breast and upon
his belly, being filled with meat and drink like cormorants, being inflamed by an insatiable cupidity,
and being incontinent in their voracity and devouring of flesh, so that whatever relates to food is in
every article something earthly, on which account he is said to eat the dust. But desire has
naturally a quarrel with the outward sense, which Moses here symbolically calls the woman; but
where the passions appear to be as it were guardians and champions in behalf of the senses,
nevertheless they are beyond all question still more clearly flatterers forming devices against them
like so many enemies; and it is the custom of those who are contending with one another to
perpetrate greater evils by means of those things which they concede. Forsooth they turn the eyes
to the ruin of the sight, the ears to hearing what is unwelcome; and the rest of the outward senses
to insensibility. Moreover they cause dissolution and paralysis to the entire body, taking away from
it all soundness, and foolishly building up instead a great number of most mischievous diseases.

(49) Why the curse pronounced against the woman is the multiplication of her sadness and
groans, that she shall bring forth children in sorrow, and that her desire shall be to her husband,
and that she shall be ruled over by him? (\#Ge 3:16). Every woman who is the companion for
life of a husband suffers all those things, not indeed as a curse but as necessary evils. But
speaking figuratively, the human sense is wholly subjected to severe labour and pain, being
stricken and wounded by domestic agitations. Now the following are the children in the service of
the outward senses: the sight is the servant of the eyes, hearing of the ears, smelling of the
nostrils, taste of the mouth, feeling of the touch. Since the life of the worthless and wicked man is
full of pain and want, it arises of necessity from these facts that every thing which is done in
accordance with the outward sense must be mingled with pain and fear. In respect of the mind a
conversion of the outward sense takes place towards the man not as to a companion, for it, like
the woman, is subject to authority as being depraved, but as to a master, because it has chosen
violence rather than justice.

(50) Why God, as he had pronounced a curse on the serpent and on the woman which bore a
relation to themselves and to one another, he did not pronounce a similar one upon the man, but
connected the earth with him, saying, ``Cursed is the earth for thy sake; in sorrow shalt thou eat of
it, thorns and thistles shall it bring forth unto thee, and thou shalt eat the grass of the field: in the
sweat of thy brow shalt thou eat they bread?'' (\#Ge 3:17). Since all intellect is a divine
inspiration, God did not judge it right to curse him in the manner deserved by his offence; but
converted his curse so as to fall upon the earth and his cultivation of it. But man, as a body of
co-equal nature and similar character to that of the earth and understanding, is its cultivator. When
the cultivator is endowed with virtue and diligence, then the body produces its proper fruit, namely
sanity, an excellent state of the outward senses, strength, and beauty. But if the cultivator be a
savage, then every thing is different. For the body becomes liable to a curse, since it has for its
husbandman an intellect unchastised and unsound. And its fruit is nothing useful, but only thorns
and thistles, sorrow and fear, and other vices which every thought strikes down, and as it were
pierces the intellect with its darts. But grass here is symbolically used for food; since man has
changed himself from a rational animal into a brute beast, having neglected all divine food, which is
given by philosophy, by means of distinct words and laws to regulate the will.

(51) What is the meaning of the expression, ``Until thou returnest to the earth from which thou wast
taken;'' for man was not created out of the earth alone, but also of the divine Spirit? (Genesis
3:18). In the first place it is clear, that the first man who was formed out of the earth was made up
both of earth and heaven; but because he did not continue uncorrupt, but despised the
commandment of God, fleeing from the most excellent part, namely, from heaven, he gave himself
up wholly as a slave to the earth, the denser and heavier element. In the second place, if any one
burns with a desire of virtue, which makes the soul immortal, he, beyond all question, attains to a
heavenly inheritance; but because he was covetous of pleasure, by which spiritual death is
engendered, he again gives himself over a second time to the earth, on which account it is said to
him, ``Dust thou art, and unto dust shalt thou return;'' therefore the earth, as it is the beginning of a
wicked and depraved man, so also it is his end; but heaven is the beginning and end of him who is
endowed with virtue.

(52) Why Adam called his wife Life, and affirmed to her, ``Thou art the mother of all living?''
(\#Ge 3:20). In the first place, Adam gave to the first created woman that familiar name of Life,
inasmuch as she was destined to be the fountain of all the generations which should ever arise
upon the earth after their time. In the second place, he called her by this name because she did not
derive the existence of her substance out of the earth, but out of a living creature, namely, out of
one part of the man, that is to say, out of his rib, which was formed into a woman, and on that
account she was called ``life,'' because she was first made out of a living creature, and because
the first beings who were endowed with reason were to be generated from her. Nevertheless, it is
possible that this may have been a metaphorical expression; for is not the outward sense, which is
a figurative emblem of the woman, called with peculiar propriety ``life?'' because it is by the
possession of these senses that the living being is above all other means distinguished from that
which is not alive, as it is by that that the imaginations and impulses of the soul are set in motion,
for the senses are the causes of each; and, in real truth the outward sense is the mother of all
living creatures, for as there could be no generation without a mother, so also there could be no
living creature without sense.

(53) Why God made garments of skins for Adam and his wife, and clothed them? (\#Ge 3:21).
Perhaps some one may laugh at the expressions here used, considering the small value of the
garments thus made, as if they were not at all worthy of the labour of a Creator of such dignity and
greatness; but a man who has a proper appreciation of wisdom and virtue will rightly and
deservedly look upon this work as one very suitable for a God, that, namely, of teaching wisdom to
those who were before labouring to no purpose; and who, having but little anxiety about procuring
useful things, being seized with an insane desire for miserable honours, have given themselves up
as slaves to convenience, looking upon the study of wisdom and virtue with detestation, and being
in love with splendour of life and skill in mean and handicraft arts, which is in no way connected with
a virtuous man. And these unhappy men do not know that a frugality, which is in need of nothing,
becomes, as it were, a relation and neighbour to man, but that luxurious splendour is banished to a
distance as an enemy; therefore the garment made of skins, if one should come to a correct
judgment, deserves to be looked upon as a more noble possession than a purple robe
embroidered with various colours. Therefore this is the literal meaning of the text; but if we look to
the real meaning, then the garment of skins is a figurative expression for the natural skin, that is to
say, our body; for God, when first of all he made the intellect, called it Adam; after that he created
the outward sense, to which he gave the name of Life. In the third place, he of necessity also made
a body, calling that by a figurative expression, a garment of skins; for it was fitting that the intellect
and the outward sense should be clothed in a body as in a garment of skins; that the creature itself
might first of all appear worthy of divine virtue; since by what power can the formation of the human
body be put together more excellently, and in a more becoming manner, than by God? on which
account he did put it together, and at the same time he clothed it; when some prepare articles of
human clothing and others put them on; but this natural clothing, contemporary with the man
himself, namely, the body, belonged to the same Being both to make and to clothe the man in after
it was made.

(54) Who those beings are to whom God says, ``Behold, Adam has become as one of us, to know
good and evil?'' (\#Ge 3:22). The expression, ``one of us,'' indicates a plurality of beings; unless
indeed we are to suppose, that God is conversing with his own virtues, which he employed as
instruments, as it were, to create the universe and all that is in it; but that expression ``as,''
resembles an enigma, and a similitude, and a comparison, but is not declaratory of any
dissimilarity; for that which is intelligible and sensibly good, and likewise that which is of a contrary
character, is known to God in a different manner from that in which it is known to man; since, in the
same way in which the natures of those who inquire and those who comprehend, and the things
themselves too which are inquired into, and perceived, and comprehended, are distinguished,
virtue itself is also capable of comprehending them. But all these things are similitudes, and forms,
and images, among men; but among the gods they are prototypes, models, indications, and more
manifest examples of things which are somewhat obscure; but the unborn and uncreated Father
joins himself to no one, except with the intention of extending the honour of his virtues.

(55) What is the meaning of the words, ``Lest perchance he put forth his hand and take of the tree
of life, and eat and live for ever;'' for there is no uncertainty and no envy in God? (\#Ge 3:23). It
is quite true that God never feels either uncertainty or envy; nevertheless he often employs
ambiguous things and expressions, assenting to them as a man might do; for, as I have said
before, the supreme providence is of a twofold nature, sometimes being God, and not acting in any
respect as a man; but, on some occasions, as a man instructs his son, so likewise should the Lord
God give warning to you. Therefore the first of these circumstances belongs to his sovereign
power, and the second to his disciplinary, and to the first introduction to instruction, so as to
insinuate into man's heart a voluntary inclination, since that expression, ``lest perchance,'' is not to
be taken as a proof of any hesitation on the part of God, but in relation to man, who, by his nature,
is prone to hesitation, and is a denunciation of the inclinations which exist in him. For when any
appearance of anything whatever occurs to any man, immediately there arises within him an
impulse towards that which appears, being caused by that very thing which appears. And from this
arises the second hesitating kind of uncertainty, distracting the mind in various directions, as to
whether the thing is fit to be accepted, or acquired, or not. And very likely present circumstances
have a respect to that second feeling; for, in truth, the Divinity is incapable of any cunning, or
malevolence, or wickedness: it is absolutely impossible that God should either envy the immortality
or any other good fortune belonging to any being. And we can bring the most undeniable proof of
this; for it was not in consequence of any one's entreaties that he created the world; but, being a
merciful benefactor, rendering an essence previously untamed and unregulated, and liable to
suffering, gentle and pleasant, he did so by a vast harmony of blessings, and a regulated
arrangement of them, like a chorus; and he being himself the only sure being, planted the tree of
life by his own luminous character. Moreover, he was not influenced by the mediation or
exhortation of any other being in communicating incorruptibility to man. But while man existed as
the purest intellect, displaying no appearance either of work or of any evil discourse, he was
certain to have a fitting guide, to lead him in the paths of piety, which is undoubted and genuine
immortality. But from the time when he began to be converted to depravity, wishing for the things
which belong to mortal life, he wandered from immortality; for it is not fitting that craft and
wickedness should be rendered immortal, and moreover it would be useless to the subject; since
the longer the life is which is granted to the wicked and depraved man, the more miserable is he
than others, so that his immortality becomes a grave misfortune to him.

(56) Why now he calls the Paradise ``pleasure,'' when he is sending man forth out of it to till the
ground from which he was taken. (\#Ge 3:23). The distinction of agriculture is conspicuous,
when man in the state of paradise, practising the cultivation of wisdom as if he were employed in
the cultivation of trees, and enjoying the food of imperishable and most useful fruits, was himself
endowed with immortality likewise. After that, being expelled from the place of wisdom, he
experienced the opposite effects of ignorance, by which the body is polluted, and at the same time
the intellect is blinded, and, being exposed to a want of proper food, he wastes away, and yields to
a miserable death. On which account, now, in contempt of the foolish man, God calls the Paradise
``pleasure,'' in order to put it in opposition to a life of pain, and misery, and savageness. In truth, the
life which is passed in wisdom is a pleasure, full of liberal joy, and is the constant enjoyment of a
rational soul; but that life which is destitute of wisdom is found to be both savage and miserable,
although it is excessively deceived by the appetites, which pain both precedes and follows.

(57) Why God places a cherubim in front of the Paradise, and a flaming sword, which turned every
way, to keep the way of the tree of life? (\#Ge 3:24). The name cherubim designates the two
original virtues which belong to the Deity, namely, his creative and his royal virtues. The one of
which has the title of God, the other, or the royal virtue, that of Lord. Now the form of the creative
power is a peaceable, and gentle, and beneficent virtue; but the royal power is a legislative, and
chastising, and correcting virtue. Moreover, by the flaming sword he here symbolically intimates the
heaven: for the air is of a flaming colour, and turns itself round, revolving about the universe.
Therefore, all these things assumed to themselves the guardianship of the Paradise, because they
are the presidents over wisdom, like a mirror; since, to illustrate my meaning by an example, the
wisdom of the world is a sort of mirror of the divine virtues, in the similitude of which it was
perfected, and by which the universe and all the things in it are regulated and arranged. But the way
to wisdom is called philosophy (a word which means the love or the pursuit of wisdom). And since
the creative virtue is endued with philosophy, being both philosophical and royal, so also the world
itself is philosophical. Some persons however have fancied that it is the sun which is indicated by
the flaming sword; because, by its constant revolutions and turnings every way, it marks out the
seasons of the year, as being the guardian of human life and of every thing which serves to the life
of all men.

(58) Whether it was properly said with respect to Cain: ``I have gotten a man from the Lord?''
(\#Ge 4:1). Here there is a distinction made, as to-from some one, and out of some one, and
by some thing. Out of some one, as out of materials; from some one, as from a cause; and by
some thing, as by an instrument. But the Father and Creator of all the world is not an instrument,
but a cause. Therefore he wanders from right wisdom who says, ``That what has been made has
been made, not from God, but by God.''

(59) Why the sacred historian first describes the employment of the younger brother, Abel, saying:
``He was a keeper of sheep; but Cain was a cultivator of the earth?'' (\#Ge 4:2). Since, although
the virtuous son was in point of time younger than the wicked son, yet in point of virtue he was
older. On which account, on the present occasion, when their actions are to be compared
together, he is placed first. Therefore one of them exercises a business, and takes care of living
creatures, although they are devoid of reason, gladly taking upon himself the employment of a
shepherd, which is a princely office, and as it were a sort of rehearsal of royal power; but the other
devotes his attention to earthly and inanimate objects.

(60) Why Cain after some days offers up the first-fruits of his fruits, but when it is said that ``Abel
offered up first-fruits of the first-born of his flock and of the fat,'' ``after some days'' is not added?
(\#Ge 4:3—4). Moses here intimates the difference between a lover of himself, and one who is
thoroughly devoted to God; for the one took to himself the first-fruits of his fruits, and very
impiously looked upon God as worthy only of the secondary and inferior offerings; for the
expression, ``after some days,'' implies that he did not do so immediately; and when it is said that he
offered of the fruits, that intimates that he did not offer of the best fruits which he had, and herein
displays his iniquity. But the other, without any delay, offered up the first-born and eldest of all his
flocks, in order that in this the Father might not be treated unworthily.

(61) Why, when he had begun with Cain, he still mentions him here in the second place, when he
says: ``And God had respect unto Abel and unto his offerings; but unto Cain and unto his sacrifices
he paid no attention?'' (\#Ge 4:5). In the first place, because the good man, who is by nature
first, is not at first perceived by the outward senses of any man except in his own turn, and by
people of virtuous conduct. Secondly, because the good and the wicked man are two distinct
characters; he accepts the good man, seeing that he is a lover of what is good, and an eager
student of virtue; but he rejects and regards with aversion the wicked man, presuming that he will
be prone to that side by the order of nature. Therefore he says here with exceeding fitness, that
God had regard, not to the offerings, but to those who offered them, rather than to the gifts
themselves; for men have regard to and regulate their approbation by the abundance and richness
of offerings, but God looks at the sincerity of the soul, having no regard to ambition or illusion of
any kind.

(62) What is the meaning of the distinction here made between a gift and a sacrifice? (Genesis
4:4). The man who slays a sacrifice, after having made a division, pours the blood around the altar
and takes the flesh home; but he who offers it as a gift, offers as it should seem the whole to him
who accepts it. Therefore, the man who is a lover of self is a distributor, like Cain; but he who is a
lover of God is the giver of a free gift, as was Abel.

(63) How it was that Cain became aware that his offering had not pleased God? (\#Ge 4:5).
Perhaps he resolved his doubts, an additional cause being added, for sorrow seized upon him and
his countenance fell. Therefore, he took the sorrow which he felt as an indication that he had been
sacrificing what was not pleasing or approved of, when joy and happiness would have been suited
to one who was sacrificing with purity of heart and spirit.

(64) Why is it that the expression used is not, because you do not offer rightly; but, because (or
unless) you do not divide rightly? (\#Ge 4:7). In the first place, we must understand that right
division and improper division are nothing else but order and the want of it. And it is by order that
the universal world and its parts were made; since the Creator of the world, when he began to
arrange and regulate the previously untamed and unarranged power which was liable to suffering,
employed section and division. For he placed the heavy elements which were prone to descend
downwards by their own nature, namely, the earth and the water, in the centre of the universe; but
he placed the air and the fire at a greater altitude, as they were raised on high by reason of their
lightness. But separating and dividing the pure nature, namely heaven, he carried it round and
diffused it over the universe, so that it should be completely invisible to all men; containing within
itself the whole universe in all its parts. Again, the statement that animals and plants are produced
out of seeds, some moist and some dry, what else does it mean but the inevitable dissection and
separation of distinction? Therefore it follows inevitably, that this order and arrangement of the
universe must be imitated in all things, especially in feeling and acknowledging gratitude; by which
we are invited to requite in some degree and manner the kindnesses of those who have showered
greater benefits liberally on us. Moreover, to pay one's thanks to God is an action which is
intrinsically right in itself: and it is not to be disapproved of that he should receive the offerings due
to him at the earliest moment, and fresh gifts from the first-fruits of every thing, not being
dishonoured by any negligence on our part. Since it is not fitting that man should reserve for
himself the first and most excellent things which are created, and should offer what is only second
best to the all-wise God and Creator; for that division would be faulty and blameworthy, showing a
most preposterous and unnatural arrangement.

(65) What is the meaning of the expression: ``You have done wrongly; now rest?'' (\#Ge 4:8).
He is here giving very useful advice; since, to do no wrong at all is the greatest of all good things:
but he who sins, and who thus blushes and is overwhelmed with shame, is near akin to him, being,
if I may use such a phrase, as the younger brother to the elder; for those persons who pride
themselves on their errors as if they had not done wrong, are afflicted with a disease which is
difficult to cure, or rather which is altogether incurable.

(66) Why he seems to be giving what is good into the hand of a wicked man, when he says, ``And
unto thee shall be his desire?'' (\#Ge 4:8). He does not deliver good into his hand; but the
expression is heard with different feelings; since he is speaking, not of a pious man, but after the
action is accomplished, saying of him: The desire and respect of the impiety of this man's
wickedness will be towards you. Do not therefore talk about necessity, but about your own habits,
in order that thus he may represent the voluntary action. And again, the sentence, ``And you shall
be his ruler over him,'' has a reference to the operation. In the first place, you begin to act with
wickedness; and now behold, another iniquity follows that great and injurious iniquity. Therefore, he
both thinks and affirms that this is the principal part of all voluntary injury.

(67) Why he slew his brother in the field? (\#Ge 4:9). That as all in fecundity and sterility arises
from a neglect of sowing and planting land a second time, he may be kept continually in mind of his
wicked murder, and self-blamed for it; since the ground was not to be the same for the future, after
it was compelled, contrary to its nature, to drink of human blood, to bring forth food to that man who
imbued it with the polluted stain of blood.

(68) Why he who knows all things asks the fratricide: ``Where is thy brother Abel?'' (\#Ge 4:10).
He puts this question to him because he wishes the man to confess voluntarily and spontaneously,
of his own accord, so that he may not imagine that every thing is done out of necessity; for he who
had slain another through necessity, would have confessed unwillingly, as having done the deed
unwillingly; since that which does not depend upon ourselves does not deserve accusation; but the
man who has done wrong intentionally denies it; for those who do wrong are liable to repentance.
Therefore, he has interwoven this principle in all parts of his legislation, because the Deity himself
is never the cause of evil.

(69) Why he who had slain his brother makes answer as if he were replying to a man; and says, ``I
do not know: am I my brother's keeper?'' (\#Ge 4:9). It is the opinion of an atheist to think that
the eye of God does not penetrate through every thing, and behold all things at the same time;
piercing not only through what is visible, but also through every thing which lurks in the deepest and
bottomless unfathomable abysses. Suppose a person said to him, ``How can you be ignorant
where your brother is, and how is it that you do not know that, when as yet he is one out of the only
four human beings which exist in the world? He being one with both his parents, and you his only
brother.'' To this question the reply made is: ``I am not my brother's keeper.'' O what a beautiful
apology! And whose keeper and protector ought you to have been, rather than your brother's? But
if you have excited your diligence to give effect to violence, and injury, and fraud, and homicide,
which are the foulest and most abominable of actions, why did you consider the safety of your
brother a secondary object?''

(70) What is the meaning of the expression, ``The voice of thy brother's blood cries to me out of
the earth?'' (\#Ge 4:10). This is especially an example by which to take warning; for the Deity
listens to those who are worthy, although they be dead, knowing that they are alive as to an
incorporeal life. But he averts his countenance from the prayers of the wicked, although they are
living a flourishing life, inasmuch as he looks upon them as dead to any real life, carrying about
their bodies like a sepulchre; and having buried their miserable souls in it.

(71) Why he is said to be cursed upon the earth? (\#Ge 4:11). The earth is the last portion of
the world, therefore if that utters curses, we must consider that the other elements do likewise pour
forth adequate maledictions; for instance, the fountains, and rivers, and sea, and the air, and the
land, and the fire, and the light, and the sun, and moon, and stars, and in short the whole heaven.
For if inanimate and earthly nature, throwing off the yoke, wars against injury, why may not still
rather those natures do so which are of a purer character? But as for him, against whom the parts
of the world carry on war, what hope of safety he can have for the future, I know not.

(72) What is the meaning of the curse, ``You shall be groaning and trembling upon the earth?''
(\#Ge 4:13).\footnote{our translation is, ``My punishment is greater than I can bear.''} This also is a general
principle; for in all evils there are some things which are perceived immediately, and some which
are felt at a later period; for those which are future cause fear, and those which are felt at once
bring sorrow.

(73) What is the meaning of Cain saying, ``My punishment is too great for you to dismiss me?
(\#Ge 4:12). In truth there is not misery greater than to be deserted and despised by God; for
the anarchy of fools is cruel and very intolerable; but to be despised by the great King, and to fall
down as an abject person cast down from the government of the Supreme Power is an
indescribable affliction.

(74) What is the meaning of Cain, when he says, ``Everyone who shall find me will kill me:'' when
there was scarcely another human being in the world except his parents? (\#Ge 4:14). In the
first place he might have received injury from the parts of the world which indeed were made for the
advantage of the good and that they might partake of them, but which nevertheless, derived from
the wicked no slight degree of revenge. In the second place it may be that he said this, because he
was apprehensive of injury from beasts, and reptiles; for nature has brought forth these animals
with the express object of their being instruments of vengeance on the wicked. In the third place,
some people may imagine that he is speaking with reference to his parents, on whom he had
inflicted an unprecedented sorrow, and the first evil which had happened to them, before they knew
what death was.

(75) Why whoever should slay Cain should be liable to bear a sevenfold punishment? (Genesis
4:15). As our soul consists of eight portions, being accustomed to be divided in its rational and
irrational individuality into seven subordinate parts, namely into the five outward senses, and the
instrument of vice, and the faculty of generation; those seven parts exist among the causes of
wickedness and evil, on which account they likewise fall under judgment; but the death of the
principal and dominant portion of man, namely of the mind, is principally the wickedness which
exists in it. Whoever therefore slays the mind, mingling in it folly, and insensibility, instead of sense,
will cause dissolution also of the seven irrational parts; since, just as the principal and leading part
had a portion from virtue, in the same manner likewise are its subject divisions composed.

(76) Why a sign is put on him who had slain his brother, that no one should kill him who found him;
when it would have been natural to do the contrary, namely, to give him over to the hands of an
executioner to be put to death? (\#Ge 4:16). This is said because, in the first place, the change
of the nature of living is one kind of death; but continual sorrow and unmixed fear are destitute of
joy and devoid of all good hope, and so they bring on many terrible and various evils which are so
many sensible deaths. In the second place, the sacred historian designs at the very beginning of
his work to enunciate the law about the incorruptibility of the soul, and to confute as deceitful those
who look upon the life which is contained in this body as the only happy life; for behold one of the
two brothers is guilty of those enormous crimes which have already been mentioned, namely,
impiety and fratricide; and he is still alive, and begetting offspring, and building cities. But the other
who was praised in respect of his piety is treacherously put to death; while the voice of the Lord
not only clearly cries out that that existence which is perceptible by the outward senses is not good,
and that such a death is not evil, but also that that life which is in the flesh is not life, but that there
is another give to man free from old age, and more immortal, which the incorporeal souls have
received; for that expression of the poet about Scylla,

``That is not mortal but an endless Woe,''\footnote{the line occurs in Homer, Odyss. 12.118.}

is asserted in the same familiarity about a person who lives ill and passes a long life for many
years in the practice of wickedness. In the third place, since Cain had perpetrated this fratricide of
enormous guilt above all other crimes, he presents himself to him, quite forgetful of the injury that
he has done, imposing on all judges a most peaceful law for the first crime; not that they are not to
destroy malefactors, but that resting for a while with great patience and long suffering, they shall
study compassion rather than severity. But God himself, with the most perfect wisdom, has laid
down the rule of familiarity and intelligence with reference to the first sinner: not slaying the
homicide, but destroying him in another manner; since he scarcely permitted him to be enumerated
among the generations of his father, but shows him proscribed not only by his parents but by the
whole race of mankind, allotting him a state separate from that of others, and secluded from the
class of rational animals, as one who had been expelled, and banished, and turned into the nature
of beasts.

(77) Why Lamech, after the fifth generation, blames himself for the fratricide of his elder Cain;
saying, as the scripture reports, to his wives, Adah and Zillah; ``I have slain a man to my injury and
a young man to my hurt; since if vengeance is taken upon Cain sevenfold, it shall certainly be taken
on Lamech seventy and sevenfold?'' (\#Ge 4:23). In numerals one is before ten, both in order
and in virtue, for it is the first beginning and element and measure of all things. But the number ten
is subsequent, and is measured by the other, being inferior to it, both in order and virtue; therefore,
also, the number seven is antecedent in its origin to and more ancient than the number seventy,
but the number seventy is younger than the number seven, and contains the calculation of
generations. These premises being laid down, he who first committed sin, as if he had been really
always ignorant of evil, like the first odd number, namely, the unit, is chastised more simply; but the
second offender, because he had the first for an example, so that there cannot possibly be any
excuse made for him, is guilty of a voluntary crime, and because he did not receive honourable
wisdom from that more simple punishment, the consequence will be that he will both suffer all that
first punishment, and will, moreover, receive this second one, which is contained in the number ten.
For as in the horse-races they pay the groom who has trained the horse twice as great a reward
as they give to the driver, so some wicked men, inclined to acts of injustice, gain the miserable
triumph of victory and then are punished with a double punishment, both by the first one which is
contained in the unit, and also by the second which is contained in the number ten; besides, Cain
being the author of a homicide, when he was ignorant of the greatness of the pollution which he
was incurring, because no death had hitherto taken place in the world, suffered a more simple
punishment, namely, only a sevenfold penalty in the order of the unit; but as his imitator could not
take refuge in the same plea of ignorance, he ought to be subjected to a twofold punishment, not
only to one equal and similar to that which had been inflicted on the first offender, but also another,
which should be the seventh among the decades. In truth, according to the law, the trial which is
before the tribunal is a sevenfold one; first of all, the eyes are put on their trial, because they
beheld what was not lawful; secondly, the ears are impeached, because they heard what they
ought not to have heard; thirdly, the smell is brought into question, as having been reduced by
smoke and vapour; fourthly, the taste is accused, as being subservient to the pleasures of the
belly; fifthly, a charge is brought against the taste, by means of which, besides the operations of
the senses abovementioned, in respect of those things which prevail over the spirit, other things,
also, are superadded separately, such as the takings of cities, the captivities of men, the
destructions of those citadels of cities in which wisdom dwells; sixthly, an accusation is urged
against the tongue and other instruments of speech, for being silent as to what should be spoken
of, and speaking of what should be buried in silence; and, in the seventh place, the lower part of the
belly is impeached for inflaming and exciting the passions by immoderate lust. This is the meaning
of that expression, according to which a sevenfold vengeance was taken upon Cain, but a seventy
and sevenfold vengeance upon Lamech for the causes above mentioned, because he was the
second offender, not having been taught by the punishment of the first delinquent, and therefore he
is altogether worthy to receive his punishment, which is the more simple one, like the unit in
numerals, and, also, a manifold punishment too equal to the number ten.

(78) Why Adam, when he begat Seth, introduces him saying, ``God has raised up for me another
seed in the place of Abel whom Cain slew?'' (\#Ge 4:25). In real truth Seth is another seed and
the beginning of a second nativity of Abel, in accordance with a certain natural principle; for Abel is
like to one who comes down below from above, on which account it was that he perished
injuriously; but Seth resembles one who is proceeding upwards from below, on which account he
also increases. And in proof of this argument Abel is explained as having been brought back and
offered upwards to God. But it is not proper that everything should be raised and borne upwards,
but only that which is good, for God is in no respect whatever the cause of evil. Therefore,
whatever is indistinct and uncertain, and mingled, and in confusion and disorder, has also, very
properly, blame and praise mingled together: praise, because it honours the cause, and blame,
since as the occurrence happened fortuitously, so it is without any plans having been formed or
any gratitude expressed. Moreover, nature also separated the two sons from him; it rendered the
good one worthy of immortality, resolving him into a voice interceding with God; but the wicked one
it gave over to corruption. But the name Seth is interpreted ``watered,'' according to the variation of
plants which grow by being watered, and put forth shoots and bear fruit. But these things are the
symbols of the soul, so that it is not lawful to assert that the Divinity is the cause of all things
equally, of the bad as well as of the good, but only of the good, and that alone ought to be planted
alive.

(79) Why Enos, the son of Seth, hoped to call upon the name of the Lord God? (\#Ge 4:26).
The name Enos is interpreted ``man;'' and it is received as meaning, not the whole of the combined
man, but as the rational part of the soul, namely, the intellect, to which it is peculiarly becoming to
hope, for irrational animals are devoid of hope; but hope is a sort of presage of joy, and before joy
there is an expectation of good things.

(80) Why, after the mention of hope, Moses says, ``This is the book of the generations of men?''
(\#Ge 5:1). It is by this that he made what has been said before worthy of belief. What is man?
Man is a being which, beyond all other races of animals, has received a copious and wonderful
portion of hope; and this is as it were inscribed on his very nature, and celebrated there; for the
human intellect hopes by its own nature.

(81) Why, in the genealogy of Adam, Moses no longer mentions Cain, but only Seth, who, he says,
was according to his appearance and form; on which account he proceeds to retain the
generations which descend from him in his genealogy? (\#Ge 5:3). It can neither be lawful to
enumerate a wicked and sinful murderer either in the list of reason or in that of number; for he must
be cast out like dung, as some one said, looking upon him as one of such a character; and on this
account the sacred historian neither points him out as the successor of his father who had been
formed out of the dust, nor as the head of succeeding generations; but he distributes both these
characteristics to him who was without pollution, and names Seth, who is a drinker of water, as
having been watered by his father, and as begetting hope in his own increase and progress; on
which account it is not inconsiderately and foolishly that he says that he was born according to the
form and appearance of his father, to the reproof of his elder brother, who, on account of the
foulness of the murder which he had committed, has nothing in him resembling his father, either in
body or soul. And on this account Moses has separated him from the family, and has given his
share to his brother, being the noble privilege of the birthright of the first-born.

(82) What is the meaning of the verse, Enoch pleased God after he begat Methuselah, two
hundred years? (\#Ge 5:22). God appointed by the law the fountains of all good things to be
under the principles of generation itself. And what I mean is something of this sort. A little while
before he appointed mercy and pardon to exist, now again he decrees that penitence shall exist,
not in any degree mocking or reproaching these men, who are believed to have offended, and at
the same time giving the soul an opportunity to mount up from wickedness to virtue, like the
conversion of those who are proceeding towards a snare. For behold, the man being made a
husband and a father together with his birth, makes a beginning of honesty. And he is said to
please God, for although he does not persevere in piety from the moment that he is born,
nevertheless, all that remaining period is counted to him as having been spent in a praiseworthy
mode of life, because he pleased God for so many years. And these things are said, not because
it perhaps was, but it might perhaps have seemed different; but he approves of the order of things,
for indulgence having been exemplified, in this case of Cain, after no long interval of time, he
introduces this statement, that Enoch practised repentance, warning us by it that repentance alone
can procure indulgence.

(83) Why Enoch, who cultivated repentance, is said to have lived before his repentance a hundred
and sixty-five years, but two hundred after his repentance? (\#Ge 5:22).\footnote{this is at variance with
the statement in the Bible, which says he lived sixty-five years before the birth of Methuselah, and walked with God three
hundred years.} This number of a hundred and sixty-five is combined of the singular addition of ten
numbers from the unit to ten; as one, two, three, four, five, six, seven, eight, nine, ten, the total of
which is fiftyfive. And again, from that by the addition of ten numbers, which, removing the unit
proceed upwards by twos, and two, four, six, eight, ten, twelve, fourteen, sixteen, eighteen, twenty,
which make a hundred and ten; the combination of which, with the numbers first mentioned,
produces a hundred and sixty-five; and in this addition the even numbers amount to twice as much
as the odd numbers; for the woman is more violent than the man, in the preposterous manner in
which the wicked man rules over the virtuous man, the outward sense over the mind, the body over
the outward sense, and matter over its cause. But the number two hundred, in which repentance
was practised, is combined of two numbers of one hundred, the first hundred of which intimates a
purification from injustice, but the other indicates the plenitude of perfect virtue. In truth, before
anything else is done, the first thing is to cut off from a sick body every diseased part, and after
that means of cure are to be applied to it, for this is the first step, and the other the second.
Moreover, the number two hundred consists of fours, for it is produced as from seed, from four
triangular numbers, and from four tetragons, and from four pentagons, and from four hexagons,
and from four heptagons; and, as one may say, it fixes its step on the number seven. Now the four
triangles are these, one, three, six, ten, which make twenty; the four tetragons are one, four, nine,
sixteen, which make thirty; from the four pentagons, one, five, twelve, twenty-two, is made the
number forty. Moreover, the four hexagons, one, six, fifteen, twenty-eight, make fifty; and the four
heptagons, one, seven, eighteen, thirty-four, make sixty; and all these numbers put together make
two hundred.

(84) Why the man who lives a life of repentance is said to have lived three hundred and sixtyfive
years? (\#Ge 5:23). In the first place, the year contains three hundred and sixty-five days;
therefore, by the symbol of the solar orbit, the sacred historian here indicates the life of the
repentant man. In the second place, as the sun is the cause of day and night, performing his
revolutions by day above the hemisphere of the earth, and his course by night under the earth, so
also the life of the man of repentance consists of alternations of light and darkness; of darkness,
that is, of times of agitation and circumstances of injury; and of light, when the light of virtue and its
radiant brilliancy arises. In the third place, he has assigned to him a complete number, as the sun
is ordained to be the chief of the stars of heaven, under an appointed number, in the time which
came before the period of his repentance, to lead to the oblivion of the sins previously committed;
since, as God is good, he bestows the greatest favours most abundantly, and, at the same time,
he effaces the former offences of those who devote themselves to him, and which might deserve
chastisement, by a recollection of their virtues.

(85) Why, when Enoch died, the sacred historian adds the assertion, ``He pleased God?'' (Genesis
5:24). In the first place, he says this because, by such a statement, he implies that the soul is
immortal, inasmuch as after it is stripped of the body, it still pleases a second time. In the second
place, he honours the repentant man with praise, because he has persevered in the same
alteration of manners, and has never receded till he has arrived at complete perfection of life; for
behold, some men appear to be readily sated after they have only tasted of excellence; and after a
hope of recovery has been given to them, they relapse again into the same disease.

(86) What is the meaning of the expression, ``He was not found because God translated him?''
(\#Ge 5:24). In the first place, the end of virtuous and holy men is not death but a translation
and migration, and an approach to some other place of abode. In the second place, in this instance
something marvellous did take place; for he was supposed to be carried off in such a way as to be
invisible, for then he was not found: and a proof of this is, that he was sought for as being invisible,
not only as having been carried away from their sight, since translation into another place is
nothing else than a placing of a person in another situation; but it is here suggested, that he was
translated from a visible place, perceptible by the outward senses, into an incorporeal idea,
appreciable only to the intellect. This mercy also was bestowed on the great prophet, for his
sepulchre also was known to no one. And besides these two there was another, Elijah, who
ascended from the things of earth into heaven, according to the divine appearance which was then
presented to him, and who thus followed higher things, or, to speak with more exact propriety, was
raised up to heaven.

(87) How it was that immediately upon the nativity of Noah his father says, ``He will make us rest
from labours and sorrows, and from the earth, which the Lord God has cursed?'' (\#Ge 5:29).
The fathers of the saints did not prophesy except for grave reasons and on important occasions;
for although those who were rendered worthy of prophetical panegyric did not prophecy at all times
or on all subjects, they did so at all events on one occasion and on one subject, with which they
were acquainted. Nor is this of no importance, but it is an emblem and an example, since Noah is a
kind of surname of righteousness, of which, when the intellect is made a partaker, it causes us to
rest from all wicked works, and releases us from sorrows and from fears, and renders us secure
and joyful. It also causes us to rest from that earthly nature which has been previously laid under a
curse, which this body, when affected by pain, is connected with, especially in those persons who
give cause for it, and who wear out their lives with pleasure. Nevertheless, if we examine
attentively the events and circumstances, and compare them with the letter of the scriptures, the
prophecy which has been already produced is deceived, because, in the time of this man, there did
not arise any putting down of evils, but a more vehement obstinacy in sin and great afflictions, and
the unprecedented event of the deluge. But you must note carefully, that Noah is the tenth in
generation from the earth-born Adam.

(88) What is meant by the three sons of Noah being named Shem, Ham, and Japhet? (Genesis
5:32). These names are the symbols of three human things, what is good, what is bad, and what is
indifferent; Shem is the symbol of what is good, Ham of what is bad, and Japhet of what is
indifferent.

(89) Why from the time that the deluge drew near, the human race is said to have increased so as
to become a multitude? (\#Ge 6:1). Divine mercies do always precede judgment; since the first
work of God is to do good, and to destroy follows afterwards; but he himself (when terrible evils are
about to happen) loves to provide and is accustomed to provide that previously an abundance of
many and great blessings shall be produced. On this principle also Egypt, when there was about to
be a barrenness and famine for seven years as the prophet himself says, \footnote{\#ge 41:28.} was for
an equal number of years continuously made exceedingly fertile by the beneficent and saving
power of the Creator of the universe. And in the same way in which he showers benefits upon men,
he also teaches them to depart and to abstain from sin; that these blessings may not be turned into
the contrary. And on this account now, by the freedom of their institutions, the cities of the world
have increased in generous virtue, so that if any corruption supervenes subsequently they may
disapprove of their own acts of wickedness as extraordinary and irremediable; not at all looking
upon the divinity as the cause of them, for that has no connection with wickedness or misery, for
the task of the Deity is only to bestow blessings.

(90) What is the meaning of the expression, ``My spirit shall not always strive with man, because he
is but flesh?'' (\#Ge 6:3). An oracle is here promulgated as if it were a law; for the divine spirit is
not a motion of the air, but intellect and wisdom; just as it also flows over the man who with great
skill constructed the tabernacle of the Lord, namely upon Bezaleel, when the scripture says, ``And
he filled him with the divine spirit of wisdom and understanding.'' Therefore that spirit comes upon
men, but does not abide or persevere in them; and the Lord himself adds the reason, when he
says, ``Because they are flesh.'' For the disposition of the flesh is inconsistent with wisdom,
inasmuch as it makes a bond of alliance with desire; on which account it is evident that nothing
important can be in the way of incorporeal and light souls, or can be any hindrance to their
discerning and comprehending the condition of nature, because a pure disposition is acquired
together with constancy.

(91) Why it is said that the days of man shall be a hundred and twenty years? (\#Ge 6:4). God
appears here to fix the limit of human life by this number, indicating by it the manifold prerogative of
honour; for in the first place this number proceeds from the units, according to combination, from
the number fifteen; but the principle of the number fifteen is that of a more transparent
appearance, since it is on the fifteenth day that the moon is rendered full of light, borrowing its light
of the sun at the approach of evening, and restoring it to him again in the morning; so that during
the night of the full moon the darkness is scarcely visible, but it is all light. In the second place, the
number a hundred and twenty is a triangular number, and is the fifteenth number consisting of
triangles. Thirdly, it is so because it consists of a combination of odd and even numbers, being
contained by the power of the faculty of the concurring numbers, sixty-four and fifty-six; for the
equal number of sixty-four is compounded of the uniting of these eight odd numbers, one, three,
five, seven, nine, eleven, thirteen, fifteen; the reduction of which, by their parts into squares, makes
a sum total of sixty-four, and that is a cube, and at the same time a square number. But again from
the seven double units there arises the unequal number of fifty-six, being compounded of seven
double pairs, which generate other productions of them, two, four, six, eight, ten, twelve, fourteen;
the sum total of which is fifty-six. In the fourth place, it is compounded of four numbers, of one
triangle, namely fifteen; and of another square, namely twenty-five; and of a third quinquangular
figure, thirty-five; and of a fourth a sexangular figure forty-five, by the same analogy: for the fifth is
always received according to each appearance; for from the unity of the triangles the fifth number
becomes fifteen; again the fifth of the quadrangular number from the unit makes twenty-five; and
the fifth of the quinquangular number from the unit makes thirty-five; and the fifth of the sexangular
number from the unit makes forty-five. But every one of these numbers is a divine and sacred
number, consisting of fifteens as has been already shown; and the number twenty-five belongs to
the tribe of Levi.\footnote{see \#Nu 8:24.} And the number thirty-five comes from the double diagram of
arithmetic, geometry, and harmony; but sixteen, and eighteen, and nineteen, and twenty-one, the
combination of which numbers amounts to seventy-four, is that according to which seven months'
children are born. And forty-five consists of a triple diagram; but to this number, sixteen, nineteen,
twenty-two, and twenty-eight, belong: the combination of which makes eighty-five, according to
which nine months' children are produced. Fifthly, this diagram has fifteen parts, and a twofold
composition, peculiarly belonging to itself; forsooth when divided by two it gives sixty, the measure
of the age of all mankind; when divided by three it gives forty, the idea of prophecy; when divided by
four it gives thirty, a nation; when divided by five, it makes twenty-four, the measure of day and
night; when divided by six, it gives twenty, a beginning; when divided by eight, we have fifteen, the
moon in the fulness of brilliancy; when divided by ten, it makes twelve, the zodiac embellished with
living animals; when divided by twelve, it makes ten, holy; when divided by fifteen, it gives eight, the
first ark; when divided by twenty, it leaves six, the number of creation; when divided by twenty-four,
it makes five, the emblem of the outward sense; when divided by thirty it makes four, the beginning
of solid measure; when divided by forty, it gives three, the symbol of fulness, the beginning, the
middle, and the end; when divided by sixty, it makes two, which is woman; and when divided by the
whole number of a hundred and twenty, the product is one, or man. And every one of all these
numbers is more natural, as is proved in each of them, but the composition of them is twofold, for
the product is two hundred and forty, which is a sign that it is worthy of a twofold life; for as the
number of years is doubled, so also we may imagine that the life is doubled too; one being in
connection with the body, the other being detached from the body, according to which every holy
and perfect man may receive the gift of prophecy. Sixthly, because the fifth and sixth figures arise,
the three numbers being multiplied together, three times four times five, since three times four
times five make sixty; so in like manner the next following numbers four times five times six make a
hundred and twenty, for four times five times six make a hundred and twenty. Seventhly, when the
number twenty has been taken in, which is the beginning of the reduction of mankind, I mean
twenty, and being added to itself two or three times, so as to make twenty, forty, and sixty, these
added together make a hundred and twenty. But perhaps the number a hundred and twenty is not
the general term of human life, but only of the life of those men who existed at that time, and who
were to perish by the deluge after an interval of so many years, which their kind Benefactor
prolonged, giving them space for repentance; when, after the aforesaid term, they lived a longer
time in the subsequent ages.

(92) On what principle it was that giants were born of angels and women? (\#Ge 6:4). The
poets call those men who were born out of the earth giants, that is to say, sons of the Earth.\footnote{the
Greek name Gigas is said to be derived from ge— and gennao—, ``to bring forth.''} But Moses here uses this
appellation improperly, and he uses it too very often merely to denote the vast personal size of the
principal men, equal to that of Hajk\footnote{hajk is an addition of the Armenian translator; it is the name of a fabulous
patriarch of the Armenian nation.} or Hercules. But he relates that these giants were sprung from a
combined procreation of two natures, namely, from angels and mortal women; for the substance of
angels is spiritual; but it occurs every now and then that on emergencies occurring they have
imitated the appearance of men, and transformed themselves so as to assume the human shape;
as they did on this occasion, when forming connexions with women for the production of giants. But
if the children turn out imitators of the wickedness of their mothers, departing from the virtue of
their fathers, let them depart, according to the determination of the will of a depraved race, and
because of their proud contempt for the supreme Deity, and so be condemned as guilty of
voluntary and deliberate wickedness. But sometimes Moses styles the angels the sons of God,
inasmuch as they were not produced by any mortal, but are incorporeal, as being spirits destitute
of any body; or rather that exhorter and teacher of virtue, namely Moses, calls those men who are
very excellent and endowed with great virtue the sons of God; and the wicked and depraved men
he calls bodies, or flesh.

(93) What is the meaning of the expression: ``God considered anxiously, because he had made
man upon the earth; and he resolved the matter in his mind?'' (\#Ge 6:6).\footnote{the translation of our
Bible is, ``It repented God that he had made man upon the earth.''} Some persons imagine that it is intimated by
these words that the Deity repented; but they are very wrong to entertain such an idea, since the
Deity is unchangeable. Nor are the facts of his caring and thinking about the matter, and of his
agitating it in his mind, any proofs that he is repenting, but only indications of a kind and
determinate counsel, according to which the displays care, revolving in his mind the cause why he
had made man upon the earth. But since this earth is a place of misery, even that heavenly being,
man, who is a mixture compounded of soul and body, from the very hour of his birth to that of his
death, is nothing else but the slave of the body. That the Deity therefore should meditate and
deliberate on these matters is nothing surprising; since most men take to themselves wickedness
rather than virtue, being influenced by the twofold impulse mentioned above; namely, that of a body
by its nature corruptible, and placed in the terrible situation of earth, which is the lowest of all
places.

(94) Why God, after having threatened to destroy mankind, says that he will also destroy all the
beasts likewise; using the expression, ``from man to beast, and from creeping things to flying
creatures;'' for how could irrational animals have committed sin? (\#Ge 6:7). This is the literal
statement of the holy scripture, and it informs us that animals were not necessarily and in their
primary cause created for their own sake, but for the sake of mankind and to act as the servants
of men; and when the men were destroyed, it followed necessarily and naturally that they also
should be destroyed with them, as soon as the men, for whose sake they had been made, had
ceased to exist. But as to the hidden meaning conveyed by the statement, since man is a symbol
for the intellect which exists in us, and animals for the outward sense, when the chief creature has
first been depraved and corrupted by wickedness, all the outward sense also perishes with him,
because he had no relics whatever of virtue, which is the cause of salvation.

(95) Why God says, I am indignant that I made them? (\#Ge 6:7). In the first place, Moses is
here again relating what took place, as if he were speaking of some illustrious action of man, but,
properly speaking, God does not feel anger, but is exempt from, and superior to, all such
perturbations of spirit. Therefore Moses wishes here to point out, by an extravagant form of
expression, that the iniquities of man had grown to such a height, that they stirred up and provoked
to anger even that very Being who by his nature was incapable of anger. In the second place he
warns us, by a figure, that foolish actions are liable to punishment, but that those which proceed
from wise and deliberate counsel are praiseworthy.

(96) Why it is afterwards said, that Noah found grace in the sight of the Lord? (\#Ge 6:8). In the
first place, the time calls for a comparison; since all the rest of mankind has been rejected for their
ingratitude, he places the just man in the place of them all, asserting that he had found favour with
God, not because he alone was deserving of favour, when the whole universal body of the human
race had had benefits and mercies heaped on them, but because he alone had seemed to be
mindful of the kindnesses which he had received. In the second place, when the whole generation
had been given over to destruction, with the exception of one single family, it followed inevitably that
that remaining household should be asserted to have shown itself worthy of the divine grace, that it
might be, as it were, a seed and a spark of a new race of mankind. And what could be a greater
grace and mercy than that the man, of whom this is said, should be at the same time the end and
beginning of the family of mankind?

(97) Why does Moses enumerate the generations of Noah with reference not to his ancestors but
to his virtues? (\#Ge 6:9). He does this in the first place, because all the men of that age were
wicked: secondly, he is here imposing a law upon the will, because, to an anxious follower of virtue,
virtue itself stands in the place of a real generation, if indeed men are the means of the generation
of men, but the virtues of minds. And on this account it is that he says, he was a just man, perfect,
and one who pleased God; but justice, and perfection, and grace before God, are the greatest of
virtues.

(98) What was the meaning of Moses when he says, ``And all the earth was corrupt in the sight of
God, and the earth was filled with iniquity?'' (\#Ge 6:11). Moses himself has given us the
reason why he speaks thus, in the sentence in which he asserts that iniquity had arisen by reason
of the corruption of the earth; for deliverance from iniquity is righteousness, both in all the parts of
the world, in heaven, that is, and earth, and among men.

(99) What is the meaning of his saying, ``All flesh had corrupted his way upon the earth?'' (Genesis
6:12). In the first place, the sacred historian calls the man who is devoted to the love of himself,
flesh; therefore, when he had already said he was flesh, he introduces not the same flesh, but the
flesh of the same being, namely, of man, or perhaps he is speaking even of man abstractedly
considered; for every one who passes a life destitute of all civilisation, and bewildered by
intemperance, is flesh. In the second place, he supposes here the cause of spiritual corruption to
be, as in truth it is, the flesh, because that is the seat of desire; and from it, as from a living spring,
arise all the peculiar appetites, and passions, and other affections. In the third place, he very
naturally says, that all flesh had corrupted his way; for ``his'' is a partial case, declined from the
nominative case of the pronoun ``he, she, or it;'' for as for the being to whom we refer honour, we
scarcely dare to speak of him by his own name, but we call him He. And it is from this that the
principle of the Pythagorean philosophers was derived, who said, ``He said it,'' speaking of their
master in a glorious manner, since they feared to speak of him by name. And the same custom
has obtained in cities and in private houses; for the servants, when speaking of the arrival of their
master, say, ``Here he comes;'' and so when the prince of any individual city arrives, they use the
same form of speech, ``He comes,'' when they speak of him. But what is the purpose of this prolix
enumeration of all these instances on my part? The truth is, that I wished to show that it is the
Father of the universe who is spoken of here; since, indeed, all his good qualities, and all his
marvellous names, are widely celebrated by the praise bestowed upon the virtues; and, therefore,
out of reverence he has used that name more cautiously, because he was about to bring on the
world the destruction of the flood; but the case of the pronoun ``He'' is used by way of honour in
these phrases. ``All flesh had corrupted His ways,'' inasmuch as it is truly convicted of having
corrupted the way of the Father, in accordance with the lusts, and desires, and pleasures of the
body; for these are the enemies and opposers of the laws of continence, and parsimony, and
chastity, and fortitude, and justice; by which the road which leads to God is found out and widened,
so that it should everywhere be a beaten and plain road.

(100) What is the meaning of the statement, ``All the time of man has come against me, because
the earth is filled with iniquity?'' (\#Ge 6:13).\footnote{the version given here does not in the least resemble that
in our Bible.} Those who resist the order of fate proceed upon these and many other arguments,
especially in that of sudden death, which oftentimes produces great slaughter in a short period of
time; as, for instance, in the overthrow of houses, in conflagrations, in shipwrecks, in civil tumults,
in battles of cavalry, in wars by land and in wars by sea, and in pestilences. To all those who
advance arguments of this kind we repeat the same assertions which are here made by the
prophet, on the principle which is derived from himself. If indeed that expression, ``All the time of
man has come against me,'' has a meaning of this kind, the term which has been determined as
the period of living for all mankind, behold it is now brought to one point and terminated at once by
the deluge; and since this is the case, they will not live any longer according to the principle of fate
which has been fixed; so that the time of each separate individual is now reduced to one, and has
received its destined termination at the same time, by I know not what harmony and periodical
revolution of the stars, by which bodies the whole race of mankind is continually preserved or
destroyed. Let those, therefore, all receive these things in any manner in which they choose who
study these things, and those too who argue against them. Nevertheless we must first of all make
this statement, that nothing can be found so contrary to, so opposite to, so wholly repugnant to, the
wonderful virtue of the Deity as iniquity; therefore, after he said, ``All the time of all mankind has
come up against me,'' he adds also the reason of its contrariety to him, that the earth is filled with
iniquity. In the second place, Time, under the name of Chronos or Saturn, is looked upon as a god
by the wickedest of men, who are desirous to lose sight of the one essential Being, on which
account he says, ``The time of all mankind has come up against me,'' because in fact the heathen
make human time into a god, and oppose him to the real true God. But, however, it is now
insinuated, in other passages also of scripture, which run thus, ``Time has departed to a distance
from them, but the Lord is in Us:''\footnote{\#nu 14:9. Compare with this Isaiah 8, \#Jer 46:21—28, Psalm
80:16.} just as if he were to say, time is looked upon by wicked men as the cause of the world, but by
wise men and virtuous men time is not looked upon in this light, but God only, from whom all times
and seasons do proceed. Again, God is the cause, not of all things, but only of good things and
good men, and of those men and things which are in accordance with virtue; for as he is free from
all wickedness, so likewise he cannot be the cause of it. In the third place, by that expression which
he uses in this manner, he indicates the excess of impiety, saying, ``that the time of all mankind has
arrived,'' that is to say, that all men, in every part of the world, have agreed together, with one mind,
to work wickedness; but the other assertion which is here made, that the whole earth is filled with
iniquity, amounts to this, that there is no part of it whatever free from wickedness, and which is also
to receive and to bear righteousness. And the expression, ``against me,'' establishes the proof of
what has been said, inasmuch as it is only the judgment of divine election which is altogether firm
and lasting.

\xchapter{Questions and Answers on Genesis, II}

(1) What is the preparation of Noah? (\#Ge 6:14). If any one should wish to make an
examination of the question of that ark of Noah's on more natural principles, he will find it to have
been the preparation of the human body, as we shall see by the examination of each particular
respecting it separately.

(2) Why does he make the ark of squared pieces of wood? (\#Ge 6:16). He does this in the
first place, because the figure of a square, wherever it may be placed, is steady and firm,
consisting as it does of right angles, and it is confirmed in a purer and clearer manner by the
nature of the human body. In the second place, he does this because, although our body is an
instrument, and although every portion of it is rounded off, nevertheless the limbs which are
compounded of all these portions do, by some manner or other, evidently reduce that circular orb
to the figure of a quadrangle or square. For example, take the breast which is rather square than
circular; in the same manner take the belly, after it is swollen with food or by any natural excess, for
there are some men potbellied by nature, who are to be excepted from our present argument. But if
any one looks upon the arms and hands, and back and thighs, and feet of a man, he will find all
these limbs compounded of a mixture of the square, with the circular figure at the same time. In the
third place, a quadrangular piece of wood shows in its extension nearly every sort imaginable of
uneven distinction, inasmuch as its length is greater than its breadth, and its breadth greater than
its depth. And such also is the formation of our bodies, which are compounded of one extension
which is great, of another which is of moderate size, of another which is small, great in its length
and small in its depth.

(3) Why does God say, you shall make the ark in nests? (\#Ge 6:14).\footnote{the word in our Bible is
rooms, not nests.} He gives this order very naturally, for the human body is formed of holes like nests;
every one of which is nourished and grows like a young bird, a certain spiritual force which exists in
it from its earliest origin penetrating through it, as, for instance, some of the holes and nests are
the eyes, in which the faculty of sight has its abode; other nests are the ears, which are the place
where hearing is situated. A third class of nests are the nostrils, in which the sense of smell is
lodged. The fourth nest, which is of larger dimensions than those already mentioned, is the mouth,
which is the seat of the taste; and it has been made of large size, since, besides taste, there is
also another still more important instrument, which is that of articulate speech, reposing in it,
namely, the tongue, which, as Socrates was wont to say, by beating in every direction in various
manners, and by touching different parts, composes and forms a word, being, in truth, an
instrument under the immediate guidance of reason. And the nest is placed under the skull, and
that which is called the membrane of the brain is a certain nest, as it were, of the genius of each
man: as also the chest is a nest, in which abide the lungs and the heart, and both these things are
receptacles of other internal organs; the lungs being the place in which the power of breathing is
lodged, and the heart being the abode of both the blood and the breath, for it has two ventricles,
which are, as it were, a certain kind of nests or receptacles in the breast; blood, from which the
veins, as if they could perceive its operations, are irrigated; and a breathing-hole, which again is
extended over and irrigates the perceptive channels of respiration. And both the harder as well as
the softer parts do, like nests prepared for the purpose, nourish the bones as real nests nourish
young birds; the harder portion of which, namely, the marrow, is the nest, and the softer flesh is the
nest of pleasure and pain; and if any one should wish to investigate the other parts, he will find that,
in every respect, the nature of man has much the same foundation as the ark.

(4) Why does God command the ark to be smeared with pitch, both on the inside and on the
outside? (\#Ge 6:14). Pitch is so called by reason of its bird-lime like tenacity, because it glues
together whatever was disunited before, so as to form one indissoluble and indivisible joint. For
everything which is held together by bird-lime is immediately held to a natural union; but our body
being composed of many parts is united on the outside, and is held together by its own proper
habit, but the previous habit of connection which binds those things together is the soul, which,
being situated in the middle, penetrates through every part till it reaches the surface, and then is
turned back again from the surface to the centre, so that our spiritual nature is rolled up compactly
in a double fold, being united in a firm solidity and union. Therefore this ark is smeared with pitch,
both on the inside and the outside, for the reason here given. But that ark which is placed in the
holy of holies, and as covered over with gold, is the similitude of the world appreciable only by the
intellect, as is declared in the account given of it: since just as there is a world appreciable by the
intellect incarnate in incorporeal figures existing at the same time, consisting of a union of all
figures by a certain invisible harmony; for, in proportion as gold is a more noble material than pitch,
in the same ratio is that ark, which is in the holy of holies, superior to this one of which we are now
speaking. And again, God ordained that its measure should be quadrangular, from a regard to
usefulness; but his object in the other ark was not so much that it should be useful as that it should
be exempt from all possibility of decay; for the nature of incorporeal things, appreciable only by the
intellect, is to be exempt from decay, being incorruptible and permanent. The one ark is tossed to
and fro by the winds and the waters, but the other has its station constantly in the holy of holies;
and being stable it is akin to divine nature, as the other, which is tossed about in every direction,
and moved from one place to another, is akin to and the emblem of created nature. Besides this,
that ark of the flood being, as it were, an example of corruption, is raised on high, but the other,
which is in the holy of holies, imitates the incorruptible condition of eternity.

(5) Why did God give the measures of the ark in the following manner; the length to be of three
hundred cubits, and the breadth thereof to be fifty cubits, and the height to be thirty cubits: and
above it was to be raised to a point in one cubit, being brought together gradually like an obelisk?
(\#Ge 6:15). It was necessary that so vast a work should be constructed in conformity with
literal directions, in order that so many animals, some of them of vast size, should be received into
it, as individuals of each class were introduced with the food necessary for them; but if the matter
is considered properly with reference to its symbolical meaning, then, for the comprehension of the
formation of our body, we shall require to make use not of the quantity of cubits, but of the certain
principles and proportions which are observed in them. But the proportions which are contained in
them are of sixfold, and double, and other portions are added. For three hundred is six times as
many as fifty, and ten times as many as thirty; and again fifty is by two thirds a larger number than
thirty. Such then are also the proportions of the body; for if any one should choose to investigate
the matter and inquire into it carefully in all its points, he will find that man is made in an exact
proportion of measurement, neither being too long or too little; and if a string be let down from his
head to his feet, he will find that to reach that distance it requires a string six times as long as the
width of his chest, and ten times as long as the depth of his ribs and their breadth as a second part
of depth added thereto. Such is the certain proportion, received in accordance with nature, of the
human body formed on exact measurement of the most excellently made men, who are incorrect
neither in the way of excess nor of defect. But again, it was with great wisdom and propriety that
God ordained the summit to be completed in one cubit; for the upper part of the ark imitates the
unity of the body; the head being forsooth as the citadel of the king, having for its inhabitant the
chief of all, the intellect. But those parts which are below the head are divided into separate
portions, as for instance into the hands, and in an especial degree into the lower parts, since the
thighs, and legs, and feet are all kept distinct from one another, therefore whoever should wish to
understand these matters, on the principle which I have pointed out, will easily comprehend the
analogy of the cubits as I have related it. But above all things he must not be ignorant that each of
these different numbers of cubits has separately a certain necessary proportion and principle,
beginning with the first, those in the length of the ark. Therefore in its length it is composed of three
hundred units, placed next to one another in continuation, according to the augmentation of units,
from these twenty-four numbers, one, two, three, four, five, six, seven, eight, nine, ten, eleven,
twelve, thirteen, fourteen, fifteen, sixteen, seventeen, eighteen, nineteen, twenty, twentyone,
twenty-two, twenty-three, twenty-four. But the twenty-fourth number is above all others a natural
number, being distributed among the hours of day and night, and also among the characters of
language, \footnote{he is referring to the Greek alphabet, which consists of twenty-four letters.} and literal speech; and it
is also compounded of three cubes, being complete, full, and compacted in equality. For the
number three constantly exhibits, as belonging to itself, the first equality of all, having a beginning,
and a middle, and an end, all of which are equal to one another; and eight is the first cube, because
it again has declared its first equality with the rest. But the number twenty-four has likewise a great
number of other virtues, since it is the substance of the number three hundred, as has been
already pointed out; this then is its first virtue; and it has another, since it is compounded of twelve
quadrangular figures, joined to one another by a continuous unity; and besides of two long figures,
and twelve double figures, being forsooth compounded of twos separately increased by two and
two. Therefore the angular numbers which make up together the twelve quadrangular figures are
these; one, three, five, seven, nine, eleven, thirteen, fifteen, seventeen, nineteen, twenty-one, and
twenty-three; but the quadrangular figure combines the following numbers, one, four, nine, sixteen,
twenty-five, thirty-six, forty-nine, sixty-four, eighty-one, a hundred, a hundred and twenty-one, and a
hundred and twenty-four. But those angular numbers which compose the other long figures are
these; one, four, six, eight, ten twelve, fourteen, sixteen, eighteen, twenty, twenty-two, twenty-four,
being twelve in all; and after these come the compound numbers, two, six, twelve, twenty, thirty,
forty-two, fifty-six, seventytwo, ninety, a hundred and ten, a hundred and thirty-two, and a hundred
and fifty-six; being also twelve. And if you put together the twelve quadrangular figures, you will find
a hundred and forty-four, and if you add the other twelve long figures, you will find a hundred and
fifty-six; and from the combination of the two you will get the number three hundred, and the
concord of full, and complete, and perfect nature rising up to the equal and infinite harmony; for a
complete and perfect nature is the maker of equality, according to the nature of a triangle; but the
equal and the infinite are the factors of inequality, according to the composition of the other long
figure. But the universe consists of a combination of equality and inequality, on which account the
Creator himself, even amid the destruction of all earthly things, placed a sort of fixed pattern of
stability in the ark. This then is enough to say about the number three hundred. We must now
proceed to speak of the fifty cubits, on the following principle; for in the first place it is composed of
the right angle of the quadrangular figures; for a right angle is compounded of three, four, and five;
and the square of these is nine, sixteen, and twenty-five, the sum total of which when added
together is fifty; in the second place, the perfect number fifty is composed of these four triangles
linked together, one, three, six, ten; and again of these four equal quadrangles also united
together, one, four, nine, sixteen; therefore these triangles when collected together make twenty;
and the quadrangles make thirty; and twenty and thirty added together make fifty. But if the triangle
and the quadrangle are added together, they make a heptangular figure: so that it is contained by
its virtue in the number of fifty, that divine and holy number; to which the prophet had regard when
he proclaimed the jubilee festival; and the whole of the jubilee year is free and a deliverer. The third
theorem is three triangles beginning with the unit, connected together in a continuous series, and
three cubes beginning also with the unit, and connected together in a similar manner, which
together make fifty; the examples of the first are one, four, and nine, which make fourteen; the
examples of the second are, one, eight, and twentyseven, which together make thirty-six; and the
sum total of the two when added together is fifty. Again, thirty is in an especial manner a natural
number, for as in the series of units the number three is, so is the number thirty in the series of
decimals; and that makes up the cycle of the moon, being the collection of separate months in full
delineation; secondly, it is composed of four numbers, which are united in the continual series of
these quadrangular figures, one, four, nine, and sixteen, which together make up thirty; on which
account it was not without some foundation and sufficient reason that Heraclitus called that number
``generation,'' when he said: a man in thirty years from the time of his birth can become a
grandfather, inasmuch as he arrives at the age of puberty in his fourteenth year, at which age he is
capable of becoming a father; and at the end of the year his offspring arrives at the birth, and
again in fifteen years more begets another offspring like himself; and out of these names of
grandfathers, fathers, and sons, as also out of the names of grandmothers, mothers, and
daughters, a generation complete in its offspring is produced.

(6) What is the meaning of a door in the side: for he says, ``Thou shalt make a door in the side?''
(\#Ge 6:16). That door in the side very plainly betokens a human building, which he has
becomingly indicated by calling it, ``in the side,'' by which door all the excrements of dung are cast
out. In truth, as Socrates says, whether because he learnt it from Moses or because he was
influenced by the facts themselves, the Creator, having due regard to the decency of our body, has
placed the exit and passage of the different ducts of the body back out of the reach of the sense,
in order that while getting rid of the fetid portions of bile, we might not be disgusted by beholding the
full appearance of our excrements. Therefore he has surrounded that passage by the back and
posteriors, which project out like hills, as also the buttocks are made soft for other objects.

(7) Why has he said that the lower part of the ark was to be made with two and with three stories?
(\#Ge 6:16). He has here admirably indicated the receptacles of food, calling them the inner
parts of the house; since food is corruptible, and what is corruptible belongs to the inward part,
because it is borne downwards, since some small portions of meat and drink which we take are
borne upwards, but the greater part is secreted and cast out into dung; and the intestines have
been made in two and in three stories, because of the providence of the Creator in order to supply
abundant support to his creatures; for if he had made the receptacles of food and its passage
having a direct communication between the bowels and the buttocks, some awkward
circumstances must have taken place; in the first place there would have been a frequent
deficiency, and want and hunger, and sudden evacuations also arising from divers unseasonable
events; in the second place there would have been an immense hunger, for when the receptacles
are emptied, it is inevitable that hunger and thirst must immediately supervene, like absolute
mistresses in difficulty from pregnancy, and then it follows also that the pleasant appetite for food
must be perverted into greediness and into an unphilosophical state; for nothing is so very
inconvenient as for the belly to be empty. And in the third place, there will be death waiting at the
door; since those persons must speedily be overtaken by death who the very moment that they
have done eating begin again to be hungry, and the moment that they have drunk are again thirsty,
and who before they are thoroughly filled are again evacuated and oppressed by hunger; but owing
to the long coils and windings of the bowels we are delivered from all feelings of hunger, from all
greediness, and from being prematurely overtaken by death; for while the food which has been
taken remains within us, not for such a time only as the distance to be passed requires, but for so
long as was necessary for us, a change in it is effected; since by the pressure to which it is
subjected, the strength of the food is extracted in the first instance in the belly; then it is armed in
the liver, and drawn out; after that whatever predominant flavour there was is emitted upwards to
the separate parts, in the case of boys in order to contribute to their growth, and in the case of
fullgrown men to add to their strength; and then nature, collecting the remaining portions into dung
and excrement, casts them out. Therefore a great deal of time is necessarily required for the
arrangement of so many and such important affairs, nature effecting its operations without
difficulty by perseverance. Moreover the ark itself appears to me to be very fitly compared to the
human body; for as nature is exceedingly prolific of living creatures, for that very reason it has
prepared an opposite receptacle similar to the earth for the creatures corrupted and destroyed by
the flood; for whatever was alive and supported on the earth, the ark now bore within itself in a
more general manner, and on that account God ordained it, being borne upon the waters as it was,
to be as it were like the earth, a mother and a nurse, and to exhibit the fathers of the subsequent
race as if pregnant with it, together with the sun and moon, and the remaining multitude of the stars,
and all the host of heaven; because men beholding by means of that which was made by art, a
comparison and analogy to the human body, might in that manner be more manifestly taught, for
this was the cause of the various disputes among mankind; since there is nothing which has so
much contributed to keep man in a servile condition as the essential humours of the body, and the
defects which arise in consequence of them, and most especially the vicious pleasures and
desires.

(8) Why does he say that the deluge will be to the corrupting of all flesh in which there is the breath
of life beneath the heaven? (\#Ge 6:17). One may almost say that what he had previously
spoken in riddles he has now made plain; for there was no other cause for the corruption of
mankind, except that, being slaves to pleasure and to desire, they did everything, and were anxious
about everything for that reason only; moreover they passed a life of extreme misery. But he
added also, in a very natural manner, the place where the breath of life is, using the expression,
``under heaven,'' because forsooth there are living beings also in heaven; for a happy body has not
been made out of a heavenly substance, as if in truth it had received some peculiar and admirable
condition, superior to that of other living creatures, but heaven appears to have been made
especially worthy of and for the sake of these admirable and divine living beings, all of which are
intellectual spirits; so that they give a share and participation in themselves and in the essence of
vitality even to the creatures which exist upon the earth, and give life to all those which are capable
of receiving it.

(9) Why does he say, all things which existed upon the earth shall be consumed; for what sin can
the beasts commit? (\#Ge 6:13). In the first place, as, when a sovereign is slain in battle the
military valour of the kingdom is also crushed, so also he now has thought it reasonable that when
the whole human race, bearing analogy to a sovereign, is destroyed, he should also destroy
simultaneously with it the species of beasts likewise, on which account also in pestilences the
beasts die first, and especially those which are bred up with and associate with men, such as dogs
and similar animals, and afterwards the men die too. In the second place, as, when the head is cut
off, no one blames nature if the other portions of the body also, numerous and important as they
are, are destroyed along with it, so too now no one can find fault with anything, since man is as it
were the head and chief of all animals, and when he is destroyed it is not at all strange if all the rest
of the beasts are destroyed also along with him. In the third place, animals were originally made,
not for their own sakes, as has been said by the philosophers, but in order to do service to
mankind, and for their use and glory; therefore it is very reasonable that when those beings are
destroyed for the sake of whom they had their existence, they also should be deprived of life, and
this is the reason of this assertion in its literal sense; but with respect to its hidden meaning we
may say, when the soul is exposed to a deluge from the overflow of vices, and is in a manner
stifled by them, those portions also which are on the earth, the earthly parts I mean of the body,
must of necessity likewise perish along with it; for life passed in wickedness is death; the eyes
though they see perish, inasmuch as they see wrongly; the ears also though they hear perish,
inasmuch as they hear wrongly; and the whole body of the senses perishes, inasmuch as they are
all exercised wrongly.

(10) What is the meaning of the expression, I will set up my treaty with you? (\#Ge 6:9). In the
first place, he here warns us that no man is the inheritor of the divine substance, except him who is
endowed with virtue; since the inheritance of men is possessed when they themselves are no
longer in existence, but when they are dead; but as God is everlasting he grants a participation in
his inheritance to wise men, rejoicing at their entering into possession of it; for he who has entered
into possession of everything is in want of nothing, but they who are in distress from a want of all
things are in the possession of no portion of truth. And on this account God, showing himself
favourable to the virtuous, benefits them, bestowing on the those things of which they have need. In
the second place, he bestows on the wise man a certain and more ample inheritance; for he does
not say, I will set up my treaty for you, but with you; that is to say, you are yourself a just and true
treaty, which I will set up for the race endued with reason, who have need of virtue, for a
possession and a glory to them.

(11) Why does he say: ``Enter thou and all thy house into the ark, because I have seen that thou art
a just man before me in that generation?'' (\#Ge 7:1). In the first place, certain faith receives
approbation, inasmuch as for the sake of one man who is just and worthy many men are saved by
reason of their relationship to him; as is the case too with sailors and armies, when the one have a
good captain and the others an excellent and skilful general. In the second place, he extols the just
man with praise, who thus acquires virtues, not for himself alone, but also for his whole family,
which in this way deserves safety. And it is with peculiar propriety that this expression is added,
namely, ``I have seen that thou art a just man before me;'' for men approve of the life of any one
upon one principle, and God on quite a different one; for they judge by what is visible, but he
derives his tests from the invisible designs of the soul. Moreover, that is a very remarkable
expression which is added as an insertion, namely, the one which says, ``I have seen that thou art
a just man in this generation;'' that he might not appear to condemn those who had gone before,
nor cut off the future hope of coming generations. This is the sense of the passage taken
according to the letter. But if we look at its inward meaning, when God will save the intellect of the
soul, which is the principal part of the man, that is to say, the head of the family, then also he will
save the whole family along with him; I mean all the parts, and all those who bear an analogy to the
parts, and to the word which is uttered, and to the circumstances of the body; for what the intellect
is in the soul, that also is the soul in the body. All the parts of the soul are in good condition, owing
to the result of counsels, and all its family derives the benefit along with it. But when the whole soul
is in a good condition, then also its habitation is again found to be benefited by purity of morals and
sobriety, those overstrained desires which are the causes of diseases being cut off.

(12) Why does he order seven of each of the clean animals, male and female, to be taken into the
ark, but of the unclean animals only two, male and female, in order to preserve seed upon all the
earth? (\#Ge 7:2). By divine ordinance he has asserted the number seven to be clean, and the
number two to be unclean; since the number seven is clean by nature, inasmuch as that is a virgin
number, free from all admixture, and without any parent. Nor does it generate any thing, nor is it
generated, as each of those numbers which are below the number ten, on account of their
similitude to the unit, because it is uncreated and unbegotten, and nothing is generated by it,
although it is itself the cause of creation and generation; because it rouses the virtues of all things
which are well-arranged, for the generation of created beings. But the number two is not clean. In
the first place, because it is empty, not solid; and because it is not full, therefore neither is it clean;
because it is likewise the beginning of infinite immensity by reason of its materiality. It also labours
under inequality on account of the other long numbers; for all the other numbers after two which are
increased in a twofold proportion are long numbers. But that which is unequal is not clean, as
neither is that which is material; but that which proceeds from such is fallible and inelegant, being
destitute of the purity of reason to conduct it to completeness and perfection; and it conducts it to
such by its own intrinsic power, and by songs of harmony and equality. This is enough to say on the
physical part of the subject; it remains for us to speak of its moral bearings. The irrational parts of
our soul which are destitute of intellect are divided into seven; that is to say, into the five senses,
and the vocal organ, and the seminal organ. Now these in a man endued with virtue are all clean,
and by nature feminine, inasmuch as they belong to the irrational species; but to a man who has
come into full possession of his inheritance they are masculine; for men endued with virtue are
also the parents of the virtue of counsel to themselves, the best part of them not permitting them to
come to the external senses in a precipitate and unbridled manner, but repressing them and
leading them back to right reason. But in the wicked man there exists a twofold wickedness; since
the injust man is full of doubts and perplexities, as a hesitating person, mingling things which ought
not to be mixed, and connecting them with one another, confounding those things which may very
easily be kept separate. Such are those passions which imbue the soul with some particular
colour, like a man spotted and leprous in body, the originally sound counsel being infected and
contaminated by that which is destructive and fatal. But the principle of the entrance and of the
custody of animals is added in a natural manner; for he says, ``for the sake of nourishing seed.'' If
we take the expression according to the letter, inasmuch as, although particular individuals may be
destroyed, still at least a race is preserved to be the seed of future generations; forsooth that the
intention of God, conceived at the formation of the world, might remain for ever and ever
unextinguishable, the different races of creatures being preserved. But if we regard the inward
meaning of the words, it is necessary that in the irrational parts of the soul, likewise, there should
be motions which are clean, as certain seminal principles, although the animals themselves are not
clean; since the nature of mankind is capable of admitting contrarieties, for instance, virtue and
wickedness; each of which he delineated at the creation of the world, by the tree bearing the name
of the tree of knowledge of good and evil. Forsooth our intellect, in which there is both knowledge
and intelligence, comprehends both good and evil; but good is akin to the number seven, and evil is
the brother of duality. Moreover, the law of wisdom, which abounds in beauty, says expressly and
carefully, that seed is to be nourished, not in one place only, but in all the earth, both naturally, in
the first instance, and also morally, in its peculiar sense; because it is very natural, and suitable to
the character of God, to cause that which in all parts and divisions of the world is said again to be
the seed of living beings, to fill places which have been evacuated a second time with similar
creatures, by a repeated generation; and not altogether to desert our body, inasmuch as it is an
earthly substance, as if it were a thing deserted by and void of all principle of life. Since, if we
practise the drinking of wines and the eating of meats, and indulge in the ardent desire of the
female, and in short practise in all things a delicate and luxurious life, we are then only the bearers
of a corpse in the body; but if God, taking compassion on us, turns away the overflow of vices and
renders the soul dry, he will then begin to make the body living, and to animate it with a purer soul,
the governing principle of which is wisdom.

(13) Why, after the entrance of Noah into the ark, did seven days elapse, after which the deluge
came? (\#Ge 7:10). The kind Saviour of the world allows a space for the repentance of
sinners, in order that when they see the ark placed in front of them as a sort of type, made with
respect to the then present time, and when they see all the different kinds of living creatures shut
up in it which the earth used to bear on its surface, according to its parts adapted to the different
species, they might believe the predictions of the deluge which had been made to them, so that,
fearing total destruction above all things, they might be speedily converted, destroying and
eradicating all their iniquity and wickedness. In the second place, this language is a most manifest
representation of the exceeding great abundance of the kind mercy of the beneficent Saviour, by
destroying the wickedness of many years, which from the time of their birth to old age has
extended itself over their conduct in those persons who practise penitence for a few days, for the
divine nature forgets all evil and is a lover of virtue. When therefore it beholds faithful virtue in the
soul, it gives it honour in a wonderful degree, in order, in the first place, to take away all kinds of evil
which impend over it from its sins. In the third place, the number of seven days after the entrance
of Noah into the ark, during which the command of God kept off the flood, is a recollection of the
creation of the world, the birthday festival of which is kept on the seventh day, showing manifestly
the authority of the Father; just as if he were to say, ``I am the Creator of the world, commanding
things to exist which have no existence; and at the same time I am he who am now about to
destroy the world with a great flood. But the original cause of the creation of the world was the
goodness which is in me, and my kindness; and the cause of its impending destruction is the
ingratitude and impiety of those persons who have been loaded by benefits by me.'' Therefore he
causes an interval of seven days, in order that the unbelieving may remember, and that those who
have abandoned their faith in the Parent of the world may in a suppliant spirit return to the Creator
of all things, and so may entreat him again that his works may be everlasting; and that they may
offer their entreaty, not with mouth and tongue, but rather with the heart of amendment and
penitence.

(14) Why did the rain of the deluge last forty days and an equal number of nights? (\#Ge 7:4).
In the first place, the word day is used in a double sense. The one meaning that time which is from
morning to evening, that is to say, from the first rising of the sun in the east to his sinking in the
west. Therefore they who make definitions, say, ``That is day, as long as the sun shines on the
earth.'' In another sense, the word day is used of the day and night together. And in this sense we
say that a month consists of thirty days, combining together and computing the period of night in
the same calculation. These premises having been first laid down thus, I say that the word now
spoken of has not been incorrectly employed, inasmuch as it implies forty days and forty nights; but
is also so used in order to suggest a double number determined for the generation of mankind,
namely, forty and eighty, as many men skilled in medicine, and indeed also in physical science,
have suggested; but it is especially described in the sacred law, which was to them also the first
principle of natural science. Since therefore destruction was on the point of overwhelming all men
and women every where on account of the excessive combination of iniquities and quarrels, the
Judge of all considered it becoming to allot an equal time to their destruction to that which he had
consumed in the original creation of nature and to the work of giving life to the world; for the
principle of procreation is the perseverance of seed in the different parts; but it was necessary to
honour the male creature with pure light, which knows not the shade; but the woman had a mixture
in her body of night and darkness. Therefore, in the creation of the whole world, the excess of the
male or the unequal number, being composed of unity, becomes the parent of square numbers; but
the woman who is an unequal number, being compounded of duality, becomes the parent of other
long numbers. Moreover, the square is splendour and light combining together by the equality of
the sides; but the other numbers being long necessarily exhibit night and darkness by reason of
their inequality, since that which is in excess throws a shade over that which lies beneath the
excess. In the second place, the number forty is the produce of many virtues, as has been
suggested in another place. It is also often used for the judgment of legislation, both with reference
to those persons who have done any thing rightly deserving of praise and honour, and also with
reference to those who on account of their sins meet with reproaches and punishments; so that it
is superfluous to adduce proofs to demonstrate what is evident.

(15) What is the meaning of the expression, ``I will destroy every living substance that I have made
form off the face of the earth?'' (\#Ge 7:4). Do you not all shrink back in astonishment when
you hear these words, by reason of the beauty of the sentence? for he has not said, ``destroy from
the earth,'' but ``from the face of the earth,'' that is to say from its surface; in order, that is, that in
the lowest depth of the earth the vital efficacy of all seeds might be preserved unhurt, and free
from all injury which could possibly bring damage to it; since the Creator was not forgetful of his
original design, but destroys those only who come in his way, and who move only on the surface of
the earth, but leaves the roots in the depth, in order to produce the generation of other causes.
Moreover, that expression, ``I will destroy,'' was also written by divine inspiration; for it happens that
if we remove the letters which require to be removed, the whole table for the reception of letters
remains the same. By which he proves that he will destroy the fickle generation on account of their
impiety, but the conversation and essence of the human race he will preserve for ever and ever to
be the seed of future generations. And what follows agrees with this, since to the expression, ``I will
destroy,'' this other is also added, all natural existence, every thing which exists, or rises upon the
earth; but existence is the destruction of the opposite characteristics; and that which is dissolved
loses quality, but retains body and materiality. This is the letter of what is said. But in the inward
meaning, the flood is symbolically representative of spiritual dissolution. When therefore by the
grace of the Father we desire to throw away and to wash off all sensible and corporeal qualities by
which the intellect was infected as by swelling sores, then the muddy slime is got rid of as by a
deluge, sweet waters and wholesome fountains supervening.

(16) Why does he say: ``Noah did every thing which the Lord commanded (or ordered) him?''
(\#Ge 7:5). A noble panegyric for the just man. In the first place, because with an ingenuous
mind and a purpose full of affection towards God he performed, not a part of what he had been
commanded, but the whole of God's commands. But the second is the more true expression,
because he does not choose so much to command as to order him; for masters command their
slaves, but friends order friends, and especially elder friends order younger ones. Therefore it is a
marvellous gift to be found even in the rank of servants, and in the list of ministers of God; and it is
a superabundant excess of kindness for any one to be a beloved friend to the most glorious
Uncreated Essence. Moreover, the sacred writer has here carefully employed both names, the
Lord God, as declaratory of his superior powers of destroying and benefiting, using the word Lord
first, and placing the name God, giving the idea of beneficence, second; since it was a time of
judgment let the name which is the indication of his destroying power come first. But still, as he is a
kind and merciful king, he leaves as relics the seminal elements by which the vacant places may be
replenished, for which reason, at the fist beginning of the account of the creation, the expression,
``Let there be,'' was not an exterminating act of power, but a beneficent one. Therefore, at the
creation, he changed the appellations and use of names; but as the name God is an indication of
his beneficent power, the sacred writer has more frequently employed that in his account of the
creation of the universe, but after everything was perfected then he called him Lord, in reference
to the creation itself, for this name betokens royal power and the ability to destroy; since, where the
act of generation is God is used first in order, but when punishment is spoken of the name Lord is
placed before the name God.

(17) Why did the deluge take place in the six hundredth year of the life of Noah, and in the seventh
month, and on the twenty-seventh day of the month? (\#Ge 7:11). Perhaps it happened that the
just man was born at the beginning of the month, at the first beginning of the commencement of that
very year which they are accustomed to call the sacred year, out of honour, otherwise the sacred
historian would not have been so carefully accurate in fixing the day and month when the deluge
began to the seventh month and the twenty-seventh day of the month. But, perhaps, by this
minuteness he intended manifestly to indicate the precise time of the vernal equinox, for that
always occurs on the twentyseventy day of the seventh month. But why was it that the deluge fell on
the day of the vernal equinox? Because about that time the birth and increase of everything take
place, whether living creatures or plants; therefore the vengeance and punishment inflicted brings
with it the more terrible and dreadful threats, as happening at the period of plenty and fertility of the
shaves of corn, and indeed, in the very midst of that productiveness, and bringing the evil of utter
destruction as a reproof of the impiety of those who are exposed to the punishment. For behold,
says he, all nature contains its own productions within itself in the greatest abundance, namely,
wheat and barley, and everything else which is produced from seed, brought on to complete
generation, as, also, it begins to generate the fruits of trees; but you, like mortals, corrupt its
mercies, perverting the divine gifts, and purposes, and mysteries. But if the deluge had taken place
at the autumnal equinox, when there was nothing growing on the earth, but when all the crops were
collected into their proper storehouses, it would not have, in any degree, been looked upon as a
punishment, but rather as a benefit, as the water would have cleansed the plains and the
mountains. But as the first man who was produced out of the earth was also created at the same
season of the year, he whom the divine writer calls Adam, because in fact it was on every account
proper that the grandfather, or original parent, or father of the human race, or by whatever name
we may choose to designate that original founder of our kind, should be created at the season of
the vernal equinox, when all earthly productions were full of their fruit; but the vernal equinox takes
place in the seventh month, which is also called the first in other passages, with reference to a
different idea. Since, therefore, the first beginning of the generation of our race, after the
destruction caused by the deluge, commenced with Noah, men being again sown and procreated,
therefore he also is recognised as resembling the first man born of the earth, as far as such
resemblance or recognition is possible. And the six hundredth year has for its origin the number
six; and the world was created under the number six, therefore, by this same number does he
reprove the wicked, putting them to shame because he would, unquestionably, never, after he had
created the universe by means of the number six, have destroyed all the men who lived on the
earth under the form of six, if it had not been for the preposterous excess of their iniquities. For the
third power of six and the minor power is the number six hundred, and the mean between both is
sixty, since the number ten more evidently represents the likeness of unity, and the number a
hundred represents the minor power.

(18) What is the meaning of the expression, ``And the fountains of the deep were broken, and the
springs of heaven were opened?'' (\#Ge 7:11). The literal meaning is plain enough, for it
suggests the two extremities of the universe, the heaven and the earth, to have met together for
the destruction of mortals deserving of condemnation, the waters running forth to meet one
another from all quarters, for part of them bubbled up from out of the earth, and part descended
downwards from heaven; and in truth, that expression is very explicit, ``The fountains were broken
up,'' for when a rupture is effected then the thing confined rushes forth without any hindrance. But
with reference to the interior meaning of the expression we may as well say this: the heaven is
symbolically the human intellect, and the earth is the sense and body, therefore there is great
distress and calamity when neither remains, but when each threatens a secret attack. But what is
the exact meaning of my words? If often happens that acuteness of intellect exhibits cunning and
wickedness, and bears itself with bitterness in every respect when the lusts of the body are
restrained and bridled; but the contrary fact often prevails, and the lusts rejoice in their
opportunities and proceed onward, gaining strength from luxury and abundance of means;
therefore, the gate of these lusts is the outward sense combined with the body; but when the
intellect, neglecting outward circumstances, is consistent with itself, then the senses lie harmless,
as if completely abandoned; but when both are united, the intellect in exerting all cunning and
wickedness, and the body irrigated with all the senses and gorged with every kind of vice to satiety,
then we are exposed to a deluge; and this is in fact a great deluge, when the streams of the
intellect are opened by iniquity, and folly, and greedy desire, and injustice, and arrogance, and
impiety, and when the fountains of the body are opened by lust, and desire, and intemperance, and
obscenity, and gluttony, and lasciviousness, with relations and sisters, and all irremediable
diseases.

(19) What is the meaning of the expression, ``And the Lord shut him in, closing the doors of the
ark?'' (\#Ge 7:16). Since we have said that the structure of the human body is symbolically
indicated by the ark, we must take notice, also, that on the outside this body is enclosed by a hard
and dense skin, to be a covering to all its parts; for nature has made this as a sort of coat, to
prevent either cold or heat from being able to do man injury. The literal meaning of the expression
is plain enough, for the door of the ark is carefully shut by divine virtue for the sake of security, lest
the water should enter in at any part, as it was to be tossed about by the waves for an entire year.

(20) What is the meaning of the expression, ``And the water was greatly increased, and bore up the
ark which floated upon the water?'' (\#Ge 7:17). The literal meaning is plain enough, but it
contains an allegorical reference to our bodies, which ought to be borne up as if on the water, and
by fluctuating with our necessities to subdue hunger and thirst, cold and heat, by which it is
agitated, disturbed, and kept in motion.

(21) Why did the water overflow fifteen cubits above all the highest mountains? (\#Ge 7:19).
With respect to the literal statement we must remark that the excess was not merely one of fifteen
cubits above all high mountains but above those which were a great deal more lofty and high than
some others; therefore it was a great deal more than that height above the lower ones. But we
must interpret this statement allegorically; for the loftier mountains shadow forth the senses in our
body, because it has been permitted to them to occupy the abode of stability in the lofty region of
our head. And there are five numbers of these, each to be considered separately, so that they
amount in all to fifteen.

As, there is the faculty of sight, the thing which is visible, the act of seeing.

The faculty of hearing, the thing which is audible, and the act of hearing.

The faculty of smelling, the thing which can be smelled, and the act of smelling.

The faculty of taste, the thing which can be tasted, and the act of tasting.

The faculty of touch, the thing which can be touched, and the act of touching.

These are the fifteen cubits in excess; for they also are overwhelmed by the overflow, being
destroyed by the unseasonable influx of infinite vices and evils.

(22) What is the meaning of the expression, ``And all flesh capable of motion perished?'' (Genesis
7:21). It is with especial propriety, and strictly in accordance with natural truth, that the sacred
historian has here pronounced all flesh capable of motion devoted to destruction; for flesh excites
pleasures, and is excited by pleasures; and such affections are the causes of the destruction of
souls, as one the other hand sobriety and patience are the causes of safety.

(23) What is the meaning of the expression, ``And everything which was on the dry land died?''
(\#Ge 7:22). The literal meaning is notorious, because in that great deluge everything which
was upon the earth was destroyed and perished; but with respect to the secret meaning, as, since
the material of timber, when it is parched and dry, is readily consumed by fire, so, likewise, when
the soul is not mingled with wisdom, and justice, and piety, and the other enduring virtues, which
alone are able to impart real joy to the thoughts, then it, being parched up and dried like a plant
which is deprived of any power of budding or producing seed, or like a withered trunk, dies, being
handed over to the mercy of the overwhelming overflow of the body.

(24) What is the meaning of the words, ``It destroyed every living substance which was on the face
of the earth?'' (\#Ge 7:23). The literal meaning of these words only announces a plain
statement of a fact, but it may be turned into an allegory in this manner. It is not without reason that
the sacred historian has used the words ``a living substance,'' for that is characteristic of ambition
and pride, which lead men to despise both divine and human laws; but ambition and arrogance do
rather appear on the face of our earthly and corporeal nature with an elated countenance and
contracted brows. Since there are some persons who come nearer to one with their feet, but with
their chests, and necks, and heads lean back, and are actually borne backwards and bend away
like a balance, so that with one half of their body, in consequence of the position of their feet, they
project forward, but backward with the upper portion of their chests, drawing themselves back like
those persons whose muscles and nerves are in pain, by which they are prevented from stooping
in a natural manner. But men of this kind it was determined to put an end to, as one may see from
the records of the Lord and the divine history of the scriptures.

(25) What is the meaning of the words, ``Noah remained alone, and they who were with him in the
ark?'' (\#Ge 7:24). The literal meaning of this is evident; but with respect to its concealed sense
we may advance an opinion, that the intellect which is desirous of studying justice and wisdom
does, like a tree, discard all noxious shoots which bud forth about it, and rejects all extravagant
humours of superfluous vigour, I mean immoderate excess of the affections, and wickedness, and
all the effects of such. Therefore he is here said to have been left alone with those of his own
kindred, and his kindred are properly all those designs and thoughts of each individual, which are
regulated in accordance with virtue, on which account the statement is added ``And he remained
alone, and they who were with him,'' in order to reveal a more genuine joy; but he remained in the
ark, that is to say, in the body, because it was purified from every vice and spiritual disease, as the
intellect was not yet put in such a condition as to be wholly incorporeal. And on this account also,
we must render thanks to the merciful Father, because he received his consort and colleague no
longer as one endowed with superior power, but to be subordinate to his own power, on which
account also the body is not submerged in the deluge, but rising above the flood is not at all
destroyed by the eddies of the cataracts, which a crapulous, libidinous, vanity-loving will,
overflowing all things, raises to an eminence.

(26) Why is it that the sacred writer says, ``And God was mindful of Noah, and of the beasts, and of
the cattle,'' but does not add that he remembered his wife and children? (\#Ge 8:1). As the
husband agrees with and is equal to his wife, and as the father is equal to his sons, there is no
need of mentioning more names than one, but one, the first, is sufficient; therefore, by naming
Noah he, in effect, names all those who were with him of his family; for when husband, and wife,
and children, and relations are all agitated by discord, then it is no longer possible for such to be
called one family, but instead of being one they are many; but when harmony exists then one family
is exhibited by one superior of the house, and all are seen to depend upon that one, like the
branches of a tree which shoot out from it, or the fruit upon a vine branch which does not fall off
from it. And in another part, also, the prophet has said, ``Have a regard to Abraham your father,
and to Sarah who brought you forth,'' where, because in fact it was one family, he displays the
agreement by mentioning the woman.

(27) Why is it that the sacred writer made mention first of the beasts and afterwards of the cattle,
saying that God remembered Noah, and the beasts, and the cattle? (\#Ge 8:1). In the first
place, that poetical rule has not been expressed in vain, that he led the bad into the middle;
therefore he places the beasts in the middle, between the domestic animals, that is to say the men
and the cattle, in order that they might be tamed and civilized by having an intimate association with
both. In the second place, he thought it scarcely reasonable to bestow a provident benefit on the
beasts by themselves, because he was about immediately to add a statement of the beginning of
the diminution of the deluge. This is the explanation of the statement taken literally. But with respect
to the inner meaning, that just intellect, dwelling in the body as if in the ark, possesses both beasts
and living animals, not those particular ones which bite and hurt, but, that I may use such an
expression, those general kinds which contain in themselves the principles of seed and origination;
since without these the soul cannot be manifest in the body. Moreover, the soul of the foolish man
employs all poisonous and deadly animals, but that of the wise man those only which have changed
the nature of wild beasts into that of domestic creatures.

(28) What is the meaning of the expression, ``He brought a breath over the earth, and the water
ceased?'' (\#Ge 8:2). Some people say that what is here meant by ``a breath'' is the wind, at
which the deluge ceased. But I am not aware that water is diminished by wind, but only that it is
disturbed and agitated into waves, for if it were otherwise the vast extent of the sea would have
been wholly dried up long ago. Therefore it appears to me that the sacred writer here means the
breath of the Deity, by which the whole universe obtains security at the same time with the
calamities of the world, and with those things which exist in the air, and in every mixture of plants
and animals. Since the deluge of that time was no trifling infliction of water, but an immense and
boundless overflow, extending almost beyond the pillars of Hercules and the great Mediterranean
Sea, since the whole earth and all the spaces of the mountains were covered with water; and it is
scarcely likely that such a vast space could have been cleared by a wind, but rather, as I have
said, it must have been done by some invisible and divine virtue.

(29) What is the meaning of the expression, ``The fountains of the deep were closed, and the
cataracts of heaven?'' (\#Ge 8:2). In the first place, it is agreed upon by all that in the first
period of forty days the waters of punishment fell uninterruptedly, the lowest fountains of the earth
being burst asunder; and from above, the cataracts of heaven being opened, and pouring down
until all places, both level and mountainous, were covered with the inundation; and for another
period of a hundred and fifty entire days the waters did not cease to fall, nor did the streams cease
to flow, nor the springs to burst up, though still in milder quantities, not so as to increase the
existing flood, but only so as to secure the duration of the existence of the deluge, which was also
assisted from on high; and this is what is indicated in the meantime by the statement that after a
hundred and fifty days the fountains and the cataracts were closed up; therefore, while as yet they
were not closed up it is plain that they were in action. In the second place, it was necessary that
that which afforded the excessive supply of waters for the deluge, namely, the double reservoir of
water, the one from the fountains of the earth, and the other from the pourings forth of heaven,
should be both closed, for the more the stores from which any material is supplied fail, the more it
is consumed by itself, especially when divine virtue has given the command. This is the literal
meaning of the expression. But with respect to the inner sense of the passage, since the deluge of
the mind arises from two things, for it arises partly from counsel, as if from heaven, and in another
degree also from the body and from sense, as if from earth, the vices being reciprocally introduced
by the passions and the passions by the vices, it was inevitably necessary that the word of the
divine physician entering in as a salutary visitation for the purpose of healing the disease, should
prevent both kinds of overflow for the future; for it is the first principle of the medical art to drive
away the cause of the infirmity and to leave no longer any materials for disease; and the scripture
teaches this, also, in the case of the leper, for when the leprosy is checked and is prevented from
extending further, it then fixes the station and abode of the leprous man in the same place by a law,
because the character of being stationary implies cleanliness, for that which is moved contrary to
nature is unclean.

(30) What is the meaning of the statement that after a hundred and fifty days the water began to
abate? (\#Ge 8:3). We must here inquire whether those hundred and fifty days, during which
the water was abating, are to be distinguished from the four months, or whether they have a
reference to the days previously mentioned, during which the deluge went on unceasingly, as still
increasing.

(31) Why does he say, ``The ark settled in the seventh month on the seven and twentieth day of the
month?'' (\#Ge 8:4). It is reasonable here to consider how the beginning of the deluge
commences in the seventh month, on the twenty-seventh day of the month, and how the diminution,
when the ark rested on the top of the mountains, again took place in the seventh month and on the
twenty-seventh day of the month; therefore we must say, that there is here an homonymy of
months and days, for the beginning of the flood took place in the seventh month, beginning at the
birthday of the just man, near the time of the vernal equinox, and its diminution took place in the
seventh month, beginning from the highest point of the flood at the autumnal equinox, since the two
equinoxes are separated from one another by seven months, having an interval of five months
between them. For the seventh month of the equinox is also by its virtue the first month, because
the creation of the world took place in it, on account of the abundance of all things at that season.
And, in like manner, the seventh month of the autumnal equinox, which, according to time, is the
first in dignity, having its principle of that number seven derived from the air; therefore, the deluge
took place in the seventh month, not according to time but according to nature, having for its
principle and commencement the spring season.

(32) Why does he say, ``In the tenth month, on the first day of the month, the heads of the
mountains appeared?'' (\#Ge 8:4). As in numerals the number ten is the extreme bound of the
units, being a definitive and perfect number, so too it is the cycle and end of the units, and also the
beginning and cycle of the decades, and of infinity of numbers; thus the Creator, on the cessation
of the deluge, condescended that the tops of the mountains should appear in the number of the
decade, being a definitive and perfect number.

(33) Why was it after forty days that the just man opened the window of the ark? (\#Ge 8:5).
We must observe carefully that the divine historian uses the same number in speaking of the influx
of the deluge and in mentioning the cessation and complete removal of the evil; forsooth on the
twenty-seventh day of the seventh month in the six hundredth year of the life of Noah, that is to say
in the six hundredth year after his nativity, the deluge began at the spring season; but on the
twenty-seventh day of the seventh month, the ark rested on the top of the mountains at the vernal
equinox. But it is plain from these circumstances that the deluge became invisible in the six
hundred and first year of Noah's life, again on the seventh month and the twenty-seventh day of it,
so that after the lapse of an entire year, it again settled and established the earth as it was at the
moment of its destruction, in the spring season, budding forth and covered with verdure and full of
all kinds of fruits. But again in a similar manner the overflow of the deluge took place for forty days,
the cataracts of heaven being opened and fountains bursting upwards from the lowest depths of
the earth; and again a hope of renewal took place at intervals of forty days after a sufficient
cessation of the rains, when he opened the window; and again the duration of the permanent
deluge lasted for a period of a hundred and fifty days, as also its gradual diminution occupied a
period of a hundred and fifty days; so that we may well admire the equality of the arrangement, for
the evil increased and ceased according to the same number, like the moon, which from its first
rise proceeds in its increase according to an equal number, going onward to its perfect fulness of
light, and then again with an equal number in its decrease, returning back to its original state, after
having been previously full; and in like manner in the case of divine chastisements, the Creator
preserves a regular order, banishing all irregularity from the divine borders.

(34) What is the window of the ark, which the just man opens? (\#Ge 8:6). The literal statement
scarcely admits of any difficulty or doubt, inasmuch as it is plain; but with reference to the inner
meaning we have this to say: each separate part of the senses has imitated the windows of the
body, since it is through them as through windows that the comprehension of sensible objects
enters into the intellect, and again it is through them that the intellect stretches forth as if escaping;
but a portion of these windows, the senses, the more noble portion too, I say, is the sight;
inasmuch as that above all the rest is akin to the soul, and it is intimately acquainted with light, the
most beautiful of the essences, and it is the minister of sacred things; moreover that is the one
which first laid open the road to philosophy. For beholding the regular motion of the sun, and of the
moon, and the erratic course of the other planets, and the unerring circular motion of the whole
heaven, and the order and harmony there existing beyond all calculation, as if it were the one real
creator of the whole world, it by itself related to its one chief counsellor and director all that it saw:
and then intellect, seeing those things with its acute eye, and by those things discerning superior
demonstrative ideas, and the cause of all those things, immediately perceived that there was a
God at the same moment that it arrived at the conception of generation and providence, because
forsooth it was plain that this visible nature was not created by itself: for it was impossible that such
a harmony, and order, and reason, and most consistent analogy, and that a concord of such a
character and extent, and that such true and perfect felicity should exist by its own power: but it
was necessary that there must be some Creator and parent of it acting like a governor and
director, who generated these things, and then having generated them preserves them safe and
sound.

(35) Why did he send out a raven first? (\#Ge 8:6). If we look to the literal statement, the raven
is said to be an animal particularly set apart for being sent on messages and employed in offices;
for to this very day many people watch its mode of flight and its chattering, judging that it gives
some intimation of unknown facts; but with respect to the hidden meaning, as a raven is a black,
and arrogant, and speedy animal, it is a sign of wickedness, which brings night and darkness over
the soul, and it is also swift to meet all the things of the world in its flight. And also that it is very
bold, so as at times to cause the destruction of those who seek to catch it, since pride produces
also rash impudence, the opposite of which is virtue, which is consistent with the brilliancy of light,
and is by nature decorated with a modest bashfulness; therefore it is quite natural that if there was
any darkness remaining behind in the intellect, darkness which exists in accordance with folly, he
should expel that and send it out beyond his borders.

(36) Why did the raven after it had gone forth not return, when there was not yet any part of the
earth dried? (\#Ge 8:7). This passage admits of an allegorical interpretation since injustice is
contrary to the light of justice; so that in comparison of the admirable actions of the man endued
with virtue, it thinks it more desirable to rejoice with its kinsman the deluge; for injustice is a lover of
confusion and corruption.

(37) Why does he speak here in an incorrect manner, ``Till the water was dried up from the earth;''
when it was not the water which was dried up from the earth, but the earth which was dried from the
water? (\#Ge 8:7). He uses this expression in an allegorical sense, indicating by the fall of the
waters the immensity of vices, by which when saturated and vigorous the soul is corrupted, but
when they are dried up and withered, it is preserved; for then they cannot inflict any mischief upon
it, since they are become impotent and dead.

(38) Why does he in the second place send forth the dove, and why does he send it forth from
himself to see whether the water had ceased, when he uses no such expressions about the
raven? (\#Ge 8:8). In the first place, the dove is a clean animal, and in the second place it is
tame, civilised, and one which associates with mankind, on which account also the honour has
been allotted to it of being offered up upon the altar in sacrifices; and on this account the sacred
writer, sanctioning this honour and adding the weight of his assertion, has said, he sent it forth from
himself, declaring by this expression that it was to see whether the water was abated, he displays
the common anxiety felt by both. But those birds, the raven and the dove, are symbols of
wickedness and virtue: for the one, whether it is wickedness or the raven, has no house, nor
habitation, nor city, being an insolent unsociable bird; but the other, namely virtue, has a regard to
humanity, and to the public good: and so the man endowed with virtue sends that bird forth as his
ambassador for desirable and salutary objects, wishing to receive from it desirable information;
and she, like an ambassador, brings us back genuine pleasure, so that what is hurtful may be
guarded against, and what is useful may be diligently and carefully admitted.

(39) Why did the dove, when it found no rest for its feet, return to Noah? (\#Ge 8:9). Is not the
reason of this evident, and is it not a plain proof that wickedness and virtue are symbolically
indicated by the raven and the dove? For behold the dove, which is the last sent out, finds no rest.
How, then, could the raven, who departed previously, while the calamity of the deluge was still
prevailing, find any place, and make a settlement? For the raven was neither a swan nor an ibis,
nor did he belong to the class of aquatic birds. But the sacred writer here points out in an
enigmatical manner, that wickedness, when it has gone forth out of doors, to the swelling whirlpools
of the vices and passions which overflow and corrupt the soul and life, joyfully admits them, and
dwells with and takes up its abode with them, as with its nearest friends and relations; but virtue,
turning away with loathing from even the first sight of them, at once springs back, and does not
return, scarcely finding rest for its feet; finding, in fact, no standing ground anywhere, and no place
worthy of itself. For what other greater evil can there be than this, that virtue should not be able to
find in the soul any place ever so small for rest and for abiding in?

(40) What is the meaning of the statement, ``Putting forth his hand, he received her, and brought
her in to himself?'' (\#Ge 8:9). The literal meaning is plain, but with respect to the hidden sense
we must elicit the truth carefully. The wise man employs truth as an overseer of and ambassador in
important affairs, which, when it perceives that those natures are worthy of it, abides among and
dwells with them, correcting them, and making them better, since wisdom is a very common, and
equal, and useful thing. But when, with reference to the opposite natures, it sees that in some
points they are preposterously redundant and in others altogether deficient, it returns to its proper
place; and the man endowed with virtue admits it in word, putting forth his hand to take it, and in fact
opening all his intellect for its reception, and unfolding it by the perfect number, full and equal, with
all imaginable promptitude. Nor even then, when he had sent her forth from himself to examine the
natures of other things, had he separated it from himself, but had only acted like the sun, which
sends forth his beams to give light to all things, because it is not at all consistent with the character
of his boundless light to be separated at all.

(41) Why did he, after waiting yet seven other days, send forth the dove a second time? (Genesis
8:10). This is an excellent example for life, since although it will behold natures obstinate at first,
still the hope of changing them into better natures is scarcely allowed to drop; and as a prudent
physician does not in a moment apply a perfect cure to a disease, or effect a complete restoration
to health, but employs salutary medicines after he has given nature an opportunity of first opening
the way to recovery, so too the man endowed with virtue behaves with respect to the employment
of the word which is in accordance with the law of wisdom. But the number seven is the sacred and
dominical number, according to which the Father of the universe, when he made the world, is said
to have looked upon his work. And the contemplation of the world, and of all the things contained in
it, is nothing else but philosophy, and that excellent and select portion of it which wisdom contains,
comprehending within itself also a work still more necessary to be seen.

(42) What is the meaning of the expression, ``The dove returned a second time to him about
evening, having in her mouth a leaf and a thin branch of olive?'' (\#Ge 8:11). All these separate
points are selected and approved signs-the returning, the returning about evening, the having an
olive-leaf and a thin branch of that tree, and oil, and the having it in her mouth; but yet every one of
these signs can be examined with a certainty beyond certainty, for the return is distinct from its
previous return, for that one bore with it an announcement of nature being wholly corrupted and
rebellious, and being wholly destroyed by the deluge, that is to say, by great ignorance and
insolence; but this second return brings the news of the world beginning to repent, but to find
repentance is not an easy task, but is a difficult and laborious business. And it is on this account
that the dove arrived in the evening, having passed the whole day from morning to evening in its
visitations; in word, indeed, examining places, but in fact investigating the different parts of nature
itself by continual visitation, and seeing them all clearly from beginning to end, for the evening is the
indication of the end. The third sign, again, is its bringing a leaf; but a leaf is a small part of a tree,
still it does not exist without a tree. And the beginning of displaying repentance is somewhat
corresponding to this, since the beginning of correction has some slight indications about it, which
we may call a leaf, by which it appears to receive guardianship, but can easily be shaken off; so
that the hope shall in that case not be great of attaining the desired improvement, which is typified
by the leaf of no other tree but of the olive alone, and oil is the material of light. For wickedness, as
I have said before, is profound darkness, but virtue is luminous brilliancy, and repentance is the
beginning of light. But you must not yet suppose that the beginning of repentance is only visible in
branches just germinating and beginning to look green, but that it exists too while they are still dry,
and while the seminal principle is dry and quiescent. And it is on this account that the fifth sign is
shown, that, namely, of the dove when it comes bearing a slender branch. And the sixth sign is that
this slender branch was in its mouth, for the number six is the first perfect number, since virtue
bears in its mouth, that is to say in its conversation, the seeds of wisdom and justice, or, in one
word, of honesty of the soul; and not only bears this, but gives some portion of participation in it
even to the foolish, by drawing up water for their souls, and irrigating them with the desire of
repentance for their sins.

(43) Why is it said, ``And Noah knew that the waters had ceased from off the earth?'' (Genesis
8:11). The literal statement is plain, since if the leaf had been taken up from off the water it would
have been wet and soaked, but now he says that it was dry and slender, as if it had become dry by
being on the earth which was dried. But with reference to its inward meaning, the wise man takes it
as a symbol of repentance, and wishes to check the calamities of excessive obstinacy by taking
the leaf, since it was not yet green, but slender, for the reason which has been already mentioned.
At the same time we may admire the Father on account of his exceeding kindness, for although
corruption had prevailed over all the men who lived on the earth from the excess of their iniquities,
still there remained some relics of antiquity and of that which was from the beginning, and a slight
seed of previous virtues; by which it is intimated nevertheless that the memory of all the good
deeds that have been done from the beginning is not wholly destroyed. On which account a certain
prophet, the kinsman and friend of Moses, uttered an oracle of this kind, ``If the omnipotent Lord
had not left us a seed, we should have been like blind and barren People,''\footnote{\#isa 1:9.} able neither
to know the truth nor to generate it. And the Chaldaeans in their native language call blindness and
sterility Sodom and Gomorrah.

(44) Why, in the third place, after seven other days, did he again send forth the dove, which did not
again return to him? (\#Ge 8:12). According to the word, the dove made no more return to him;
but what in fact is meant is virtue, which, however, is not an indication of alienation, since, as I
have said before, she was not separated from him at that time, but sent forth like a sun-beam to
pay a visit of examination to the natures of others, but then, not finding any one to listen to her
precepts of correction, she returns, and properly comes to him alone. But this time she is no longer
the possession of one single individual, but is rather a common good to all those who have been
willing to receive the emanations of wisdom as if coming up from the earth, those persons, that is,
who from the very beginning have laboured under a great thirst of perfect wisdom.

(45) Why in the six hundred and first year of the life of Noah, and on the first day of the first month,
did the waters of the deluge cease from off the earth? (\#Ge 8:13). The word first, according to
the defect of time, is spoken of with reference either to the month or the man, and each
interpretation has reason to support it; for if we are bound to maintain that the water began to
abate in the first month, we are equally obliged to consider that the sacred historian intended also
to speak of the seventh month, that is, of that month which is the second equinox, since the same
month is both the first and the seventh; that is to say, the first as respects nature and virtue, and
the seventh in point of time. Therefore in another place he says, \footnote{\#ex 12:2.} ``This month is unto
you the beginning of months, the first among the months of the year;'' calling that the first which is
so in respect of nature and virtue, and which as to number is in time the seventh month, since the
equinox has its appointed order in regular series, and in point of time is assigned the better season
of the year. But if you take that word first to have reference to the man, then it will be used with
more truth, and with strict propriety, for the just man was truly and properly the first, as in a vessel
the captain is the first man, and in a state the prince. But he is first not only in virtue, but also in
order, inasmuch as in the very circumstances of the regeneration of the second sowing of the
human race he was the beginning and the first. Moreover, it is very admirably considered with
reference to this passage, that the deluge took place during the life of the first man, and that again,
when it abated, things returned to their former steadiness, since after the deluge took place he had
to live by himself with his whole family, and after that evil was removed he alone was found upon
the earth during the latter period of his life until the regeneration of mankind began. But it is not to
no purpose that this testimony is given both of the preceding portion of his life, and also of the later
period, for he alone burnt with a desire for that genuine life which is in accordance with virtue, while
all the rest of the world were hastening on to death by reason of their fatal wickednesses.
Therefore of necessity the evil ceased on the six hundred and first year of his life, since in truth
the destruction came with reference to the sixth number, and safety was restored in unity since
unity is more a generativeness of the soul, and is the best for giving life, wherefore also a
deficiency of water in the sea takes place at the new moon, in order that the units may be
preferred in dignity both among months and years, when God saves those things which are upon
the earth; since the man who cultivates just habits is called by the Hebrews in their native language
Noah, but by the Greeks he is named Dikaios; however, he is not exempted from the laws
affecting the body. For although he is not subordinate to the power of others, but is a prince, yet
still, because he is nevertheless devoted to death, as he is dead, the principle of that number six is
connected with unity; since it was not in one year taken separately that the deluge ceased, but
together with the number six (as contained in the number six hundred), which is connected with it
according to corporeality and inequality; since the other being a long number is in the first place six
(that is to say, six hundred); on which account it is said, in the six hundred and first year. But the
just man is so in his generation, not in that which is general, nor again in that in which he is just by
comparison with the general corruption, but according to some especial generation; for his
generation bears with it a certain comparison. But that man also is deserving of praise whom God
selected beyond all other generations as being considered worthy of life, placing a limit to that life,
and to him as being about to be both the end and the beginning of each generation and of each
age; the end of that which is corruptible, the beginning of that which is to follow. And truly it is much
more proper to praise him who, bending upwards with his whole body, looked up by reason of his
friendship with God.

(46) What is the meaning of the expression, ``And Noah opened the roof of the ark?'' (Genesis
8:13). The text stands in need of no explanation. But with reference to its meaning, because the
ark is symbolically our body, we must consider that that is spoken of as the roof of our body, which
covers it and for a long time preserves its strength; such is concupiscence, by which the body is
preserved and made to last, in a moderate degree, that is, and in accordance with the law of
nature; as also it is dissolved by pain. When therefore the intellect is attracted by a desire for
heavenly things it wishes to spring upwards, and in that way it bursts asunder every appearance of
concupiscence; so that that thing being as it were removed which threw a veil of shade over it and
obscured it, it might be able to apply its senses to undisguised and incorporeal natures.

(47) Why is it that the earth was dried up in the seventh month, and on the twenty-seventh day?
(\#Ge 8:14). Do you not see that he here calls that month the seventh, which a little while before
he styles the first? for the seventh, as far as related to time, is the same, as I have said before, as
that which is the first in nature, being the beginning of the equinox. But it is with great propriety that
the beginning of the deluge is fixed to the seventh month, and the twenty-seventh day of the month;
and again, the end and cessation of the deluge is fixed to the same seventh month and the same
day; for, both the deluge and the removal of life took place at the equinox; the principle of which we
have indicated a little time ago; for the seventh month is found to be synonymous with months and
days of this time, and then again, the twentyseventh day occurs with the same meaning, when the
ark rested on the mountains. This is the month which by nature is the seventh, but in point of time
the first, which in fact is the month of the equinox. Therefore, at the equinoxes a power of selection
is given for seven months and twenty-seven days; for the deluge took place in the seventh month,
on which the vernal equinox takes place; so that it is in time the seventh, but in nature the first. And
the cessation of the deluge and the display of mercy belong to the same measure, when the ark
rested on the tops of the mountains; again in truth in the seventh month, but not the same month,
but in that in which the autumnal equinox occurs; that is to say, the seventh by nature, but the first
in point of time. But the most perfect cure, the fact of the evil being wholly dried up, is again fixed to
the seventh month and the twenty-seventh day of the vernal season; in order that both the
beginning and the end of the deluge might find its boundary at the same season; and that the
middle season when human life is repaired, is fixed to the intermediate season. In the meantime
that expression is more certainly to be observed, namely, that the whole year, by a strict
computation of days, made the deluge equal to the exact time of the remedy; for it began in the six
hundredth year of Noah's age, in the seventh month, and on the twenty-seventh day; so that the
whole space of the intermediate time completed a perfect year, the beginning being placed at the
vernal equinox, and the flood also ending equally at the same epoch of the vernal equinox. And in
this manner, after all things on earth, things full of fruit, had undergone destruction, as I have said
before, now that the persons who used the fruits were also destroyed, the earth being wholly
relieved of all evil was again found full of seeds and fruit-bearing trees, according to the production
of spring; for he thought it reasonable that, as the earth after it had suffered the deluge was in a
similar condition when dried again to that in which it was before, so it should now show itself, and
pay the debt which it owed to nature. Nor ought any one to wonder that in one day the earth when
left to itself produced every thing by divine virtue, both seeds and trees, all complete, entirely and
suddenly, with perfect and excellent herbs, and grain, and plants, and fruits; since in the creation of
the world on one day of the six he finished and brought to perfection the whole generation of plants.
But the present fruits were already perfect in themselves, and produced all kinds of fruits in a
manner suitable and corresponding to the season of spring; for all things are possible to God, who
scarcely requires time to effect any thing.

(48) Why was it that after the earth was dried, Noah did not depart out of the ark, before he had
received a fresh command from God, for God said to Noah: ``Go forth, thou, and thy wife, and thy
sons, and thy sons' wives, together with all the rest of the living creatures?'' (\#Ge 8:16).
Justice is commonly inspired with fear, as on the other hand injustice is rash and self-confident.
But the proof of a fear of God is the not giving up more to, or guiding one's self more by one's own
reason than by God. And above all other men it was natural for that man who had seen the whole
earth suddenly become an immense sea, to suspect that it might be possible that the same
misfortune would again return. Besides this, he also gave a thought to the corresponding
consequence, namely, that as he had entered into the ark at the command of God, so it was fitting
that he should also leave it at the command of the same being; for let no one believe that he can
ever do any thing perfectly unless God himself guides him by his preventing precepts.

(49) Why, when they entered into the ark was the order as follows: first himself and his sons, and
after them his wife and his sons' wives, but when they went forth the order was changed, for the
sacred historian says, ``Noah and his wife went forth, and after them his sons and his sons'
wives?'' (\#Ge 8:18). By the literal statement the sacred writer gives an obscure intimation, in
the order in which they entered, that the propagation of seed was taken away, but by the order of
their egress he implies the continuance of the process of generation; since, while they are
entering, the sons are mentioned together which their father, and the daughtersin-law with their
mother-in-law, but when they are going forth the wives are all mated again, the father being
accompanied by his wife, and each of his sons also by his wife, since he chose to show by fact
rather than by words everything which it is fitting for his friends to do. Moreover he had in express
words, and not by any vague intimations, commanded the men, as they were about to enter into the
ark, that they while there were to keep themselves from connection with women; but now that they
were about to depart from it, he plainly intimates to them that offspring is to be begotten in
accordance with nature, by the order in which he appoints their going forth; nor did he employ
words only, in order to make his proclamation about the state of the ark, saying, ``After a
destruction of all things on earth, of such a character and of such extent, do not indulge in
pleasures, for that is not decorous. It is sufficient, however, for you to have received your lives; but
while you are actually in the ark, to ascend up into the marriage bed with your wives would be a
proof of your being devoted to lasciviousness.'' And, indeed, it was natural for them, as being
relations to those who were being destroyed, to be moved with compassion for the perishing
human race, especially because they themselves also were still in doubt whether, from some
quarter or other, calamity might not fall also upon themselves; and besides these considerations it
was absurd, while those who were alive were perishing, for those in the ark to be contriving that
others who did not exist should be born, being warm at an unreasonable time, and burning with an
inopportune desire. But after the anger of God ceased, then he commanded those who had been
delivered from the calamity, when they had again gone forth out of the ark in order, to apply
themselves to the procreation of a succeeding generation, when he tells us, that the men did not go
forth with the men, nor the women with the women, but the wives with their husbands. But with
respect to the inner meaning of this fact, we must say this, that when the mind is about to wash off
and cleanse away its sins, then it is fit for male to live with male, that is to say, for the intellect, the
chief part of the man, to be as a father, united to each separate thought, as a father to his sons,
without any admixture of the female race, which is in accordance with the outward sense; since it is
a time of battle, in which it is necessary to keep the order of the cohort distinct, and to preserve it
strictly in order, that the soldiers may not be mingled in confusion, and so, instead of gaining a
victory over the enemy, be conquered themselves; but when the purification is completed, and
when the soul is dried up from all ignorance, and when a complete deliverance from everything
pernicious has taken place, then it becomes the man to collect his scattered forces together, not in
order that masculine counsels may be rendered effeminate by softness, but that the female race,
that is to say, the outward senses, may clothe themselves with the vigour of the male, attaining to
masculine counsels, and from their receiving seed for the production of a generation; so that, from
this time forth they may cherish, in all things, sentiments of wisdom, and honour, and justice, and
courage, and, in one word, of virtue. But, besides this, it will be reasonable also to take notice, that
when once a confusion, in the similitude of a deluge, has overwhelmed the intellect, and when the
different senses, being perplexed by the affairs of this world, like so many bulwarks erected
against them, begin to quarrel, it is utterly impossible that any one should be able, either to sow, or
to conceive, or to generate any good thing. But when all the hostile attacks of various agitations
and passions are checked, and when the ceaseless invasions of lawless counsels are repressed,
then the soul produces virtue and excellent works, as the most fertile portion of the earth, when
dried, produces fruits.

(50) Why did Noah build an altar without having been commanded to do so? (\#Ge 8:20). The
requital of gratitude which is due to God ought to be offered to him without command, and without
any delay or hesitation, showing the mind to be free from vices; for it becomes that man, who has
been endued with blessings by God, to offer him his thanks with a grateful and willing mind; but he
who delays to do so, waiting for an express command, is ungrateful, being as it were compelled by
necessity honour his benefactor.

(51) Why is he said to have built an altar to God, and not to the Lord? (\#Ge 8:20). In passages
of beneficence and regeneration, as at the creation of the world, the sacred writer only refers to
the beneficent virtue of the Creator, by which he makes everything in its integrity, and he implies
this by concealing the royal name of Lord, as one which bears with it supreme authority; therefore
now also, since what he is describing is the beginning of the renewed generation of mankind, he
borrows for his description the beneficent virtue, which bears the name of God; for he used the
kingly attribute, which declares his imperial power, by which he is called Lord, when he was
describing the punishment inflicted by the flood.

(52) What is the meaning of the statement, ``He took of the cattle and of the flying animals, and he
offered whole burnt offerings on the altar?'' (\#Ge 8:20). All this is said with reference to an
inward meaning, both because he received everything from God as a favour and gift; and also
because he took of the clean sorts of animals, and burnt those which were unpolluted and clean,
as entire and pure first fruits; for they are proper victims for good men to offer, and are themselves
entire, being full of integrity; and they may be classed as fruits, for fruit is the end, for the sake of
which the plant exists. This indeed is the literal statement; but with respect to the inner meaning, the
clean cattle and the clean birds are the outward senses and intellect of the wise man, with the
thoughts which are received in his mind; all which things it is reasonable to offer in their integrity as
entire and perfect fruit, in the way of a display of gratitude to the Father, and to offer them to him
as an unpolluted and clean oblation of a victim.

(53) Why does he offer his sacrifice to the beneficent virtue of God, but the acceptance of it takes
place by means of both the qualities of the Lord and God, for Moses says, ``And the Lord God
smelled a savour of sweetness?'' (\#Ge 8:20). He says this since, when unexpectedly, after all
hope is gone, we are preserved from dangers which are coming over us, we then, looking solely at
the beneficence of him who has preserved us, do, on account of our joy, display ingratitude, and
prefer the benefits which we have received rather to the beneficent power than to the Lord. But the
beneficent preserver himself, by means of both his attributes, looks down upon and honourably
accepts grateful minds, that he may not appear to halt in rewarding them; but he declares that such
a display of gratitude is pleasing to both attributes of the one God.\footnote{or, ``But the one God very much likes
to act by means of both his attributes.''—Note to the Latin version.}

(54) What is the meaning of the words, ``And the Lord God said, repenting him, I will not again
proceed to curse the earth for the works of man, for the thoughts of the mind of man are toward,
and are diligently and ceaselessly exercised in, wickedness from his youth up; therefore I will not
now proceed to smite all living flesh as I have done at other times?'' (\#Ge 8:21). The reasons
alleged appear to indicate a change of purpose, which is an affection not usual nor akin to the
divine virtue; for the dispositions of mankind are variable and inconstant, so that all affairs among
them are altogether uncertain; but with God nothing is uncertain, nothing incomprehensible, for he
is a being of mighty and consistent determination; how then, when reasons of the same kind are
present to him, because he was forsooth aware from the very beginning that the mind of man was
deliberately inclined to wallow in wickedness from his youth on, could he have originally intended to
destroy the human race by a flood; and yet afterwards say, that he did not intend to destroy it any
more, when the same evils still exist in the mind? But we must think that every kind of expression of
this sort is, by law, connected with learning and the utility of instruction rather than with the nature of
truth, since there are, as it were, two kinds which occur in the whole course of the law; in the first
place, as it is said, ``Not as a man;'' and in the second place, as it is said, ``As a man,'' the one God
is believed to instruct his son. That first expression relates to the actual truth; for, in real fact, God
is not as a man, nor again, as the sun, nor as the heaven, nor as the world, which is perceptible by
the outward senses, but as God, if it is justifiable to assert that also; since that most happy and
blessed being will not endure similitude, or comparison, or enigmatical description; nay, rather he
surpasses even blessedness and felicity itself, and whatever can be imagined as better than and
preferable to them. But the second expression relates to instruction and direction, I mean the
express words, ``As a man,'' in order that it may be observed, that he is willing to impress us beings,
born of the earth, lest perchance we should unceasingly incur his anger and his chastisement by
our implacable hostility to him, without any peace; for it is sufficient for him to be roused and
embittered against us once, and once to exact vengeance against sinners; but to inflict punishment
over and over again for the same thing is the conduct of a savage and ferocious disposition:
since, says he, ``when I shall inflict deserved retribution, as is possible, on every one, I will cause a
burning recollection of my design to be preserved.'' Therefore behold, the sacred historian has
excellently expressed himself, saying, ``That God observed in his mind,'' for his mind and
disposition rejoice in a superior degree of constancy; but our wills are found to be inconsistent and
vacillating, on which account we cannot be properly said to observe and think with our minds, since
it is by the thoughts that the passage of the mind is allowed to take place, \footnote{``if you connect the Armenian
words in a different manner, the sense will be 'meditation is the purification of the course of the mind,' and this is perhaps
better.''—Note by the Latin Translator.} but the human intellect is unable to be extended over everything,
since it is incapable of penetrating all things in a perfect and suitable manner. But that expression,
``I will not proceed any more to curse the earth,'' is used with great propriety, for it is not becoming
to add more curses to what has already been done, because the evils that have been inflicted are
already complete; because, although they are in some sense imperfect, inasmuch as the Father is
kind and merciful, and most humane, still he is rather inclined to alleviate the evil than to add to
men's misery. But that is as it were the same thing, according to a common proverb, to wash a
brick, or to draw water properly, and wholly to eradicate wickedness, with all its deeply imprinted
tokens from the mind of man; for if it is implanted in it at first, it does not exist accidentally, but is
engraven deeply on it and clings to it. But since the mind is a potential and principal part of the soul,
he introduces that word ``diligently;'' but that which has been weighed with diligence and care is
exquisite thought, examined more certainly than certainty itself. But this diligence does not tend to
any one evil, but as is plain, to mischief, and to all mischief; nor does it exist in a perfunctory
manner; but man is devoted to it from his youth, not only in a manner, but from his very cradle, as if
he were in some degree united to, and nourished, and bred up with sin. But yet God says, ``I will not
any more smite all flesh;'' giving notice that he will not, at any future time, destroy every portion of
mankind altogether, but only single individuals, in ever such great numbers, who perpetrate
unspeakable wickednesses; for he does not leave wickedness unpunished, nor does he grant it
liberty or impunity, but indulging his care for the human race on account of his original design, he of
necessity fixes destruction as a punishment for sinners.

(55) What is the meaning of the expression, ``Sowing-time and harvest, cold and heat, summer and
spring, shall not cease day nor night?'' (\#Ge 8:22). If taken literally this expression signifies the
continuation of the duration of the annual seasons, and that the earthly temperature adapted to
animals and plants is not again to be destroyed; since indeed, if the weather is corrupted it would
corrupt them likewise, and if it is preserved in its existing state it would preserve them also safe
and sound; for it is according to the weather and temperature that all animals and plants are
preserved safe and sound, without any infirmity, being accustomed, in some measure, to be
produced separately, in an admirable way, and to grow up together. But nature is like a harmony,
composed of opposite sounds, both flat and sharp; for thus, also, the world is compounded of
opposite qualities, for when, in the first place, the mortal commixtures of cold and heat, of moisture
and dryness, preserve their natural order, without any confusion, they are themselves a cause
which prevents destruction from overwhelming everything upon the earth. But if we regard the
inward sense of the passage, the seed time is the beginning and the harvest time is the end, and
both the beginning and the end are concurrent causes of safety, for either thing alone is by itself
imperfect, because the beginning requires an end, and the end has a natural inclination for the
beginning; but cold and heat bring round winter and autumn; for the autumn is fiery, but only in such
a degree as succeeding in its annual revolution to cool the fiery summer. And, symbolically, with
reference to the mind, cold indicates fear, since it causes terror and trepidation; but heat indicates
anger, because an angry disposition bears in itself a resemblance to flame and fire; for it is
necessary that those things should always exist and always remain among created and corruptible
beings; since summer and spring have been instituted for the production of fruits; spring for the
perfection of the seeds, and summer for the perfecting of fruits and the buds of trees. These
things indeed are discerned symbolically in addition to the inward sense of the words, producing a
double fruit; what is necessary being computed in the season of spring, and what is superfluous in
the summer. Therefore necessary food is for the most part for the body, being whatever is
produced freely from seeds; as virtues are necessary for the soul. But as many fruits as come by
way of excess from trees in summer, besides the advantage which they are to the body, do also
bring corporeal goods to the mind, as external advantages: for these external advantages are
subservient to the body, and the body is subservient to the mind, and the mind to God. But day and
night are the measures of times and numbers; and time and number exist without interruption. Day
indicates lucid wisdom, and night betokens obscure folly.

(56) Why was it that God, blessing Noah and his sons, said, ``Increase, and multiply, and replenish
the earth, and rule over it; and let your fear and the dread of you be upon all beasts, and upon flying
fowls, and upon reptiles, and upon the fishes which I have placed under your hand?'' (\#Ge 9:1).
This devotion of the inferior animals to man, God also at the beginning of the creation bestowed on
the sixth day upon man, after he had created him in his own image; for the scripture saith, ``And
God made man; in the image of God created he him; male and female created he them. And God
blessed them, and said, Be fruitful, and multiply, and replenish the earth; and be ye lords over it,
and be ye rulers of the fishes, and of the flying fowls, and of every creeping thing that creepeth
upon the earth.'' And did he not by these words evidently intimate that Noah, at the beginning of
what we may call the second creation of mankind, was found equal in honour to that creature who
in the first instance was made as to his form in the likeness of himself? Therefore he equally
assigned both to the one and to the other the principality and power over all the creatures that live
upon the earth. But do thou diligently take notice that he showed this man, who at the time of the
deluge was the only just man and the king of all the creatures which live upon the earth, to be equal
in honour, not to the identical man who was first created and formed out of the earth, but to that
one who was made according to the likeness and form of the true incorporeal entity, to whom also
he gives power, making him a king, not the very created man (or the man formed out of the earth),
but him who is according to his form and similitude, that is to say, incorporeal. Wherefore also the
creation of that man, who as to his form is incorporeal, was marked to have taken place on the
sixth day, in accordance with the perfect number six; but the creation of that man who was created
after the completion of the world and subsequent to the generation of all animals on the seventh
day, because it is after that that the manly figure was fashioned out of clay. Therefore after the
days of generation he says, ``on the seventh day of the world;'' for God had not yet rained upon the
earth, and no man did exist who could cultivate the earth. And then he proceeds to say, ``But God
formed a man out of the clay of the earth, and breathed into his face the breath of life, and man
became a living soul.'' Therefore how he can be made worthy of the same kingly power according
to the image of the man thus formed, he, I mean, who is the beginning of the second creation of
mankind, is indicated by the letter of the history that relates these events. But with reference to the
inward sense of the passage we must give an explanation in the following manner. God wills that
the souls of wise men should increase in the magnitude and multitude of the beauty of their virtues,
and should fill the mind as if it were the earth with those beauties, leaving no portion empty and void
so as to become occupied by folly. And he wills also that they should rule over, and strike terror
into, and inflict alarm upon all beasts; that is to say, he wills that all wickedness should be subdued
by their will, since wickedness is of an untamed and savage nature. Also he willed that they should
be lords over all flying fowls, which by reason of their lightness are raised on high, being armed with
courage and empty pride, and which thus cause the greatest mischief, being scarcely controlled at
all by fear. Moreover, he made them rulers over all creeping things, which are the symbols of
destructive vices, for they creep through the whole soul, namely, concupiscence, desire, sadness,
and cowardice, striking and goading; as also they are indicated by the fishes, which eagerly
cultivate a moist and delicate life, but one which is far from being sober, wise, or lasting.

(57) Why does God say, ``Every creeping thing which lives shall be to you for food?'' (\#Ge 9:3).
Creeping things are of a twofold nature; some being venomous, and others domestic. The
venomous ones are serpents, which, instead of feet, use their bellies and breasts, creeping upon
the earth; but the domestic ones are those which have legs above their feet. This is the literal
meaning of the statement. But if we look to the inward sense of it, then the creeping things
represent the foul vices, but the clean ones represent joy; for in connexion with the passion of
concupiscence there will exist joy and pleasure; and in connexion with desire there will be will and
counsel, and in connexion with sorrow goading and compunction, and in connexion with avidity
there will be fear. Therefore such disordered perturbations of the passions threaten souls with
death and destruction; but the joys do really live, as he himself has warned us in an allegory; and
they also give life to those who possess them.

(58) What is the meaning of the expression, ``As the green herb I have given you all things?''
(\#Ge 9:3). Some persons say that by this expression, ``As the green herb I have given you all
things,'' the eating of flesh was permitted. But I say that even though God had intended to give that
permission, still that before all things he must have intended to establish by law the necessary use
of herbs, that is to say of vegetables. And under the general name of herb he includes all the other
additional descriptions of green food, without mentioning them expressly in the law. But now the
power of this command is adapted not to one nation alone among all the select nations of the earth
which are desirous of wisdom, among which religious continence is honoured, but to all mankind,
who cannot possibly be universally prohibited from eating flesh. Nevertheless, perhaps the present
expression has no reference to eating food, but rather to the possession of the power to do so; for
in fact every herb is not necessarily good to eat, nor again is it the uniform and invariable food of
all uniform living animals; since God said that some herbs were poisonous and deadly, and yet they
are included in the number all. Perhaps therefore, I say, he means to express this, that all brute
beasts are subjected to the power of man, as we sow herbs and take care of them by the
cultivation of the land.

(59) What is the meaning of the expression, ``You shall not eat flesh in the blood of its life?''
(\#Ge 9:4). God appears by this command to indicate that the blood is the substance of the
soul; I mean of that soul which exists by the external senses and by vitality, not of that which is
spoken of with a certain especial pre-eminence, being the rational and intellectual soul; for there
are three parts of the human soul; one the nutritive part, another that which is connected with the
external senses, and the third that which exists in reason. Therefore the rational part is the
substance of the divine spirit according to the sacred writer Moses: for in his account of the
creation of the world, he says, ``God breathed into his face the breath of life,'' as being what was to
constitute his life. But of that part of the soul which is connected with the external senses and with
vitality, blood is the substance; for he says in another place, ``The blood exists in every breath of
flesh.'' It is with great propriety in fact that he has called the blood the breath of all flesh, because
there are in the flesh senses and passions, but not intellect nor thoughts. But again by the
expression ``the spirit of blood,'' he intimates that the spirit is one thing and the blood another; so
that the essence of the soul is truly and beyond all possible question spirit. But that spirit has a
place not by itself separately, apart from the blood in the body; but it is interwoven and mingled with
the blood. As also the veins which exhibit a pulse, as if they were vessels to convey breathing,
bear with them most unmixed and pure air, but blood likewise, though perhaps in a less degree; for
there are two vessels, the veins and the breathing channels; but the veins have more blood than
breath, and the breathing channels have more breath than blood. Therefore the proper admixture
in each vessel is distinct, as the greater and the lesser proportion. This is the meaning of these
words when taken literally; but if we look to their inner meaning, he calls the blood of the soul that
warm and fiery virtue belonging to it which we name courage. And he who is full of this wisdom
despises all food, and every pleasure of the belly, and of those parts which are below the belly. But
if any one adopts a profligate life, and becomes a wanderer like the wind, and gradually inactive
from laziness and a luxurious life, he in fact does nothing else but fall upon his belly, as a reptile
creeping upon the earth, and greedily licking up earthly things, closing his life without ever tasting of
that heavenly food which the souls which are desirous of wisdom receive.

(60) What is the meaning of the expression, ``The blood of your souls will I require from every
beast, and from the hand of man's brother will I require the life of man?'' (\#Ge 9:5). The
multitude of creatures which do injury is twofold; some being beasts, and others men. But beasts
are rather the least injurious of the two, because they have no actual familiarity with those whom
they wish to injure, principally because they do not fall under their power, but destroy those who
have properly power over them. But when he speaks of brothers, he means men who are
murderers, intimating these three things. First of all, that all we men are akin to one another, and
are brothers, being connected with one another according to the relation of the highest kind of
kindred; for we have received a lot, as being the children of one and the same mother, rational
nature. In the second place, he intimates that very commonly numerous and terrible quarrels arise,
and acts of treachery take place, between relations, and rather between brothers, on account of
the division of their inheritance, or on account of some superiority of dignity in the household; since
a quarrel between those of the same family is worse and altogether unseemly, because brothers
who are really so by the ties of nature meet in contest with a great knowledge of one another's
internal circumstances; being therefore well aware what kind of attack they must employ in their
present warfare. But, in the third place, as it appears to me, he employs the appellation of brothers
in order to warn men of the implacable and severe punishment which is reserved for murderers;
that they, without meeting any compassion, shall suffer what they have inflicted; for they have not
slain strangers, but their own brothers in blood. It is with exceeding great propriety that he calls
God the protector and overseer of those who are slain by man; for although men despise the
revenge, yet let them not behave negligently, but although impure men of savage disposition
escape for the moment from danger, still let them know that they are already caught and brought
before the greater tribunal of justice, namely, before the divine judgment-seat, which rises up to
inflict vengeance on the wicked for the defence of those who have received shameful and
unworthy treatment. This is the literal meaning of the words; but if we look to the inward sense of
them they have a regard to the merit of the purity of the soul, to which it is suitable to avoid
unceasing destruction brought in from outward parts; which merit, that propitious and beneficent
being, the most merciful and only Saviour, does not despise; but he expels and destroys all its
enemies who stand around it, calling them beasts, and men brothers; for beasts are a symbolical
expression for furious men threatening calamitous death; but men and brothers are both separate
individual thoughts, and words uttered by mouth and tongue, because they are akin to them, and, by
consequence, they bring on great and destructive evils, leaving no stone unturned, no work or word
omitted to do injury.

(61) What is the meaning of the expression, ``Whoso sheddeth man's blood by man shall his blood
be shed?'' (\#Ge 9:6). There is no excess in this declaration, but rather an indication of a still
more formidable denunciation, because he says, ``He himself shall be poured out like blood who
pours out blood.'' For that which is poured out flows forth and is lost, so that it has no longer any
power or substance. And by this he shadows forth the fact that the souls of those who perpetrate
unworthy actions imitate the mortal body in its corruption, as far as corruption is accustomed to
come upon individuals; for the body is then dissolved into those parts of which it was composed,
returning into its proper elements. But the miserable soul, labouring under distresses, is borne
hither and thither by the overflow of a lascivious life; and the very evils which have grown up along
with it are accustomed to suffer the same overflow, in the manner of the parts of the limbs.

(62) Why is it that he speaks as if of some other god, saying that he made man after the image of
God, and not that he made him after his own image? (\#Ge 9:6). Very appropriately and without
any falsehood was this oracular sentence uttered by God, for no mortal thing could have been
formed on the similitude of the supreme Father of the universe, but only after the pattern of the
second deity, who is the Word of the supreme Being; since it is fitting that the rational soul of man
should bear it the type of the divine Word; since in his first Word God is superior to the most
rational possible nature. But he who is superior to the Word holds his rank in a better and most
singular pre-eminence, and how could the creature possibly exhibit a likeness of him in himself?
Nevertheless he also wished to intimate this fact, that God does rightly and correctly require
vengeance, in order to the defence of virtuous and consistent men, because such bear in
themselves a familiar acquaintance with his Word, of which the human mind is the similitude and
form.

(63) What is the meaning of the words, ``There shall not again be a deluge to destroy all the earth?''
(\#Ge 9:11). By his last saying he declares sufficiently that there may be various inundations,
but that there shall never be one of such a character as to be able to change the whole earth into a
lake or sea. This is the literal meaning of this saying. But if we look to its inward sense, there a
divine kindness is intimated, according to which, although it is not every part of the soul which is
allowed to make proficiency in every virtue, still some are adorned in a considerable degree. So
that, supposing any one is not able to display excellence in his whole body, he still may labour with
all diligence to acquire all the means in his power to display excellence; and that exertion is within
his reach. And it does not follow that if any one is less highly endowed, or is unable to make every
portion of his life altogether perfect, that he is on that account to despair of those things which he
is able to do and to attain to. Since as there is power in every individual, he who does not exert
himself in accordance with it is both idle and ungrateful; idle because of his laziness, and ungrateful
because, though he has received most excellent means, he still sets himself in opposition to the
essential qualities of things.

(64) Why does God say that, as a sign that he will never again bring a deluge over the whole earth,
he will place his bow in the clouds? (\#Ge 9:13). Some persons imagine that by the bow he
means that thing which by some is called Jupiter's belt, from its figure, dwelling on its continual
similitude to the rainbow; but I do not perceive that that has been positively asserted. In the first
place, because the bow aforesaid ought to have a peculiar and essential nature of its own,
because it is called the bow of God; for he says, ``I will set my bow in the clouds.'' But that which
belongs to God and is said to have been set in any place as his, indicates plainly that it is not
devoid of essence or of substance. But the belt of Jupiter has not, properly speaking, any separate
nature of its own, but is merely an appearance of the solar rays on a wet cloud, all the
phaenomena of which are non-existent and incorporeal. And moreover, this is a further proof of
that, that it is never seen at night, though clouds exist by night as well as by day. In the second
place, we must also say that even in the day-time, when clouds obscure the whole face of heaven,
the belt of Jupiter is never at all seen in them. But what remains may also be affirmed without any
falsehood, when the Maker of the law says, ``I will set my bow in the clouds;'' for, behold, while
clouds are present there is no appearance of the belt of Jupiter visible. But he said, ``Where there
is a collection of clouds let there be a bow seen in the clouds.'' Still it often happens, when the
clouds are collected and when the air is obscured and thickened, that no appearance of a rainbow
is seen anywhere. We must consider, therefore, whether haply the sacred historian indicates
something else by this mention of the bow, namely, that in the very exercise of the mercy of God,
and also in the moment of his bitterness towards men on earth, there still shall not be any ultimate
destruction of them, in the fashion of a bow, which is too soft and unfit for such a purpose, nor shall
there be any violence added, so as to cause a rapid destruction, but there shall be a moderate
determination, each attribute being carefully measured; for the great deluge took place with a
breaking asunder and disruption of the clouds and of all things; as he himself asserts, when he
says, ``The fountains of the deep were broken up.'' And yet it was not an unmeasured vehemence.
Moreover, a bow is not itself a weapon, but only an instrument for the use of weapons, namely, for
the arrow which strikes; and the arrow being sent forth by means of the bow strikes a part which is
at a distance, while the parts which are nearest to it remain unhurt. And this is given as a proof that
the whole earth shall never for the future suffer any deluge, since no one arrow ever hits all places,
but only those which are at a distance. Therefore the divine virtue, being invisible, is symbolically
indicated by the bow in the cloud; being in truth dissolved according to the figure of tranquillity, and
condensed in accordance with a cloud; so that it does not permit all the clouds to be altogether
dissolved into water, so that the earth may not be made a lake by an inundation, which it carefully
forbids, and arranges the condensation of air, checking it as by a bridle, though it is at that time the
more accustomed to exhibit itself as rebellious by reason of its excessive fulness. For by reason
of the clouds it also shows itself to be replenished, dripping, and saturated.

(65) Why is it that after the sons of the just man have been named Shem, Ham, and Japhet, he
relates only the generations of the middle one, saying, ``And Ham was the Father of Canaan;'' and
afterwards he adds, ``These are the three sons of Noah?'' (\#Ge 9:18). Mentioning four men,
Noah and his sons, he says that these were obedient. Because the grandson Canaan was in his
habits like his Father who begat him, on that account, instead of mentioning only one, he includes
both in his enumeration, so that they are four in number, three in virtue. But in the meantime in the
scripture he mentions only the generations of the middle one, on account of the just man whom he
is going to speak of subsequently, because although he was his father, since Ham is the Father of
Canaan, still he does not mention the father with blame, but with respect to the man with whom he
thought it fair that the son should be a partaker, he yet did not give the father a participation with
him. In the second place, perhaps he thus gives a premonitory warning also to those persons who
by the acuteness of their mental vision can see a long way off what is at a distance, namely, that
he designs to take away the land of the Canaanites from them after the lapse of many ages, and to
give it to his chosen people who are thoroughly devoted to God. Therefore he chooses to
designate the chief inhabitant of that region, namely Canaan, and to show that he both practised
singular and peculiar wickedness of his own, and also all the wickedness of his father, so that in
every part he might be convicted of an ignoble slavery and submission. This is the literal meaning
of these words. But if we have a regard to the inward sense, he does not say that Ham had a son
named Canaan, but he predicates offspring of him alone, saying, ``Ham was the Father of
Canaan.'' Since such a disposition as that of Ham is always the Father of such designs as those
of Canaan, and that the very names themselves intimate this. For if we translate them into another
language, Ham means heat or hot; and Canaan means merchants, or buyers, or causes, or
recipients. Accordingly, he is not now speaking manifestly of generations, nor is he saying that one
man is the Father or the son of another man, but he is evidently demonstrating the connection
between one counsel and another, by reason of its alienation from all familiarity with virtue.

ABOUT THE CULTIVATION OF THE EARTH

(66) What is the meaning of the statement, ``Noah began to be a cultivator of the earth?'' (Genesis
9:20). He is here comparing Noah to the first created man who was formed out of the earth; for in
that manner also does he speak of him when he came forth out of the ark; since both then and now
there took place a first beginning of the cultivation of the land, each being after a deluge. For also,
at the time of the original creation of the world the earth was, as it were, a lake, being covered by
an inundation of water, for the sacred historian could not tell us that God said, ``Let the waters be
gathered together into one body, and let the dry land appear,'' unless it had previously been
inundated with waters which now returned into certain depths of the earth. Nor again is the
expression a purposeless one, ``He began to be a tiller of the earth,'' for in the second generation
he was himself the beginning of men, and also of seed, and of the cultivation of the land, and of the
life of all other things. This is the literal meaning of the words. But if we look to their inner sense, a
distinction is made between being a cultivator of the earth and a tiller of it; as the murderer of his
brother is represented as tilling the earth, but not as cultivating it. For by the earth our body is
symbolically represented, which is by its nature earthly, and which the unjust and wicked man tills
like a lazy hireling, but which the man endued with virtue cultivates like a skilful manager of plants
and an agriculturist of good works appointed to superintend it. Because the workman of the body,
the mind, as being carnal, procures carnal pleasures; but the cultivator of the earth is careful to
produce useful fruits, those, namely, which are to be obtained by the study of continence, and
modesty, and sound wisdom; and he prunes away all superfluous excesses and bad habits which
spring up around, like the thin and misplaced branches of trees.

(67) Why does the just man first plant a vineyard? (\#Ge 9:20). It was very natural for it to be a
subject of anxiety and doubt to him in what quarter he was to find any plants after the deluge, when
everything upon the earth was destroyed. Therefore it appeared natural, as was said a little while
ago, that the earth was made dry in the spring season; therefore when the spring produced the
buds of trees, the roots and stems of the vine could easily be found by the just man still alive, and
might thus be collected by him. But we have to consider why the first thing he did was to plant a
vineyard, and why he did not rather sow wheat and barley, since the latter are necessary
productions of the earth, without which life cannot be supported, but the former is only a material
for superfluous pleasure. The answer is that Noah, adopting a salutary design, consecrated and
offered up to God those things which are necessary to support life and which require no
co-operation for the production of the fruit; but the superfluous plants he devoted to men; for the
use of wine is superfluous and not necessary. As therefore God ordered fountains of water fit to
drink to burst up from the earth without the cooperation of man, so he also of his own accord
granted to man in a similar manner wheat and barley, in order that he himself might be the sole
giver of each kind of food which serves for necessary eating and drinking. But he did not take
away the power nor grudge them providing for themselves by their own industry those things which
contribute to pleasure.

(68) What is the meaning of the statement, ``He drank of the wine and was drunken?'' (Genesis
9:21). In the first place, the just man did not drink the wine, but a portion of the wine, not the whole
of it; in which case an incontinent and debauched man does not quit his means of debauchery, till
he has first swallowed all the wine that there is before him; but by the religious and sober man
everything necessary for food is used in a moderate degree. And the expression, ``he was
drunken,'' is here to be taken simply as equivalent to ``he used the wine.'' But there are two modes
of getting drunk, the one is that of an intemperate sottishness which misuses wine, and this
offence is peculiar to the depraved and wicked man; the other is the use of wine, and this belongs
to the wise. It is therefore in the second of these meanings that the consistent and wise Noah is
here called drunken, not as having misused but as having used wine.

(69) What is the meaning of the statement, ``He was naked in his house?'' (\#Ge 9:21). This is
a praise of the wise man both in the literal sense of the words, and also in their hidden meaning,
that his exhibition of nakedness took place not out of doors but in his house, being concealed by
the roof and walls of his house; for the nakedness of the body is concealed by a house which is
made of stones and beams of wood: but the covering and clothing of the soul is the discipline of
wisdom. Therefore there are two kinds of nakedness, one which takes place by accident, which is
the result of an involuntary offence, because the just man, using, if I may say so, his honesty as if it
were a garment with which he is clothed, stumbles out of his own accord like men who are
intoxicated, or who are afflicted with insanity; for in such men their offences are not deliberately
committed: but it is his task and pleasing duty to clothe himself, as with a garment, with the
discipline and study of honesty. There is also another kind of nakedness of the soul which is
caused by perfect virtue, which expels from itself the whole carnal weight of the body, as if it were
flying from a tomb, as indeed it has long been buried in it as in a tomb; as also it avoids pleasures,
and also a great number of miseries arising from the different passions and many anxieties arising
from misfortunes, and indeed all the evil effects of these different circumstances. He therefore,
who has been able with distinction to pass through such various and great dangers, and to escape
such injuries, and to emancipate himself from such evils, has attained to the destiny of happiness,
without any stain or disgrace; for I should pronounce this to be the ornament and badge of beauty
in those individuals who have been rendered worthy to pass their existence in an incorporeal
manner.

(70) Why is it that the sacred writer has not simply said, Ham saw his nakedness, but Ham the
father of Canaan saw the nakedness of his father? (\#Ge 9:22). By stating the fact thus, he
both blames the son in the father and the father in the son, as performing together in common the
deed of folly, and iniquity, and impiety, and every other kind of wickedness. This is the literal
meaning of the statement; and as to the inner sense, we must look at that in the same manner in
which we have hitherto treated these subjects.

(71) What is the meaning of the statement, ``He told it to his two brothers out of doors?'' (Genesis
9:22). The sacred historian is here adding to the gravity of the transaction. In the first place,
because he did not report the involuntary evil of his father to one brother only, but to both of them;
and no doubt if he had had any more he would have told it to them all, as he did in fact to every one
he could; and he did so with ridicule in his very words, making a jest of what ought not to have been
treated with laughter and derision, but rather with shame and fear mingled with reverence. In the
second place, when the historian says he told it them, not in the house but out of the house, he
evidently points out that he displayed his father when naked, not only to his brothers, but also to the
bystanders with whom they were, both men and women. This is the literal information conveyed by
the words. But if we look to their inward meaning, then we shall see that a depraved and malignant
habit of life is full of derision and contempt: and it is a bad thing to judge of the miseries of others
even by one's self like a chastising judge. But in this case what has happened is worse than this,
for any man with a joyful mind to ridicule the involuntary misfortune of a devoted disciple of wisdom,
and to make a song of and proclaim abroad his misery, is the part of a thoroughly hostile accuser,
who ought rather to have pardoned such an occurrence than to have added accusation or
vituperation to it. Moreover, because these three things are, as I have said before, as it were
brothers together; namely, good, bad, and indifferent, being all the offspring of one parent thought:
in accordance with each of these principles, they have been found to be overseers, some
celebrating virtues with praise, others upholding acts of malignity, and others supporting riches and
honours and other good things which, however, are not attached to and which are external to the
body. The overseers who emulate wickedness rejoice at the fall of the wise man, and ridicule and
disparage him, as if he had done no good by the part which he adopts and to which he applies
himself as better for the mind, or for his body, or for his external circumstances, to his internal
virtues or to any of the good things which are around and exterior to his body. Unless indeed that
man alone is eminently able to attain his object, who applies himself to iniquity, as that alone is
accustomed to confer advantages on human life. Pronouncing these and similar precepts, those
who are overseers of iniquity ridicule those who devote themselves to virtue, and to those things by
which virtue is produced and consolidated: as some look upon those things to be which are around
the body, and outside it, and which may be regarded in the light of instruments serving to that end.

(72) What is the meaning of the statement, ``Shem and Japhet, taking a garment, laid it upon both
their shoulders and went backwards, and covered the nakedness of their father, and they
themselves did not see it?'' (\#Ge 9:23). The literal meaning of the statement is evident; but with
respect to the inner sense contained in it, we must say that the light man who is in too great haste
only sees those things which are before his eyes and exposed to his sight: but that the evil man
also sees those things which are at his back, that is to say, the future. And since what is posterior
is postponed to what is anterior, so is what is future to what is present, the sight of which is
peculiar to the virtuous and wise man, who in truth is a second Lynceus, being according to the
fables gifted with eyes in every part. Therefore every wise man, who is not so much man as actual
intellect, walks backward, that is to say, he sees what is behind him or future, as if it were placed in
brilliant light; and seeing every thing on all sides of him with a perfect sight, and looking all around
him, he is found to be armed, and protected, and fortified, so that no part of his soul is ever found
naked or in an unseemly plight, on account of any accidents which occur unfortunately.

(73) What is the meaning of the statement, ``And Noah became sober after the wine?'' (Genesis
9:24). The literal meaning is too notorious. Therefore we need only here speak of what concerns
the inner sense of the words. When the intellect is strengthened, it is able by its soberness to
discern with a certain accuracy all things, both before and behind it, both present, I mean, and
future; but the man who can see neither what is present nor what is future with accuracy, is
afflicted by blindness; but he who sees the present, but who cannot also foresee the future, and is
not at all cautious, such a man is overcome by drunkenness and intoxication; and he, lastly, who is
found to be able to look all around him, and to see, and discern, and comprehend the different
natures of things, both present and future, the watchfulness of sobriety is in that man.

(74) Why is it that after the sacred historian has enumerated Ham in the middle of the offspring of
Noah, or has placed him in the middle between his brethren, he nevertheless points out that he was
the younger, saying, ``Noah saw what his younger son had done to him?'' (\#Ge 9:25). This is a
manifest allegory, because he here takes as the younger, not him who was so in age and in point
of time, but him who was younger in mind; since wickedness is unable to attain to a perception of
the learning which is proper to the elder; but the elder thoughts belong to a will which is truly growing
old, not indeed in body, but in mind.

(75) Why did Noah when praying for Shem speak thus: ``Blessed is the Lord God, the God of
Shem: and Canaan shall be his servant?'' (\#Ge 9:23). The names Lord and God are here
used together on account of his principal attributes, both of benevolence and of kingly power by
which the world was created; for as king he created the world according to his beneficence; but
after he had completed it then the world was arranged and set in order by his attribute of kingly
power. Therefore he at that time rendered the wise man worthy of a common honour, which the
whole world also received, all the parts of the world being formed in an admirable manner with the
attributes of the Lord and God, doing so by his especial prerogative, munificently pouring forth the
favour and liberality of his beneficent power. And it is on this account that the beneficent power of
God is mentioned twice. Once, as has been already stated, being placed in opposition to his kingly
power; and a second time without any such connexion, in order, forsooth, that the wise man having
been rendered worthy of his gifts, both such as are common to him with others and such as are
peculiar to himself, he might also be rendered acceptable both to the world and to God; to the world
on account of the excellence imparted to him in common with it, and to God for such as was
peculiar to himself.

(76) Why, when Noah prayed for Japhet, did he say, ``God shall enlarge Japhet, and bid him to
dwell in the house of Shem: and Canaan shall be their servant?'' (\#Ge 9:27). Without
examining the literal statement, for the meaning of that is plain, we had better approach the inner
sense contained in it, and examine that, in which the second and third blessings mentioned are
capable of an enlarged and ample extension. As, for instance, good health, and a vigorous state of
the outward senses, and beauty, and strength, and opulence, and nobleness of birth, and friends,
and the power of a prince, and numbers of other things. And on this account he said, ``God shall
enlarge,'' etc. Because taken separately, the abundant possession of such numerous and great
blessings has of itself been injurious to many persons who have scarcely dwelt with justice, or
wisdom, or any other virtues, the complete possession of which dispenses to man in an admirable
manner the advantages which are external to and which surround the body; but the deprivation or
absence of them leaves him without the enjoyment or use of them; and man, if deprived of all good
protectors, and of the use of these enjoyments, is exposed to as much suffering as he is capable
of. Therefore he prays on behalf of he man who has those things which are around and exterior to
the body, that he may dwell in the house of the wise man; so that attending to the rules of all good
men he may see and regulate his own course by their example.

(77) Why because Ham had sinned did God pronounce that his son Canaan should be the servant
of Ham and Japhet? (\#Ge 9:27). In the first place, God pronounced this sentence because
both father and son had displayed the same wickedness, being both united together and not
separated, and both indulging in the same disposition. But in the second place, he did so because
the father would be exceedingly afflicted at the curse thus laid upon the son, being sufficiently
conscious that he was punished not so much for his own sake as for that of his father. And so the
leader and master of the two suffered the punishment of his wicked counsels, and words, and
actions. This is the literal meaning of the statement. But if we look to its inward meaning, then in
reality they are no more two different men than two different dispositions. And this is made plain by
the names given to them, which manifestly denote the nature of the facts; for Ham being
interpreted means heat or hot; and Canaan means merchants of causes.

(78) Why was it that Noah lived after the deluge three hundred and fifty years? (\#Ge 9:28). It
is now declared that in two periods of seven years the form of the world was originally created and
now renewed under Noah. But the wise man lives for a period of fourteen quarters of a century;
and fourteen times twenty-five is equal to seven times fifty, or fifty times seven. And it is the
principle of the seventh year and also of the fiftieth, which has an especial order of its own
explained and ordained in Leviticus.

(79) Why among the three sons of Noah does Ham appear always to occupy the middle place, but
the two extremities are varied; for when their birth is mentioned, Shem is placed in the first rank, in
this manner, Shem, Ham, and Japhet; but when they are spoken of as fathers, then Japhet is
mentioned first, and the beginning of the enumeration of the nations is derived from Japhet
himself? (\#Ge 10:1). Those who inquire into the literal nature of the divine writings think thus of
the order in which these men are mentioned, looking upon him who is the first named, that is Shem,
as the younger; and upon him who is named the last, that is Japhet, as the elder. However they
may choose to think of this let them, being guided by the principle of mere opinion. But we who look
to the real meaning of these statements think that there is here a reference to the three things,
good, bad, and indifferent; which last are called secondary goods; and we must therefore think that
the sacred writer always puts the bad in the middle, so that being confined at either extremity it may
be subdued on one side by the one, and on the other side by the other; so that, being confined, it
may be kept in and subdued. But the good and the indifferent, or secondary good, change the
order with one another; for when there is such great evil present, and yet not wholly and altogether,
the good rejoices in the first place, having the position of the dispenser and chief of the whole. But
when it is placed in the position of the will in a state of conspiracy, and injustice remains not only in
the intellect but is also conducted to its end by unjust works, then that first good is changed from its
original order into another place, together with all the good habits which depend upon it, rejecting all
education and all arrangement, as being wholly unable to attain its proposed end, just as a
physician does when he sees an incurable disease. But the elder good manages that virtue which
is around the body and exterior to it; therefore, by observing the extremities with greater caution,
and closing in the beast within its toils, it is sufficiently demonstrated that it does not dare to bite or
injure any more. But while it feels that it has done no injury, it is transferred into a more secure and
more permanent position, and then, a higher and better fortified place being assigned to it, it easily
retains the lower position too as one easy to be preserved; for, in consequence of the superior
power of its guardian, it is always practicable to watch it closely, since nothing is more mighty than
virtue.

(80) Why do the people of Ceos, and of Rhodes, and the isles of the Gentiles, spring from Japhet?
(\#Ge 16:4—5). Since he has the name denoting breadth (namely Japhet), being expanded in
his growth and increase, that part of the things of the world which have been assigned by nature for
the use of mankind, that is to say, the earth, can no longer hold him, therefore he passes over into
the other part, that is to say, the sea and the islands belonging to it. This is the literal meaning of
the statement. But if we look to its inner sense, all the external blessings which are bestowed by
nature, such as riches, and honour, and principalities, are lavished and poured forth in every
direction on those men into whose hands they come, and are also extended widely to others who
are not so much within reach; so that in a greater, or at all events, in no less a degree do they
surround and hem the man in, in accordance with the greediness of the lovers of riches and glory,
since they are eager for principalities, and are never satisfied because of their insatiable desires.

(81) Why the eldest son of Ham is Chus. (\#Ge 10:6). The sacred historian has here produced
a word most completely in accordance with nature, saying that Chus was the elder son of evil,
Chus being the dissolved and loose nature of the earth, for the earth, when dense and fertile, and
moist, is full of herbs, and hills, and trees, and is well arranged for the production of different fruits;
but when dissolved and reduced to dust and dry, it is unfruitful and barren; and besides it is tossed
about in the air, when it is raised from the ground by the wind, by its dust making the air all alive.
Such as this is the first origin and the first shoots of evil being destitute of the generation of good
pursuits, and the cause of barrenness to the soul and to all its parts.

(82) Why was Chus the father of Nimrod, who began to be a giant and a hunter before the Lord: on
which account they said, ``Like Nimrod the mighty hunter before the Lord?'' (\#Ge 10:8). The
father in this case, having a nature truly dissolute, does not at all keep fast the spiritual bond of the
soul, nor of nature, nor of consistency of manners, but rather like a giant born of the earth, prefers
earthly to heavenly things, and thus appears to verify the ancient fable of the giants and Titans; for
in truth he who is an emulator of earthly and corruptible things is always engaged in a conflict with
heavenly and admirable natures, raising up earth as a bulwark against heaven; and those things
which are below are adverse to those which are above. On which account there is much propriety
in the expression, he was a giant against God, which thus declares the opposition of such beings
to the deity; for a wicked man is nothing else than an enemy, contending against God: on which
account it has become a proverb that every one who sins greatly ought to be referred to him as the
original and chief of sinners, being spoken of ``as a second Nimrod.'' Therefore his very name is
an indication of his character, for it is interpreted Aethiopian, and his art is that of hunting, both of
which things are detestable: an Aethiopian because unmitigated wickedness has no participation in
light, but imitates night and darkness: and the practice of the huntsman is as much as possible at
variance with rational nature, for he who lives among wild beasts wishes to live the life of a beast,
and to be equal to the brutes in the vices of wickedness.

\xchapter{Questions and Answers on Genesis, III}

(1) What is the meaning of the expression, ``I am the Lord thy God who brought thee out of the land
of the Chaldaeans to give thee this land for an inheritance?'' (\#Ge 15:7). As the literal
statement is plain enough, we need only consider the inner meaning, which was meant to be
interpreted in this manner. The law of the Chaldaeans taken symbolically is mathematical
speculation, one part of which is recognised to be astronomy, which the Chaldaeans study with
great industry and with great success. Therefore God is here honouring the wise man with a gift; in
the first place, by taking men out of the sect of the astrologers, that is to say, away from the
hallucinations of the Chaldaeans, which, as they are difficult to detect and refute, are found to be
the cause of great evils and wickedness, since they ascribe the attributes of the Creator to
created things, and persuade men to worship and to venerate the works of the world as God. In the
second place, God honours him by granting to him the wisdom which bears fruit, which he has here
symbolically called the earth; but the Father of the universe shows that wisdom and virtue are
invariable and immutable, since it is not consistent with his character that God should show to any
one that which can undergo any variation or change, for that which is shown by the being who is
immutable and consistent must be so too; but that which is liable to change, as being incessantly in
the habit of suffering variation, admits of no proper or divine demonstration.

(2) Why does he say, ``Lord, by what shall I know that I shall inherit it?'' (\#Ge 15:8). He here is
seeking a sign for a ratification of the promise; but two things only are described deserving of
study; one that which is an affection of the mind, namely, the belief in God according to his literal
word; the other a being borne on with the most exceeding desire not to be left in want of some
signs, by which the hearer may feel, to the conviction of his outer senses, a confirmation of the
promise: and to him who has given the promise he offers worthy veneration by the appellation,
``Lord.'' For by this title he says, I know thee to be the Lord and prince of all things, who art also able
to do all things, and there is no disability with thee. But in truth, if I have already given credence to
thy promise, still I nevertheless wish to obtain speedily if not a completion of it, yet at all events
some evident signs by which its consummation may be indicated; in truth I am thy creature, and
even if I were to arrive at the highest degree of excellence, I am not always able to restrain the
violence of my desire, so as not, when I have seen or heard anything good, to be contented with
obtaining it slowly and not immediately; therefore I entreat that thou wilt give me some means of
knowledge, by which I may comprehend those future events.

(3) Why is it that he says, ``Take for me a heifer of three years old, and a goat of three years old,
and a raven of three years old, and a turtle dove and a pigeon?'' (\#Ge 15:9). He here mentions
five animals, which are offered on the sacred altar; for these are divided into classes of victims,
three kinds of terrestrial animals, the ox, the goat, and the sheep; and two kinds of birds, the turtle
dove and the pigeon; for the sacred writer constantly tells us that the everlasting reverence of
victims derived its origin from the patriarch, who was also the origin of the race: but instead of the
expression, ``Bring to me,'' he has very admirably used the words, ``Take for me;'' since there is
nothing especially and peculiarly belonging to the creature, but everything is the gift of and blessing
bestowed by God, who is altogether willing that when any one has received anything he should
offer thanks for it with all his heart. But he orders him to take every animal at the age of three
years; since three is a full and perfect number, consisting of a beginning, a middle, and an end; but
still we may raise the question, why of these three animals, he takes two females, the heifer, and
the she-goat, and one male, the ram; may it not be perhaps because the heifer and the she-goat
are offered as an atonement for sin; but the sheep is not, as sin arises from frailty, and the female
is frail? This much I have thought fit to say with especial appositeness to this question; but I am not
however ignorant that all things of this kind offer a handle to those who wish to cavil, to disparage
the sacred scriptures; therefore in this instance they say that there is nothing here described and
indicated but a command to sacrifice, by the division of the animals and an examination of their
entrails; and what is visible in them they affirm to be an indication of what is convenient, and of the
similitude which arises from things visible. But those men, as it appears to me, are of that class
which forms a part alone from a judgment of the whole, but which on the contrary does not from a
judgment of a part from the whole, which last is the better way of coming to an opinion, as being
that by which both the name and the fact are altogether established. Therefore the giving of the
law, that is to say the sacred scriptures, that I may so express myself, is a sort of living unity, the
whole of which one ought to examine carefully with all one's eyes, and so discern with truth, and
certainty, and clearness, the universal intention of the whole of the scripture without dissecting or
lacerating its harmony, or disuniting its unity; by any other mode everything would appear utterly
inconsistent and absurd, being dissociated from all community or equity. What then is the intention
of the delivery of the law as exhibited to us? It is scientific, and so is everything which describes
scientific species; since the offering of sacrifice and all science admits of a consistent usage, and
of expression well adapted to them, and of various opinions, by which not only the footsteps of truth
are occupied, but sometimes are even darkened, as affection is by flattery; but in such way that
the very things which are genuine and established by experiment are perverted by things which are
both inconsistent and unproved. And the natures of the animals above mentioned have an intimate
connection with the parts of the universe; the ox is connected with the earth, as being an animal
employed in drawing the plough and in tilling the earth; the goat again is connected with the water (it
is called in Greek and Armenian aix, or ajx), being an animal deriving its name from driving and
rushing on (from ago— or aisso—); since water is an impetuous thing, and the course of rivers, and
the extent of the breadth of the sea, and the sea itself agitated as it is by its ebb and flow, are
witnesses of the propriety of the name and of the closeness of the connection. And the ram (aries)
is connected with the air, as being a very violent and vivacious animal, on which account too the
ram is more useful to mankind than any other animal as affording them raiment. Therefore, on
account of these reasons, as I think, God orders him first to take these two female animals, the
cow and the she-goat; since both these elements, earth and water, are material, and for the most
part feminine. But the third he will have a male, namely the ram; because the air or wind has been
explained as masculine; since the natures of all things are divided into bodies or into earth and
water, and female animals exist by nature. But that which exhibits a similitude to the soul is
arranged under the head of air and the breath of life. And this, as I have said, is masculine. If
therefore we are to call that masculine which is the moving and active cause we must call that
feminine which is moved and passive. But the whole heaven is found to be familiarly connected
with flying birds such as the pigeon and turtle dove, being distributed as it is into the rotatory path of
the planets and fixed stars. Therefore he dedicates the pigeon to the planets, for that is a tame
and domestic animal, as also the planets are more familiarly connected with us as being nearer to
the earth, and as having sympathies with us; but he consecrates the turtle dove to the fixed stars,
for that animal is a lover of solitude, and flees from the conversation of the multitude, and from all
connection of every kind. And so also the globe itself is remote, and a thing which wanders into the
furthest extremities of the world. Therefore both the species of these two birds are assimilated to
the divine attributes, since as Plato, the disciple of Socrates, says it is fitting that the heaven
should have a swift chariot by reason of its very swift rotatory motion, which in fact surpasses
even the birds themselves in the velocity of their course. But the birds above mentioned are
singers; the prophet indicating by an enigmatical expression that perfect music which exists in
heaven harmoniously adapted from the motion of the stars, since it is a proof of human art when
the corresponding music of the voices of animals and of living instruments is adapted together by
the industry of genius. But this heavenly music has been abundantly extended over the earth by the
Creator, as he has also extended the rays of the sun, being always prompt to exercise his
beneficent care for the human race. For such music excites frenzy in the ears, and brings
unrestrained pleasure to the mind; and so causes men to forget even their meat and drink, and
even when hunger brings death to the door to be willing even to die out of a desire to hear music.
And if the song of the Sirens, \footnote{he alludes here to the description in Homer, Od. 12.39—47 (as translated by
Pope)—``next, where the Sirens dwell, you plough the seas; / Their song is death, and makes destruction please. / Unblest
the man, whom music wins to stay / Nigh the curst shore, and listen to the lay; / No more that wretch shall view the joys of
life, / His blooming offspring, or his beauteous wife! / In verdant meads they sport, and wide around / Lie human bones
that whiten all the ground; / The ground polluted floats with human gore, / And human carnage taints the dreadful shore.''
And further on in the same book, the poet describes the effect of these songs upon Ulysses, Od. 12.183—194 (as translated
by Pope)—``o stay, O pride of Greece! Ulysses, stay! / O cease thy course and listen our to lay! / Blest is the man ordained
our voice to hear, / The song instructs the soul, and charms the ear. / Approach! thy soul shall into raptures rise! /
Approach! and learn new wisdom from the wise! / We know whate'er the kings of mighty name / Achieved at Ilion in the
field of Fame; / Whate'er beneath the sun's bright journey lies. / O stay and learn new wisdom from the wise!''} as
Homer tells us, invites the heathen so forcibly, that they forget while listening to it, their country,
their houses, their friends, and necessary food; how much more must that most perfect and
consummate music, so truly heavenly and endowed with the highest degree of harmony, when it
touches the organs of the ear, compel men to go mad and to yield to rapture. But the reason on
account of which every one of the animals to be offered is to be three years of age has already
been explained; and we must now discuss it under another form of mystery, since it has been seen
that every one of those things which were called into existence and subsequently to the moon,
such as the earth, water, and air, rejoice in an order connected with the number three. In the
divisions of earth there is a vast quantity of dry continent, islands and peninsulas. Water is divided
into sea, rivers, and lakes; and the air into the two equinoxes, the vernal and the autumnal; and
they may be taken as one, for they have an equal proportion of day and night, and accordingly the
equinoxes are neither hot nor cold. Add to these the changes of summer and winter, for the sun is
borne through those three circles into the seasons of summer, winter, and the equinoxes.
Therefore, in the first place, the natural arrangement will be of this kind; and the moral arrangement
is properly thus. In every one of us there are three things: flesh, the outward sense, and reason;
therefore the calf exhibits a familiarity with the corporeal substance, since our flesh is subdued by,
and kept in subservience to, and in connection with the ministrations of life; also their nature is
female according to matter, being calculated rather to be passive and to be subject rather than the
be active. But the similitude of the she-goat is connected with the communion of the outward
senses, either because all the objects of those outward senses are each borne towards their
appropriate sensation, or because each impulse and motion of the soul takes place in
consequence of an imagination formed of the objects received through the medium of the external
senses. And this is followed, in the first place, by a certain inflexion or alienation, which by some is
called an occasion, that is to say, an impulse affecting each kind of sense. But since the female is
the outward sense, as being passive on consequence of what is subjected to the outward senses,
therefore God has adapted to it a female animal, the she-goat. But the ram is akin to the word, or
to reason. In the first place, because it is a male animal; secondly, because it is a working animal;
and thirdly, because it is the cause of the world, and of the firmament; that is to say, the ram is so
by means of the clothing which it supplies; and reason, or the word, is so in the arrangement of life;
for whatever is not irregular and absurd immediately exhibits reason. And there are two species of
reason; the one derived from that nature by which the affairs of the world subjected to the outward
senses are finished; the other from that of those things which are called incorporeal species, by
which the affairs of that world which is the object of the intellect are brought to their
accomplishment. Therefore the pigeon and the turtle dove are found to resemble these. The
pigeon, forsooth, resembles speculation in natural philosophy; for it is a more familiar bird, as the
objects of the outward sense are exceedingly familiar to the sight: and the soul of the inquirer into
natural science flies upward as if it were furnished with wings; and being borne aloft is carried
round the heaven, discerning every part of every thing, and the principles of every separate thing;
for the turtle dove imitates that species which is the subject of intellect and incorporeal; for as that
animal is fond of solitude, so it is superior to the violent species which come under the outward
sense, associating itself as it does with the invisible species by its essence.

(4) Why does he say, ``And he took unto him all these things?'' (\#Ge 15:10). He has added
also that expression, ``And he took unto him,'' with especial propriety; for it is the sign of a soul
thoroughly imbued with the love of God to ascribe whatever good and noble theories and feelings it
receives, not unto itself, but wholly to God who is the giver of all benefits.

(5) What is the meaning of, ``He divided them in the middle and laid the pieces opposite to one
another?'' (\#Ge 15:10). Also the whole structure of the body, as of flesh, is to be looked at in
such a light as this according to its whole creation; for the parts are brothers; not as they are
divided and placed opposite to one another; but, being naturally inclined to one another, and having
a mutual regard to one another, on account of their natural co-operation; the original Creator who
gave them life making this division for the sake of usefulness, so that one part should be opposed
to the other part, and again that both should reciprocally seek one another in all necessary
ministrations. In this way he has directly separated the sense of sight, distributing it equally to two
eyes by placing the nose between them and thus turning each eye to the other; for the pupils, if I
may so say, lean both in one direction so as mutually to behold the same thing, scarcely ever
straying beyond the position in which they are placed, but only looking towards one another,
especially when anything comes across their sight. And in similar manner the faculty of hearing is
distributed between the two ears, which are both reciprocally turned to one another, both tending to
one and the same operation. And the sense of smell is divided between the two nostrils, being
turned towards the two tubes of the nostrils, which are not revolving around or inclined towards the
cheeks, so as being drawn in two different directions to look the one towards the right and the
other towards the left, but being both collected together and turned inwards they await all smells
with a common action. So also the hands are not made of an appearance contrary to that of one
another, but being like brothers and like divisible parts, looking to one another mutually, and being
prepared by nature for an operation and employment suitable to them, they thus act in the
operations of receiving, giving, and working. And the feet are not constituted differently from the
hands; as each of them behaves in such a manner that they both yield the one to the other, and
progress is effected by the motion of both together, so that nothing can be accomplished by one
alone. Nor is it only the feet and shins, but also the legs and knee-pans, and hips, and the breasts,
and in fact every part on the right or left of the body, being divided in a similar manner, indicate one
general harmony and correspondence and union as it were of connatural parts; that is to say, of all
of those different members enumerated according to their separate species. And generally,
whoever considers together and in an equal manner all the above mentioned parts thus subdivided,
in reference to their joint operation, will find one nature combined of the two parts. As the hands,
united and connected together with the fingers, are seen when in union with them to exhibit a
harmony; and the feet, when re-united in operation, are seen to tend to union; and the ears, when
similarly combined in the figure of an amphitheatre, are seen to unite themselves, in effect
extending across the space which separates them. Therefore our nature, continually making in this
manner a division of those parts which exist in us according to each separate species, has first of
all separated and arranged the different sections, placing them as it were opposite to one another
in the same way in which it has arranged the world; and it has also arranged them with reference to
the easy discharge of their several duties. And again it has combined each of these members
according to each species into one action, and into the same operation, collecting together all of
them when considered generally. Nor is it only the parts of the body which any one may see thus
united and in pairs, separated in their union, and again united in their division, but the parts of the
soul are so too. But since the two superior sections of this are so many separate classes, namely
the rational and the irrational, so also the separate parts of each section have their own
appropriate division; as for instance, the rational part is divided into the intention and into the
uttered word; and that part which exists in accordance with the outward senses is divided into the
four senses; for the fifth sense, touch, is common to the other four, two of which, those with which
we see and hear, are philosophical senses, so that it is by means of them that the power of living
well is acquired for us; the others are nonphilosophical, namely smell and taste, but are servile,
being created only for living; for the sense of smell, by means of its exercise, contains many things
which awaken it, and receives a continual breathing which is as it were the continual food of living
creatures; therefore smell and taste support this mortal body, but sight and hearing afford service
to the immortal soul. Therefore these divisions of our members, according to our body and soul,
were made and separated by the Creator; however, we must know that the parts of the world also
are arranged in two divisions and are placed opposite to one another; the earth being divided into
mountainous and champaign districts; the water into sweet and salt, sweet being that which is
supplied by springs and rivers, and salt being that which comes from the sea; as also the
atmosphere is divided into summer and winter, and also into spring and autumn. And it is on this
account that Heraclitus wrote his books about nature, having borrowed his theory of contraries
from our sacred historian, with the addition of an infinite number of laborious arguments.

(6) Why is it said, ``But he did not divide the birds?'' (\#Ge 15:10). He is shadowing forth a fifth
and periodical nature, from which the ancients say that the heaven was made; for the four
elements are mixtures rather than elements: by which he subdivides those things which are already
divided into those materials of which they were originally composed, as the earth includes within
itself a portion of the elements of water, and also of air, and also of fire, which however obtains the
appellation not so much in accordance with our apprehension of it, as with our sight; and again the
water is not so clear or pure, as not to have some participation in wind and earth; and so also in
each of the other elements there is a certain tempering and combination; but the fifth substance is
the only one which has been made unmixed and pure, on which account it was not accustomed to
be mentioned at all. Therefore it is well said, he did not divide the birds; since the heavenly nature,
both of the planets and also of the fixed stars, is raised on high like that of birds, in the similitude of
both kinds, that is to say, of clean birds, the turtle dove and the pigeon, which scarcely admit of
being divided or cut up; for the indivisible nature is of a fifth essence, more unmixed and pure than
the others, and therefore it more closely resembles unity.

(7) What is the meaning of, ``And the birds descended on the bodies which were divided?''
(\#Ge 15:11). Since the three animals, the heifer, and the shegoat, and the ram, were divided in
a symbolical manner, they are signs, as we have already said, of the earth, and water, and air; still
it is necessary to give now a reason for this, examining the truth carefully under the mystery of a
similitude. Perhaps therefore he designs and intimates by the descent of the birds on the cut
pieces an invasion of enemies; for all the nature of the world beneath the moon is full of battles and
ill will, both domestic and external; and the birds in truth appear to fly down on the divided bodies for
the sake of meat and drink; naturally indeed it is the stronger which descend upon the weaker
animals, as upon dead bodies, attacking them in general unexpectedly, but they do not fly down on
the turtle dove and pigeon, since the heavenly bodies are free from desires and unconscious of
suffering wrong.

(8) Why is it that he says, ``Abraham passed over and sat upon them?'' (\#Ge 15:11). Those
who think that sacrifice is indicated by the matters about which we are at present speaking will say
that the virtuous man, sitting as it were in a synagogue, has examined into the entrails of the
divided animals, as if that were looked upon as an unerring symbol for the declaration of the truth;
but we, who adhere to Moses and who are thoroughly acquainted with the views of that teacher,
one who, turning away his face from every sophistical appearance and prognostic, trusted in God
alone, will rather say, that he has here introduced the just man who is endued with virtue with the
birds themselves, who were congregated together and flying about over him, intending to denote
nothing else by this parabolical presentation, but that he is desirous of hindering injustice and
covetousness, and is most hostile to quarrels and wars, and a lover of consistency and peace; for
he himself is truly a guardian of peace. Since no one state has ever rested in tranquillity owing to
the conduct of the wicked, but kingdoms have become fixed steadily when one or two men endued
with virtue have arisen, whose virtue has put an end to civil disturbances, God granting to those
who are earnest in the pursuit of virtue good habits calculated to procure them honour; and not to
them only, but to those also who approach near to the production of general advantage.

(9) What is the meaning of the words, ``About the time of the setting of the sun a trance fell upon
Abraham; and lo, a great horror of darkness came over him?'' (\#Ge 15:12). A certain divine
excess was suddenly rendered calm to the man endued with virtue; for the trance, or ecstacy as
the word itself evidently points out, is nothing else than a departure of the mind wandering beyond
itself.\footnote{ekstasis, derived from existamai, in 2nd aor. act. exeste—n, ``I was beside myself.''} But the class of
prophets loves to be subject to such influences; for when it is divining, and when the intellect is
inspired with divine things, it no longer exists in itself, since it receives the divine spirit within and
permits it to dwell with itself; or rather, as he himself has expressed it, as spirit falls upon him; since
it does not come slowly over him, but rushes down upon him suddenly. Moreover, that which he has
added afterwards applies admirably, that a great horror of darkness fell upon him. For all these
things are ecstacies of the mind; for he also who is in a state of alarm is not in himself; but
darkness is a hindrance to his sight; and in proportion as the horror is greater, so also do his
powers of seeing and understanding become more obscured. And this is not said without reason:
but as an indication of the evident knowledge of prophecy by which oracles and laws are given
from God.

(10) Why was it said to him, ``Thou shalt know to a certainty that thy seed shall be a stranger in a
land that is not theirs, and shall be reduced to slavery, and shall be grievously afflicted for four
hundred years?'' (\#Ge 15:13). That expression is admirably used, ``It was said to him,'' since a
prophet is supposed to utter something, but yet he is not pronouncing any command of his own, but
is only the interpreter of another who sends something into his mind; and moreover whatever he
does utter and deliver in words is all true and divine. And in the first place, he declares that a family
of the human race is to dwell in a land belonging to another; for all things which are beneath the
heaven are the possession of God, and those living creatures which exist on the earth may more
properly and truly be said to be sojourners in a foreign land than to be dwelling in a country of their
own which by nature they have not got. In the second place, he thus declares to us that every
mortal is a slave after his kind. But no man is found to be free, but every one has many masters
who vex and afflict him both within and without; for instance, without there are the winter which
affects him with the cold, the summer which scorches him with heat, and hunger, and thirst, and
many other calamities; and within there are pleasures and concupiscences, and sorrows, and
fears. But his servitude is limited to a period of four hundred years, during which the aforesaid
pleasures shall rise up against him. On which account it has been said above, that Abraham
passed over and sat upon them, hindering and repelling them; as far as the literal words go,
repelling those carnivorous birds which were hovering over the divided animals, but in fact repelling
the afflictions which come upon men. Since a man who is in his own proper nature a lover of, and
also by diligent practice a studier of virtue, is a most humane physician of our race, and a true
protector of it, and guardian of it from evil. For all these things have an allegorical reference to the
soul. For while the soul of the wise man, descending from above from the sky, comes down upon
and enters a mortal and is sown in the field of the body, it is truly sojourning in a land which is not
his own. Since the earthly nature of the body is wholly alien from pure intellect, and tends to subdue
it and to drag it downwards into slavery, bringing every kind of affliction upon it, until the sorrow,
bringing the attractive multitude of vices to judgment, condemns them; and thus at last the soul is
restored to freedom. And it is on this account that he subsequently adds the sentence,
``Nevertheless the nation which they shall send I will judge: and afterwards they shall go forth with
great substance;'' namely, with the same measure, and still better. Because then the mind is
released from its mischievous colleague, departing out of the body and being transferred not only
with freedom but also with much substance; so as to leave nothing good or useful behind to its
enemies. Since every rational soul is productive, but he who thinks himself loaded and endued with
virtue in his own counsel, is unable to preserve his fruit unto the end. For it becomes a virtuous
man to attain to the objects which he has intended of his own accord, as also the counsels of
wisdom correspond to those objects. Since, as some trees, although they appear productive at the
first season of the budding of their fruits, are yet unable to bring them to maturity, so that the whole
fruit before it becomes ripe is shaken off by every trifling cause; in the same manner the souls of
inconstant men feel many influences which contribute to their productiveness, but nevertheless are
unable to keep them sound till they arrive at perfection, as a man studious of virtue ought to do in
order eventually to gather them as his own possessions.

(11) What is the meaning of, ``But thou shalt go to thy fathers in peace, being nourished in a fair old
age?'' (\#Ge 15:15). He here clearly indicates the incorruptibility of the soul: when it transfers
itself out of the abode of the mortal body and returns as it were to the metropolis of its native
country, from which it originally emigrated into the body. Since to say to a dead man, ``Thou shalt go
to thy fathers,'' what else is this but to propose to him and set before him a second existence apart
from the body as far as it is proper for the soul of the wise man to dwell by itself? But when he says
this he does not mean by the fathers of Abraham his father, and his grandfather, and his
great-grandfathers after the flesh, for they were not all deserving of praise so as to be by any
possibility any honour to him who arrived at the succession of the same order, but he appears by
this expression to be assigning to him for his fathers, according to the opinion of many
commentators, all the elements into which the mortal man when deceased is resolved. But to me
he appears to intend to indicate the incorporeal substances and inhabiters of the divine world,
whom in other passages he is accustomed to call angels. Moreover the words which follow are not
by any means without an object, that he is nourished in peace and in a fair old age. For the wicked
and depraved man is nourished in battle, and lives and departs in a very bad old age. But the good
man, in both phases of existence, both in that which is in connection with the body and in that which
is apart from the body, cultivates peace, and is alone completely virtuous, such as no foolish
person is found to be, even though he should live longer than an elephant; on which account he
here carefully said, ``Thou shalt go to thy fathers, being nourished-not in an advanced old age,
but-in a fair old age.'' For many foolish persons also have their lives extended to a greatly
lengthened period, but it is only the man who is desirous of virtue who enjoys a good old age and
one endued with virtue.

(12) Why is it that he says, ``And in the fourth generation they shall return again hither?'' (Genesis
15:16). The number four is more fit than any other number, for this reason, that as it is more
perfect, and is the root and foundation of the perfect number ten; and it is according to the principle
of the number four that all collected are to return hither, as he himself has said. But as he by
himself is perfect, so also those of whom he is the father are evidently perfect. But what is it that I
am saying? In the generation of animals the sowing of the seed has the first place; in the second
place, comes the fact of each instrument being, in some manner, impressed by something akin to
nature; thirdly, there is the growth after the first formation of the creature; fourthly, after everything
else comes the perfection, that is to say, the birth. And the same principle and order prevail in
plants; the seed is cast into the earth, then it pushes its way both upwards and downwards, partly in
roots and partly in branches; after that it increases; and fourthly, it produces fruit; and in the same
manner again the trees, when made, first of all produce fruit, which subsequently grows; then, as it
becomes ripe, it changes colour; and, fourthly, and this is the last operation, it completes and
perfects its work, the consequence of which is the use and enjoyment of it by men.

(13) What is the meaning of, ``For the sins of the Amorites were not as yet completed?'' (Genesis
15:16). Some persons have said, that by this expression of the principle of Moses fate is
expressly introduced, as if, in truth, everything was to be accomplished according to some
particular hour and appointed period of time.

(14) What is the meaning of, ``And when the sun was in the west a flame arose?'' (\#Ge 15:17).
It means either that the sun himself appeared in the west in the similitude of a flame, or that some
other flame appeared at eventide, not lightning, but some fire like it, which descended from above.
The manifest interpretation of the oracle is this; but we must now discuss that which regards the
inner sense.

(15) What is the meaning of the expression, ``Behold there was a smoking furnace and torches of
fires, which passed through the middle of those divisions?'' (\#Ge 15:18). The literal meaning of
the statement is plain, for the fountain or root of the divine word will have the victims consumed, not
by that fire which is given for our use, but by that which descends from above, out of heaven, in
order that the purity of the essence of heaven may bear witness to the sanctity of the victims. But if
we regard the inward meaning of the words, all things which are done beneath the moon are here
compared to a smoking furnace, on account of the vapour which rises up out of the earth and
water. As also the divisions of nature are, as has been already shown, every portion of the world
being divided into two parts; and by these there are kindled, as it were, torches of fire, being
powers which are more rapid in motion and more efficacious, being burning, in truth, like divine
fiery discourses, at one time keeping the whole universe in a state of integrity reciprocally with
themselves, and at another cleansing away the superfluous darkness. But the following
interpretation may also be given with propriety in a more familiar manner. Human life is like unto a
smoking furnace, because it has not a pure fire and an unalloyed brilliancy, but a great deal of
smoke, smoking darkly through the flame, which causes mist and darkness, and an obscuration,
not of the body but of the soul, so that this last cannot discern things clearly, until God the
redeemer commands the heavenly lamps to arise, I mean those more pure and more holy
radiations which unite those parts previously divided in two, on the right hand and on the left, and,
at the same time, illuminate them, being the causes of harmony and of lucid clearness.

(16) Why did he say, ``On that day, God made a covenant with Abraham, saying, To thy seed will I
give this land, from the river of Egypt to the great river Euphrates?'' (\#Ge 15:19). The literal
expression describes the boundaries of the space which lies in the middle, between the two rivers
Egyptus and Euphrates, for anciently the river was also called by the same name as the district,
Egypt, as the poet also testifies when he says-

``And in the river Egypt did I fix

My double-oared Ships.''\footnote{the line is in Odyssey 14.258.}

But if we look to the inner meaning of the expression, it intimates happiness, which is the perfect
fulness of three good things, namely, of spiritual, and corporeal, and external blessings, as some
of those men describe it in their panegyrics, who were afterwards called philosophers, such as
Aristotle and the Peripatetics; nevertheless, such a giving of the law as this is called Pythagorean.
Therefore the Egypt is a symbol of corporeal and external blessings, and the Euphrates of spiritual
advantages, in which alone, it is plain, their real joy consists, which has wisdom and all the other
virtues for its foundation; and the boundaries of this happiness are very rightly described as
beginning with the Egyptus and ending with the Euphrates; for the things affecting the soul come at
the end, which we usually approach with difficulty after we have passed through corporeal and
external things, in such a manner that, by this progress, we have felt our unity, the integrity of our
outward senses, and the beauty and strength which existed in our youth, advance, increase, and
come to maturity. And in a similar manner, those things which relate to acquiring gain and to
trafficking, as the management of ships, and agriculture, and commerce; for it is well said, that all
things, especially those above-mentioned, become a young man.

(17) Who are the Kenites, the Kenezites, and the Kadmonites, and the Hittites, and the Perizzites,
and the Rephaims, and the Amorites, the Canaanites, the Girgashites, and the Jebusites?
(\#Ge 15:20). Ten nations of wickedness are here enumerated, which he here destroys
because of their neighbourhood, since the number ten, when false and improperly stamped, is very
near to that which is good and an object of affection; but the complete perfection of the number ten
is exceedingly fit, as being the measure of infinite numbers, since the world is arranged in
accordance with it, and so likewise is the mind of the wise man, the substance of which,
nevertheless, wickedness perverts and overthrows, despising all very necessary powers, so that
that alone remains which the sacred writer has said, namely, that the pursuit of virtue is a blessing,
for the wicked man is such that he embraces vague opinions rather than truth, and of such is
Ishmael, though the seed of the prophets.

(18) Why it was that Sarah, the wife of Abraham, bore him no children? (\#Ge 16:1). The
mother of opinion is here spoken of as barren. In the first place in order that the son of generation
might appear more wonderful, as being born by a miracle. In the second place in order that his
conception and nativity might appear to be owing not more to the marriage of the man than to
divine providence. For it is not owing to the faculty of conception that a barren woman should bear
a son, but rather to the operation of divine power. This is the literal meaning of the statement. But if
we look to its inward sense, then we shall say, in the first place, that to bring forth is peculiar to the
female sex, as to beget is the office of the male: therefore God wills in the first place to render the
mind, which is filled with virtue, like to the male sex rather than to the female, thinking it suited to its
character to be active, not passive. In the second place both do generate, both the virtuous mind
and the wicked one: but they generate in a different manner, and they produce contrary offspring,
the virtuous mind producing good and useful things, but the depraved or wicked mind producing
base and useless things. In the third place he who is still advancing and making progress is to be
incited to the summit itself, and is near to the light which by some persons is said to be delivered to
oblivion, and to be made unknown. He therefore, as he is making progress, does not generate bad
things, nor yet good things, because he is not yet perfect; but he resembles that man who is neither
sick nor yet thoroughly well, but who, after a long sickness, is at last proceeding to convalescence.

(19) What is the meaning of the statement, she had an Egyptian handmaid whose name was
Hagar? (\#Ge 16:1). Hagar is interpreted travelling, and she is the servant of a more perfect
nature, being by nature an Egyptian less naturally; for the study of encyclical learning loves an
abundance of knowledge, and abundant knowledge is, as it were, the handmaid of virtue, since the
whole course and connection of sciences and arts is subservient to his use who is able to profit by
their acquisition so as to attain to virtue, for virtue has the soul for its abode; but the course of arts
and sciences stands in need of bodily instruments. But the body is symbolically Egypt; therefore
the sacred writer here properly asserts the likeness of encyclical knowledge to Egypt.
Nevertheless he has also given it a name by reason of its travelling abroad, since sophistry is a
foreign thing, unconnected with the acquisition of that wisdom which alone is native, and which
alone is necessary, which is the mistress of intermediate wisdom, and which conducts itself in a
beautiful course through the guidance of encyclical studies.

(20) Why did Sarah say to Abraham, Behold the Lord has shut me up so that I shall not bring forth:
go in now unto my handmaid so as to beget a son by her? (\#Ge 16:2). In the actual letter of
this statement it is the same thing to feel no envy, and also to provide for the welfare of the wise
man who is her husband and her genuine brother; so that she, wishing to find a remedy for her own
barrenness by means of her handmaid of whom she was mistress, gives her as a concubine to her
husband. But there is a still greater abundance of her affection towards her husband indicated by
this; for as she herself was accounted barren, she did not think it reasonable that the family of her
husband should be left entirely without offspring, but preferred his advantage to her own dignity.
This is what is indicated by this statement taken literally. But if we look to the inner sense of the
passage it bears such an interpretation as this: it becomes those persons who are unable in
respect of their virtue to bring forth beautiful works deserving of praise, to apply themselves to the
intermediate kind of study, and, if I may so express myself, to procure themselves children from
the encyclical branches of knowledge; for an abundance of knowledge is as it were the whetstone
of the mind and of the intellect. And it is with great propriety that she says, The Lord has shut me
up; for that which is shut up is generally opened again at a seasonable time. Therefore she was
not destitute of hope, nor was her wisdom fixed in the belief that she should be for ever without
offspring, but she knew that some day or other she should bring forth. Nevertheless she will not
bring forth at present, but when the soul displays the purity of its perfection. But inasmuch as it is at
present imperfect it is satisfied with using a milder kind of learning, such as is attainable by
encyclical studies. On which account it is not without a purpose that in the sacred contests at
Olympia also, those who are unable to attain to the first prize of victory are contented to be thought
worthy of the second; for there is offered to the competitors a first, and a second, and a third prize
by the presidents of the games, who are representatives of nature. So now to her the sacred writer
attributes the first prize of virtues, and the second prize of encyclical study.

(21) Why has he called Abraham's wife Sarah, for he says, Sarah the wife of Abraham, taking her
handmaid Hagar the Egyptian, gave her into his hand? (\#Ge 16:2). The sacred writer here
sums up with his approbation the marriage of the good on account of those who are incontinent
and lascivious; for those persons despise their wise wives for the sake of concubines, whom they
love with a frantic passion: on which account he here introduces the man endued with virtue, the
constant husband of one wife, at that time in which it was lawful for him to make use of her
handmaid; and his wife in fact indicates that he is wise, that is to say temperate, when he enters
into the bed of another woman, since his connection with his concubine was only a connection of
the body for the sake of propagating children; but his union with his wife was that of two souls
joined together in harmony by heavenly affection. This is the literal effect of the statement. But if
we look to the inner meaning of it, then he who has truly entrusted all his secret wishes to wisdom,
and justice, and the other virtues, when once he has received the counsel of wisdom, and has
tasted the joys of a matrimonial connection with it, remains constant to it as the partner and
companion of his life; although encyclical education would lead him in a beautiful course, since
when the man eminent in virtue has become master of the sciences of geometry, and arithmetic,
and grammar, and rhetoric, and the other exercises of the mind, he is not the less on that account
mindful of the pursuit of honesty, but is borne on towards the one as to a necessary aim, to the
other as an accessory. But it is altogether fair that that fact also should meet with our
approbation, -the fact I mean of his calling his handmaid also by the name of wife, because he went
up to her bed out of complacency to and at the exhortation of his real wife, and not of his own
genuine inclination; on which account he no longer calls her his handmaid, that even if it were not
wholly deserved still his handmaid having been given to him to wife might at least obtain the same
title. But those who study allegory may be allowed to say that the exercise of the middle disciplines
also stands in the place of a concubine, having nevertheless the shape and ornaments of a wife,
for all encyclical learning re-produces in itself and imitates genuine virtue.

(22) What is the meaning of, ``When she saw that she had conceived her mistress was despised
before her?'' (\#Ge 16:4). The sacred writer now carefully calls Sarah the mistress when it
might else have been thought that her dignity was diminished, and that she was surpassed by her
handmaid, that she, that is, who had no children, was surpassed by her who was gifted with
offspring. But this kind of language is extended to nearly all the necessary affairs of human life: for
a poor man who is wise is more approved of and is superior in authority to a rich man who is
destitute of wisdom and reputation, or than a boasting man; and even a sick man who is wise is
better than a foolish man who is well; for whatever is united with wisdom is genuine, and is endued
with an authority of its own, but whatever is combined with folly is found to be slavish and
inconstant. But it has been excellently said not that she despised her mistress, but that her
mistress was despised; for the one statement would imply an accusation of the person, but the
other contains only a declaration of an event. The scripture forsooth does not intend here to
impute blame to any one while praising another, but only to hand down in an intelligible manner the
pure truth of the facts. This is what is indicated by the literal statement. But if we seek the inner
meaning of the words, whoever honours and embraces rank before genius and wisdom, and
whoever esteems and considers the external senses of more importance than prudence and
counsel, is departing from the real character of things, thinking that they have brought forth much
offspring, and that having produced a great generation of visible things they are great and perfect
goods, and in a singular degree noble, but that barrenness in this respect is evil, and deserving of
disapprobation, because they do not see that invisible seed and that offspring which is appreciable
only by the intellect, which the mind is accustomed to generate in itself and by itself.

(23) Why does Sarah as it were repent of what she has done, saying to Abraham, I am receiving
injury from you: I gave my handmaid into your bosom, and now, because she sees that she has
conceived, I am despised before her? (\#Ge 16:5). This language indicates her anxiety and
hesitation; displaying them first in the expression, ``since,'' that is to say from the time that I gave
my handmaiden, and in the second place it betokens a regard to the person of whom complaint is
made, for she says, ``I am receiving injury from you,'' a statement which in fact is a reproof, since
she thinks that her husband ought always to be preserved without any stain, or any liability to
blame, always virtuous and true, and in no respect forgetful of her, for she always introduces him,
honouring him with all possible veneration, and calling him lord. Nevertheless, the first fact stated
by her is true; for from the time that she gave her handmaiden to him to be his concubine, she
herself was looked upon as despised. This is the literal meaning of her words. But if we look to
their inner sense, when any one bestows on another the handmaid of wisdom, she being
influenced by the counsels of sophistry, will, because she is ignorant of propriety, despise her
mistress; for as she herself possesses encyclical knowledge, and is delighted with its brilliancy,
where every one of the separate branches of education is by itself very attractive to the soul, as if
it possesses the power of drawing it by force to itself, then she, the handmaiden, can no longer
agree with her mistress, that is to say, with the image of wisdom and its glorious and admirable
beauty, until that acute judge of all things, the word of God, coming in, separates and distinguishes
what is probable from what is true, and the middle from the extremities, and what is second from
what is placed in the first rank. On which account Sarah says, at the end of her remonstrance, ``Let
God judge between me and thee.''

(24) Why does Abraham say, ``Behold thy handmaid is in thy hand, do unto her what seems good to
thee?'' (\#Ge 16:6). The literal expression used by the wise man contains a panegyric; for he
does not call the woman who had conceived by himself, his wife, or his concubine, but the
handmaiden of his wife. But since he saw that she also was a mother, he did not indulge in anger
and embitter the feelings of her mind, but rather tranquillised her, and made her prudent. But the
passage contains an allegory in the expression, ``In thy hand:'' as if, if I may so say, sophistry lives
under the dominion of wisdom, which indeed does spring forth from the same fountain, but only in
one part, and not directly; nor does it preserve the whole of its emanations pure, but draws up with
its waters many fetid things, and many others of a similar character. Since, therefore, it is in thy
hand and in thy power (for to whomsoever wisdom belongs, he is possessed also of all the
branches of encyclical learning), do with it whatsoever pleases thee, for I am quite persuaded that
you will judge with not more severity than justice; because that very thing is especially agreeable to
you: I mean the distributing to every one according to his deserts, and giving to no one more than
is just, either in the way of honouring or despising him.

(25) Why does he say, Sarah afflicted her? (\#Ge 16:6). The literal meaning of the words is
plain: but if we look to the inner sense of them, they contain a principle of this kind. It is not every
affliction that is injurious, but there are even some occasions when they are salutary; and this is
experienced by sick men at the hands of physicians, and by boys under their tutors, and by foolish
people from those who correct them so as to bring them to wisdom. And this I can by no means
consent to call affliction, but rather the salvation and benefit of both soul and body. Now a part of
such benefit wisdom affords to the circle of encyclical knowledge; rightly admonishing the soul
which is devoted to an abundance of discipline, and which is pregnant with sophism, not to rebel as
if it had acquired some great and excellent good, but to acquiesce and venerate that superior and
more excellent nature as its genuine mistress, in whose power is constancy itself, and authority
over all things.

(26) Why did Hagar flee from her face? (\#Ge 16:6). It is not every soul which is capable of
proper respect and of submitting to salutary discipline, but the mind which is gentle, and
good-tempered, and consistent loves reproof, and becomes more and more attached to those who
correct it. But the stubborn soul becomes malignant and hates them, and turns from them, and
flees away from them, preferring those discourses which are agreeable rather than those which
tend to his advantage, and looking upon them as more excellent.

(27) What is the meaning of the statement, ``The angel of the Lord found her sitting by a fountain of
water in the desert in the way to Sur?'' (\#Ge 16:7). All these statements are as it were symbols
by which the sacred writer indicates that the wellinstructed soul, which is the possession of virtue,
is nevertheless not yet able to discern the beauty of her mistress. They are, I say, symbols; I mean
the statements that she was found, and that she was found by an angel in the desert, and in no
other way than that leading to Sur. However we must begin with what is plain. Now the too subtle
sophist and the real lover of disputations is commonly unable to be detected by reason of his
artifices and sophistical persuasions, with which he is accustomed to deceive and perplex men.
But he who, being free from bad habits, has only an eager desire for obtaining instruction by the
course of encyclical training, although he is difficult to be detected, is yet not altogether incapable
of being so; for perdition is near at hand to him who cannot be detected, but safety to him who can
be discovered, especially when he is sought for and found by a more holy and more excellent spirit.
And who is more holy and more excellent than the angel of the Lord? For it is to him that it has
been entrusted to seek out the erring soul, the soul which, on account of its presumed erudition, is
continually ignorant of her whom it ought to respect, but still she could be susceptible of correction
and amendment; for which object she was sought out. Nor was she found imperfect, but ready to
the hand, since the soul was found which had fled from perfect virtue, not being able to submit to
discipline. But the third symbol takes place after she is found and after the discovery has been
made by an angel, namely, in the fact of her being found by a fountain, that is to say, by nature; for
it is nature which bestows on clever people abilities in proportion to the industry of each individual,
effacing unseasonable learning, which is no learning at all: and praise is implied in the very place in
which the soul is found, which is thirsting after genius and after its placid law, wishing to draw water
while in the society of those who drink wine; for thus it associates with those who feed upon and
are delighted with the exercise of proper training, where nature itself affords sufficient
nourishment, namely, education and instruction as if from a fountain. The fourth symbol is
contained in the fact that the discovery took place in the desert; since difficulty coming over each
of the outward senses, together with an influx of each separate desire, represses the mind, and
does not permit it to drink pure water: but when it cannot avoid these things as in the desert, it
acquiesces, and, abandoning the thoughts which agitated and perplexed it, it becomes
convalescent, so as to receive a hope not only of life, but even of eternal life. The fifth symbol is
contained in the fact that she was found in the way; for dispositions which are incorrigible are led
by devious paths; but that one which can be changed for the better, lo! it proceeds along the road
which leads to virtue, and that road is like a fortified wall and guardian to the souls which are
capable of being saved, for Sur means a fortified wall. Do you not see, then, that the whole is a
symbolical, or indeed a legitimate, figure of an improving soul? And, in fact, the soul which is
improving does not perish as one which is wholly foolish does; for if the divine word be found by it,
then again it seeks it; and he who is not pure and clean in his habits and disposition, flees from the
divine word; but yet he has a fountain of water in which he washes away his vices and
wickednesses, drawing from thence the fertility of the law. Besides this, it loves the desert, to which
it has fled from its vices and wickednesses, and when it has once beheld the way of virtue it
returns from the devious paths of wickedness. And all these things are fortified walls and bulwarks
to it, so as to protect it from being ever injured by any words of circumstances which attack it, and
from suffering any damage.

(28) Why did the angel say to her, ``Hagar, the handmaid of Sarah, whence comest thou, and
whither goest thou?'' (\#Ge 16:8). The plain letter of the question requires no explanation, for it
is exceedingly clear; but with reference to the inner meaning contained in it, there is come asperity
expressed; since the divine word is full of instruction, and is a physician of the infirmity of the soul.
Therefore the angel says to her, ``Whence comest thou?'' knowest thou not what good thou has
abandoned? Art thou not altogether lame and blind? For thou dost not see at all; and though
endowed with the outward senses, dost not feel, and dost not appear to me to have any portion
whatever of intellect, as if thou wert quite senseless. But ``whither goest thou?'' From what
excellence to what misery? Why have you so erred as to cast away the blessings which you had in
your power, and to pursue good things which are more remote? Do not, do not, I say, act thus; but,
quitting your insane impetuosity, go back again, and return into the same way as before, looking
upon wisdom as thy mistress, her whom you had before as your governess and directress in all the
things which you did.

(29) What is the meaning of the answer, ``I am fleeing from the face of Sarah, my mistress?''
(\#Ge 16:8). It is reasonable to praise a sincere disposition, and to think it friendly to truth. And
moreover it is reasonable now to admit the veracity of a mind which confesses what it has
suffered; for she says, ``I am fleeing from the face,'' that is to say, I have recoiled at the outward
appearance of wisdom and virtue; since, beholding its royal and imperial presence, she trembled,
not being able to endure to look upon its majesty and sublimity, but rather thinking it an object of
avoidance; for there are some people who do not turn from virtue from any hatred of it, but from a
reverential modesty, looking upon themselves as unworthy to live with such a mistress.

(30) Why did the angel say to her, ``Return to thy mistress and be humbled beneath her hands?''
(\#Ge 16:8). As the letter is plain, we must rather investigate its inner meaning. The word of
God corrects that soul which is able to be lured, and instructs it, and converts it, leading it to
wisdom as its mistress, that it may not, through being abandoned by its mistress, rush at once into
absurd folly. But it warns it, not only to return to virtue, but also to be humbled beneath its hands,
that is to say, beneath its several excellencies. But there are two kinds of humiliation; one, in
accordance with defect, which arises from spiritual infirmity, which it is easy to overcome, seize
upon, and reprove. But there is another kind which the word of the Lord enjoins, proceeding from
reverence and modesty; such as that humility which children exhibit to their fathers, pupils to their
masters, and young men to the aged; since it is very advantageous to be obedient, and to be
subject to those who are better than one's self; for he who has learnt to be under authority is in a
moment imbued with a power which he alone may exercise; for, although any one were to be
clothed with the authority of all the earth and sea, yet he would not be able to possess the royal
supremacy of virtue, unless he had first been instructed and taught to obey.

(31) Why did the angel say to her, ``I will multiply thy seed, and it shall not be numbered for
multitude?'' (\#Ge 16:9). It is the honour of the docile mind not to be presumptuous or rebellious
on account of its progress in knowledge, or because of the very useful seed which it has received
from various kinds of erudition; for it does not any more, as wordcatchers and cavillers do, employ
all the arguments of encyclical learning to establish any whimsical object, but to prove the truth
which is contained in them. And when it has begun to prosecute that by diligent investigation, it is
then rendered worthy to behold the sight of its mistress, free from all acceptance of persons, and
from all reproof.

(32) What is the meaning of the statement, ``The angel said to her, Behold, thou hast conceived,
and thou shalt bring forth a son, and shalt call his name Ishmael, because the Lord has heard the
voice of thy affliction?'' (\#Ge 16:11). The literal sense of the words admits of no question
except this allegorical explanation. Erudition, which is acquired and trained by the dispensation of
virtue as a mistress, is found not to be barren, but it has conceived the seed of wisdom; and when
it has conceived it brings forth; but it brings forth a work which is not perfect but imperfect, like an
infant which has need of care, and ailment, and nourishment; for in truth, it is quite plain that the
offspring of a perfect soul is perfect, that is to say, its words and works; but that of the soul of the
second class, which is still lying in servitude and subordination, is more imperfect. On which
account it has a certain name given to it, Ishmael, which is interpreted ``the hearing of God.'' But
hearing is honoured with the second dignity among the outward senses, being next to sight; for
nature has arranged a succession of ranks in the contests of the senses, giving the first place to
the eyes, the second to the ears, the third to the nostrils, and the fourth to that sense by which we
taste.

(33) What is the meaning of the statement, ``He shall be a wild man; his hand shall be upon every
one, and every one's hand shall be upon him, and he shall dwell over against all his brethren?''
(\#Ge 16:12). If we look to the letter of the statement, up to this time Ishmael has not any
brothers, for he was the first child of his parents. But the sacred writer is here figuring a certain
nature, too secret to be thoroughly investigated; for he has set forth the figure of his future
character. And such a figure evidently represents the sophist whose mother is erudition or wisdom.
But the sophist himself is a man of wild opinions; since the wise man as being civilized is fitted for
living in cities, and for urbanity, or for statesmanlike and political companionship; but he who is wild
and a man of wild opinions is immediately also quarrelsome. And it is on this account that the
sacred writer makes an addition, saying, ``His hand shall be upon every one, and every one's hand
shall be upon him;'' for the abundance of science and the use of erudition is able to contradict all
men. As those men of the present day who are called academicians and inquirers, consistently
setting no bounds to the determinations of their will and resolution, and among the different
opinions which they investigate preferring neither this nor that one, admit those men to be
philosophers who attack the opinions of every sect; and those whom it has been usual to call
opposers of will, as if they called them thele—machoi or thele—mamachoi, because they in the first
place raise contentions and declare themselves the champions of their national sect, not to be
convinced or put down by those who oppose them. But they are all kinsmen, and as it were
brothers of the same womb, being the offspring of one mother, namely, of philosophy. And it is on
this account that he says, ``And he shall dwell over against all his brethren;'' for in good truth the
academician and the inquirer are diametrically opposed to sects, finding fault in each of them with
their certain limitation of the resolution.

(34) Why does he say, ``But Hagar called on the name of the Lord, who spoke to her, saying, 'Thou
God who hast had regard unto me.' Because he said, 'In truth I have beheld thee appearing before
me.'''(\# Ge 16:13). In the first place, take notice carefully that the angel, after the manner of
the handmaiden of wisdom, was a minister to her on the part of God. But still why is he here called
Lord or God who ought only to have been styled his angel? It was in order to adapt the fact to the
proper person; for it was right that the Lord and chief of all the universe should appear to wisdom
as God, and that his word should appear as a minister to the handmaid and servant of wisdom. But
we may not suppose that she mistakenly looked upon the angel as God; for those who are unable
to behold the first cause may easily be deceived and look upon the second as the first; in the same
manner as he who has but weak sight, not being able to behold the sun which is in heaven in its
real appearance, thinks that the ray which falls upon the earth is the sun itself; and those who have
never seen the king attribute frequently the dignity of the supreme sovereign to his ministers. And
in truth mild and rustic men who never have beheld a city, not even from the summits of the hills
where they live, think every country house or farm-yard a mighty city, and look upon the people who
dwell there as citizens of a great city, out of ignorance of what a city really is.

(35) What is the meaning of, ``On this account she called that well the well of him whom I have seen
face to face?'' (\#Ge 16:14). The well has both a spring and depth. But the learning of the
students of encyclical science is neither all on the surface, nor is it destitute of first principles; for it
has for its source corrective discipline. Therefore it is with perfect correctness that she says that
the angel appeared before the well as God; since the erudition of the encyclical training
possessing the second rank is supposed to rejoice in the first authority, though it is in reality
separated from that first wisdom which it is permitted to wise men to behold, but not to sophists.

(36) Why is the well said to have been between Cadesh and Pharan? (\#Ge 16:14). Cadesh is
interpreted holy, but Pharan is translated hail, or corn.

(37) What is the meaning of the statement, ``Hagar brought forth a son to Abraham?'' (Genesis
16:15). It is made in perfect accordance with nature; for no habit of possession brings forth for
itself, but for him who possesses it; as grammar does for the grammarian, and music for the
musician, and mathematical science for the mathematician; because it is a part of him, and stands
in need of him. And the habit is not received as a thing in need of something, just as fire has no
need of heat, for it is heat to itself; and it gives a portion of the participation in it to those who
approach it.

(38) Why is Abraham said to have been eighty and six years old when Hagar bore Ishmael to him?
(\#Ge 16:16). Because the number which follows eighty, that is to say six, is the first perfect
number, being equal to its parts, and being the first number which is composed of the multiplication
of an odd and an even number; receiving also something from its efficient cause according to the
odd or redundant number, and from its material and effective cause according to the even number.
On which account, among the most ancient of our ancestors, some persons have called it
matrimony, and others harmony; and our sacred historian too has divided the creation of the world
into six days. But among numbers, eighty rejoices in perfect harmony, since it is composed of two
generous diameters in a double and treble proportion, according to the figure of a square of four
sides. And it contains within itself all the four inferences; the arithmetical, and the geometrical, and
the harmonious one. Being in the first place composed of double numbers, as of six, eight, nine,
twelve, the union of which makes thirty-five; in the second place of triple number, six, nine, twelve,
eighteen, the sum of which amounts to forty-five. And from these two numbers thirty-five and
forty-five, the whole number eighty is completed. Again, when the sacred historian Moses himself
began by divine inspiration to utter the oracular precepts which he was commissioned to deliver, he
was eighty years old. And the first man who existed of our nation according to the law of
circumcision, being circumcised on the eighth day, being eminent for virtue, bears that name of joy,
being called Isaac in the Chaldaic tongue, and Isaac means laughter; being naturally called so
because nature rejoices or laughs at everything, being never vexed at any thing which is done in
the world, but rather looking with complacency on every thing which occurs as being done well and
profitably.

(39) Why when he was ninety and nine years old does the sacred writer say, ``The Lord God
appeared to him and said, I am the Lord thy God?'' (\#Ge 17:1). He here makes use of both the
titles of each superior virtue, applying them in the case of his address to the wise man, because it
was by them that all things were created, and by them that the world is regulated after it had been
created. By one of them therefore the wise man, just in the same manner as the world itself, was
fashioned and made according to the likeness of God; and God is the name of creative virtue; and
by the other of them that he was made according to the Lord, as falling under his authority and
supreme power. Therefore he designs here to show that the man who is conspicuous in virtue is
both a citizen of the world, and also equal in dignity to the whole world, declaring that both the
virtues of the world, the divine and the royal attributes, are in a singular manner appointed to and
set over him as protectors. And it was with great correctness and propriety that this appearance
took place when he was about ninety and nine years old, because that number is very near the
hundred. And the number a hundred is composed of the number ten multiplied by itself, which the
sacred historian calls the holy of holies. Since the first court, the first ten, is simply called holy, and
that is permitted to be entered by the sweepers of the temple; but the ten of tens, which he again
enjoins the sweepers of the temple to pay above all things to the existing high priest, is the number
ten computed along with the number a hundred, for what else is the tenth of the tenths but the
hundredth? However the number ninety and nine has been set forth and adorned not only by its
affinity to the number a hundred, but it has also received a particular participation in a wonderful
nature, since it consists of the number fifty, and of seven times seven. For the fiftieth year, as the
year of Pentecost or the Jubilee, is called remission in the giving forth of the law, as then all things
are given their liberty, whether living or inanimate. And the mystery of the seventh year is one of
quiet and profound peace to both body and soul. For the seventh is the recollection of all the good
things which come of their own accord without industry of labour, which at the first creation of the
world nature produced of herself; but the number forty-nine, consisting as it does of seven times
seven, indicates no trifling blessings, but rather those which have virtue and wisdom, in such a
degree as to contribute to invincible and mighty constancy.

(40) What is the meaning of, ``Do thou please me, and keep thyself from stain, and I will make my
treaty between me and thee, and I will multiply thee exceedingly?'' (\#Ge 17:1). God here lays
down a law for the human race in a somewhat familiar manner; for he who has no participation in
wickedness and is free from evil, will be perfectly good, which is peculiar to incorporeal natures.
But those who are in the body are called good in proportion to the measure in which wickedness
and the practice of sin are removed from them. Therefore the life of those men has appeared
honourable, not that of those who have been free from sickness from the beginning to the end, but
that of those who from a state of infirmity have advanced to sanity; on which account he says
directly and plainly, ``Keep thyself free from stain,'' for it is sufficient to conduct a mortal nature to
felicity not to be blamed, and neither to do nor say anything deserving of reproof; and such conduct
is at once pleasing to the Father. Therefore it is that he said, ``Do thou please me, and keep thyself
free from stain.'' Where the form of expression implies a mutual conversion; since the habits which
please God do not deserve reproof, and he who keeps himself free from stain and avoids reproof
in all things is altogether pleasing to God. Therefore he promises to bestow a double blessing on
him who keeps himself free from all reproof; in the first place, to make him the guardian of the
deposits of the divine covenant: and in the second place to cause him to increase to a multitude
without any limit. For that expression, ``I will make my treaty, or covenant, between me and thee,''
shows the office of guardianship of the truth which is entrusted to an honest man; for the whole
treaty of God is the incorporeal word; which is the form and measure of the universe according to
which this world was made. And then repeating the expression, ``I will multiply thee exceedingly,''
twice manifestly shows the immense numbers to which the multitude promised shall grow, I mean
the increase which shall take place in the people, not in human virtue.

(41) What is the meaning of, ``Abraham fell on his face?'' (\#Ge 17:3). The present expression
is the interpretation of what has already been promised; for God had said, ``Keep thyself free from
stain,'' but there is not other cause of a man leading a life which is disapproved but the outward
sense, because that is the origin and source of the passions; on which account he rightly and
properly falls on his face, that is to say, the offences caused by the outward senses fall to the
bottom, showing that the man is now devoted to all good works. This is enough to say in the first
place, But in the second place we must say that he was so struck by the manifest appearance of
the living God that he was scarcely able to behold him through fear, but fell to the ground and
offered adoration, being overwhelmed with awe at the appearance which presented itself to him. In
the third place, he fell to the ground on account of the revelation thus made to him, at the form of
his appearance by the living God who exists alone, whom he knew and regarded as truth opposed
to created nature; since the one exists in unvarying constancy and the other vacillates and falls
into its proper place, that is to say, to the earth.

(42) What is the meaning of, ``And God conversed with him, saying, And I, behold, my covenant is
with thee, and thou shalt be the father of a multitude of nations?'' (\#Ge 17:4). Since he had
previously used the expression, ``treaty,'' he now proceeds to say, do not seek that treaty in letters,
since I myself, in accordance with what has been said before, am myself the genuine and true
covenant. For after he has shown himself and said, ``I,'' he makes an addition, saying, ``Behold, my
covenant,'' which is nothing but I myself; for I am myself my covenant, according to which my treaty
and agreement are made and agreed to, and according to which again all things are properly
distributed and arranged. Now the form of this prototypal treating is put together from the ideas and
incorporeal measures and forms in accordance with which this world was made. Is it not therefore
a climax to the benefits which the Father bestowed on the wise man, to raise him up and conduct
him not only from earth to heaven, nor only from heaven to the incorporeal world appreciable only
by the intellect, but also to draw him up from this world to himself, showing himself to him, not as he
is in himself, for that is not possible but as far as the visual organs of the beholder who beholds
virtue herself as appreciable by the intellect are able to attain to. And it is on this account that he
says, ``Be no more a son but a father; and the father, not of one individual but of a multitude; and of
a multitude, not according to a part, but of all nations;'' therefore of the revealed promises two
admit of a literal interpretation, but the third of one which is rather spiritual. One of those which
admit of a literal interpretation is to be construed in this way: in truth thou shalt be the father of
nations, and shalt beget nations, that is to say, each individual among thy sons shall be the founder
of a nation. But the second is of this kind; like a father you shall be clothed with power over, and
authority to rule, many nations; for a lover of God is necessarily and at once also a lover of men;
so that he will diligently devote his attention, not only to his relations but also to all mankind, and
especially to those who are able to go through the discipline of strict attention, and who are of a
disposition the reverse of anything cruel or hard, but of one which easily submits to virtue, and
willingly gives obedience to right reason. But the third we may explain under this allegory: the
multitude of nations spoken of indicates as it were the multifarious inclination of the will in each of
our minds, both those inclinations which it is accustomed to form with reference to itself, and also
those others which it admits by the agency of the senses, as they enter clandestinely through the
intervention of the imagination, and if the mind possesses the supreme authority over all these, it,
like a common father, turns them to better objects, cherishing their infant opinions, as it were, with
milk, exhorting those which are older and more mature, though still imperfect, to improvement, and
honouring with commendation those which perform their duty aright; and again, putting a bridle, by
means of discipline and reproof, on those which rebel and act rashly; since, wishing to imitate the
Deity, it receives a twofold influx from the virtues of that same being, one from his beneficent
attributes and another from his avenging might, as if from two sources; therefore the docile receive
his kindness, and towards the rebellious he uses reproof; so that some are led to improvement by
praise and others by chastisement: in truth, he who is eminent for virtue is able to be of great, and
extensive, and just service to all, according to his power.

(43) What is the meaning of, ``Thy name shall not be called Abram, but Abraham shall thy name
be?'' (\#Ge 17:5). Some of those who are destitute of all knowledge of music and dancing,
some indeed being wholly foolish and keeping aloof from the divine company, mock the one
existing or only wise Being, immaculate by nature, saying, in a tone of vituperation, ``Oh the great
gift, the governor and Lord of the whole universe has given one letter, by which the name of the
patriarch was to be increased and become of great importance, so as to be made a trisyllable
instead of a dissyllable!'' Oh the great misery, and wickedness, and impiety, of such men! If some
persons dare, in any respect, to endeavour to detract from God, being deceived by the outward
appearance of a name, when they ought rather to thrust their minds down into the depths, and
inquire into the things themselves more closely, on account of the real magnitude and importance
of the possession. Besides this, why do ye not think the concession of one letter, although a small
and easy gift, nevertheless an act of providence? and why do ye not weigh its value? since, above
all things, the very first element of language, as expressed in letters, is A, both in order and in
virtue. In the second place, it is also a vowel, and the very first of vowels, being placed above them
as their head. In the third place, because it does not belong to long properties, nor to short
properties, but it is of the number of those which comprise each characteristic, for it is extended
into greater length, and then again it is recalled into shortness, by reason of its softness,
resembling wax, and being figured into many shapes, and afterwards figuring words, according to
infinite numbers; besides all this it is a cause, for it is the brother of unity, from which all things
begin and in which all things terminate. Therefore, when any one sees such great beauty, and a
letter set forth with such great importance and necessity, how can he accuse it as if he had not
seen this? for if he has seen it, he then shows himself to be a person of insulting disposition and a
hater of what is good; and if he has not seen a fact, which is so easy to comprehend, how does he
presume to ridicule and despise that which he does not understand as if he did understand it? But
however these things may be said by the way, as I stated before. But we must now examine into its
necessary and most important task. The addition of the letter A, by one single element, changed
and reformed the whole character of the mind, causing it, instead of the sublime knowledge and
learning of sublime things, that is to say, instead of astronomy, to acquire a comprehension of
wisdom, since it is by the knowledge of things above that the faculty is acquired of mounting up to
one portion of the world, that is to say, to heaven, and to the periodical revolutions and motions of
the stars; but wisdom has reference to the nature of all things, both such as are visible to the
outward senses, and such as are appreciable only by the intellect, for the intellect is the wisdom
which gives a knowledge of divine and human things and of their principles. Therefore, in divine
things there is something which is visible, and something else which is invisible, and a
demonstrative idea. And in human affairs there are some things which are corporeal and some
which are incorporeal; to attain to the right comprehension of which is a great task, and a real
employment for the abilities and courage of man. But to be able, not only to behold the substances
and natures of the universe, but also the principles which regulate each separate fact, indicates a
virtue more perfect than that which is allotted to mankind; for it is necessary for the mind, which
perceives so many and such great things, to be altogether and wholly eye, and to dispense with
sleep, passing its whole existence in the world in a state of incessant wakefulness, and being
surrounded by a light which knows no darkness, and which exhibits the appearance of light itself,
as by an ever-flashing lightning, taking God for its leader and guide, to the comprehension of the
knowledge of those things which are, and to the faculty of explaining their principles. Therefore the
dissyllabic name Abram is explained as meaning ``excellent father,'' on account of his affinity to the
knowledge of sublime wisdom, that is, astronomy and mathematics. But the trisyllabic name
Abraham is interpreted ``the father of elect sound,'' being the name of a really wise man; for what
else is sound in us, except the utterance of a pronounced word? for which object we have an
instrument constructed by nature, passing through the thick tube of the throat, and united with the
mouth and tongue; and the father of such a sound is our intellect, and elect intellect is endued with
virtue. But if we are to keep to exact propriety, then it is plain that the mind is the familiar and
natural father of the uttered word, because it is the especial property of the father to beget, and the
word is born from the mind; and it will be a certain proof of this if we recollect that when it is set in
motion by counsels it sounds, and when they are absent it ceases to sound: and the evidences of
this are the rhetoricians and philosophers who demonstrate its habit by objects; for whenever the
mind publishes abroad different heads of designs, and in the manner of a mother about to bring
forth produces each individual means previously stored up in itself, then also the word, flowing forth
like a fountain, is borne to the ears of the bystander as to its appropriate receptacles: but when
those are wanting, then it also is unable to publish itself further, and rests, and the sound is inactive
as being struck by no one. Now therefore, O ye men, full and crammed with superfluous loquacity,
ye men devoid of wisdom, does not the gift of one single element appear to you to have been such
that by the intervention of a single letter the wise man is rendered worthy of the divine attribute of
wisdom, than which there is nothing more excellent in our nature? because instead of the sublime
erudition of astronomy he gave him intellect, that is to say, instead of a small part of wisdom, he
gave him the whole and perfect blessing of entire wisdom, since a knowledge of things above is
included and comprehended in wisdom, as a part is included in the whole; for mathematics are only
a part. But it becomes you, O men, to consider this point also, that the man who is well instructed
and skilful in the investigation of the nature of things above may by possibility be a man of
depraved and wicked habits; but the wise man is altogether approved as virtuous. Shall we then
now any longer ridicule this gift, than which nothing more excellent can be found? For what is more
shameful than wickedness or more excellent than virtue? Can anything be found here not good,
and is it not wholly opposed to evil? Or can this gift be compared to riches, or honour, or liberty, or
health, or to any other superfluous possession of any kind around or exterior to the body? For the
whole of philosophy is thus added to our life as a sort of college of medicine to the soul, in order
from thence to dispense to it freedom from suffering and immunity from disease; but in truth it is
noble to be a philosopher, and that wonderful knowledge is truly noble; and the end is even more
admirable, on account of which the act is called into existence. Here therefore is wisdom, and that
the best kind of wisdom, which God called in the Chaldaic dialect Abraham, namely the father of
elect sound, giving as it were a definition of a wise man; for as the definition of man is a mortal
animal endowed with reason, so also the mysterious definition of a wise man is the father of elect
sound.

(44) What is the meaning of, ``I will greatly increase thee, and set thee among the nations, and
kings shall proceed from thee?'' (\#Ge 17:6). That expression, ``I will greatly increase thee,'' was
used to the wise man with exceeding propriety; since every wicked or bad man does increase and
advance, not to improvement but towards deficiency; as withering flowers advance not towards life
but towards death; but the man whose life is extended long and is greatly increased is like a
passing cloud, or like the continually flowing stream of a river, because as it increases it is
extended more and more out of doors, as its wisdom also is divine. And that expression, ``I will set
thee among the nations,'' was used in order that God might the more evidently demonstrate that he
was making him worthy to be as a foundation and firm support to the nations through his wisdom,
not only to his own nation, but also to all other peoples who in various manners are in want in
respect of their minds, as has been said before; since the wise man is the redeemer of nations
and intercessor for them before God, and since it is he who implores pardon for the sins of his
relations. Last of all, the promise, ``Kings shall come forth from thee,'' is again used with especial
propriety; for everything which relates to wisdom is a royal seed; the offspring of the chief and
master according to nature: but the wise man has no seed or fruit of his own, but is fertile and
abundant in the seed which proceeds from the great cause himself.

(45) What is the meaning of, ``I will give this land to thee and to thy seed after thee, in which thou
hast sojourned, namely all the land of Canaan, for an everlasting possession?'' (\#Ge 17:8).
The letter of the promise is so clear that the language does not stand in need of any explanation
whatever; but with respect to the inward meaning of it we must have recourse to an allegory of this
kind. The mind which is endowed with virtue is rather a sojourner in the corporeal space allotted to
it, than a regular inhabitant of it; for its real country is the air and the heaven; and the earth, and the
earthly body in which it is said to sojourn, is only a colony; therefore the Father, conferring a
benefit upon it, gives to it the sovereign authority over all the things of the earth for ever and ever,
as he says himself, for an everlasting possession; so that it for the future shall not be governed by
the body, but shall always be its master and ruler, having the body for its servant and attendant.

(46) What is the meaning of, ``And every male of you shall be circumcised, and you shall
circumcise, or you shall be circumcised, in, the flesh of your foreskin?'' (\#Ge 17:10). I see
here a twofold circumcision, one of the male creature, and the other of the flesh; that which is of
the flesh takes place in the genitals, but that which is of the male creature takes place, as it seems
to me, in respect of his thoughts. Since that which is, properly speaking, masculine in us is the
intellect, the superfluous shoots of which it is necessary to prune away and to cast off, so that it,
becoming clean and pure from all wickedness and vile, may worship God as his priest. This
therefore is what is designated by the second circumcision, where God says by an express law,
``Circumcise the hardness of your hearts,'' that is to say, your hard and rebellious thoughts and
ambition, which when they are cut away and removed from you, your most important part will be
rendered free.

(47) Why orders he the males only to be circumcised? (\#Ge 17:11). For in the first place, the
Egyptians, in accordance with the national customs of their country, in the fourteenth year of their
age, when the male begins to have the power of propagating his species, and when the female
arrives at the age of puberty, circumcise both bride and bridegroom. But the divine legislator
appoints circumcision to take place in the case of the male alone for many reasons: the first of
which is, that the male creature feels venereal pleasures and desires matrimonial connexions
more than the female, on which account the female is properly omitted here, while he checks the
superfluous impetuosity of the male by the sign of circumcision. But the second reason is, that the
material of the female is supplied to the son from what remains over of the eruption of blood, while
the immediate maker and cause of the son is the male. Because therefore the male supplies the
most indispensable part in the fact of generation, God deservedly represses his pride by the figure
of circumcision, but the material or feminine cause, as being inactive, does not display ambition in
the same degree. And this is enough to say on this head. But afterwards we must note this
likewise, that the intellect in us is endued with the power of sight, therefore it is necessary to cut
away its superfluous shoots. And these superfluous shoots are empty opinions, and all the actions
which are done in accordance with them. So that the intellect after circumcision may only bear
about with itself what is necessary and useful; and that whatever causes pride to increase may be
cut away; with which also the eyes are circumcised as if they did not see.

(48) Why did he say, ``And let the child, every male child, be circumcised at eight days old?''
(\#Ge 17:12). He orders the freeborn to be circumcised, which, in the first place, was permitted
on account of diseases that might arise; for it is more difficult to heal a disease in the genitals, and
it is commonly done by burning by fire those parts over which a membrane grows, but this rarely
affects those who have been circumcised. And in truth, if it were possible that other infirmities also
could be avoided by amputating any member or any part of the body, so that though it was
amputated still the operation of each necessary part would not be hindered, then without the
knowledge of mortal man he would be transmitted into immortality. But that here it was thought fit
that man should be circumcised out of a provident care for his mind without any previous infirmity is
plain, since not the Jews alone, but also the Egyptians, and Arabians, and Ethiopians, and nearly all
the nations who live in the southern parts of the world, down to the torrid zone are circumcised.
What then is the chief reason of this fact? except that in those districts, and especially in the
summer, when the genitals are protected with a skin, it burns and is injured by inflammation, but
when that covering is laid bare by circumcision it is cooled, and the disease is repelled; and on this
account the northern nations and others, to whom the cooler portion of the habitable earth has
been allotted, are not circumcised, for not only is the solar heat moderate in those regions, but so
is also all inflammatory disease which affects the membranes of the members. Let every one take
a firm judgment, and from that time when the disease comes in more vigorously; for it never comes
at all in the winter, but in the summer it shows itself and flourishes and ripens; for it loves, if I may
so say, like fire to burn in those parts. In the second place, it was not only from a regard to sound
health that our ancestors diligently employed this method of cure, but also from a regard to the
multiplication of the human race, seeing that nature was very vivacious and too eager to propagate
the human species. Therefore they knew, like wise men, how the seed when poured over the folds
of the membrane is often accustomed to be wasted and so to become unfruitful; but if no
impediment arises then it would easily be able to arrive at the situation suited to receive it. On
which account also those nations which adopt the practice of circumcision have grown into an
exceedingly numerous population: and our legislator, weighing the consequences also,
commanded the circumcision of infants to be performed at an earlier age, keeping in view the
same effect of circumcision with regard to the population. Therefore it is in truth, as it seems to
me, that the Egyptians also in the fourteenth year of their children's age, in which the desire to
propagate the species usually begins, have said that it is suitable to circumcise them, with the view
of increasing the population; but it was better and more carefully done in our nation, where the
circumcision of infants was ordained, since perhaps the man when grown up would delay the
operation out of fear, because he then has a will of his own. In the third place, he says this with a
view to cleanliness in the sacred oblations; for in truth those who enter into the courts of the
temples are made clean by sprinkling and ablutions. Moreover the Egyptians scrape the whole
body, removing all the hairs which cover and envelop the body, so as to appear white all over; but
the circumcision of the skin is no small assistance towards cleanliness, otherwise everyone would
abhor it when he beheld it as it is in itself. In the fourth place, there are in us two generative
principles, one in the soul and one in the body; the generative principle of the soul is the intellect,
and that of the body is the corporeal organ; therefore the ancients chose to refer the generative
principle of the body to an imitation of the intellect which is rather the generative principle of the
heart. And in truth there is nothing to which it is found more like than the circumcision of the heart;
these therefore are real facts like the celebrated reasons for things which have been investigated.
But we must now speak of those which have greater symbols belonging to them and which exhibit a
certain principle. Therefore the circumcision of the skin is said to be a symbol, but as one
indicating that it is proper to cut away all superfluous and extravagant desires, by studying
continence and religion; for as the skin of the prepuce is quite superfluous for generation, and is
moreover especially injurious by reason of the disease of inflammation which burns within it, so
also an over abundance of desire is as superfluous as it is pernicious, superfluous because it is
not necessary, and pernicious because it is the cause of diseases to both body and soul; and by
the greater desire he also warns us that all the other desires are likewise to be cut off. And that is
called the greater desire which has a regard to the matrimonial connexion of the male and the
female; since it is the beginning of a great thing, namely, of generation; and since it creates a great
affection on the part of the father towards her who is to bring forth; for it is natural for them both to
be influenced by love and affection for their offspring. Therefore, he here warns us to cut away not
only all the superfluous desires, but also pride, as being a great wickedness and an associate of
wickedness. For pride, as the language of the ancients tells us, is what keeps men back and
hinders them in their improvement; since it will not exhibit that honesty which it really possesses,
thinking that it is itself an adequate cause for anything. Moreover it naturally influences those who
think themselves the causes of generation; so that they scarcely ever turn their minds at all to
behold the true Father of the universe. For he is in truth the one real and genuine Father of all; and
we, who are called fathers, are only instruments of his, serving to generation; since, as in a
wonderful resemblance, all things which are represented in appearance are yet in reality inanimate,
but that which strengthens the nerves is invisible, and yet is itself the cause of virtue, and of
motion, and of sight. So, in like manner, from everlasting and invisible space there extends the
Creator of the universe, and we, like so many puppets, are strengthened by him with nerves for the
purpose which belongs to us, namely, sowing seed and raising a generation; unless we choose to
fancy that a flute is blown by itself, and is not made by an artist in a way adapted for the production
of harmony, by whom it was constructed as an instrument for service and for its own necessary
end.

(49) Why does he order circumcision to be performed on the eighth day? (\#Ge 17:12). The
number eight has many beauties in it; for it is, in the first place, a cubic number. Secondly, it has
beauties, because it everywhere contains in itself the form of equality, because longitude, and
breadth, and depth, which are all equal to one another, are indicated by the first number eight. In
the third place, the composition of the number eight produces agreement, namely, the number
thirty-six, which the Pythagoreans call agreement, since that is the first number in which odd
numbers being added together agree with even numbers. If, indeed, four odd numbers from the unit
are separately taken and added together and four even numbers beginning with two, they united
make thirty-six. Now the odd numbers are these: one, three, five, seven, which make sixteen. And
the even numbers are these: two, four, six, eight, which make twenty. And the addition of the two
together makes thirty-six, which is in truth a more fertile number. Since it is a square, having each
side composed of the number six; the first of which is both odd and even; which some persons
most correctly call harmony or matrimony; and it was by the employment of this number that the
Creator of the universe made the world, as the holy and admirable book of Moses relates. In the
fourth place, the idea of eight produces sixty-four, which is the first number, which is a cube and
also a square, being the type of incorporeal substance appreciable only by the intellect and
invisible, and also of corporeal substance. Of incorporeal substance, inasmuch as it produces
superficies according to the square; and of corporeal substance, as producing a solid according to
the cube. In the fifth place, it is always a kindred number to the virgin number seven, for seven
makes up the parts of eight; because four is the half of it, two is the fourth part of it, one the eighth
of it, and four, two, and one, added together, make seven. In the sixth place, the power of eight is
sixtyfour, which we call the first number, being both a cube and a square. In the seventh place,
taken separately from the units by these doubled numbers, one, two, four, eight, sixteen, thirty-two,
the sum makes sixtyfour. And the number eight has also other more distinguished virtues still,
which we have enumerated in another place; but now it seems better to explain the principle which
corresponds to the present question, and which depends on the grounds now laid down. But in the
first place we must premise this: that nation to whom it is enjoined, having the commandments give
to it, that it should be circumcised on the eighth day, is called in the Chaldaic language Israel, that
is to say, ``he that seeth by day.'' Therefore God wills, in the first place, that he should be a partaker
both of his own just rights, and also of those which exist according to election, and according to the
principle of Genesis (or creation), by that first number six, which immediately followed the creation.
This number, in fact, the Father and Creator of all things evidently exhibited to the world as the
festival of generation, completing the world on the sixth day. And the other number, that which is
according to election, he exhibited by the number eight, which is the beginning of the second
seven; as eight is seven and one, so the race which has been honoured is always a race receiving
that number also in addition, so that it should be elect, both by nature and in accordance with the
decree of the Father. In the second place, the number eight exhibits equality everywhere, showing
that all its separations are equal, as has been already said, I mean its length, and breadth, and
depth. And equality it is which is the parent of equity and justice, by which he shows that the nation
which loves God is adorned with equity or justice, and has advanced to complete possession. In
the third place, eight is not only a measure of complete equity in all its dimensions, but is the very
first number that is so, for it is the first cube; since the number eight indicates equality, and so it
has the second and not the first rank: therefore it demonstrates in a symbolical manner that that
nature was the first which was ever completely furnished with consummate and perfect equity and
justice, and that it is the first nature of the human race, not in point of creation or of time, but in the
dignity of virtue, as if justice united with equality were a connatural part of it. In the fourth place,
since there are four elements, the appearances of earth, water, air, and fire; fire has received for
its figure a shape becoming a similar name, a pyramid; \footnote{pyramis, resembling the word pyr, ``fire.''} and air
has received for its figure an eight-sided one; water, a twenty-sided one; and the earth, a cube.
Therefore he thought it necessary that the earth, which was to be the allotment of the race of man,
who were endowed with virtue, should participate in the cubic number, as the whole earth has been
formed in its figure. And a part of it receives the parts of that which should bring forth, because by
nature the earth is very fertile, producing all the various and distinct species of every kind of animal
and plant.

(50) Why does he order all slaves to be circumcised, those born in the family and also those who
are bought? (\#Ge 17:12). The literal meaning is plain; for it is fitting that servants should
imitate their masters on account of their necessary employment, and the services to which they
are bound in life. But with respect to the inner object of the command, those dispositions are what
may be called born in the family which are influenced by nature itself, and those are bought which
can be changed for the better by teaching and instruction. Each of these has its appropriate
employment, and requires like a plant to be cleared and pruned in order that the good and fruitful
parts may acquire constancy; for fertile plants produce many superfluous things by reason of their
fecundity, and those superfluities must be cut away; but those who are taught by instructors cut
away their ignorance.

(51) What is the meaning of, ``And it shall be my covenant (or agreement) in your flesh?'' (Genesis
17:13). God is willing to do good, not only to the man who is endued with virtue, but he wishes that
the divine word should regulate not only his soul but his body also, as if it had become its physician.
And it must be its care to prune away all excesses of seeing, and hearing, and taste, and smell,
and touch, and also those of the instrument of voice and articulation, and also all the redundant and
pernicious impulses of the genitals, as also of the whole body, the effect of which is, that at times
we are delighted by our passions and at times pained by them.

(52) Why is it that he pronounces a sentence of death on an infant, saying, ``Every male child who
is not circumcised, who has not been circumcised (or, as the Greek has it, who shall not be
circumcised) in the flesh of his foreskin on the eighth day, that soul shall be cut off from his
generation?'' (\#Ge 17:14). The law never declares a man guilty for any unintentional offence;
since even those who have committed an unintentional homicide are pardoned by it, cities being
set apart into which such men may flee and there find security; for whoever escapes to them is
rendered secure and free from danger; and no one has the power to drag him forth, or to cite him
before the tribunal of the judge for the deed. Therefore, if a boy is not circumcised on the eighth
day after his birth, what offence will he have committed that he is to be held guilty, and suffer the
penalty of death? Some persons may perhaps say that the form of the command points to the
parents themselves, for they look upon them as despisers of the command of the law. But others
say that it has here exerted excessive severity against infants, as it seems, imposing this heavy
penalty in order that grown up persons who break the law may thus be irrevocably subjected to
most severe punishment. This is the literal effect of the words. But if we look to their inward
meaning, then what is male in us is most especially the intellect, and that God here commands to
be circumcised on the eighth day, for the reason previously stated, not in any other part, but in the
flesh of the foreskin, by this expression symbolically indicating those parts which in the flesh do
subsequently become the organs of pleasure and impulse. And on this account it is that he
introduces a legitimate reason, warning men that the intellect, which is not circumcised and cleared
away from the flesh and the vices of the flesh, is corrupt and cannot be saved. But that this
language is not to be applied to the man, but to the intellect, which is thus put in a sound condition,
he tells us in the subsequent words, saying, ``that soul shall be cut off,'' not that human body, or that
man, but that soul and mind. Cut off from what? From its generation; for the whole generation is
incorrupt. Therefore the wicked man is removed from incorruption to corruption.

(53) Why does God say, ``Sara thy wife shall not be called Sara, but Sarra shall be her name?''
(\#Ge 17:15). Here again some foolish persons may laugh at the addition of one single letter,
that is to say, of a hundred, for in Greek characters the letter r means a hundred; but if they jest in
this way they are foolish, as being unwilling to behold the inward merits of things and to cleave to
the footsteps of truth; for that element, r, which is here thought of merely as the addition of one
letter, is the parent of all harmony, making things great instead of small, general instead of
particular, and mortal instead of immortal; since Sara, when called Sara with one r, is interpreted
``thy princedom,'' but with two r's, Sarra, ``princess.'' Let us then be careful, and see how these two
names are distinguished from one another. In me wisdom (or prudence), integrity (or temperance),
justice, and fortitude have only a prince-like power and are mortal; moreover, when I die they die
too. But this wisdom is herself a princess, and justice is a prince too, and each separate one of
these virtues is not the principal or princely part in me, but is itself a mistress and a queen, an
everlasting monarchy and sovereignty. Do you not now see the magnitude of the gift? By this slight
change, God changes the part into the whole, the species into the genus, the corruptible into the
incorruptible. And all these things are previously dispensed on account of the impending birth of a
more perfect joy than all joys, whose name is Isaac.

(54) Why does he say, ``And from her I will give thee children, and I will bless him, and he shall be
over the nations, and kings of the nations shall come forth from him?'' (\#Ge 17:16). It is
scarcely proper to inquire why he has said children in the plural number, when he meant their only
and beloved son; for the intention of God's words applies to his offspring, from which nations and
kings should arise. This is the literal meaning of the words. But if we look to their more inward
sense, when the soul possesses that virtue, small and mortal as it is, which is only particular, she is
still barren. But from the time that it acquires a share of the divine and incorruptible virtue, it begins
to conceive and to bring forth varieties of nations, namely, of all other holy and sacred persons; for
ever one of the everlasting virtues is subject to an immense number of voluntary laws, which bear
in themselves a similarity to nations and kingdoms; for virtue and the generations of virtue are
royal things, being previously instructed by nature what it is which rejoices in princely power, and
has no knowledge of a servile condition.

(55) Why did Abraham fall on his face and laugh? (\#Ge 17:17). Two things are indicated by his
falling on his face. One an act of adoration on account of the excess of his divine ecstacy; the
other that it corresponds to and is suitable to the aforesaid harmony, by which the intellect has
confessed that God alone exists in a continual and unvarying existence. But those creatures which
owe their existence to creation and generation, all are subject to changes in time; for they fall to a
certain extent, inasmuch as they are accustomed to rise up, and to be corrected in accordance
with their original appearance. And it was very natural for Abraham to laugh at the promise, as he
was then filled with the great hope that the things which he expected should be accomplished,
especially because he had received a manifest revelation from that appearance, by which he
became more thoroughly acquainted with him who exists for everlasting without variation, and with
him also who is continually stooping and falling.

(56) Why did Abraham appear to hesitate about the promise, for the sacred writer says, He said in
his mind, shall there be a son to one who is a hundred years old; and shall Sarra, who is ninety
years old, bring forth a child? (\#Ge 22:18). This expression, ``he said in his mind,'' is not added
without an object or gratuitously, for words which are articulated in the tongue and the mouth incur
guilt, and become liable to punishment, but those which are restrained within the mind are not liable
to punishment, because the mind without any intention on its part is led away by irregularities, all
kinds of passions being introduced from different quarters, which it for a while resists, being
indignant at them, and wishing to keep aloof from their representations. But perhaps we should not
say that he hesitated, but rather that he was struck by wonderment at the amazing nature of the
gift, and so said, ``Behold my body is advanced in years, and has passed the age of generation;
nevertheless all things are possible to God, so that he may transmute old age into youth, and lead
those who have no seed nor fruit to fertility and generation: and if a man who is a hundred years
and a woman who is ninety years old become parents, all commonplace occurrences and all
regularity of nature will be done away, and it will be clearly seen that it is only the power and the
grace of God.''

But what virtue the number one hundred has must now be explained.

In the first place, a hundred is the power of the number ten.

In the second place, the number ten thousand is the power of this number a hundred, and ten
thousand is the brother of the unit, for as one times one is one, so ten thousand times one is ten
thousand.

In the third place, every part of the number a hundred is honourable.

In the fourth place, this number consists of thirty-six and sixty-four, which is a cube, and at the
same time a triangle.

In the fifth place, it is composed of all these separate odd numbers: one, three, five, seven, nine,
eleven, thirteen, fifteen, seventeen, nineteen, which added together make a hundred.

In the sixth place, it is composed of these four numbers: one and its double, and four and its
double; as one, two, four, eight, which make fifteen, and of these four numbers also added
together, one, four, fifteen, sixty-four, which make eightyfive. And the principle of doubling
pervades all these numbers, containing that principle which is by fours and by fives: and the
principle of four times and twice pervades them all.

In the seventh place, it is composed of five numbers taken simply, one, two, three, four, which
make ten; and of five triangular numbers, one, three, six ten, which make twenty; and of five
quadrangular numbers, one, four, nine, sixteen, which make thirty; and of five quinquangular
numbers, one, five, twelve, twenty-two, which make forty: and all these added together make a
hundred.

In the eighth place, it is composed of four cubes taken simply beginning with the unit, for after
giving one, two, three, four, their cubes one, eight, twenty-seven, sixty-four, make a hundred.

In the ninth place, it is divided into forty and sixty, each of which is a very natural number; and in
accordance with the first order of decimals up to ten thousand in a quinquangular figure the number
a hundred holds the middle place; for instance: one, ten, a hundred, one thousand, ten thousand,
where a hundred is the middle number of one, ten, a hundred, one thousand, and ten thousand.

But we ought also not to pass over in silence the number ninety as far as it concerns the visible
characters. As it seems to me the number ninety is second only to the number a hundred,
inasmuch as the tenth part of it, that is to say ten, is taken away, since I see that in the law
two-tenths of the firstfruits were set apart, first a tenth of the whole, secondly a tenth of the
remainder, for when a tenth of the fruits of the earth, of corn, or wine, or oil, is taken, another tenth
is also taken from the remainder; therefore of these two that which is the first and principal one is
honoured with the greater share; and in the second place that which follows it, since the number a
hundred of the years of the wise man comprises both the first-fruits with which it is consecrated,
both the first and the second kind; but the number ninety of the years of the female parent,
comprehends the second and lesser first-fruits, namely, the remainder of the first, which is the
great one among the sacred numbers. This therefore may be called the first vision in the sacred
law which is familiar; and the other has a general character, for the number ninety is fertile; on
which account also it happens that the woman begins to bring forth in the ninth month; but the tenth
is the sacred and perfect number; and when the two numbers nine and ten are multiplied together
ninety is made, as being the virtue of the sacred birth, receiving a fertile generation according to
the number nine, and a holy one according to ten.

(57) Why did Abraham say to God, O may this my son Ishmael live before thee? (\#Ge 17:18).
In the first place, I do not despair, says he, O Lord, of a better generation, but I believe thy
promise: nevertheless, it would be a sufficient blessing for me for this son to live who in the
meantime is a living son, standing visibly, even though he be not so according to the legitimate
blood, but is only born of a concubine. In the second place, that blessing which he is now asking for
is an additional one; for he does not entreat for life alone for his sons, but for an especial life in
God; and we must suppose that there is nothing more perfect than the rejoicing in the presence of
God with a salutary soundness of mind, which is equal to immortality. In the third place, he by a
conjecture intimates that the divine law, when heard, ought not to be considered enough if merely
heard, but that it ought also to enter more deeply into the inward man, and to form his principal part;
for that life is worthy of being beheld by the Deity which is formed in accordance with his word.

(58) Why does the divine oracle, in the way of intimation, say to Abraham, Yes, be it so: behold
Sarah thy wife shall also bring forth a son unto thee? (\#Ge 17:19). The meaning of this
sentence is as follows: that confession and admission, says God, is on my part an admission of
thy wish, being manifestly full of unadulterated joy; and your faith is not doubtful, but without any
hesitation it has a share of modest awe and reverence; therefore that which thou hast received
before, as to be done unto thee on account of thy faith in me, shall certainly be done; for this is
what is meant by yes.

(59) Why does he say, But behold I will also listen to thee concerning Ishmael, and I will bless him,
and he shall become the father of twelve nations? (\#Ge 17:20). God says, I will grant to thee
both the first and the second blessing, that is to say, both the blessings of nature and the blessings
of instruction; by nature that which is according to the legitimate course of nature, that is Isaac,
and by instruction that which is according to Ishmael, who is not legitimate: for hearing, when
compared with sight, is like the illegitimate compared with the legitimate, and what is brought about
by instruction is not of the same class with that which owes its existence to nature; and the man
who is desirous of encyclical wisdom becomes the father of nations, for the encyclical number is a
period of twelve days and years.

(60) Why does he say, But I will set up my covenant with Isaac, whom Sarah shall bring forth about
this time in the succeeding year? (\#Ge 17:21). As in men's wills some persons are set down
as heirs, and some are entered as worthy of gifts which they are to receive from the heirs, so also
in the divine testament that man is set down as the heir who is by nature a worthy disciple of God
being adorned with all perfect virtues; but he who is introduced by learning, and is made subject to
the law of wisdom, and partakes in encyclical instruction, is not at all an heir, but only a receiver of
gifts gratuitously given. But it is said with great wisdom and propriety that his mother shall bring
forth Isaac in the succeeding year, since this birth unto life does not belong to the present time, but
to another great and holy time; and that which is divine rejoices in excessive abundance, and is by
no means like the nations of this world.

(61) Why does he say, Abraham was ninety and nine years old when he was circumcised, and
Ishmael his son was thirteen years old? (\#Ge 17:24). The number of ninety and nine years is
arranged here as approximating to the number a hundred. And it is in accordance with this number
that it is arranged that the seed of the perfect man becomes the beginning of generation, which
appears more evidently in the number a hundred; but the number thirteen is composed of the first
square numbers of four and nine, the odd and even numbers; so that the even number has for its
sides a twofold material form; and the odd number has an operative form, from all which a triple
number is made, which is the greatest and most perfect of the festival victims which the
examinations of the sacred scriptures contain. This is one reason. A second also it may be
allowed to us to mention, that the age namely of thirteen years is very near to and a partaker with
the fourteenth year, in which the motions of seed towards generation begin to have life. In order,
therefore, that no foreign seed should be sown, he arranged that the first generations should be
kept pure, figuring the instrument of generating under the figure of generation. In the third place, he
teaches that he who is about to go through the operations of matrimony ought by all means first of
all to cut away concupiscence, reproving all lascivious and effeminate persons as those who bring
together superfluous mixtures which were not for the sake of the generation of children but to
gratify incontinent desires.

(62) Why did Abraham also circumcise strangers? (\#Ge 17:27). The wise man is as useful as
the humane man, who saves and invites to himself not only his relations and neighbours, but also
strangers and men of another family, giving them a share of his own habit of patient and religious
continence; for these are the foundations of constancy, which is the object of all virtue, and the
point at which it rests.

\xchapter{Appendix 1: Concerning the World}

APPENDICES

A TREATISE CONCERNING THE WORLD--

\emph{[----Yonge's edition includes this treatise not found in Cohn-Wendland (Loeb). In a note, Yonge asserts that it is virtually
identical to the Loeb treatise, On the Eternity of the World (which Yonge titled, On the Incorruptibility of the World).]}

I. There is no existing thing equal in honour to God, but he is the one Ruler, and Governor, and
King, to whom alone it is lawful to govern and regulate everything; for the verse-

``A multitude of masters is not good,

Let there one sovereign be, one king of All,''\footnote{hom. Il. 2.204.}

is not more appropriate to be said with respect to cities and men than to the world and God, for it
follows inevitably that there must be one Creator and Master of one world; and this position having
been laid down and conceded as a preliminary, it is only consistent with sense to connect with it
what follows from it of necessity. Let us now, therefore, consider what inferences these are. God
being one being, has two supreme powers of the greatest importance. By means of these powers
the incorporeal world, appreciable only by the intellect, was put together, which is the archetypal
model of this world which is visible to us, being formed in such a manner as to be perceptible to our
invisible conceptions just as the other is to our eyes. Therefore some persons, marvelling at the
nature of both these worlds, have not only worshipped them in their entirety as gods, but have also
deified the most beautiful parts of them, I mean for instance the sun, and the moon, and the whole
heaven, which, without any fear or reverence, they called gods. And Moses, perceiving the ideas
which they entertained, says, ``O Lord, King of all Gods,''\footnote{\#de 10:17.} in order to point out the
great superiority of the Ruler to his subjects. And the original founder of the Jewish nation was a
Chaldaean by birth, being the son of a father who was much devoted to the study of astronomy,
and being among people who were great studiers of mathematical science, who think the stars,
and the whole heaven, and the whole world gods; and they say that both good and evil result from
their speculations and belief, since they do not believe in anything as a cause which is apart from
those things which are visible to the outward senses. But what can be worse than this, or more
calculated to display the want of true nobility existing in the soul, than the notion of causes in
general being secondary and created causes, combined with an ignorance of the one first cause,
the uncreated God, the Creator of the universe, who for these and innumerable other reasons is
most excellent, reasons which because of their magnitude human intellect is unable to apprehend?
but this founder of the Jewish nation having conceived an idea of him in his mind, and looking upon
him as the true God, forsook his native country and his family, and his father's house, knowing that
if he remained, the deceits of the polytheistic doctrine also remaining in his soul would render his
intellect incapable of discovering the nature of the one God, who is alone everlasting, and the
father of everything else, whether appreciable only by the intellect or perceptible to the outward
senses; but if he departed and emigrated, then he saw that deceit would also depart from his mind,
which would then change its erroneous opinions into truth. And at the same time the oracular
commands of God, which had been given to him, did further excite the desire which he felt to
become acquainted with the living God. And he went forth like a man under the immediate guidance
of others, with the most unhesitating promptness, to search after the knowledge of the one God;
and he did not relax in his search till he had arrived at a more accurate and correct perception, not
indeed of his essences (for that is impossible), but of his existence and of his providence; on
which account he is the first man of whom it is said that he believed in God, since he was the first
who had an accurate and positive notion of him, believing that there is one supreme cause of all
things, which by his providence takes care of the world and of all things that are therein. Since the
Creator of the world bringing an essence previously without any order and in complete confusion,
into distinct order and regularity, began to arrange and adorn the earth and the water, and
established them in the middle of the world, and the trees, and air, and fire he drew up from the
middle to the higher regions, and he fixed the regions of the aether all around, placing it as a
boundary to and a preserver of the things which were inside, from which also it derived its name of
Heaven.\footnote{horos is the Greek word for boundary, from which Philo thinks that ouranos, ``heaven,'' is derived.} And
these things, then, were the perfect seeds of the whole universe, but the great and all productive
tree raised from this seed is this world, of which the aforesaid branches are the offshoots.

II. Where, then, God placed the roots, and what foundation it has upon which it is so firmly fixed like
a statue, we must now consider. It is not natural that any body which is left behind should wander
out of its limits, since God has made and arranged in its proper place, the materials of the whole
universe. For it was fitting that the greatest of all works, being also the most perfect, should be
created by the greatest of all workmen. And it would not have been completely perfect if it had not
been completed in perfect parts. So that if this world consists of every kind of material, nothing
being beyond, and not even the most insignificant thing being omitted, it follows of necessity that
whatever is outside the world must either be vacuum or nothing at all. If it be a vacuum, then how
can it be found to balance the world, which is full and closely packed, and the heaviest of all things,
when there is nothing solid to support it? from which consideration it would appear to resemble a
vision. Since the mind is always looking for a corporeal basis, it is natural to suppose that one
whole should have such a thing if it happens to be put in motion, and the world above all things,
inasmuch as it is the greatest of bodies, and as it embraces in its bosom a multitude of other
bodies as its own appropriate parts. Therefore, if any one wishes to escape the perplexities which
arise in treating of doubtful matters, let him speak his mind freely, and affirm that there is no
material so strong as to be able to support the weight of the world. But the eternal law of the
everlasting God is the strong and lasting support of the universe. This law being extended from the
centre of the world to its furthest extremities, and again back from its extremities to the centre,
moves on in the unwearied irresistible course of nature, uniting and binding together all the parts of
the universe. For the Father who established it made it to be the indissoluble bond of the universe.
Therefore we are naturally led to conclude that the whole earth will not be dissolved by water, which
its bosoms contain; nor again will fire be extinguished by the air, nor again will the air be burnt up
and consumed by fire, since the divine law has placed itself as a boundary to keep all these
elements distinct from one another. As yet the allproductive plant was not rooted, and had not the
power which was to be derived from being rooted. But of the subordinate, particular, and less
important plants, some were moveable in such a way as easily to change their places, and some,
without being liable to any change of places, were made as if they were to stand for ever in the
same position. Those therefore which are exposed to a motion which involves a change of place,
which we call animals, were added to the most entire and perfect parts of the universe. The earth
receiving the terrestrial animals, the water the aquatic animals; the air those creatures which fly;
and the stars being assigned to the heaven.

III. But the Creator created two different kinds both in the earth, and in the water, and in the air. In
the air he placed those animals which fly, and other powers also which cannot by any means, or on
any occasion, be comprehended by the outward senses. Thus the company of incorporeal souls is
arranged in regular order according to their nature. For it is said that some of them are separated
off and assigned to mortal bodies, and that, at certain definite and predetermined periods, they
again depart from them. But that others of a more divine nature are utterly regardless of any
situation in earth, but are raised to a great height, and placed in the aether itself, being of the
purest possible character, which those among the Greeks who have studied philosophy call heroes
and daemones, and which Moses, giving them a most felicitous appellation, calls angels, acting as
they do the part of ambassadors and messengers, announcing to the subjects all kinds of
blessings from their rulers, and acting as servants to the king to whom they are subject; and they,
descending into the body as into a river, at one time are carried away by the violence of a most
irresistible current and swallowed up, and at other times, being wholly unable to resist the powers of
destruction, at first, indeed, raise their heads above the flood, and afterwards sink down again to
the place from whence they have started. These, then, are the souls of those who devote
themselves to the vigorous study of philosophy respecting divine subjects, from the beginning to
the end of their existence studying things which may concern them after the life has left the body,
that thus they may enjoy an incorporeal and endless life in the presence of the uncreated and
immortal God. But those souls of other men which I have spoken of as being overwhelmed, being
such as have disregarded wisdom, giving themselves up to uncertain circumstances, such as
depend wholly on chance, of which none have any reference to the soul or to the intellect, but all to
the body, which is but a corpse to which we are joined, or to other things even more inanimate and
insensible than that; I mean such things as glory, and riches, and power, and honour, and all such
other things as through the deceitfulness of false opinions are looked upon as real and living
objects by people who do not see what is really beautiful. Therefore, if you look upon souls, and
daemones, and angels, as things differing indeed in name, but as meaning in reality one and the
same thing, you will thus get rid of the heaviest of all evils, superstition. For as people in general
speak of good daemones and bad daemones, in the same manner also do they speak of good and
bad souls; and so they speak of some angels as being by their title worthy ambassadors from men
to God, and from God to men, being sacred and inviolable guardians on account of their blameless
and most excellent service which they have allotted to them. And, again, if you look upon others as
unholy and unworthy of any such appellation, you will not err. And the Psalmist himself is a witness
in favour of what I have here asserted, where he speaks as follows: ``He sent among them the
anger of his wrath, by the operation of evil Angels.''\footnote{psalm 77:49.} Again, all animals that swim and
zoophytes are allotted to the water, and all terrestrial animals and plants to the land. And the plants
he placed with their heads downwards, fixing their heads in the deepest parts of the earth; but the
heads of the irrational animals he dragged up from the earth and placed upon a lofty neck, placing
the fore-feet beneath them as a kind of pedestal. But man has had a separate formation of a
higher character; for in the case of other animals, God has placed their eyes in the side of their
heads and bent them down to the ground, on which account they are all inclined downwards to the
earth. But the eyes of man, on the other hand, he has raised up, that he might behold the heavens,
being not a terrestrial but a celestial plant, as the old saying Is.\footnote{this is in accordance with the idea of
Ovid, who says (as may be Translated)—``and while all other creatures from their birth, / With downcast eyes gaze on their
kindred earth, / Man walks erect, and proudly scans the heaven / From which he sprung, to which his hopes are given.''}
But the other class, who affirm that our intellect is a portion of ethereal nature, connect man in a
relationship with the air. Accordingly, the great Moses has not spoken of the rational soul as it
resembled in its species any created thing, but he has called it the image of the divine and invisible
God, looking upon it to be a glorious and carefully wrought image, the seal of God, the character of
which is the everlasting Word; for, says he, ``God breathed into his face the breath of
Life.''\footnote{\#ge 2:7.} So that it follows inevitably that he who received it must be made in the image
and likeness of him who sent it. On which account he also says that man was created in the image
of God, and not in the likeness of any created thing.

IV. But, taking up our discourse again at the beginning for the sake of clearness, let us say that of
bodies some have put on habit, and others nature, and others soul, and others a rational soul.
Therefore those stocks and stones which are torn from any intimate connection, have made for
themselves that strongest of all forms, namely habit, and that is a breath returning constantly to
itself; for it begins at the centre and extends to the furthest extremities, and when it has touched
the outermost circumference it turns back again until it arrives at the same place from which it
originally started. This it he continued course of habit over which it runs and returns. And he has
allotted a nature of their own to plants, having combined it of many powers, especially the nutritious
and the generative power. And the Creator has made the soul different from nature in three
particulars-the outward sense, and fancy, and impulse. Now plants have no participation in any of
these things, but every living animal has a share in all of them. Therefore the outward sense, as its
very name in my opinion shows, is a certain imposition which represents to the intellect the things
which have appeared to it. And it represents to the fancy a sort of outline in the soul, being, as it
were, a kind of representation of light; for those things which each of the outward senses has
introduced, like a ring as it were or a seal, it impresses on them its own character, or else it
preserves the impression which has been made until the rival of memory, forgetfulness, having
softened the impression, at least makes it very dim, or else entirely effaces it; and what has
appeared to have been impressed upon it disposes the soul at one time as if it belonged to it, and
at another time as if it belonged to some other: and this feeling is called impulse, which those who
have attempted to give accurate definitions have called the primary motion of the soul.

V. In such important particulars are animals superior to plants. Let us now therefore see in what
man is superior to the other animals. He now has received as an especial and pre-eminent
honour, the gift of intellect, by which he is accustomed to comprehend the natures of all things,
whether they be bodies or things; for as the predominant part in the body is the sight, and as the
nature of light is the most important part of the universe, so in the same manner the most important
and influential of all the parts in us is the mind; for this is the light of the soul, being irradiated and
enlightened by its own beams, by which that dense and profound darkness, ignorance of facts
which was shed around it, is dissipated. And this portion of the soul is not composed of the same
elements as those of which the other parts are made, but it has a pure and more excellent
essence, from which the divine natures were made; on which account the intellect alone, of all the
parts within us, appears very reasonably and naturally to be imperishable, for that is the only
portion which the Father who generated it has thought worthy of freedom, and loosing the bands of
necessity, he has allowed it to roam at large without restraint, having endowed it with a share of his
own most glorious and becoming attribute, freewill, the highest present which it was able to receive.
For the other animals in whose souls there does not exist that intellect which is thus especially
appropriated to freedom, have been given up to them to submit to their yoke and to receive their
bridle in their mouths, so as to serve them as servants obey their masters. But man having a
spontaneous will, subject to no promptings but those of his own nature, and exerting his energies in
accordance with his own deliberate purpose, is very properly subject to blame for whatever unjust
actions he commits from deliberate intention, and to praise for all the good deeds which he
intentionally performs; for as he has received from God a power of voluntary motion, and as he is
in this respect like unto God himself, being delivered from all subservience to that most severe and
grievous mistress, necessity, he very properly is open to accusation when he does not pay
worship to that being who has thus delivered him. Therefore he will most justly in such a case
suffer the punishment which has been inexorably pronounced against ungrateful people who do not
deserve freedom. On which account also, the body being raised up towards the purest portion of
the universe, the heaven, raises its eyes upwards, that so by an observation of what is visible, it
may arrive at an adequate comprehension of what is invisible. Since, therefore, it would be
impossible to behold the attraction of the intellect towards the living God, excepting as far as those
who are attracted towards him can themselves perceive it, for each man in an individual and
especial degree knows what happens to himself, he has made a visible image of the invisible eye,
namely, the eyes of the body which are thus able to look towards the sky. For when the eyes, which
are made of perishable materials, have gone to such heights as even to soar upwards to the
heaven which is removed to such an immense distance from the regions of the earth, and to touch
its borders, to how great a distance must we not suppose that the eyes of the soul can reach?
which, being excited by a vehement desire to see the one Being clearly and distinctly, stretch
forward not only to the furthest extremity of the sky, but, leaving beneath them the boundaries of
the universal world, hasten onwards to the uncreated.

VI. Having now, therefore, gone through the whole question of the more important plants in the
world, let us see in what manner also the all-wise God has fashioned the trees which exist in man,
that lesser world. Therefore immediately having taken our body as a region of fertile soil, he has
made in it the outward senses as so many channels; and then he has carefully trained each of
those outward senses as a plant susceptible of cultivation and of the greatest use, implanting the
sense of hearing in the ears, and that of seeing in the eyes, and that of smell in the nostrils, and all
the other senses in the places akin to and appropriate to them. And I have a witness in favour of
this my argument in that god-like man who speaks thus in the Psalms: ``He that planted the ear,
shall he not hear? and he that fashioned the eye, must not he See?''\footnote{\#ps 93:9.} Moreover, those
other faculties which reside apart from the main body, being situated in the legs and hands and the
other parts of the body, whether within or without, all these faculties, I say, are noble and excellent
offshoots. And the more excellent and more perfect parts he very appropriately stationed near the
dominant portion of the whole, as being in the centre, and able preeminently to bring forth fruit, as
being the lord of the whole. And these faculties are perception and comprehension, and felicity of
conjecture, and study, and recollection, and habit, and disposition, and every variety of art, and
certainty of knowledge, and an ever-mindful apprehension of the speculations of every kind of
virtue. Now, no one can properly or sufficiently cultivate any one of these within, but the one
uncreated Maker of them, and who has not merely created them, but who also makes all these
plants to correspond to everything which takes place; he alone can manage them and perfect them
as they should be perfected.

VII. And the way in which Paradise was planted is in strict conformity with what has been here said;
for we read that ``God planted a Paradise in Eden, towards the east, and there he placed the man
whom he had Made.''\footnote{\#ge 2:8.} Now, to think that this means that God planted vines and olives,
and trees of apples and pomegranates, and things of that kind, is great and incurable folly. But in
order that no one might imagine that the Creator had need of anything that he had created, Moses
has made a most important declaration when he says, ``The Lord, the King of ages, for ever and
Ever.''\footnote{\#ex 15:18.} Accordingly, God is both the Father, and the Creator, and the Governor, in
reality and truth, of all the things that are in heaven and in the whole world. And, indeed, the future is
concealed and separated from the present moment at one time by a brief, and at another time by a
long interval. But God is also the Creator of time, for he is the Father of that which is the father of
time; and the father of time is the world, which proves that its own birth is the motion of time. But
nothing is future to God, because he is in possession of and the author of the boundaries of time;
for it is not time, but rather the archetype and model of time. But in eternity nothing is passed,
nothing is future, but everything is at the present moment.

VIII. Having now, then, discussed these matters at sufficient length, we must proceed to investigate
its imperishableness. Now, there are three opinions in vogue among the philosophers on this
subject: some affirming it is everlasting, and uncreated, and free from all liability to destruction;
others, on the contrary, that it is created and perishable. There is also a sect which has adopted
some portions of the doctrine of each of the beforementioned parties, taking from the latter sect
the doctrine that it is created, and from the former the idea that it is imperishable; and thus they
have left a mixed opinion, looking upon it as at the same time created and yet imperishable.
Therefore Democritus, and Epicurus, and the chief body of the philosophers of the Stoic school,
believe the generation and also the destructibility of the world; but they do not all do so in the same
manner. For some give a sketch of many worlds, the creation of which they attribute to the
concourse and conflicting combination of atoms, and their destruction they attribute to the
repercussion and shattering of what has been thus formed. But the Stoics affirm that there is one
world, and that God is the cause of its creation, but that God is not the cause of its destruction; but
that the power which is contained in existing things, in the long periods of never-ending time,
attracts everything to itself, from which again a regeneration of the world is caused by the
prudence of the Creator. But Aristotle pronounced the world to be both uncreated and
imperishable, and he affirmed that those who maintained a contrary doctrine were guilty of terrible
impiety, as they considered that so great a work of God was in no respect superior to things made
by the hand of men. And they say too that it has been proved to be both uncreated and
imperishable by Plato in his Timaeus. But some persons interpret Plato's words sophistically, and
think that he affirms that the world was created, not inasmuch as it has had a beginning of creation,
but inasmuch as if it had been created it could not possibly have existed in any other manner than
that in which it actually does exist as has been described, or else because it is in its creation and
change that the parts are seen. But the forementioned opinion is better and truer, not only because
throughout the whole treatise he affirms that the Creator of the gods is also the father and creator
and maker of everything, and that the world is a most beautiful work of his and his offspring, being
an imitation visible to the outward senses of an archetypal model appreciable only by the intellect,
comprehending in itself as many objects of the outward senses as the model does objects of the
intellect, since it is a most perfect impression of a most perfect model, and is addressed to the
outward sense as the other is to the intellect. But also because Aristotle bears witness to this fact
in the case of Plato, who, from his great reverence for philosophy, would never have spoken
falsely. But some persons think that the father of the Platonic theory was the poet Hesiod, as they
conceive that the world is spoken of by him as created and indestructible; as created, when he
says, -

``First did Chaos rule;

Then the broad-chested earth was brought to light,

Foundation firm and lasting for whatever

Exists among Mankind;''\footnote{hesiod, Theogon. 116.}

and as indestructible, because he has given no hint of its dissolution or destruction. Now Chaos
was conceived by Aristotle to be a place, because it is absolutely necessary that a place to
receive them must be in existence before bodies. But some of the Stoics think that it is water,
imagining that its name has been derived from Effusion.\footnote{chysis, as if chaos were derived from cheo—, ``to
pour.''} But however that may be, it is exceedingly plain that the world is spoken of by Hesiod as
having been created: and a very long time before him Moses, the lawgiver of the Jews, had said in
his sacred volumes that the world was both created and indestructible, and the number of the
books is five. The first of which he entitled Genesis, in which he begins in the following manner: ``In
the beginning God created the heaven and the earth; and the earth was invisible and without form.''

IX. But we must place those arguments first which make out the world to be uncreated and
indestructible, because of our respect for that which is visible, employing an appropriate
commencement. To all things which are liable to destruction there are two causes of that
destruction, one being internal and the other external; therefore you may find iron, and brass, and
all other substances of that kind destroyed by themselves when rust, like a creeping disease,
overruns and devours them; and by external causes when, if a house or a city is burnt, they also
are consumed in the conflagration, being melted by the violent impetuosity of the fire. A similar end
also befalls animals, partly when they are sick of diseases arising internally, and partly when they
are destroyed by external causes, being sacrificed, or stoned, or burnt, or when they endure an
unclean death by hanging. And if the world also is destroyed, then it must of necessity be so either
by some external cause, or else by some one of the powers which exist within itself; and both
these alternatives are impossible, for there is nothing whatever outside of the world, since all things
are brought together in order to make it complete and full, for it is in this way that it will be one, and
whole, and free from old age; it will be one, because if anything were left outside of it, then another
world might be created resembling that which exists now; and whole, because the whole of its
essence is expended on itself; and exempt from old age and from all disease, since those bodies
which are liable to be destroyed by disease or old age are violently overthrown by external causes,
such as heat, and cold, and other contrary qualities, no power of which is able to escape so as to
surround and attack the world, all those being entirely enclosed within, without any part whatever
being separated from the rest. But if indeed there is any external thing it must by all means be a
vacuum, or else a nature absolutely impossible, which it would be impossible should either suffer or
do anything. And again, it will also not be dissolved by any cause existing within itself; first of all
because, if it were, then the part would be greater and more powerful than the whole, which is the
greatest possible absurdity, for the world, enjoying an unsurpassable power, influences all its parts,
and is not itself influenced or moved by any one of them; in the second place because, since there
are two causes of corruption, the one being internal and the other external, those things which are
competent to admit the one must also by all means be liable to the other; and a proof of this may be
found in oxen, and horses, and men, and other animals of similar kinds, because it is their nature
to be destroyed by the sword, or to be liable to die by disease.

X. Since, therefore, the arrangement of the world is such as I have endeavoured to describe it, so
that there is no part whatever left out, so as for any force to be applied, it has now been proved
that the world will not be destroyed by any external thing, because in fact nothing whatever external
has been left at all; nor will it be destroyed by anything in itself on account of the proof which has
already been considered and stated, according to which that which was obnoxious to the power of
one of those causes was also naturally susceptible of the influence of the other. And there are
testimonies also in the Timaeus to the fact of the world being exempt from disease and not liable to
destruction, such as these: ``Accordingly, of the four elements the constitution of the world receives
each in all its integrity; for he who compounded it made it to consist of the whole of fire, and the
whole of water, and the whole of air, and the whole of earth, not leaving any portion or any power of
any one of them outside, from the following intentions:-in the first place, in order that the whole
might be as far as possible a perfect animal made up of perfect parts. And besides all these things,
he ordained that it should be one, inasmuch as there is nothing left out of which another similar
world could be composed. Moreover, he willed that it should be exempt from old age, and free from
all disease, considering that those things which in the body are hot or cold, or which have mighty
powers, if standing all around and falling upon it unseasonably, would be likely to dissolve it, and, by
introducing diseases and old age, cause it to decay and perish. For this cause, and because of
this reason, God made the whole universe to consist of entire and perfect elements, and exempt
from old age and free from disease.''

XI. Let this be taken as a testimony delivered by Plato to the imperishable nature of the world. Its
uncreated character follows from the truth of natural philosophy; for dissolution must of necessity
attend everything which is born, and incorruptibility must inevitably belong to everything which is
unborn; since the poet who wrote the following iambic verse,

``All that is born must surely Die,''\footnote{timaeus, p. 32.}

appears to have spoken very correctly when he asserted this connection of destructibility with
birth. The argument may be stated in a different way as follows. All compound things which are
destroyed are dissolved into the elements of which they were compounded; accordingly,
dissolution is nothing else but a return of everything to its original constituent parts; just as, on the
contrary, composition is that which compels the things combined to come together in a manner
contrary to their nature; and indeed, this appears to be the most exact truth; for men are
composed of the four elements which together make up the whole of the universe, the heaven, the
earth, the air, and fire, borrowing a few parts of each in a manner at first sight hardly consistent
with nature. But the things which are thus combined together are necessarily deprived of a motion
in accordance with nature; for instance, warmth is deprived of its upward motion, and coldness of
its downward tendency, the earthy and somewhat weighty substance being lightened and assuming
the higher place, which the most earth-like of our own parts, the head, has obtained in us. But of all
bonds, that is the worst which is forged by violence, and which, being violent, is also short-lived; for
it is speedily broken by those who are bound in it, since they become restive from their desire for a
motion in accordance with nature, to which they hasten; for as the tragic poet says, -

``And for things sprung from earth, they must

Return unto their parent dust,

While those from heavenly seed which rise

Are borne uplifted to the skies.

Nought that has once existed dies,

Though often what has been combined

Before, we separated find,

Invested with another Form.''\footnote{a fragment from the Chrysippus of Euripides.}

And this law and ordinance is established with reference to everything which is destroyed, that
wherever composite things are existing in combination they are thrown into disorder instead of into
the order in accordance with nature, which they previously enjoyed, and they are removed to
situations opposite to those in which they were previously placed, so that they seem in a manner to
be sojourners; and when they are dissolved again, then they return to the appropriate parts allotted
to them by nature. But since the world has no participation in that irregularity which exists in the
things which I have just been mentioning, let us stop awhile and consider this point. If the world
were liable to corruption and destruction, it follows of necessity that all its parts would at present be
arranged in a position not in accordance with nature: but it is impious even to imagine such a thing
as this; for all the parts of the world have received the most excellent position possible, and an
arrangement of the purest symmetry and harmony; so that each individual part, being content with
its place as a native country to it, does not seek any change for the better. On this account it is that
the most central position of all has been assigned to the earth, to which all things belonging to it
adhere, and to which they descend again even if you throw them into the air: and this is a proof that
their place is in accordance with nature; for wherever anything is borne without any violence, and
where it then remains firm and stationary, that is clearly its natural place. And then, in the second
place, water was poured over the earth, and air and fire have gone from the central to the upper
part, air having received for its portion the region which is on the borders between air and fire, and
fire having received the highest place of all: on which account, if you light a torch and press it down
towards the ground, nevertheless the flame will still turn in a contrary direction, and lightening itself
in accordance with the natural motion of fire, will rise upwards: if, then, motion contrary to nature is
the cause of corruptibility and destruction in the case of other animals, but if in the case of the
world every one of its parts is arranged in complete accordance with nature, having had
appropriate positions allotted to each of them, then surely the world must most justly be
pronounced incorruptible and imperishable.

XII. Moreover this point is manifest to every one, that every nature is desirous to keep and
preserve, and if it were possible to make immortal, everything of which it is the nature; the nature of
trees, for instance, desires to preserve trees, and the nature of animals desires to preserve each
individual animal. But particular nature is of necessity unable to conduct what it belongs to to
eternity; for want, or heat, or cold, or innumerable other ordinary circumstances, when they affect
particular things, shake them and dissolve the bond which previously held them together, and at last
break them to pieces; but if nothing resembling any of these things were lying in wait outside, then
in that case nature itself, as far as it is possible, would preserve everything both great and small
free from old age. It follows therefore of necessity, that the nature of the world must desire the
durability of the universe; for it is not worse than particular natures, so that it should run away and
desert its proper duties, and attempt to produce disease instead of health, and corruption and
destruction instead of complete safety, since,

``High over all she lifts her beauteous face,

And towers above her nymphs with heavenly grace,

Fair as they all Appear.''\footnote{homer, Odyssey 6.107, where the lines quoted are applied to Latona among her
nymphs.}

But if this be true, then the world cannot be capable of destruction. Why so? Because the nature
which holds it together is itself invincible by reason of its exceeding strength and power, by which it
gets the mastery over every thing else which might be likely to injure it; wherefore Plato has well
Said:\footnote{timaeus, p. 33.} ``For nothing ever departed from it, nor did anything ever come to it from any
quarter; for that was not possible; for there was nothing in existence which could come; for since it
supplies itself with nutriment out of its own consumption, it also does everything and suffers
everything in itself and by itself, and is compounded with the most consummate art. For he who
created it thought that it would be better if wholly selfsufficient, than if in continual need of
accessories from other quarters.''

XIII. However, this argument also is a most demonstrative one, on which I know that vast numbers
of philosophers pride themselves as one most accurately worked out, and altogether irresistible;
for they inquire what reason there is for God's destroying the world. For if he destroys it at all he
must do so either with the intention of never making a world again, or with the object of creating a
second fresh one; now the former idea is inconsistent with the character of God; for it is proper to
change disorder into order, and not order into disorder: in the second place, it is so because it
would give rise to repentance, which is an affliction and a disease of the soul. For he ought either
never to have created a world at all, or else, if he judged that it was a fitting employment for him, he
ought to have been pleased with it after it was made. But the second reason deserves no
superficial examination; for if he were intending to make another world instead of that which exists
at present, then of necessity this second world that would be made, in that case, would be either
worse than, or similar to, or better than the first; everyone of which ideas is inadmissible; for if the
new world is to be worse than the former, then the maker must be also worse: but all the works of
God are without blemish, beyond all reproach and wholly faultless, inasmuch as they are wrought
with the most consummate skill and knowledge; for as the proverb says; -

``For e'en a woman's wisdom's not so coarse

As to despise the good and choose the worse.''

But it is consistent with the character of, and becoming to God to give form to what is shapeless,
and to invest what is most ugly with admirable beauty. Again, if the new world is to be exactly like
the old one, then the maker is only wasting his labour, and differs in no respect from infant children
who, very often while playing on the sea shore raise up little mounds of sand, and then pull them
down again with their hands and destroy them; for it would have been much better than making
another world exactly like the former, neither to take anything from, nor to add anything to, nor to
change either for the better or for the worse, what existed originally, but to let it remain just as it
was. If, on the other hand, he is about to make a world better than the former one, then the maker
too must be better than the maker of the former world, so that when he made the former world he
was inferior both in his skill and in his intellect, which is impious even to imagine, for God is at all
times equal and similar to himself, being neither capable of any relaxation which can make him
worse, nor of any extension which can make him better. Men, indeed, do admit of such inequalities
in either direction, being naturally liable to alter either for the better or for the worse, and continually
admitting of increase, and advance, and improvement, and everything contrary to these states;
and besides this, the works of us who are but mortal men may very appropriately be perishable, but
the works of the immortal must in all consistency and reason be likewise imperishable, for it is
natural that what is made should resemble the nature of the maker.

XIV. But Boethus adduces the most convincing arguments, which we shall proceed to mention
immediately; for if, says he, the world was created and is liable to destruction, then something will
be made out of nothing, which appears to be most absurd even to the Stoics. Why so? Because it
is not possible to discover any cause of destruction either within or without, which will destroy the
world. For on the outside there is nothing except perhaps a vacuum, inasmuch as all the elements
in their integrity are collected and contained within it, and within there is no imperfection so great as
to be the cause of dissolution to so great a thing. Again, if it is destroyed without any cause, then it
is plain that from something which has no existence will arise the engendering of destruction, which
is an idea quite inadmissible by reason; and, indeed, they say that there are altogether three
generic manners of corruption, one which arises from division, another which proceeds from a
destruction of the distinctive quality which holds the thing together, and the third from confusion;
therefore the things which consist of a union of separate members, such as flocks of goats, herds
of oxen, choruses, armies; or, again, bodies which are compounded of limbs joined together, are
dissolved by disjunction and separation. But wax, when stamped with a new impression, or
softened before being remodelled so as to present a new and different appearance, is corrupted
by a destruction of the distinctive quality which previously held it together. Other things are
corrupted by confusion, as the medicine which the physicians call tetrapharmacon, for the powers
of the drugs brought together and combined were destroyed in such a manner as to produce one
perfect medicine of especial virtue. By which, then, of these modes of corruption is it becoming to
say that the world is destroyed? By that which is caused by separation? No, for it is not
compounded of separate members so that its different parts can be dispersed, nor of portions
joined together so that they can be dissolved; nor is it united together in a similar manner to our
own bodies, for they have the seeds of decay in themselves, and they are subject to influence of a
great variety of things by which they are at times injured; but the power of the world is invincible,
since by its great superiority to other things it has dominion over everything. Is it then destroyed by
a complete destruction of its distinctive qualities? This again is impossible, for there remains, as
the adversaries affirm, a quality of arrangement which by the process of conflagration is only
diminished to a lesser substance ... Is it destroyed then by confusion? Away with such an idea, for
in that case it would be necessary to confess that the corruption of a body can be reduced to a
state of non-existence. Why so? Because if each of the particular elements were destroyed
separately, it would be possible for it to become changed into another; but if they are altogether
destroyed at one and the same moment by confusion, then it would be necessary to imagine what
is absolutely impossible.

XV. Is it not however worth while to examine this question, in what manner there can be a
regeneration of all those things which have been destroyed by fire, and resolved into fire? for when
their substance has been wholly destroyed by the fire, it follows of necessity that the fire itself must
also be extinguished as no longer having any nourishment. Therefore, as long as it remained the
seminal principle of arrangement was likewise preserved, but when it is destroyed that principle is
destroyed with it. But it would be impious, and an impiety of double dye, not only to attribute
destruction to the world, but also to take away the possibility of its regeneration; as if God delighted
in disorder, and irregularity, and all kinds of evil things. But we must examine this question more
accurately, in the following manner. There are three species in fire; the coal, and the flame, and the
light. Now coal is the fire in its earthy substance, which, like a sort of spiritual habit, couches and
lies hid in a sort of cavern, pervading it all to its very extremities. And the flame is that part which,
being raised on high, is lifted up from its fuel. And the light is that which is emitted from the flame,
so as to co-operate with the eyes, in order to enable them to comprehend what is seen. And the
flame occupies the middle position between the coal and the light; for when it is extinguished it ends
in coal, and when it is kindled it excites the light, which, being deprived of its burning power, blazes.
If therefore, we affirm that the world is dissolved by conflagration, it would not be coal, because, in
that case there will be a great deal of the earthy substance left behind, in which also fire must
necessarily be contained. But we must agree, that none of the other bodies subsist any longer, but
that earth, and water, and air, are all dissolved into unmixed fire. Nor, again, would it become
flame; for that can only exist in connexion with nourishment; and, if nothing is left behind, being
deprived of all nourishment it will immediately be extinguished. It follows from all this, that it cannot
become light either; for light by itself has no substance at all, but flows from the things before
mentioned, coal and flame, not in a great degree from the coal, but very much from the flame; for it
is diffused over a very great space indeed. But if, as has been already proved, those things had no
existence from the conflagration of all things, then there could not be any light either. So that it is
impossible for the world to be susceptible of any regeneration, inasmuch as there is no spermatic
principle smouldering beneath; from which consideration it is plain that it is uncreated, and that it will
be for ever imperishable.

XVI. However, besides what has been here said, any one may use this argument also in
corroboration of his opinion, which will certainly convince all those who are not determined to be
obstinate beyond all bounds; of those things which in pairs are exactly contrary to one another it is
impossible that one thing should be, and that the other should not be; for since there is white it
follows as a matter of absolute necessity that there must also be black, and since there is a great
there must likewise be a little; since there is an odd there must inevitably be an even; since there is
a sweet there must be a bitter; since there is day there must be night; and so on in an infinite
number of similar cases; but if a conflagration should take place, then something would ensue
which is impossible; for then, of things in a pair, the one will happen and the other will not. Come,
now, let us consider the matter thus: if everything is resolved into fire, there is then something light,
and rare, and warm; for all these are the especial properties of fire; but there can be nothing
heavy, or cold, or thick, which are the opposites of the qualities which I have just enumerated. How
then can any one more completely overturn the idea of the universal disorder which would be
involved in such a conflagration than by showing that those things which by a law of nature must
exist together, are by this process separated from their natural conjunction? And the separation
has extended to such a degree, that those who maintain this doctrine attribute eternal durability to
the one and deny any existence at all to the other. Again, there is this assertion made by some of
those who diligently employ themselves in investigating truth which appears to me to be a
sufficiently felicitous one; if the world is destroyed it will either be destroyed by some other efficient
cause, or by God; now there is certainly nothing else whatever from which it can receive its
destruction, for there is nothing whatever which it does not surround and contain; but that which is
surrounded and confined within something else is manifestly inferior in power to that which
surrounds and confines it, by which it is therefore mastered; on the other hand, to say that it is
destroyed by God is the most impious of all possible assertions; for God is the cause not of
disorder, and irregularity, and destruction, but of order, and beautiful regularity, and life, and of
every good thing, as is confessed by all those whose opinions are based on truth.

XVII. But some of those persons who have fancied that the world is everlasting, inventing a variety
of new arguments, employ also such a system of reasoning as this to establish their point: they
affirm that there are four principal manners in which corruption is brought about, addition, taking
away, transposition, and alteration; accordingly, the number two is by the addition of the unit
corrupted so as to become the number three, and no longer remains the number two; and the
number four by the taking away of the unit is corrupted so as to become the number three; again,
by transposition the letter Z becomes the letter H when the parallel lines which were previously
horizontal (3/43/4) are placed perpendicularly (1/2 1/2), and when the line which did before pass
upwards, so as to connect the two is now made horizontal, and still extended between them so as
to join them. And by alteration the word oinos, wine, becomes oxos, vinegar. But of the manner of
corruption thus mentioned there is not one which is in the least degree whatever applicable to the
world, since otherwise what could we say? Could we affirm that anything is added to the world so
as to cause its destruction? But there is nothing whatever outside of the world which is not a
portion of it as the whole, for everything is surrounded, and contained, and mastered by it. Again,
can we say that anything is taken from the world so as to have that effect? In the first place that
which would be taken away would again be a world of smaller dimensions than the existing one, and
in the second place it is impossible that any body could be separated from the composite fabric of
the whole world so as to be completely dispersed. Again, are we to say that the constituent parts of
the world are transposed? But at all events they remain in their original positions without any
change of place, for never at any time shall the whole earth be raised up above the water, nor the
water above the air, nor the air above the fire. But those things which are by nature heavy, namely
the earth and the water, will have the middle place, the earth supporting everything like a solid
foundation, and the water being above it; and the air and the fire, which are by nature light, will have
the higher position, but not equally, for the air is the vehicle of the fire; and that which is carried by
anything is of necessity above that which carries it. Once more: we must not imagine that the world
is destroyed by alteration, for the change of any elements is equipollent, and that which is
equipollent is the cause of unvarying steadiness, and of untroubled durability, inasmuch as it
neither seeks any advantage itself, and is not subject to the inroads of other things which seek
advantages at its expense; so that this retribution and compensation of these powers is equalized
by the rules of proportion, being the produce of health and endless preservation, by all which
considerations the world is demonstrated to be eternal.

XVIII. Theophrastus, moreover, says that those men who attribute a beginning and destructibility to
the world are deceived by four particulars of the greatest importance, the inequalities of the earth,
the retreat of the sea, the dissolution of each of the parts of the universe, and the destruction of
different terrestrial animals in their kinds; and he proceeds to establish the first point thus: if the
earth had never had any beginning of its creation, then there would have been no portion of it rising
above the rest so as to be conspicuous, but all the mountains would have been level, and all the
pieces of rising ground would have been even with the plain. For as there are such vast showers
falling from heaven throughout all ages, it would be natural that of any places which were originally
raised on high some would be broken down and washed away by torrents, and others would
subside of their own accord and so become lowered, and that every place everywhere would be
smoothed; but now, as things are, the constant inequalities which exist, and the vast heights of
many mountains, reaching up even to the sky, are so many proofs that the earth is not eternal. For
otherwise, as I have said before, all the earth would long since have been rendered level from one
extremity to the other by the vast rains which would have fallen from the eternal commencement of
time; for it is the character of the nature of water, and especially of such as descends in a heavy
fall from lofty places, to push some things away by force, and to cut out hollow others places by its
continual dropping, and in this manner to operate on the hard, rugged, stony ground not less than
men digging. And again, the sea, as they affirm, is already somewhat diminished, and for proof of
this fact we can appeal to the most celebrated islands, Rhodes and Delos, for these were in
ancient times invisible, being overwhelmed by and sunk under the sea, but by lapse of time, as the
sea gradually diminished, they by slow degrees rose above it and came into sight, as the histories
which are written concerning them record. And they used to call Delos Anaphe, confirming the
account here given by both names, since when it appeared above the Waters\footnote{the Greek word is
anaphaneisa, from which Anaphe— is derived.} it became evident, \footnote{de—le—, from which De—los is derived.} having
been formerly invisible and unseen. And in addition to these arguments they adduce the facts that
many great and deep bays and gulfs of vast seas have been dried up, and have become land, and
have so turned out no insignificant addition to the adjacent country when sown and planted, and on
that soil there is still left plenty of proof of such spots having formerly been sea, in the pebbles, and
shells, and other things which are commonly washed up on the sea-shore being found in them. But
if the sea is gradually being diminished then the earth also will be diminished; and in long
revolutions of years every one of the elements will be entirely consumed and destroyed; and the
whole air will be consumed, being diminished by little and little; and all things will be absorbed and
dissolved into the one substance of fire. And for the purpose of establishing the third alternative of
this question they use the following argument: beyond all question that thing is destroyed all the
parts of which are liable to destruction; but all the parts of the world are liable to destruction,
therefore the world also is liable to destruction. But we must now proceed to consider the question
which we postponed till the present time. What sort of a part of the earth is that, that we may begin
from this, whether it is greater or less, that is not dissolved by time? Do not the very hardest and
strongest stones become hard and decayed through the weakness of their conformation (and this
conformation is a sort of course of a highly strained spirit, a bond not indissoluble, but only very
difficult to unloose), in consequence of which they are broken up and made fluid, so that they are
dissolved first of all into a thin dust, and afterwards are wholly wasted away and destroyed? Again,
if the water were never agitated by the winds, but were left immoveable for ever, would it not from
inaction and tranquillity become dead? at all events it is changed by such stagnation, and becomes
very foetid and foul-smelling, like an animal deprived of life. And so also the corruptions of the air
are plain to everyone, for it is the nature of the atmosphere to become sick and to decay, and, as
one may say, in a manner to die; since what else is it which a man, who is not aiming at selecting
plausible language, but only at truth, would call a plague except a death of the atmosphere, which
diffuses its own disease and suffering to the destruction of everything which is endowed with life?
And why need I speak at great length concerning fire? for if it is deprived of nourishment it is
immediately extinguished. If then, each of the separate parts of the world awaits utter destruction, it
is plain that the world which is compounded of these can not be itself exempt from destruction. We
must now consider with accuracy the fourth and remaining argument. Thus they argue: if the world
were eternal then the animals also would be eternal, and much more the human race, in proportion
as that is more excellent than the other animals; but, on the contrary, those who take delight in
investigating the mysteries of nature consider that man has only been created in the late ages of
the world; for it is likely, or I should rather say it is inevitably true, that the arts co-exist with man, so
as to be exactly co-eval with him, not only because methodical proceedings are appropriate to a
rational nature, but also because it is not possible to live without them; let us therefore examine the
dates of each of these, disregarding the fables invented by the tragedians about the gods; but if
man is not eternal then neither is any other animal, so that then neither are the places which
receive them, the earth, or the water, or the air; from all which considerations it is plain that this
world is liable to destruction.

XIX. But it is necessary to encounter such quibbling arguments as these, lest some persons of too
little experience should yield to and be led away by them; and we must begin our refutation of them
from the same point from which the Sophists begin their deceit. They say, ``There could no longer
be any inequalities existing on the earth, if the world were eternal.'' Why not, my most excellent
friends? For other persons will come up and say that the natures of trees are in no respect
different from mountains; but just as they at certain seasons lose their leaves, and again at certain
seasons recover their verdure again; (on which account there is admirable truth in those lines of
the poet:-

``Like leaves on trees the race of man is found,

Now green in youth, now withering on the ground;

Another race the following spring supplies;

They fall successive and successive Rise.'')\footnote{homer, Il. 6.147.}

And so in like manner some portions of the mountains are broken off, and others grow in their
stead: but after a long lapse of time the additional growth becomes conspicuous, because the
trees having a more rapid nature display their increase with great rapidity; but mountains have a
slower character, on which account it happens that the additions which take place in their case are
not perceptible by the outward senses except after a long time. And these men appear to be
ignorant of the manner in which they are produced, since if they had not been, perhaps they would
have been silent out of shame; but still there is no reason why we should not teach them; but there
is nothing new in what is now said, neither are they our words but the ancient sayings of wise men,
by whom nothing which was necessary for knowledge has been left uninvestigated; when the fiery
principle which is contained beneath, in the earth, is thrust upwards by the natural power of fire, it
proceeds to its own appropriate place; and if it meets with any respite or relaxation, though ever so
slight, it draws up with it a large portion of the earthy substance, as much as it can; and when it has
emerged from the earth it proceeds more slowly; but the earthy substance being compelled to
follow it for a long time, being at last raised to an immense height, is contracted at the top, and at
last comes to end on a sharp point imitating the general appearance of the flame of fire; for there
arises then a most violent contention between two things which are natural adversaries, the lightest
and the heaviest of things, each of them pressing onwards to reach its own place, and each
striving against the violent efforts of the other; accordingly the fire, which is drawing up the earth
with it, is compelled to sink down by its descending power; and the earth naturally inclining to the
lowest point is nevertheless to a certain degree made light, and lifted up by the upward tendencies
of fire, and so is raised on high, and being at last overpowered by the more influential power which
lightens it is thrust upwards towards the natural seat of fire, and established on high. Why then
need we wonder if the mountains are not entirely washed away by the impetuosity of the rains,
when so great a power, which keeps them together, and by which they are raised up, is very firmly
and steadfastly connected with them? For if they were released from the bond which holds them
together, it would be natural for them to be entirely dissolved and to be dispersed by the water; but
since they are bound together by this power of fire, they resist the impetuosity of the rains more
surely.

XX. These things, then, may be said by us with respect to the argument that the inequalities of the
surface of the earth are no proof of the world having been created and being liable to destruction;
but with respect to that argument which was endeavoured to be established by the diminution of the
sea, we may reasonably adduce this statement in opposition to it: ``Do not look only at the islands
which have risen up out of the sea, nor at any portions of land which, having been formerly buried
by the waters, have in subsequent times become dry land; for obstinate contention is very
unfavourable to the consideration of natural philosophy, which considers the search after truth to
be the chief object of rational desire; but look rather at the contrary effects: consider how many
districts on the main-land, not only such as were near the coast, but even such as were completely
inland, have been swallowed up by the waters; and consider how great a portion of land has
become sea and is now sailed over by innumerable ships.'' Are you ignorant of the celebrated
account which is given of that most sacred Sicilian strait, which in old times joined Sicily to the
continent of Italy? and where vast seas on each side being excited by violent storms met together,
coming from opposite directions, the land between them was overwhelmed and broken away; from
which circumstance the city built in the neighbourhood was called Rhegium, \footnote{rhe—gion, from
rho—gnymi, ``to break.''} and the result was quite different from what any one would have expected; for
the seas which had formerly been separated now flowed together and were united in one expanse;
and the land which had previously united was now separated into two portions by the strait which
intersected it, in consequence of which Sicily, which had previously formed a part of the mainland,
was now compelled to be an island. XXI. And it is said that many other cities also have
disappeared, having been swallowed up by the sea which overwhelmed them; since they speak of
three in Peloponnesus-

``Aegira and fair Bura's walls,

And Helica's lofty halls,

And many a once renowned town,

With wreck and seaweed overgrown,''

as having been formerly prosperous, but now overwhelmed by the violent influx of the sea. And the
island of Atalantes which was greater than Africa and Asia, as Plato says in the Timaeus, in one
day and night was overwhelmed beneath the sea in consequence of an extraordinary earthquake
and inundation and suddenly disappeared, becoming sea, not indeed navigable, but full of gulfs and
eddies. Therefore that imaginary and fictitious diminution of the sea has no connection with the
destruction or durability of the world; for in fact it appears to recede indeed from some parts, but to
rise higher in others; and it would have been proper rather not to look at only one of these results
but at both together, and so to form one's opinion, since in all the disputed questions which arise in
human life, a wise and honest judge will not deliver his opinion before he has heard the arguments
of the advocates on both sides.

XXII. And as for the third argument, it is convicted by itself, as being derived only from an unsound
system of questioning proceeding from the assertions originally made; for in truth it does not
necessarily follow that a thing, all the parts of which are liable to corruption, is likewise perishable
itself; but this is only inevitably true of that thing of which all the parts are perishable when taken
collectively and together in the same place and at the same time, since in the case of a person
who has the tip of his finger cut off, he is not disabled from living, but if he had the whole collection
of all his parts and limbs cut off at once, he would die immediately. Therefore in the same manner,
if all the elements of the world together were all to disappear at one and the same moment, then it
would be necessary to admit that the world was liable to corruption and destruction; but if each of
these elements separately only changes its nature so as to assimilate to that of its nature, it is then
rendered immortal rather than destroyed, according to the philosophical statement of the tragic
poet-

``Nought that has once existed dies,

Though often what has been combined

Before, we separated find,

Invested with another form.''

For it is the greatest folly imaginable to estimate the antiquity of the human race from the state of
art; for if any one were to follow the absurdity of such a system of reasoning as this, he will prove
the world to be very young indeed, and to have been made scarcely a thousand years, since all
those men whom we have heard of traditionally as the discoverers in different branches of science
do not go back to a greater number of years than that which I have mentioned. But if we must
speak of the arts as co-eval with the race of mankind, then we must speak, drawing our arguments
from natural history, and not inconsiderately or carelessly. And what is this history? The
destruction of the things on the earth, not all together, but of the greatest number of them, is
attributed to two principal causes, the indescribable violence and power of fire and water. And they
say that each of these elements attacks them in its turn, after very long periods of revolving years.
When, therefore, a conflagration seizes upon things, a stream of ethereal fire being poured down
from above is frequently diffused over them, overrunning many districts of the habitable world; and
when a deluge draws down the whole of the rainy nature of water, the regular rivers and torrents
overflowing, and not only that, but even far exceeding the ordinary measure of a common flood.
Accordingly, when the greater part of mankind is destroyed in the manners above mentioned,
besides an infinity of other ways of less power and importance, it follows of necessity that the arts
also must fail, for it cannot be possible to discuss science by itself without some one to reduce it to
method and practice. But when those common pestilences relax their fury, and when the human
race begins again to recover vigour and to flourish, descending from those who have not been
previously destroyed by the evils which pressed upon them, then the arts also begin again to exist,
not indeed as they were at first, but in thinner numbers from the diminution of the numbers of those
who practise them. I have now then set forth to the best of my ability what I have been able to learn
or to understand concerning the indestructibility of the world.

\xchapter{Appendix 2: Fragments}

FRAGMENTS

EXTRACTED FROM THE PARALLELS OF JOHN OF DAMASCUS

About the unstable and changeable condition of human affairs.

Page 326. C. If one is to tell the plain truth, man is without real power in anything, never taking a
firm hold of anything. I do not mean merely of common things, but not even of those which concern
himself; neither of health, nor of a good condition of the outward senses, nor of soundness in
respect of the other parts of his body, nor of his voice, nor of his presence of mind; for as to
wealth, or glory, or friends, or power, or all the other things which depend on fortune, who is there
who does not know how thoroughly unstable they are? So that we must of necessity confess that
the supreme power over everything belongs to one being alone, the true Lord of all existing things.

About impious men, sinners, etc.

Page 341. D. If you wish to be governed under God as your king, take care not to sin; but if you
commit sin, how can you be under the government of God as your king?

About those people who have renounced such and such a line of conduct, and then turning

back again, have adopted that very line which they had renounced.

Page 343. D. Some men, making improvement, have returned back to virtue before coming to the
end, the ancient principle of oligarchy having destroyed the principle of aristocracy lately
engendered in the soul, which having been quiet for a little while, has subsequently come up over
again with greater power than before.

Page 343. D. When a man rightly establishes himself in a virtuous life, with meditation, and
practice, and good government, and when having been known by all men as a pious man and one
who fears God, he falls into sin, that is a great fall, for he has ascended up to the height of heaven,
and fallen down into the abyss of hell.

About resurrection and judgment.

Page 349. A. It is not possible with God that a wicked man should lose his good reward for a single
good thing which he may have done among a great number of evil actions; nor, on the other hand,
that a good man should escape punishment, and not suffer it, if among many good actions he has
done wickedly in anything, for it is infallibly certain that God distributes everything according to a
just weight and balance.

Page 349. B. The mind is the witness to each individual of the things which they have planned in
secret, and conscience is an incorruptible judge, and the most unerring of all judges.

About those who are ruled.

Page 359. A. He who has learnt how to submit to be ruled, immediately learns how to rule others;
for even if a man were invested with the supreme power over all the earth and all the sea, he would
not be a true ruler unless he had also learnt and been previously taught to submit to the rule of
others.

About anarchy.

Page 359. D. Alas, how many and great evils are produced by anarchy! Famine, war, the
devastation of lands, the deprivation of money, abductions, fears of slavery, and death.

About the foolish and senseless man, etc.

Page 362. E. No wicked man is rich, not even though he should be the owner of all the mines in the
whole world; but all foolish men are poor. Every foolish man is straitened, being oppressed by
covetousness, and ambition, and a love of pleasure, and things of that sort, which do not permit the
mind to dwell at ease or to enjoy plenty of room.

Page 363. A. There is no greater evil to a man than folly, and the being deprived of the proper use
of his reasoning powers and intellect.

Page 363. A. Ignorance is the cause of disease and destruction.

About deceit affecting the management of a household.

Page 367. D. Every stratagem is not blameable, since guardians of the night appear to act
properly when they lie in wait for robbers, and generals when they form ambuscades against the
enemy, whom they cannot catch without a stratagem; and the same principle is applicable to what
are called manoeuvres, and to the artifices practised in the contests of wrestlers, for in such
cases deceit is accounted honourable.

About impossible things.

Page 370. B. It is as impossible that the love of the world can co-exist with the love of God, as for
light and darkness to co-exist at the same time with one another.

About holy men.

Page 372. E. The happy nature is that which rejoices on every occasion, and which is not
discontented with anything whatever which exists in the world, but is pleased with whatever
happens, as being good, and beautiful, and expedient.

About leisure and quiet.

Page 376. A. The wise man endeavours to secure quiet and leisure, and periods of rest from work,
that he may devote himself peacefully to the meditations on divine matters.

About evil-speaking.

Page 369. D. Foul speakers and random accusers, who seek to make a display of their art with
vain words, being slow to learn what is good, are very quick and ready at learning what is of the
opposite character.

About counsel.

Page 397. D. Everything which is not done with reason is discreditable, just as what is done with
reason is beautiful.

About old men.

Page 404. C. Old age is an unruffled harbour.

Page 404. C. Old age is the time when the vigour of the body is passed by; the period when the
passions can be checked.

About gymnasia.

Page 405. D. Continued practice makes knowledge firm, just as want of practice engenders
ignorance. And, again, practice in any matter increases experience.

Page 405. D. Study is the nurse of knowledge.

About calumny.

Page 436. D. Calumniators and men discarded from the divine grace, who are afflicted with the
same evil disposition of calumny with him, are in all respects hated and detested by God, and
removed to a distance from all happiness.

Page 436. D. What can be worse than calumny? for it seduces the ears and perplexes the minds
of those who listen to it, and it makes them brutal and always on a watch for evil, like men engaged
in hunting; but those who are well ballasted and restrained by prudent reason, hate the man who
utters calumnies more than him against whom they are uttered, reproving and seeking to check all
desire of blaming others until it be either proved by evidence or demonstrated by undeniable proof.

About justice and virtue.

Page 438. D. If any one embraces all the virtues with earnestness and sobriety, he is a king, even
though he may be in a private station.

About voluntary and involuntary sins.

Page 526. B. As to sin intentionally is unjust, so to sin unintentionally and out of ignorance is not at
once justifiable, but perhaps it is something between the two, that is between righteousness and
unrighteousness, and is of what some persons call an indifferent character, for no sin can be an
act of righteousness.

About initiation into divine mysteries.

Page 533. C. It is not lawful to speak of the sacred mysteries to the uninitiated.

About the sea.

Page 551. D. It is proper to marvel at the sea, by means of which countries requite one another for
the good things which they receive from each other, and by which they receive what they are in
need of, and export what they have a superfluity of.

About equality.

Page 556. D. To give equal things to unequal people is an action of the greatest injustice.

About physicians and medical science.

A good physician would not be inclined to apply every kind of salutary medicine at once and on the
same day to a patient, as he would know that by such a course he would be doing him more harm
than good, but he would measure out the proper opportunities, and then give saving medicines in a
seasonable manner; and he would apply different remedies at different times, and so he would
bring about the patient's restoration to health by gentle degrees.

About opportunity.

Page 563. C. Say what is right, and at the time when it is right, and you will not hear what is not
right.

Page 563. C. It is well to economise time.

About mysteries.

Page 576. D. Chatterers divulging what ought to be kept buried in silence, do in a manner from a
disease of the tongue pour forth into people's ears things which are not worthy of being heard.

About people who are in a state of pupillage.

Page 613. D. To inquire and put questions is the most useful of habits with a view to acquiring
instruction.

Page 613. D. He who hungers and thirsts after knowledge, and who is eager to learn what he does
not know, abandoning all other objects of care, is eager to become a disciple, and day and night
watches at the doors of the houses of wise men.

Page 613. D. For any one to know that he is ignorant is a piece of wisdom, just as to know that one
has done wrong is a piece of righteousness.

About reproach.

Page 630. C. Never reproach any one with misfortune, for nature is impartial, and the future is
uncertain; lest if you yourself should fall into similar misfortunes, you should be found to be
convicted and condemned by your own conscience.

About a proper constitution.

Page 657. C. It is advantageous to submit to one's betters.

About a blameable constitution.

From the fifth book of the Essays on Genesis.

Page 658. E. A shameless look, and a high head, and a continual rolling of the eyes, and a
pompous strut in walking, and a habit of blushing at nothing, however discreditable, are signs of a
most infamous soul, which stamps the obscure topics of the reproaches which belong to itself upon
the visible body.

About familiarity and habituation.

Page 681. D. A change of all kinds of circumstances at once to the opposite direction is very
harsh, especially when the existing powers are established by the length of time that they have
lasted.

About correction.

Page 683. D. It is useful to be warned by the misfortunes of others.

Page 683. D. Punishment very often warns and corrects those who do wrong; but if it fails to do so
to them, at all events it corrects the bystanders, for the punishments of others improve most
people, from fear lest they should suffer similar evils.

About associating with wicked men.

Page 692. A. Associations with wicked men are mischievous, and very often the soul against its
will receives the impression of the insane wickedness of one's associates.

About wisdom.

Page 693. E. Every wise man is a friend of God.

About haughty men.

Page 693. E. Self-conceit, as the proverb of the ancients has it, is the eradication of all
improvement, for the man who is full of self-conceit is incapable of improvement.

Self-conceit is by nature an unclean thing.

About natural things.

Page 711. C. As it is difficult to inoculate anything in a manner contrary to nature, and to introduce
anything into nature which does not belong to it, so likewise is it hard to change things which are of
such and such a nature from that nature, and to restrain them; for it has been well said by some
one, everything is vain if nature sets herself against it.

About man.

From the Questions arising in Genesis.

Page 748. A. What is the meaning of the expression, ``Until''\footnote{\#ge 3:19.} thou return to the dust
from which thou wast taken? For man was not formed of the dust alone, but also of the divine
Spirit; but since he did not continue in an unchanged condition, he neglected the divine command,
and cutting off that constitution which imitated the heaven from his better part, he made himself
over wholly to the earth; for if he had been a lover of virtue, which is immortal, he would beyond all
question have received heaven for his inheritance, but since what he sought was pleasure, by
means of which the death of the soul is brought upon mankind, he became appropriated to the
earth.

About Adam.

From the Questions arising in Genesis.

Page 748. B. ``And God brought all the animals to Adam, to see what he would call Them;''\footnote{genesis
2:19.} for God does not doubt, but since he has given mind to man, the first born and most excellent
of his creatures, according to which he, being endowed with knowledge, is by nature enabled to
reason; he excites him, as an instructor excites his pupil, to a display of his powers, and he
contemplates the most excellent offspring of his soul. And, again, he visibly by the example of this
man gives an outline of all that is voluntary in us, looking with disfavour on those who affirm that
everything happens through necessity, by which some men must be influenced, he on that account
commanded man to take upon himself the regulation of these things. And this is an employment
peculiarly fitting for man, as being endowed with a very high degree of knowledge and most
surpassing prudence, the giving of names to the animals being suited to him not only as being wise,
but also as being the first nobly born creature.

For it was fitting that he should be the founder of the human race, and also the king of everything
that is born of the earth, and that he should have this as an especial honour of his own, that, as he
was the first who had any acquaintance with the animals, he might also be the first inventor and
pronouncer of their names; for it would have been absurd for them to be left without names, and
subsequently to have names given to them by some younger man, to the honour and glory of the
elder.

And when Adam saw the figure of his wife, as the prophet says, and that it had been produced not
by any connexion, nor out of a woman, as human beings in after times were produced, but that she
was as it were a nature on the borders between these two kinds, like a graft from a shoot of
another vine taken off and grafted into a second one, on which account he says, ``For this cause a
man shall leave his father and his mother, and shall cleave to his wife, and they two shall become
one Flesh;''\footnote{\#ge 2:24.} in saying which he used a most gentle expression, which was at the
same time most perfectly true, meaning that they would be united by sympathy in their griefs and
joys.

From the same book, or else from the last book of the Questions arising in Exodus.

Truly the divine place is inaccessible, and one which is hard to be approached, nor is it given even
to the purest intellect to be able to ascend to such a height as to touch it. It is impossible for human
nature to behold the face of the living God; but the word ``face'' is not used here in its literal
meaning, but it is a metaphorical expression, here intended to manifest the purest and simplest
form of the living God, since man is not recognized more by anything than by his face, according to
his peculiar distinctive qualities and form. For God does not say, ``I am not visible in my nature.'' But
who, in fact, is more visible than he who is the Father of all visible things? And being such with
regard to being seen, I am, says he, seen by no mortal man; and the reason of this is the inability
of the created man to behold him.

And that I may not become prolix while weaving in all kinds of arguments, it is inevitable that God
must first be created (which is not possible), in order for any one to be able to comprehend God.
But if any one dies as to this mortal life, but still lives, having received in exchange a life of
immortality, perhaps he will see what he never saw before.

All the different philosophical sects which have flourished in Greece, and in the countries of the
barbarians, when investigating the secrets of nature, have never been able to arrive at a clear
perception of even the most trivial circumstances; and a clear proof of this assertion may be found
in the disagreements, and dissensions, and contentions of those of each sect who are seeking to
establish their own opinions, and to overthrow those of their adversaries. And the households of
those who have been contending for the predominance of this and that sect, have been the causes
of universal wars, blinding the human mind by their contradictory quarrels, which might otherwise
have been able to see the truth, and fighting hard about what doctrines ought to be abandoned and
what ought to be preserved.

Now he who desires to form to himself a conception of the most excellent of all beings, ought in the
first place to stand firm in his mind, being steadfastly fixed in one opinion, and not varying or
wandering in different directions. And in the next place, he ought to take his stand upon nature, and
upon solid grounds, and to abandon all barren and corruptible things, for if anything of a somewhat
effeminate character approach him, he will be disappointed of his object, and he will be unable,
even if he exert the most acute faculties of sight imaginable, to behold the uncreated God; so that
he will become blind before he sees him, on account of the brilliancy of his beams and the flood of
light which distils therefrom. Do you not see that the power of fire in the case of those who stand at
a measured distance from it affords light to them, but it burns those who approach too near? Take
care that you do not suffer such an injury as this in your mind, and lest an extravagant desire of an
impossible object destroy you.

About those who are governed.

Out of the first book of the Questions in Genesis.

Page 749. E. As pillars support whole houses, so also the power of God supports the whole world,
and the best and most God-loving section of the human race.

Out of the Questions in Genesis.

Page 750. C. If any one is either in any house, or village, or city, or nation, who is a lover of
wisdom, it is absolutely inevitable that that house or city should be the better for his existence in it,
for a virtuous man is a common good to all men, bestowing on them advantages proceeding from
himself as from a prepared store.

About people who carry news, and act as intermediate bearers of answers.

From the Questions arising in Exodus.

Page 751. B. The influx of evils agitates and disturbs the soul, enveloping it in a giddiness which
darkens its perceptions, and compels it to suffer that power of sight which by nature was
preeminent, but which by habit has become blinded, to be obscured.

Page 751. B. There is nothing so opposite to and inconsistent with the most holy powers of God as
injustice.

About the sinner and offender.

From the Questions arising in Genesis.

Page 751. C. Never to err in any point whatever is the greatest blessing; but when one has erred,
to repent is next akin to it, as a younger good, if one may say so, by the side of an elder, for there
are some persons who exult in the offences which they have committed as if they had done good
actions, though they are in reality afflicted with a disease difficult to be cured, or I should rather say
incurable.

About its being impossible to escape from God.

From the last book of Questions arising in Exodus.

Page 752. A. He contains all things, while yet he is himself contained by nothing; for as place is
that which contains bodies, and that to which they flee for refuge, so also the divine reason
contains the universe and is that which has completed it.

About truth and faithful evidence.

From the second book of the Questions in Exodus.

Page 754. C. By some lawgivers the practice of giving hearsay evidence has been forbidden, on
the ground that the truth is established by the eyesight, but falsehood by hearing.

About quiet and ease.

From the fourth book of the Questions in Genesis.

The wise man is desirous of peace and leisure, that he may have time for meditation on heavenly
things.

From the fifth book.

Page 754. E. For thus the lover of wisdom never unites with any rash person, even though he may
be closely united to him by blood; nor does he ever consent to dwell with a wicked man, being
separated from the multitude by his reasoning powers, on account of which he is said not to be a
fellow voyager, or a fellow citizen, or a companion of such men.

Page 754. E. The wise man is a sojourner and a settler, having come as an emigrant from a life of
confusion and disorder to one suitable to peaceful and happy men.

About the fearful expulsion.

From the first book of the Questions in Genesis.

Page 772. B. But the essence of the angels is spiritual, but they are very often made to resemble
the appearance of men, being transformed on any emergencies which arise.

From the second book of the same Questions.

Page 772. B. All the powers of God are winged, being always eager and striving for the higher path
which leads to the Father.

About heretics.

From the first book of the Questions in Exodus.

Page 774. B. All those who have stumbled, being unable to proceed with upright feet, go on slowly,
being fatigued a long time before they come to their journey's end; so also the soul is hindered
from proceeding successfully on the path which leads to piety if it has previously fallen in with any
of the byroads of wickedness, for they are great hindrances to it, and the causes of its stumbling,
by means of which the mind becoming lame, proceeds too slowly on the road, according to nature;
and this road, according to nature, is that which ends at the Father of the universe.

From the same book.

The contentious investigations which men enter into about the virtues of God, improve the intellect
and train it in most pleasant labours, which are also most beneficial to it, and especially when men
do not (as those of the present day do) disguise themselves under a false appellation, and contend
for the doctrines in appearance only, but do, in an honest and true heart, seek out truth in
connection with knowledge.

From the second book of the same treatise.

... not being more anxious to display melody and harmony in their voices than in their minds; the
eloquence of the wise man does not display its beauty in words only, but in the matters which it
proves by its words.

From the last book of the Questions in Exodus.

Those men who apply themselves to the study of the holy scriptures ought not to cavil and quibble
at syllables, but ought first to look at the spirit and meaning of the nouns and verbs used, and at the
occasions on which and the manners in which each expression is used; for it often happens that
the same expressions are applied to different things at different times; and, on the contrary,
opposite expressions are at different times applied to the same thing with perfect consistency.

From the Questions in Genesis.

Those men act absurdly who judge of the whole from a part, instead of, on the contrary, forming
their estimate of a part from their knowledge of the whole; for this is the more proper way to form
one's opinion of anything, whether it be a body or a doctrine; therefore the divine code of laws is, in
a manner, a united creature, which one must regard in all its parts and members at once with all
one's eyes, and one must contemplate the meaning and sense of the whole scripture with
accuracy and clearness, not disturbing its harmony nor dissevering its unity; for the parts will have
a very different appearance and character if they are once deprived of their union.

From the fourth book of the same treatise.

Let there then be a law against all those who profess to look on what is venerable and divine, in
any other than a respectful and holy spirit, inflicting punishment on their blindness.

From the second book of the Questions in Exodus.

Page 775. There is nothing either more pleasant or more deserving of respect than to serve God,
whose power is superior to that of the mightiest sovereign; and it appears to me that the greatest
kings have also been chief priests, showing, by their actions, that it is right for those who are the
masters of other men nevertheless to serve as servants of God.

About a king not being greatly respected.

From the first book of the Questions in Genesis.

Page 775. E. No foolish man is a king even though he be invested with supreme power by sea and
land, but he only is a king who is a virtuous and God-loving man, even though he may be deprived
of those supplies and revenues, by means of which kings in general are strengthened in their
sovereignty; for as a rudder, or a collection of drugs, or a flute, or a harp, are all superfluities to a
man who has no knowledge of the art of steering, or medicine, or music, because he is not able to
employ any one of them to the purpose for which it is made, while they may be said to be
excellently adapted to and to be very seasonable for a pilot, or a physician, or a musician; so also,
since kingcraft is an art, and the best of arts, we must look upon him who does not know how to
exert it as a private individual; and as the man who does know how to exert it well as the only king.

About the stable and unstable man.

From the Questions in Genesis.

Page 776. E. A facility of change must of necessity belong to man, by reason of the unsteadiness
of external circumstances. Accordingly we thus oftentimes, after we have chosen friends, and
have associated with them for some time, though we have nothing to accuse them of, turn away
from them with aversion as enemies.

About those who change their minds and blame themselves.

Page 776. E. These are the words of Philo:-

Gaius, as he was ignorant of the greatness of the cause, that he should never fall into death,
suffered a more simple punishment; but his imitator, not being able to take refuge in the plea of
ignorance, is subjected to a double punishment; on which account Lamech shall be avenged
seventy and seven fold, for the reason above mentioned, according to which he was the second
offender who had not thought fit to take warning from the punishment of him who had offended
before, and he clearly receives his punishment, being a more simple one; as in numbers the units
have a highly multiplied power, resembling that of the decades, such as now Lamech, changing his
mind, denounces against himself.

From the same book of the same author.

Page 777. To be aware of what one has done amiss, and to blame one's self, is the part of a
righteous man; but to be insensible to such things causes still more grievous evils to the soul, and
the conduct of wicked men.

About the courage of a woman.

From Philo, from the Questions arising in Exodus.

Page 777. B. It is said by men who have applied themselves to the study of natural philosophy, that
the female is nothing else but an imperfect male.

About the oracles of God.

The words of Philo, out of the second book of his Questions arising in Genesis.

Page 782. A. It is not lawful to divulge the sacred mysteries to the uninitiated until they are purified
by a perfect purification; for the man who is not initiated, or who is of moderate capacity, being
unable either to hear or to see that nature which is incorporeal and appreciable only by the intellect,
being deceived by the visible sight, will blame what ought not to be blamed. Now, to divulge sacred
mysteries to uninitiated people, is the act of a person who violates the laws of the privileges
belonging to the priesthood.

From the same author.

Page 782. B. It is absurd that there should be a law in cities that it is not lawful to divulge sacred
mysteries to the uninitiated, but that one may speak of the true rites and ceremonies which lead to
piety and holiness to ears full of folly. All men must not partake of all things, nor of all discourses,
above all, of such as are sacred; for those that desire to be admitted to a participation in such
things, ought to have many qualifications beforehand. In the first place, what is the greatest and
most important, they ought to have deep feelings of piety towards the only true and living God, and
correct notions of holiness, avoiding all inextricable errors which perplex so many about images
and statues, and in fact about any erections whatever, and about unlawful ceremonies, or illicit
mysteries.

In the second place, they must be purified with all holy purifications, both in soul and body, as far as
it is allowed by their national laws and customs. In the third place, they must give credible evidence
of their entering into the common joy, so that they may not, after having partaken of the sacred
food, like intemperate youths, be changed by satiety and overabundance, becoming like drunken
men; which is not lawful.

About evil-doers.

The words of Philo, out of the Questions arising in Exodus.

Page 782. D. The man who lives in wickedness, bears about destruction within him, since he has
living with him that which is both treacherous, designing, and hostile to him. For the conscience of
the wicked man is alone a sufficient punishment to him, inflicting cowardice on his soul from its own
inmost feelings, as it feared blows.

From the same author.

Page 782. D. The life of the wicked man is subject to pain and sorrow, and full of fear; and in
everything which it does according to the outward senses, it is mingled with fear and grief.

About monks who break their vows.

The words of Philo, from the Questions arising in Exodus.

Page 784. C. The reasoning of some persons is very rapidly satiated, who, though they have been
borne upwards on wings for a little while, yet do presently return back again; not so much flying
upwards, says Philo, as being dragged down again to the lowest depths of hell. But happy are they
who do not draw back.

From the same author.

Page 784. C. Before now, some persons who have tasted happiness, being very speedily satiated,
after they have given hopes of their being in health, have fallen back into the same disease as
before.

From the same author, out of the Questions arising in Genesis.

Page 784. D. To commit perjury is impious and mischievous.

About good friends.

The words of Philo, out of the first book of the Questions arising in Exodus.

Page 788. I. We ought to look upon those men as our friends who are inclined to assist us, and to
requite our kindnesses with kindness, even if they are destitute of power; for friendship is a thing
which is seen more in moments of necessity, than in a steady conjunction or union of dispositions.
So that in the case of each person who unites with another in an association of friendship, one
may apply the expression of Pythagoras to him, and say, ``A friend is a second I.''

About the mercies of God.

The words of Philo, out of the first book of the Questions arising in Exodus.

Page 789. A. When the fruits of these crops which are raised from seed are in a state of
perfection, they receive the beginnings of the generation of trees in order that the mercies of God
may last for ever, and then that one continually succeeding the other, and connecting ends with
beginnings and beginnings with ends, they may be in reality never ending.

From the second book of the same treatise.

Page 789. A. The mercies of God give us not only what is necessary, but also all such things as
conduce to a more excessive and liberal enjoyment of life.

FRAGMENTS FROM A MONKISH MANUSCRIPT

About man: to show that God when he made him endowed him with free will.

It is said to you, O noble man, who live in obedience to the divine precepts, endeavour with all thy
might not only to preserve the gifts which you have received unimpaired and unalloyed, but also
think them worthy of all imaginable honour and regard, as being endowed with free will and
independent power, so that he who has committed them to your charge may have no reason to find
fault with you for having neglected to take proper care of them; and the Creator of the world has
entrusted to your care to employ them according to your own deliberative purpose, a soul, and
speech, and the outward senses. Therefore, those men who receive these gifts in a proper spirit,
and who preserve them for him who has bestowed them on them, have kept their intellect carefully
in such a way that it shall never think of anything else than of God and his virtues; and their speech
in such a manner that with unwearied mouth it shall honour the Father of the universe with praises
and hymns; and their outward senses in such a way that after they have represented to
themselves the whole of the world which is perceptible to those senses, namely, the heaven and
the earth, and the natures which are between those two, they may relate what they have been in a
pure and guileless manner to the soul.

About people who are governed.

The words of Philo, from the fourth book of his Allegorical Interpretation of the Sacred Laws.

If you take away their resources of wealth from politicians, you will find nothing left but empty
arrogance devoid of sense, for as long as there is an abundant supply of external good things,
wisdom and presence of mind appear also to attend them, but when that plenty is taken away all
appearance of wisdom is taken away at the same time.

About the best men.

From the same author, in his Treatise on Drunkenness.

Good men, to speak somewhat metaphorically, are of more value than whole nations, since they
support cities and constitutions as buttresses support large houses.

From the same author.

If it depended on wicked men, no city would ever enjoy tranquillity; but states continue free from
seditious troubles on account of the righteousness of one or two men who live in them, whose
virtue is a remedy for the diseases of war, because God, who loves mankind, grants this effect as
a reward to those who are virtuous and honourable, so that they should not only benefit
themselves, but all who are near them.

From the same author.

There is no place upon earth more sacred than the mind of a wise man, while all the virtues hover
around like so many stars.

About things which are uncertain and unknown to us.

The words of Philo.

The comprehension of the future does not belong to the nature of man.

From the same author.

All things are not known to the mortal race.

From the same author.

God alone is acquainted with the ultimate results of things.

About evil report.

Quiet, which is free from danger, is better than words, the object of which is only to give pleasure.

About self-satisfied people, etc.

The words of Philo.

The lawgiver says, ``You shall not do all the things which we will do here this day, \footnote{deuteronomy
12:28.} every one doing that which is pleasant in his own sight,'' by which words he declares as
loudly as possible that there is no evil which may not be produced by selfishness and
self-sufficiency, which must be eradicated from the mind as unholy feelings. Let no one embrace
that which is pleasing to himself rather than that which is agreeable to nature, for the one is found
to be the cause of mischief and the other the cause of benefit.

From the same author.

Those who do everything for their own sake alone practise selfishness, which is the greatest of
evils, which produces unsociability, want of fellowship, unfriendliness, injustice, impiety, for nature
has made man not like those beasts which love solitude, but like the gregarious beasts which live
together like the most sociable of all creatures, that he may live not to himself alone, but also to his
father, and to his mother, and to his brethren, and to his wife, and to his children, and to all his
other relations and friends, and to those of the same borough as himself, and to those of the same
tribe, and to his native country, and to his fellow countrymen, and to all mankind, and moreover to
the different parts of the universe, and to the whole world, and much more to the Father and
Creator of the world, for he must be (if at least he is really endowed with reason) sociable, loving
the world, and loving God, that he may also be beloved by God.

About God being incomprehensible.

From the first book of the Questions arising in Exodus.

There are thousands and thousands, I do not say only of important matters, but also of those which
appear to be most trivial, which escape the human intellect.

From the same author.

No one may so far yield to unreasonable folly as to boast that he has seen the invisible God.

About the doctrine that God has made angels to be guardians of us.

The words of Philo, from the first book of the Questions arising in Genesis.

As pillars support whole houses, so also do the divine powers support the whole world, and that
most excellent and God-loving race of mankind.

About avoiding sin.

From the treatise on the Giants.

I think it absolutely impossible that no part of the soul should become tainted, not even the outer
most and lowest parts of it, even if the man appears to be perfect among men.

About slowness of counsel.

Slow counsel is profitless, and change of purpose in extremities is mischievous.

About heretical teachers, etc.

From the same book.

A teacher of a good and virtuous disposition, even if he sees his pupils at first stiff-necked by
nature, does not despair of producing in them a change for the better; but, like a good physician,
he does not apply a remedy at once at the first moment of the disease attacking the patient, but he
gives nature time that it may recede a little, so that he may first make ready the path to safety, and
then apply healthful and salutary remedies. And in the same manner does the virtuous man apply
the arguments and doctrines of philosophy.

If, when a pupil is first introduced to you, and first comes to learn of you, you hasten to eradicate all
his ignorance at once, and attempt to introduce every kind of knowledge in a lump, you will produce
the contrary effect to that which you desire, for it will not be likely that such an eradication, having
taken place all in a moment, will continue effectual, nor that the pupil will be able at once to contain
such an abundant influx and overflow of instruction; but being exceedingly perplexed and troubled,
he will resist both these operations, that of eradicating one thing and that of introducing another;
but the system of taking away his ignorance with gentleness and moderation, and of, in the same
manner, gently instilling wisdom into the mind, will be the causes of admitted advantage.

About people who meditate and design mischief.

The words of Philo, from his treatise on Things Improperly Named.

The ordinary production or wickedness enslaves the mind, even if it has not as yet produced any
perfect fruit; for it is, as the proverb says, washing a brick, or taking up water in a net, to try and
eradicate wickedness out of the soul of man. For ``behold,'' says Moses, ``with what designs the
minds of all men are Impressed.''\footnote{\#ge 8:21.} And he speaks truly, for he does not say, what
designs are attached to and adapted to it, but that which has been considered with care and
deliberation is also explained with accuracy, and this too not slowly and with difficulty, but from
man's earliest youth, or as one may almost say, from his very cradle, as if it were a part of him,
kept in continual exercise.

About cowardly and wavering people.

Those who are unmanly from an innate effeminacy, falling down of their own accord before they
meet with any opposition, are a disgrace and ridicule to themselves.

From the same author.

Wickedness in a foolish man has a twin offspring, for the foolish man is wavering and hesitating,
mingling considerations together which ought not to be mingled, and humbling and confusing what
ought to be kept distinct, having as many colours in his soul as a viper has in his body, and polluting
even his sound thoughts with those which cause trouble and death.

From the same author.

The thoughts of a bad man are one thing, and his words another, and his actions indeed are many,
but they are all inconsistent and at variance one with another, for he does not say what he thinks,
and he has decided on the contrary of what he affirms, and he does things which are not
consistent with his original designs, so that, to speak truly, one may say that the life of the wicked
man is a life of enmity.

About distinctness.

The words of Philo.

That which is not distinct is unsuited to a free man, being the most shameful product of folly and
haughtiness; for as distinctness in everything that is to be done is a mark of acuteness and
wisdom, and deserves honour and praise, so also an absence of shame is a sign of folly and
infamy, on which account the other definition which you disregard, classifies a man who is afflicted
with this disease thus, saying, he is impious who does not know how to respect the face of an
honourable man, nor to rise up in the presence of an elder, \footnote{juvenal speaks of this as a custom of the
ancient Romans.} nor to guide his own steps in the right way.

About those who serve God.

The servants of virtuous men submit to voluntary obedience to God, for they are not servants to
human caprices, but to wise men; and he who is the servant of wisdom may justly be said to be
also the servant of God.

About just men.

The words of Philo.

An irreconcileable and endless war is carried on by the atheists against the godly, so that they
threaten them even with slavery.

About justice.

The words of Philo.

Justice, above all things, conduces to the safety both of mankind and of the parts of the world,
earth and heaven.

About the judgments of God.

From the same author.

It is good to begin every day with divine and holy employments, and after that to proceed to the
necessary duties of life. On this account God has commanded Us\footnote{\#de 6:7.} to take care to
obey his commandments, and especially at the first moment of the dawn, as soon as we are risen,
to pay our adoration to Him, that their offerings to God may precede every human occupation,
having the recollection of God for their prompter and leader.

From the same author.

Every soul which piety fertilises with its own mysteries is necessarily awake for all holy services,
and eager for the contemplation of those things which are worth being seen, for this is the feeling
of the soul at the great festival, and this is the true season of joy.

About the difference between God and man.

The words of Philo.

The things of creation are far removed from the uncreated God, even though they are brought into
close proximity following the attractive mercies of the Saviour.

About bold and brave men.

The words of Philo, from his treatise about the Giants.

It is a sign of courage not to be easily alarmed by the terrors of death, and to be full of cheerful
confidence in dangers, and to be of valiant hardiness amid disasters, and to prefer dying with
honour to being saved disgracefully, and to wish to be the cause of victory; and a happy boldness,
and a cheerfulness of soul, and fortitude, are the attendants on a manly spirit.

About equality.

The words of Philo.

As an equality of measurement is the cause of the most perfect blessings, so also a want of
measure is the cause of the greatest evils, as it dissolves that most useful bond of equality.

About drunkenness.

From the same author.

Inequality is a grievous thing and the cause of differences, just as equality is free from all
annoyances and contributes to unite men for advantageous ends.

From the same author.

Obedience to the law and equality are the seeds of peace, and the causes of safety and continued
durability; but inequality and covetousness are excitements to war, and dissolvers of all existing
things.

About evil-doers.

The words of Philo.

Those things which chastise the first, are, if men are wise, preventatives of the second.

About the eye and sight.

The words of Philo, from the treatise about the Creation of the World.

The outward senses resemble windows; for through them, as through windows, the comprehension
of the objects of the outward senses enters into the mind, and again through them the mind goes
out to investigate such objects. But the sight is a part of these windows, that is to say, of the
outward senses, since above all others it is akin to the soul, because it is nearly connected with the
most beautiful of all things, namely light, and is a servant of divine things; and, indeed, that is the
sense which first opened the way to philosophy. For when the eye had beheld the motions of the
sun and moon, and the periodical revolutions of the stars, and the unvarying motions of the whole
host of heaven, and the indescribable order and harmony of the whole universe, and the one
unerring Creator of the world, it then related what it had seen to reason, as having the supreme
authority; and reason, having beheld with a still more acutely piercing eye both these things, and
things of a still more sublime character in their appearance and species, and the great cause of all
things, it then immediately arrived at a due conception of God, and of creation, and of providence;
considering that the whole nature of all things was not brought into existence of its own accord, but
that of necessity it had a creator, and a father, and a guide, and a governor, who also created it,
and who also preserves everything which he has created.

About contentment.

The words of the same author.

If you have a great deal of wealth, take care and do not be carried away by its overflow; but
endeavour to take hold of some dry ground, in order to establish your mind with proper firmness;
and this will be the proper exertion of justice and fairness. And if you should have abundant
supplies of all the things requisite for the indulgence of those passions which lie beneath the belly,
be not carried away by such plenty, but oppose to them a saving degree of contentedness, taking
in this way dry ground to stand upon instead of an absorbing quicksand.

By the same author.

One should practise being contented with a little, for this is being near God; but the contrary habit is
being very far from him.

About faith in and piety towards God.

The words of Philo.

What can be a real sacrifice except the piety of a soul devoted to the love of God? whose grateful
feelings are made immortal by God, having conferred on them an immortal duration like that of the
sun and the moon, and the whole world.

About wicked and impious men.

From the same author.

The hopes of wicked men are unstable, as they expect a good fate, but suffer a contrary destiny of
which they are worthy.

About a bad conscience.

The words of Philo, from his treatise on Men and Things which are Improperly Named.

Who is there who does wrong who is not convicted by his own conscience as if he were in a court
of justice, even though no man correct him?

About advisers.

The words of Philo, from the Questions in Genesis.

Since the mind of those who have not studied philosophy is blind with respect to many of the
circumstances of life, one must take those who do see the character of affairs for one's guides.

About hasty talkers.

The words of Philo.

He who has not shame or fear for his companions, has an unbridled mouth and a licentious
tongue.

About perfection.

The words of Philo.

Perfection and an absence of deficiency are found in God alone. But deficiency and imperfection
exist in every man. For man is taught, even if he be the wisest of his race, by some other man, and
he knows nothing without being taught by his own nature. And if one man has more knowledge than
another, still he has it not naturally, but because of instruction which he has received.

About those who think lowlily of themselves.

The words of Philo.

These things are proved to be most completely natural, that the descent of the soul is its elation by
means of self-conceit, and that its ascent and elevation is its return from arrogance.

From the same author.

It is desirable to eradicate self-conceit, which is the friend of endurance, and prudence, and
justice, \footnote{it is evident that there is great corruption in this and the next sentence.} and also to destroy
overbearing pride; for it is no small proof and exercise of folly to study virtue in an illegitimate
manner.

From the same author.

If you are puffed up by glory and authority so as to desire great things, then remember, like the
skilful pilot of a ship, to take in your sails, that you may not be carried away into absurd conduct.

About sleep.

The words of Philo.

Sleep, according to the prophet, is a trance, not indeed in accordance with insanity, but proceeding
from a relaxation of the outward senses and the retreat of reason; for at that time the outward
senses cease from attaching themselves to their proper objects, and the mind is quiet, neither
being any longer under the influence nor affording any motion to them, and they, being in
consequence cut off from any energy because they are separated from the objects which are
perceptible to them, are dissolved in a state of motionless inactivity.

From the same author.

Very naturally some who have been wise enough to arrive at correct notions of the truth, have
described sleep as a thing to teach us to meditate upon death, and a shadow and outline of the
resurrection which is hereafter to follow, for it bears in itself visible images of both conditions, for it
removes the same man from his state of perfection and brings him back to it.

About promises, etc.

The words of Philo.

It is better absolutely never to make any promise at all than not to assist another willingly, for no
blame attaches to the one, but great dislike on the part of those who are less powerful, and intense
hatred and long enduring punishment from those who are more powerful, is the result of the other
line of conduct.

About haughty men, etc.

From the first book of the Sacred Allegory of the Holy Laws.

Some persons say that the last thing which the wise man puts off is the tunic of vain glory, for even
if a man gets the mastery over his other passions, still he is inclined by nature to be influenced by
glory and the praises of the multitude.

From the same author.

Self-conceit is an impure thing by nature.

About promises, etc.

The words of Philo.

To give thanks to God is intrinsically right, but not to do so to him in the first place, and not to begin
with the first reasons for gratitude, is blameable, for it is not right to give the chief honour to the
creation, and the inferior honour to God, who is the giver of all things in the creation; and indeed
that is a most culpable division, inasmuch as it is laying down a certain disorder of order.

About envy.

The words of Philo.

Envy naturally attaches itself to whatever is great.

About industrious people.

The words of the same author.

The most perfect and greatest of all good things are usually the result of laborious exercise and
energetic vigorous labour.

From the same author.

It is absurd for a man who is in the pursuit of honours to flee from labours by which honours are
acquired.

About the soul and the mind.

From the same author.

What is the meaning of the expression, ``You shall not eat the flesh in the blood of the
Soul?''\footnote{\#ge 9:4.} God appears by this expression to intend to show that the blood is the
essence of the soul, that is to say, of the soul endowed with the outward senses, not the soul
spoken of in the most excellent sense of the word, that is to say, as far as it is endued with reason
and intellect; for there are three divisions of the soul, one part being nutritious, a second being
endued with the outward senses, and the third being endued with reason. Accordingly the divine
Spirit is the essence of the rational portion, according to the sacred historian of the creation of the
world, for he says that ``God breathed into his face the breath of Life.''\footnote{\#ge 2:9.} But of that part
which is endued with the outward senses, and which has the revivifying power, blood is the
essence, for he says in another place that ``the soul of all flesh is the Blood;''\footnote{\#de 12:23.}
but what is connected with the flesh is the outward sense and the passions, and not the mind and
the intellect; not but what that expression, ``in the blood of the soul,'' also indicates that the soul is
one thing and the body another. So that in real truth the breath is the essence of the soul, but it has
not any place of itself independently of the blood, but it resembles and is combined with blood.

About the assistance of God.

The words of Philo, from the fourth book of his treatise on the Allegories contained in the Sacred

Laws.

The extremity of happiness is the assistance of God, for there can be no such thing as want when
God gives his aid.

About the creation of the world.

From the same author, from the first book of the Questions arising in Genesis.

It is impossible that the harmony, and arrangement, and reason, and analogy, and that all the great
accord and real happiness which we see existing in the world can have been originated by
themselves, for it follows inevitably that these things must have had a creator, and a father, and a
regulator and governor, who generated them in the first place, and who now preserves what he has
generated.

About the church of God.

From the same author.

God wishing to send down from heaven to the earth an image of his divine virtue, out of his
compassion for our race, that it might not be destitute of a more excellent portion, and that he
might thus wash off the pollutions which defile our miserable existence, so full of all dishonour,
established his church among us.

About seeking God.

From the same author, from the last book of the Questions arising in Exodus.

The one most powerful relaxation of the soul leads to the sacred love of the one living God,
teaching mankind to take God as its guide in all their plans, and words, and actions.

From the same author.

The extremity of happiness is to rest unchangeably and immovably on God alone.

About the last day.

The words of Philo, from the second book of the Questions arising in Exodus.

The stars are turned round and revolve in a regular circle, some proceeding on in the same
manner through the whole heaven, and others have special eccentric motions of their own.

About the detestation of wickedness felt by God.

The words of Philo, from the second book of the Questions arising in Exodus.

Some men think that repentance appears at times to take possession of God on account of the
oaths which he has sworn, but they do not form correct notions; for apart from the fact that the
Deity does not change, neither the expression, ``God repented,'' nor that, ``And it grieved him at the
Heart,''\footnote{\#ge 6:6.} is indicative of repentance, for the Deity is unchangeable; but they only show
the character of the pure intellect which is now deeply meditating on the cause for which he
created man upon the earth.

By the same author, from the same book.

There is no hesitation and no envy in God; but he often uses expressions indicative of hesitation
or of uncertainty from a reference to man, who is susceptible of such feelings; for as I have often
said, there are altogether two supreme sources; in the one case God does not speak as man
speaks, in the other he instructs man as a man instructs his son, the former being a sign of his
power, the second of the way in which he teaches and guides man.

About promises.

The words of Philo, from the last book of the Questions arising in Exodus.

He who does not offer to God first fruits of his own free will does not really offer first fruits at all,
even if he brings everything which is great, with a most royal abundance of treasure; for the real
first fruits consist not in the things offered, but in the pious disposition of him who offers them.

About the mildness of God and his love for mankind.

The words of Philo, from the Questions arising in Exodus.

The mercies of God do always outstrip justice, for the work which he has chosen for himself is that
of doing good, and the task of punishing follows that; and it is common, when great evils are about
to arise, for an abundance of great and numerous blessings to happen first.

FRAGMENTS PRESERVED BY ANTONIUS

SER. I.

The virtues alone know how to regulate the affairs of men.

The contemplation of virtue is exceedingly beautiful, and actions according to it, and the exercise
of it, are desirable above all things.

SER. II.

If you wish to have a good reputation in a twofold manner, then honour exceedingly those who are
doing well, and reprove those who are doing ill.

SER. VIII.

When you are entreated to pardon offences, pardon willingly those who have offended against
you, because indulgence given in requital for indulgence, and reconciliation with our fellow
servants, is a means of averting the divine anger.

SER. IX.

The virtuous man is a lover of his race, and he is merciful and inclined to pardon, and never bears
ill will towards any man whatever, but thinks it right to surpass in doing good rather than in injuring.

What is beautiful is then beautiful, when a man has no need of the assistance of another, but when
he contains in himself all the signs of excellence as his own.

SER. X.

It is well that the worse should always follow the better, on account of the hope of improvement.

SER. XI.

One ought to call a city, and a country, and a house, happy, when they contain a virtuous man; and
one ought to call those miserable, when they have no such man within them.

SER. XVI.

Those who are tyrannical in their natures, but without power, make their designs succeed by
treachery.

SER. XX.

The friendships of the wicked are mischievous, and very often the soul of such men, being
influenced by such associations, takes the impressions of downright insanity.

It is not the country which makes men bad, or the city which makes them good, but the habits of
living with such and such men.

SER. XXVIII.

One need not dread the blow of a weak man, nor the threat of a fool.

Light-minded men, like empty vessels, may easily be taken and moved by their ears.

SER. XXX.

Nothing that is done can be beautiful without scientific contemplation, for knowledge is the offspring
of counsel, but folly is the source of all evils.

Every argument on behalf of justice is superfluous, when those who listen are unanimous in a bad
object.

SER. XXXVIII.

The wicked man disturbs the city, and is eager for the confusion and the disorder of all men and all
things within the city; for a desire of interference, and covetousness, and the acts of a
demagogue, and the influence with the populace, are looked upon as honours by such a man, and
quiet he looks upon with disdain.

Excellence is a thing difficult to find, or rather is absolutely undiscoverable in a troubled life.

SER. XLIII.

There is nothing so calculated to cause good will as kind words, on account of good actions.

SER. XLVII.

It is sufficient not to bear witness one's self, but that which stands in need of the advocacy of
another is inadequate to bring conciliation to the mind.

SER. LII.

Reject with aversion the deceitful words of flatterers, for they, obscuring reason, do not contribute
to the truth of things; for either they praise actions which are deserving of blame, or else they often
blame things beyond all praise.

SER. LVI.

Peace is the greatest blessing which no man is able to afford, since this is a divine action.

SER. LVII.

Behave to your servants in the same manner in which you desire that God should behave to you;
for as we hear them we shall be heard by him, and as we regard them we shall be regarded by him.
Let us therefore let our compassion outrun compassion, that we may receive a like requital from
him for our mercy to them.

SER. LXIX.

How great a relief of nature is sleep, it is the image of death, and the rest of the outward senses.

Sleep is one thing only, but the desire of it has many reasons and causes; I mean from nature,
from food, from fate, and perhaps also from excessive and intense fasting, by means of which the
flesh, becoming unnerved and deprived of strength, wishes to recover itself for subsequent actions
by means of sleep.

As much drinking is called a habit, so is much sleep, and it is difficult to get rid of an inveterate
habit.

SER. LXXIV.

Pardon is apt to engender repentance.

SER. LXXIX.

Shamelessness is the characteristic of a worthless man, and modesty of a virtuous man, but never
to feel either ashamed or bold is a mark of one who is slow of comprehension, and who is without
the power of giving assent.

SER. LXXXII.

Since God penetrates invisibly in the region of the soul, let us prepare that region in the best
manner that we are able to, or rather that it may be a habitation fit for God, otherwise, without our
being aware of it, God will depart and remove to some other abode.

The mind of a wise man is the house of God, and he is called, in an especial manner, the God of all
mankind, as the prophet says when speaking of the mind of a wise man, he calls it ``that in which
God Walks,''\footnote{\#le 26:12.} as in a palace.

What is visible and actually before us is comprehended by the eyes, but the pure faculty of reason
extends even to what is unseen and future.

SER. LXXXVII.

God who is merciful by nature will never exonerate from guilt the man who swears falsely for an
unrighteous object, as such a man is impure and defiled, even though he may escape the
punishments inflicted by men.

SER. XCIX.

Those things which are kept in the dark for a while by envy, are at last released and brought to
light.

SER. CIV.

In his essential character a king is equal to every man, but in the power of his authority and rank he
is equal to God who ruleth over all things; for there is nothing on earth that is higher than he.
Therefore it becomes him as being a mortal not to be too much elated, and as being a kind of God
not to yield too much to passion; for if he is honoured as being of the likeness of God,
nevertheless he is in some degree entangled in terrestrial and vile dust, by means of which he
should learn simplicity and meekness towards all men.

SER. CXVI.

A severe master is best for untractable and foolish servants; for they, fearing his threats and
punishments, though against their will, are made to do right by fear.

SER. CXVIII.

It is the greatest praise of a servant to neglect nothing which his master commands, but to attempt
with an honest heart to perform in a proper and successful manner, even if it be beyond his power,
all that is commanded him with energy and without hesitation.

SER. CXXIII.

When once the wife of Philo was asked in an assembly of many women why she alone of all her
sex did not wear any golden ornaments, she replied: ``The virtue of a husband is a sufficient
ornament for his wife.''

SER. CXXX.

The virtues of children are the glory of their fathers.

Those who are well acquainted with what is honourable and virtuous, are happy in their children.

SER. CXXXV.

To drink poison out of a golden goblet, and to take advice from a foolish friend, is the same thing.

New vessels are better than old ones, but old friendship is better than new.

The fruits produced by the earth come once a year; but those which we derive from friendship are
to be gathered on every occasion. Many men select for their friends not those who are the most
virtuous, but those who are rich.

Many who appear to be friends are not so, and many who do not appear to be such are so in
reality; but it is the part of a wise man to discern both these classes.

SER. CLII.

Youth which is not willing to work is laying up misfortunes for old age.

SER. CLVI.

What is bad is, not being punished here, but being worthy of punishment hereafter.

SER. CXXXV.

God has implanted hope in the human race that, having a comfort innate in them, those who have
committed errors which are not irremediable may feel their sorrows lightened.

SER. CLXXXII.

Pleasure appears to be an equable kind of motion, but in reality it both is and is found to be rough.

THE FOLLOWING FRAGMENTS ARE FROM AN ANONYMOUS COLLECTION IN THE

BODLEIAN LIBRARY AT OXFORD

EXTRACTS FROM PHILO

About friends.

A steadiness towards one's friends is a sign of a general stability of disposition, on which account one
ought not to form friendship till one has carefully tested the characters of those with whom he proposes
to form it; for not only is the forming of such friendship pleasant, but so also is the feeling that one has
not to bear by one's self burdens which oppress the soul, and not to depart from the association; for he
who is the cause of differences in friendship is not known to the generality of men, but he is accustomed
to bring common blame upon both parties, and very commonly on the innocent party more than on the
guilty one.

Of secret things, you may share with mean persons those which increase your virtue; but as to those
which deteriorate your mind, you must not pursue them yourself, nor impute them to your friends.

The life of man is like a sea, it is liable to every description of agitation and change, even in the height
of prosperity; for nothing earth-born is firmly established, but all such things are carried about to and
fro, like a vessel which is driven about in the sea by contrary winds.

About sin.

Let us fear not the diseases which come upon us from without, but those offences on which account
diseases come, diseases of the soul rather than of the body.

About pain.

Every foolish man is in a strait, being oppressed by covetousness, and love of glory, and desire of
pleasure, and things of that sort, which do not allow the mind freedom of motion.

About gluttony.

The sons of the physicians have laid it down as a maxim that regularity is the parent of a healthy
condition of the body, paying but little attention to the health of the soul; but we lay it down that
regularity is not only destructive of all diseases of the body, but much more do we recognise the fact
that the truest health is that which destroys the passions which injure the soul.

About custom and familiarity.

An inveterate habit is more powerful than nature, and little things, if they are not hindered, grow up and
increase and become of a large size.

THE FOLLOWING FRAGMENTS ARE FROM AN UNPUBLISHED MANUSCRIPT IN THE LIBRARY

OF THE FRENCH KING

From the works of the Hebrew Philo, on \#Ge 6:7.

Why is it that God, when he threatens to extirpate mankind, does also destroy the irrational animals?
Because the irrational animals were not originally created designedly for their own sakes, but for the
sake of man, and to perform services of which he might be in need; and when man was destroyed it
followed naturally that they should also be destroyed at the same time, when the beings for whose sake
they had been created were no longer in existence.

From the same author, on \#Ge 17:14.

The law does not treat any action done involuntarily as guilty, since it even pardons a man who has
committed murder unintentionally; but if a child is not circumcised eight days after its birth, what evil
has it done so as to be subjected to the punishment of death? Therefore some persons say that the
manner of the punishment is to be referred to its parents, and think that they ought to be punished
severely as having neglected the commandments of the law; and others think that it is by an excess of
indignation that God is here represented as inflicting punishment, as far as appearance goes, on the
child, in order that this inevitable punishment may be inflicted on those people of mature age who have
violated the law.

Not because the action of circumcision is important in itself, but because if that is neglected the
covenant itself is treated with contempt when the seal by which it is recognised and ratified is not made
perfect.

From the same author, on \#Ge 19:23.

Why did the sun go forth upon the land when Lot entered into Segor? And he says the very same place is
a safety for those who are making progress, and a punishment to those who are inwardly wicked. And
again the moment that the sun rises in the beginning of the day it brings with it justice; wishing to show
that the sun, and the day, and the light, and everything else in the world which is beautiful and
honourable, are given only to the virtuous and to no worthless man who embraces incurable
wickedness.

From the same author, on \#Ge 27:24—27.

Having been spies rather than friends under truce, and being prepared for either alternative; for war if
they saw that the other was weak, and for peace if they found him stronger than themselves.

From the same author, on \#Ge 26:28, etc.

These are the covenants which they made, not to be destroyed as the other nations had been, and the
Philistines were at a subsequent period by the Israelites; whom the holy scriptures call sometimes
Canaanites, and sometimes Cappadocians; but afterwards the Cappadocians emigrated.

From the same author, on \#Ge 26:30.

Not on account of praise, for the wise man is not attracted by flattery or by any other kind of
subserviency, but because he has accepted their repentance.

From the same author, on \#Ge 27:6, etc.

When he had two sons, the one good and the other guilty, he says that he will bless the guilty one, not
because he preferred him to the good one, but because he knew that the other one could do right by
himself, but that the other was convicted by his own disposition, and had no hope whatever of salvation
except in the prayers of his father; and if he did not obtain them, then he would be the most miserable of
all men.

From the same author, on \#Ge 27:11, etc.

It is proper here to admire also the good will of his mother, who confessed herself willing to take upon
herself the cause for his sake, in order that her son might have the honour to which the two were
entitled, for she is carried away by her affection for both of them; for she had feared his father, lest she
should be looked upon as imposing on him, and to be filching away the honour to which the other was
entitled; and his mother, lest he should be considered by her as disobedient to her when she urged him
vehemently; on which account he says, with great prudence and propriety, Will not my father curse me?
and I shall be bringing a curse on myself. He had confidence because of the promise of God, which said,
``The elder shall serve the younger.'' But, on the other hand, he feared as a man, lest the blessing of his
father, as a just man, should overturn the assertion of God.

From the same author, on \#Ge 27:30.

He is not so indignant at his disappointment in not obtaining the blessings, as at the fact of his brother
having been thought worthy of them; for being of an envious disposition, he regarded his want of
success as more desirable than even his own advantage, and he shows this by his great and bitter
lamentations, and by his subsequent exclamation, ``Bless me now also, O my father.''

From the same author.

But if he obtained it by fraud, a man will be inclined to say, he was not to be praised. What then does
his father say? ``And he shall be blessed.'' But he appears by what he here says to intimate, in an
enigmatical and obscure manner, that it does not follow that every stratagem is blameable, since
guardians of the night when they lie in wait for robbers, and generals when they form ambuscades for
enemies whom they would not be able to subdue by open force, appear to act rightly: and what are
called stratagems proceed on the same principle as the contests of wrestlers, for in these cases too tricks
are accounted honourable; and those who by trickery get the better of their antagonists are thought
worthy of the prize, and of the crown of victory; so that it is not a charge against a man to say, he has
done a thing by trick, but it is rather a panegyric, being equivalent to saying, he has done it skilfully, for
the virtuous man does not do anything unskilfully.

From the same author, on \#Ex 20:25.

What is the meaning of ``thy dagger,'' and what comes next? Those who by their nature venture to make
improper attempts, and who by their own private endeavours metamorphose the works of nature, defile
what ought not to be defiled, for all the things of nature are perfect and complete, and stand in need of
no addition.

From the same author, on \#Ex 22:19.

He shows most evidently that he is a proselyte, inasmuch as he is not circumcised in the flesh of his
foreskin, but in the pleasures and appetites, and all the other passions of the soul; for the Hebrew race
was not circumcised in Egypt, but being ill-treated with every imaginable circumstance of ill-treatment
by the natural cruelty of the natives of the country to strangers, it nevertheless lived among them with
fortitude and patience, and that no more from compulsion than voluntarily, because of the refuge which
it possessed in God the saviour, who, sending down his beneficent power, delivered his suppliants from
their difficult and apparently inextricable troubles. On this account Moses adds, ``For you know the soul
of a Proselyte.''\footnote{\#ex 23:10.} Now what is the mind of a proselyte? a forsaking of the opinions of
the worshippers of many gods, and a union with those who honour the one God, the Father of the
universe. In the second place, some persons call foreigners also proselytes, and those are
strangers who have come over to the truth in the same manner with those who have been
sojourners in Egypt; for the one are strangers newly arrived in the country, but the last are
strangers also to the customs and laws, but the common name of proselytes is given to both.

From the same author, on \#Ex 22:22.

It is forbidden to injure a widow and orphan, for these are under the protection of the especial
providence of God, since they are deprived of their natural protectors and guardians, for God wills
that those who enjoy natural associations should make amends to the others from their own
abundance of resources.

From the same author, on \#Ex 23:1.

He says that we must not approach folly or falsehood, either with the ears or with any other of the
outward senses, for great injuries are the result of being deceived; on which account some
lawgivers have forbidden any one to give hearsay evidence, since the truth is confirmed by
eyesight, but falsehood by hearing.

From the same author, on \#Ex 23:6.

Poverty by itself claims compassion, in order to correct its deficiencies, but when it comes to
judgment, it then has for the arbitrator the law of equity, for justice is a divine and incorruptible
thing, on which account it is expressly affirmed in another passage that the judgment of God is
Just.\footnote{\#de 32:4.}

From the same author, on \#Ex 23:18.

Instead of saying leavened bread must not come among the things which are offered, but all things
which are brought as a sacrifice or an offering must be unleavened, he intimates two most
necessary things by an obscure and symbolical expression; one being to despise pleasure, for
leaven is the seasoning of food and not food itself; and the other being that it is not right for men to
be elated, because of being puffed up by vain self-conceit; for each is a wicked state, and
pleasure and selfconceit are both the offspring of one mother, deceit.

The blood of the sacrifices is a proof of a soul making its offerings to God; and it is not in
accordance with the divine law that things which will not unite should be mingled together.

From the same author, on \#Ex 23:20.

One must suppose that the angel mentioned a little before indicated the voice of God; for the
prophet is the messenger of the Lord, who is the real speaker; for it is inevitable that he who hears
with his ears, that is to say who firmly receives what is said, must also accomplish what is said to
him by his actions; for an action is the proof of what is said; and he who is obedient to what is said,
and who performs actions corresponding to his orders, must of necessity have him who has
commanded him for his ally and champion, who in appearance indeed brings assistance to his
pupil, but in reality to his own doctrines and commandments, ... which his enemies and adversaries
seek to overthrow.

From the same author, on \#Ex 23:24.

Pillars symbolically mean the doctrines which appear to stand and to be firmly established. Now of
the doctrines established in this firm way, some are good, which ought to be stored up and to be
fixed in a most lasting manner; but others are open to blame, and such it is desirable should be
overthrown. But the expression, ``overthrowing you will overthrow, and destroying you will destroy,''
has such a meaning as the following. Some men pull down some things as if they meant to raise
them again, and destroy some things as if they meant at a future time to re-establish them. But
God wills that what has been once destroyed and pulled down shall never be raised or
re-established again, but shall be utterly destroyed and for ever, as being contrary to what is good
or beautiful.

From the same author, on \#Ex 23:28.

And we ought to consider that the wasps are a sign of unexpected power coming by the divine
mission; which, bringing down its blows from high places so as to reach the extremity of the ear,
takes a good aim with all its strokes, and regulating them well will meet with no failure whatever
itself.

From the same author, on \#Ex 23:31.

These things God announced to them, if they obeyed him and kept his commandments. But when
they were found to be transgressing and disobedient to the divine law, he then contracted his
promise from Dan to Beersheba.

From the same author, on \#Ex 24:9, 10.

The express command as uttered has a subsequent proposition evident, as all were preserved in
safety. But the real meaning is that they all were of one mind in respect of piety and differed in no
good thing.

From the same author, on \#Ex 24:10.

When he speaks of the seventy men he means those with Moses, and Aaron, and Nadab, and
Abihu. And the statement that they did not differ, rather shows that they all equally saw the place
where God had stood, than that nothing was left.

From the same author, on \#Ex 24:13.

He is most manifestly offended with those who being near thought, out of their impiety or folly, that
the motions of the Deity were those of peace, and belonging to the act of changing his abode; for
behold he says expressly, not that the God who exists in essence, and who is duly thought of in
respect of his existence, came down, but that his glory came down. And the acceptation of the
word glory may be twofold; for in one sense it may signify the presence of his powers, since the
power of his army is spoken of as the glory of a king; and in another sense it may refer to the
appearance of him alone, and to the apprehension of his divine glory; so that an idea of the actual
arrival of God may have been created in the minds of those who were present, as if he had come
in order to give a most undeniable information to the laws which were about to be given.

From the same author, on \#Ex 24:17.

But he says that the appearance of the glory of the Lord is very like unto flame, or rather not that it
is so, but that it appears like it to the beholders; since God shows what he chose to appear to be,
in order to strike the beholders with amazement without in reality being what he appeared.
Accordingly he brings him before the face of the children of Israel, affirming in the plainest
language that it was an appearance as of flame, but not a real flame. But as flame consumes
every material which is exposed to it, so also when the true conception of God once enters into the
soul, it destroys all the heterodox reasonings of impiety, and purifies and sanctifies the whole mind.

From the same author, on \#Ex 24:18.

Because the generation which had thus quitted its former abode was about to be condemned, and
to wander in a state of desolation for forty years, having received innumerable benefits, but having
displayed its ingratitude in still more countless instances.
